\documentclass[a4paper,10pt]{article}
\usepackage{amsmath,amssymb}
\usepackage{fontspec}
\setmainfont{Liberation Serif}
\usepackage{graphicx}
\usepackage{eso-pic}
\usepackage{tikz}
\usepackage[most]{tcolorbox}
\usepackage{geometry}
\usepackage{fancyhdr}
\usepackage{lastpage}
\usepackage{booktabs}
\usepackage{listings}
\usepackage{enumitem}
\usepackage{microtype}

\geometry{
    a4paper,
    top=1.5cm,
    bottom=1.8cm,
    left=1.2cm,
    right=1.2cm
}

% Persistent Unremovable Canvas Watermark on Every Page
\newsavebox{\watermarkbox}
\sbox{\watermarkbox}{%
    \begin{tabular}{c}
        \includegraphics[width=8cm]{/home/johan/GATE/GATE_CS_PYQ_DB/assets/studiva-butterfly-icon.jpg} \\[-0.5em]
        {\fontsize{18}{22}\selectfont\bfseries\scshape www.studiva.co.in}
    \end{tabular}%
}
\AddToShipoutPictureBG{%
    \begin{tikzpicture}[remember picture, overlay]
        \node[opacity=0.12] at (current page.center) {\usebox{\watermarkbox}};
    \end{tikzpicture}%
}

% Header and Footer
\pagestyle{fancy}
\fancyhf{}
\fancyhead[L]{\small\scshape GATE 2018--2025 CS/IT $\cdot$ \textbf{Multiple Choice Questions (MCQ)}}
\fancyhead[R]{\small\scshape Practice Question Bank with Solutions}
\fancyfoot[L]{\scriptsize\textsc{Copyright \copyright\ 2026} Studiva\textsuperscript{TM}}
\fancyfoot[C]{\small Page \thepage\ of \pageref{LastPage}}
\fancyfoot[R]{\scriptsize\texttt{www.studiva.co.in}}
\renewcommand{\headrulewidth}{0.4pt}
\renewcommand{\footrulewidth}{0.4pt}

% Code Listing Style (Monochrome)
\lstset{
    basicstyle=\footnotesize\ttfamily,
    numbers=none,
    frame=none,
    tabsize=2,
    breaklines=true,
    showstringspaces=false
}

\begin{document}

% ==================== PAGE 1: IMPRINT & TITLE PAGE ====================
\thispagestyle{empty}
\vspace*{0.5cm}
\begin{center}
    \includegraphics[width=2.5cm]{/home/johan/GATE/GATE_CS_PYQ_DB/assets/studiva-butterfly-icon.jpg}\\[0.6em]
    {\LARGE \textbf{GATE 2018--2025 COMPUTER SCIENCE \& IT}}\\[0.3em]
    {\Large \textbf{PREVIOUS YEAR QUESTION BANK WITH SOLUTIONS}}\\[0.4em]
    {\large \textsc{\textbf{Question Type Wise: Multiple Choice Questions (MCQ)}}}\\[0.2em]
    {\normalsize \textsc{Combined 2018--2025 Examination Papers $\cdot$ Comprehensive Edition ($\mathbf{206}$ Questions)}}\\[1.2em]
    \rule{0.6\textwidth}{0.6pt}
\end{center}

\vspace{1em}
\noindent\begin{tabular}{ll}
\textbf{Author / Compiler:} & Rupam Ghosh \\
\textbf{Publisher \& Platform:} & Studiva\textsuperscript{TM} (\texttt{www.studiva.co.in}) \\
\textbf{Official Inquiries:} & \texttt{contact@studiva.co.in} \\
\textbf{Edition:} & First Standard Edition (2026) \\
\textbf{Coverage:} & Complete GATE CS/IT Papers (2018, 2022, 2023, 2024 CS1, 2024 CS2, 2025 CS1, 2025 CS2) \\
\textbf{Syllabus Reference:} & Official GATE 2026 Syllabus (Organizing Institute: IIT Guwahati) \\
\end{tabular}

\vspace{1.5em}
\begin{tcolorbox}[enhanced, standard jigsaw, opacityback=0, colframe=black!40, arc=1mm, boxrule=0.5pt, left=10pt, right=10pt, top=8pt, bottom=8pt]
\textbf{\large COPYRIGHT \& INTELLECTUAL PROPERTY NOTICE}\\[0.6em]
Copyright \copyright\ 2026 Rupam Ghosh \& Studiva\textsuperscript{TM}. All Rights Reserved.\\[0.5em]
This educational compilation, problem organization, taxonomy classification, LaTeX typesetting, hints, and explanatory solutions are the proprietary intellectual property of \textbf{Rupam Ghosh} and \textbf{Studiva\textsuperscript{TM}}.\\[0.6em]
\textbf{Terms of Use \& Attribution Rules:}
\begin{itemize}[leftmargin=*, nosep]
    \item \textbf{Personal Educational Use:} Permitted freely for individual students and educators preparing for competitive examinations.
    \item \textbf{Mandatory Attribution:} Any reproduction, republication, hosting, or sharing (digital or physical) must prominently credit \textbf{Rupam Ghosh \& Studiva\textsuperscript{TM}} and include an active, direct link to \texttt{https://www.studiva.co.in}.
    \item \textbf{Commercial Restrictions:} Commercial resale, sublicensing, course bundling, or inclusion in paid test series without prior written permission is strictly prohibited.
    \item \textbf{Anti-Tampering Notice:} Removing, altering, or obscuring copyright notices, page watermarks, or website links from this document is strictly prohibited.
\end{itemize}
\vspace{0.6em}
\textbf{Fair Use Disclaimer:}\\
The underlying examination questions originate from the Graduate Aptitude Test in Engineering (GATE), conducted by the Indian Institutes of Technology (IITs) and Indian Institute of Science (IISc). Question texts are reproduced for non-commercial educational study under Fair Use doctrines. No copyright is claimed over the original GATE examination question texts.
\end{tcolorbox}

\vspace{1.2em}
\begin{center}
\small\itshape
Practice Recommendation: Use an opaque sheet of paper or ruler to cover the right-hand column during mock attempts. Tick your chosen answers directly in the check boxes on the left, then slide the cover to check the solution and hint.
\end{center}

\newpage
% ==================== PAGE 2 ONWARDS: QUESTION BANK ====================

\vspace{8pt}
\subsection*{General Aptitude}
\vspace{-4pt}
\noindent
\begin{tcolorbox}[
    blanker,
    sidebyside,
    sidebyside align=top,
    lefthand width=0.55\textwidth,
    top=4pt, bottom=4pt, left=0pt, right=0pt,
    middle=16pt,
    underlay={
        \draw[dashed, black!80, line width=0.8pt, dash pattern=on 4pt off 3pt]
            ([xshift=0.55\textwidth+8pt]frame.north west) --
            ([xshift=0.55\textwidth+8pt]frame.south west);
    }
]
\textbf{Q.1} \quad {\footnotesize\textsc{[MCQ $\cdot$ 1 Mark(s) $\cdot$ Neg: 0.33]}} \hfill {\footnotesize\textbf{GATE 2018 CS} \quad \tiny\texttt{2018\_CS\_GA\_Q01}}
\vspace{2pt}\hrule height 0.3pt\vspace{4pt}

``From where are they bringing their books? \_\_\_\_\_\_\_\_ bringing \_\_\_\_\_\_\_ books from \_\_\_\_\_.''
\par\medskip
The words that best fill the blanks in the above sentence are

\par\vspace{6pt}\noindent
\begin{tabular}{@{}p{0.48\linewidth}@{\hspace{0.04\linewidth}}p{0.48\linewidth}@{}}
$\square$\ \textbf{(A)} Their, they're, there & $\square$\ \textbf{(B)} They're, their, there \\
\addlinespace[3pt]
$\square$\ \textbf{(C)} There, their, they're & $\square$\ \textbf{(D)} They're, there, there
\end{tabular}

\tcblower

\textbf{Official Key:} \textbf{\large B} \hfill \textit{\small [Easy]}
\vspace{2pt}\hrule height 0.3pt\vspace{3pt}
\textbf{Topic:} Verbal Aptitude \\
\textbf{Subtopic:} \textit{Vocabulary: words, idioms, and phrases in context}
\vspace{2pt}\hrule height 0.3pt\vspace{3pt}
\textbf{Guide / Hint:} \\
{\small "They're" is the contraction for "they are", "their" is the possessive pronoun, and "there" refers to a place. Thus, the complete sentence is: "They're bringing their books from there."}

\end{tcolorbox}
\vspace{6pt}\hrule height 0.35pt\vspace{6pt}
\noindent
\begin{tcolorbox}[
    blanker,
    sidebyside,
    sidebyside align=top,
    lefthand width=0.55\textwidth,
    top=4pt, bottom=4pt, left=0pt, right=0pt,
    middle=16pt,
    underlay={
        \draw[dashed, black!80, line width=0.8pt, dash pattern=on 4pt off 3pt]
            ([xshift=0.55\textwidth+8pt]frame.north west) --
            ([xshift=0.55\textwidth+8pt]frame.south west);
    }
]
\textbf{Q.2} \quad {\footnotesize\textsc{[MCQ $\cdot$ 1 Mark(s) $\cdot$ Neg: 0.33]}} \hfill {\footnotesize\textbf{GATE 2018 CS} \quad \tiny\texttt{2018\_CS\_GA\_Q02}}
\vspace{2pt}\hrule height 0.3pt\vspace{4pt}

``A \_\_\_\_\_\_\_\_\_ investigation can sometimes yield new facts, but typically organized ones are more successful.''
\par\medskip
The word that best fills the blank in the above sentence is

\par\vspace{6pt}\noindent
\begin{tabular}{@{}p{0.48\linewidth}@{\hspace{0.04\linewidth}}p{0.48\linewidth}@{}}
$\square$\ \textbf{(A)} meandering & $\square$\ \textbf{(B)} timely \\
\addlinespace[3pt]
$\square$\ \textbf{(C)} consistent & $\square$\ \textbf{(D)} systematic
\end{tabular}

\tcblower

\textbf{Official Key:} \textbf{\large A} \hfill \textit{\small [Easy]}
\vspace{2pt}\hrule height 0.3pt\vspace{3pt}
\textbf{Topic:} Verbal Aptitude \\
\textbf{Subtopic:} \textit{Vocabulary: words, idioms, and phrases in context}
\vspace{2pt}\hrule height 0.3pt\vspace{3pt}
\textbf{Guide / Hint:} \\
{\small The sentence contrasts an erratic or wandering investigation with an "organized" one. "Meandering" means proceeding in a winding or circuitous course without a clear direction, which fits the context precisely.}

\end{tcolorbox}
\vspace{6pt}\hrule height 0.35pt\vspace{6pt}
\noindent
\begin{tcolorbox}[
    blanker,
    sidebyside,
    sidebyside align=top,
    lefthand width=0.55\textwidth,
    top=4pt, bottom=4pt, left=0pt, right=0pt,
    middle=16pt,
    underlay={
        \draw[dashed, black!80, line width=0.8pt, dash pattern=on 4pt off 3pt]
            ([xshift=0.55\textwidth+8pt]frame.north west) --
            ([xshift=0.55\textwidth+8pt]frame.south west);
    }
]
\textbf{Q.3} \quad {\footnotesize\textsc{[MCQ $\cdot$ 1 Mark(s) $\cdot$ Neg: 0.33]}} \hfill {\footnotesize\textbf{GATE 2018 CS} \quad \tiny\texttt{2018\_CS\_GA\_Q03}}
\vspace{2pt}\hrule height 0.3pt\vspace{4pt}

The area of a square is $d$. What is the area of the circle which has the diagonal of the square as its diameter?

\par\vspace{6pt}\noindent
\begin{tabular}{@{}p{0.48\linewidth}@{\hspace{0.04\linewidth}}p{0.48\linewidth}@{}}
$\square$\ \textbf{(A)} $\pi d$ & $\square$\ \textbf{(B)} $\pi d^2$ \\
\addlinespace[3pt]
$\square$\ \textbf{(C)} $\frac{1}{4}\pi d^2$ & $\square$\ \textbf{(D)} $\frac{1}{2}\pi d$
\end{tabular}

\tcblower

\textbf{Official Key:} \textbf{\large D} \hfill \textit{\small [Easy]}
\vspace{2pt}\hrule height 0.3pt\vspace{3pt}
\textbf{Topic:} Quantitative Aptitude \\
\textbf{Subtopic:} \textit{Mensuration and geometry}
\vspace{2pt}\hrule height 0.3pt\vspace{3pt}
\textbf{Guide / Hint:} \\
{\small For a square of area $d$, the side length is $s = \sqrt{d}$ and its diagonal is $s\sqrt{2} = \sqrt{2d}$. Since this diagonal is the diameter of the circle, the radius is $r = \frac{\sqrt{2d}}{2} = \sqrt{\frac{d}{2}}$. Hence, the area of the circle is $\pi r^2 = \pi \left(\frac{d}{2}\right) = \frac{1}{2}\pi d$.}

\end{tcolorbox}
\vspace{6pt}\hrule height 0.35pt\vspace{6pt}
\noindent
\begin{tcolorbox}[
    blanker,
    sidebyside,
    sidebyside align=top,
    lefthand width=0.55\textwidth,
    top=4pt, bottom=4pt, left=0pt, right=0pt,
    middle=16pt,
    underlay={
        \draw[dashed, black!80, line width=0.8pt, dash pattern=on 4pt off 3pt]
            ([xshift=0.55\textwidth+8pt]frame.north west) --
            ([xshift=0.55\textwidth+8pt]frame.south west);
    }
]
\textbf{Q.4} \quad {\footnotesize\textsc{[MCQ $\cdot$ 1 Mark(s) $\cdot$ Neg: 0.33]}} \hfill {\footnotesize\textbf{GATE 2018 CS} \quad \tiny\texttt{2018\_CS\_GA\_Q04}}
\vspace{2pt}\hrule height 0.3pt\vspace{4pt}

What would be the smallest natural number which when divided either by 20 or by 42 or by 76 leaves a remainder of 7 in each case?

\par\vspace{6pt}\noindent
\begin{tabular}{@{}p{0.48\linewidth}@{\hspace{0.04\linewidth}}p{0.48\linewidth}@{}}
$\square$\ \textbf{(A)} 3047 & $\square$\ \textbf{(B)} 6047 \\
\addlinespace[3pt]
$\square$\ \textbf{(C)} 7987 & $\square$\ \textbf{(D)} 63847
\end{tabular}

\tcblower

\textbf{Official Key:} \textbf{\large C} \hfill \textit{\small [Easy]}
\vspace{2pt}\hrule height 0.3pt\vspace{3pt}
\textbf{Topic:} Quantitative Aptitude \\
\textbf{Subtopic:} \textit{Numerical computation and estimation: ratios, percentages, powers, exponents and logarithms, permutations and combinations, series}
\vspace{2pt}\hrule height 0.3pt\vspace{3pt}
\textbf{Guide / Hint:} \\
{\small The required number is $\text{LCM}(20, 42, 76) + 7$. Calculating prime factorizations: $20 = 2^2 \times 5$, $42 = 2 \times 3 \times 7$, $76 = 2^2 \times 19$. Thus $\text{LCM} = 2^2 \times 3 \times 5 \times 7 \times 19 = 7980$. Adding the remainder 7 gives $7980 + 7 = 7987$.}

\end{tcolorbox}
\vspace{6pt}\hrule height 0.35pt\vspace{6pt}
\noindent
\begin{tcolorbox}[
    blanker,
    sidebyside,
    sidebyside align=top,
    lefthand width=0.55\textwidth,
    top=4pt, bottom=4pt, left=0pt, right=0pt,
    middle=16pt,
    underlay={
        \draw[dashed, black!80, line width=0.8pt, dash pattern=on 4pt off 3pt]
            ([xshift=0.55\textwidth+8pt]frame.north west) --
            ([xshift=0.55\textwidth+8pt]frame.south west);
    }
]
\textbf{Q.5} \quad {\footnotesize\textsc{[MCQ $\cdot$ 1 Mark(s) $\cdot$ Neg: 0.33]}} \hfill {\footnotesize\textbf{GATE 2018 CS} \quad \tiny\texttt{2018\_CS\_GA\_Q05}}
\vspace{2pt}\hrule height 0.3pt\vspace{4pt}

What is the missing number in the following sequence?
\[ 2,\; 12,\; 60,\; 240,\; 720,\; 1440,\; \_\_\_\_\_,\; 0 \]

\par\vspace{6pt}\noindent
\begin{tabular}{@{}p{0.48\linewidth}@{\hspace{0.04\linewidth}}p{0.48\linewidth}@{}}
$\square$\ \textbf{(A)} 2880 & $\square$\ \textbf{(B)} 1440 \\
\addlinespace[3pt]
$\square$\ \textbf{(C)} 720 & $\square$\ \textbf{(D)} 0
\end{tabular}

\tcblower

\textbf{Official Key:} \textbf{\large B} \hfill \textit{\small [Easy]}
\vspace{2pt}\hrule height 0.3pt\vspace{3pt}
\textbf{Topic:} Quantitative Aptitude \\
\textbf{Subtopic:} \textit{Numerical computation and estimation: ratios, percentages, powers, exponents and logarithms, permutations and combinations, series}
\vspace{2pt}\hrule height 0.3pt\vspace{3pt}
\textbf{Guide / Hint:} \\
{\small Examining consecutive multipliers: $2 \times 6 = 12$, $12 \times 5 = 60$, $60 \times 4 = 240$, $240 \times 3 = 720$, $720 \times 2 = 1440$. Following this decreasing multiplier pattern, the next term is $1440 \times 1 = 1440$, and the term after is $1440 \times 0 = 0$.}

\end{tcolorbox}
\vspace{6pt}\hrule height 0.35pt\vspace{6pt}
\noindent
\begin{tcolorbox}[
    blanker,
    sidebyside,
    sidebyside align=top,
    lefthand width=0.55\textwidth,
    top=4pt, bottom=4pt, left=0pt, right=0pt,
    middle=16pt,
    underlay={
        \draw[dashed, black!80, line width=0.8pt, dash pattern=on 4pt off 3pt]
            ([xshift=0.55\textwidth+8pt]frame.north west) --
            ([xshift=0.55\textwidth+8pt]frame.south west);
    }
]
\textbf{Q.6} \quad {\footnotesize\textsc{[MCQ $\cdot$ 2 Mark(s) $\cdot$ Neg: 0.66]}} \hfill {\footnotesize\textbf{GATE 2018 CS} \quad \tiny\texttt{2018\_CS\_GA\_Q06}}
\vspace{2pt}\hrule height 0.3pt\vspace{4pt}

An old man walks in a town to distribute old age pensions. There are 300 persons in the town who are eligible for the pension. The monthly pension is Rs. 750 for each eligible man and Rs. 1000 for each eligible woman. 8/9 of the eligible men and 2/3 of the eligible women claim the pension. The total pension claimed per month is

\par\vspace{6pt}\noindent
\begin{tabular}{@{}p{0.48\linewidth}@{\hspace{0.04\linewidth}}p{0.48\linewidth}@{}}
$\square$\ \textbf{(A)} Rs. 1,50,000 & $\square$\ \textbf{(B)} Rs. 2,00,000 \\
\addlinespace[3pt]
$\square$\ \textbf{(C)} Rs. 2,25,000 & $\square$\ \textbf{(D)} Cannot be determined
\end{tabular}

\tcblower

\textbf{Official Key:} \textbf{\large B} \hfill \textit{\small [Medium]}
\vspace{2pt}\hrule height 0.3pt\vspace{3pt}
\textbf{Topic:} Quantitative Aptitude \\
\textbf{Subtopic:} \textit{Numerical computation and estimation: ratios, percentages, powers, exponents and logarithms, permutations and combinations, series}
\vspace{2pt}\hrule height 0.3pt\vspace{3pt}
\textbf{Guide / Hint:} \\
{\small Let $M$ and $W$ be eligible men and women, with $M + W = 300$. Number claiming pension is $\frac{8}{9}M$ men and $\frac{2}{3}W$ women. Total pension claimed $= \frac{8}{9}M(750) + \frac{2}{3}W(1000) = \frac{2000}{3}M + \frac{2000}{3}W = \frac{2000}{3}(M + W) = \frac{2000}{3}(300) = 2{,}00{,}000$.}

\end{tcolorbox}
\vspace{6pt}\hrule height 0.35pt\vspace{6pt}
\noindent
\begin{tcolorbox}[
    blanker,
    sidebyside,
    sidebyside align=top,
    lefthand width=0.55\textwidth,
    top=4pt, bottom=4pt, left=0pt, right=0pt,
    middle=16pt,
    underlay={
        \draw[dashed, black!80, line width=0.8pt, dash pattern=on 4pt off 3pt]
            ([xshift=0.55\textwidth+8pt]frame.north west) --
            ([xshift=0.55\textwidth+8pt]frame.south west);
    }
]
\textbf{Q.7} \quad {\footnotesize\textsc{[MCQ $\cdot$ 2 Mark(s) $\cdot$ Neg: 0.66]}} \hfill {\footnotesize\textbf{GATE 2018 CS} \quad \tiny\texttt{2018\_CS\_GA\_Q07}}
\vspace{2pt}\hrule height 0.3pt\vspace{4pt}

If $p^{-x} = \frac{1}{q}$, $q^{-y} = \frac{1}{r}$, and $r^{-z} = \frac{1}{p}$, what is the value of $xyz$?

\par\vspace{6pt}\noindent
\begin{tabular}{@{}p{0.48\linewidth}@{\hspace{0.04\linewidth}}p{0.48\linewidth}@{}}
$\square$\ \textbf{(A)} $-1$ & $\square$\ \textbf{(B)} 0 \\
\addlinespace[3pt]
$\square$\ \textbf{(C)} 1 & $\square$\ \textbf{(D)} $pqr$
\end{tabular}

\tcblower

\textbf{Official Key:} \textbf{\large C} \hfill \textit{\small [Easy]}
\vspace{2pt}\hrule height 0.3pt\vspace{3pt}
\textbf{Topic:} Quantitative Aptitude \\
\textbf{Subtopic:} \textit{Numerical computation and estimation: ratios, percentages, powers, exponents and logarithms, permutations and combinations, series}
\vspace{2pt}\hrule height 0.3pt\vspace{3pt}
\textbf{Guide / Hint:} \\
{\small From $p^{-x} = \frac{1}{q}$, we have $q = p^x$. Substituting into $q^{-y} = \frac{1}{r}$ gives $r = q^y = (p^x)^y = p^{xy}$. Then $r^{-z} = \frac{1}{p}$ gives $p = r^z = (p^{xy})^z = p^{xyz}$. Since $p > 0$ and $p \neq 1$, it follows that $xyz = 1$.}

\end{tcolorbox}
\vspace{6pt}\hrule height 0.35pt\vspace{6pt}
\noindent
\begin{tcolorbox}[
    blanker,
    sidebyside,
    sidebyside align=top,
    lefthand width=0.55\textwidth,
    top=4pt, bottom=4pt, left=0pt, right=0pt,
    middle=16pt,
    underlay={
        \draw[dashed, black!80, line width=0.8pt, dash pattern=on 4pt off 3pt]
            ([xshift=0.55\textwidth+8pt]frame.north west) --
            ([xshift=0.55\textwidth+8pt]frame.south west);
    }
]
\textbf{Q.8} \quad {\footnotesize\textsc{[MCQ $\cdot$ 2 Mark(s) $\cdot$ Neg: 0.66]}} \hfill {\footnotesize\textbf{GATE 2018 CS} \quad \tiny\texttt{2018\_CS\_GA\_Q08}}
\vspace{2pt}\hrule height 0.3pt\vspace{4pt}

In a party, 60\% of the invited guests are male and 40\% are female. 80\% of the invited guests attend the party. If all the invited female guests attend the party, what is the ratio of male guests who attended to the female guests who attended?

\par\vspace{6pt}\noindent
\begin{tabular}{@{}p{0.48\linewidth}@{\hspace{0.04\linewidth}}p{0.48\linewidth}@{}}
$\square$\ \textbf{(A)} 2 : 3 & $\square$\ \textbf{(B)} 1 : 1 \\
\addlinespace[3pt]
$\square$\ \textbf{(C)} 3 : 2 & $\square$\ \textbf{(D)} 2 : 1
\end{tabular}

\tcblower

\textbf{Official Key:} \textbf{\large B} \hfill \textit{\small [Easy]}
\vspace{2pt}\hrule height 0.3pt\vspace{3pt}
\textbf{Topic:} Quantitative Aptitude \\
\textbf{Subtopic:} \textit{Numerical computation and estimation: ratios, percentages, powers, exponents and logarithms, permutations and combinations, series}
\vspace{2pt}\hrule height 0.3pt\vspace{3pt}
\textbf{Guide / Hint:} \\
{\small Out of 100 invited guests, 60 are male and 40 are female. Since 80 guests attend and all 40 invited females attend, the remaining $80 - 40 = 40$ attendees must be male. Thus, the ratio of attending males to attending females is $40 : 40 = 1 : 1$.}

\end{tcolorbox}
\vspace{6pt}\hrule height 0.35pt\vspace{6pt}
\noindent
\begin{tcolorbox}[
    blanker,
    sidebyside,
    sidebyside align=top,
    lefthand width=0.55\textwidth,
    top=4pt, bottom=4pt, left=0pt, right=0pt,
    middle=16pt,
    underlay={
        \draw[dashed, black!80, line width=0.8pt, dash pattern=on 4pt off 3pt]
            ([xshift=0.55\textwidth+8pt]frame.north west) --
            ([xshift=0.55\textwidth+8pt]frame.south west);
    }
]
\textbf{Q.9} \quad {\footnotesize\textsc{[MCQ $\cdot$ 2 Mark(s) $\cdot$ Neg: 0.66]}} \hfill {\footnotesize\textbf{GATE 2018 CS} \quad \tiny\texttt{2018\_CS\_GA\_Q10}}
\vspace{2pt}\hrule height 0.3pt\vspace{4pt}

A six sided unbiased die with four green faces and two red faces is rolled seven times. Which of the following combinations is the most likely outcome of the experiment?

\par\vspace{6pt}\noindent
\begin{tabular}{@{}p{0.98\linewidth}@{}}
$\square$\ \textbf{(A)} Three green faces and four red faces \\
\addlinespace[3pt]
$\square$\ \textbf{(B)} Four green faces and three red faces \\
\addlinespace[3pt]
$\square$\ \textbf{(C)} Five green faces and two red faces \\
\addlinespace[3pt]
$\square$\ \textbf{(D)} Six green faces and one red face
\end{tabular}

\tcblower

\textbf{Official Key:} \textbf{\large C} \hfill \textit{\small [Medium]}
\vspace{2pt}\hrule height 0.3pt\vspace{3pt}
\textbf{Topic:} Quantitative Aptitude \\
\textbf{Subtopic:} \textit{Elementary statistics and probability}
\vspace{2pt}\hrule height 0.3pt\vspace{3pt}
\textbf{Guide / Hint:} \\
{\small The number of green faces $X$ follows a Binomial distribution $B(n=7, p=4/6=2/3)$. The mode of a binomial distribution is $\lfloor (n+1)p \rfloor = \lfloor 8 \times \frac{2}{3} \rfloor = \lfloor 5.33 \rfloor = 5$. Thus, the most likely outcome is 5 green faces and $7 - 5 = 2$ red faces.}

\end{tcolorbox}
\vspace{6pt}\hrule height 0.35pt\vspace{6pt}
\noindent
\begin{tcolorbox}[
    blanker,
    sidebyside,
    sidebyside align=top,
    lefthand width=0.55\textwidth,
    top=4pt, bottom=4pt, left=0pt, right=0pt,
    middle=16pt,
    underlay={
        \draw[dashed, black!80, line width=0.8pt, dash pattern=on 4pt off 3pt]
            ([xshift=0.55\textwidth+8pt]frame.north west) --
            ([xshift=0.55\textwidth+8pt]frame.south west);
    }
]
\textbf{Q.10} \quad {\footnotesize\textsc{[MCQ $\cdot$ 1 Mark(s) $\cdot$ Neg: 0.33]}} \hfill {\footnotesize\textbf{GATE 2022 CS} \quad \tiny\texttt{2022\_CS\_GA\_Q01}}
\vspace{2pt}\hrule height 0.3pt\vspace{4pt}

The \_\_\_\_\_\_\_\_\_ is too high for it to be considered \_\_\_\_\_\_\_\_\_.

\par\vspace{6pt}\noindent
\begin{tabular}{@{}p{0.48\linewidth}@{\hspace{0.04\linewidth}}p{0.48\linewidth}@{}}
$\square$\ \textbf{(A)} fair / fare & $\square$\ \textbf{(B)} faer / fair \\
\addlinespace[3pt]
$\square$\ \textbf{(C)} fare / fare & $\square$\ \textbf{(D)} fare / fair
\end{tabular}

\tcblower

\textbf{Official Key:} \textbf{\large D} \hfill \textit{\small [Easy]}
\vspace{2pt}\hrule height 0.3pt\vspace{3pt}
\textbf{Topic:} Verbal Aptitude \\
\textbf{Subtopic:} \textit{Vocabulary: words, idioms, and phrases in context}
\vspace{2pt}\hrule height 0.3pt\vspace{3pt}
\textbf{Guide / Hint:} \\
{\small 'Fare' denotes the price charged for a journey or conveyance, whereas 'fair' means reasonable or just. Therefore, 'The fare is too high for it to be considered fair.'}

\end{tcolorbox}
\vspace{6pt}\hrule height 0.35pt\vspace{6pt}
\noindent
\begin{tcolorbox}[
    blanker,
    sidebyside,
    sidebyside align=top,
    lefthand width=0.55\textwidth,
    top=4pt, bottom=4pt, left=0pt, right=0pt,
    middle=16pt,
    underlay={
        \draw[dashed, black!80, line width=0.8pt, dash pattern=on 4pt off 3pt]
            ([xshift=0.55\textwidth+8pt]frame.north west) --
            ([xshift=0.55\textwidth+8pt]frame.south west);
    }
]
\textbf{Q.11} \quad {\footnotesize\textsc{[MCQ $\cdot$ 1 Mark(s) $\cdot$ Neg: 0.33]}} \hfill {\footnotesize\textbf{GATE 2022 CS} \quad \tiny\texttt{2022\_CS\_GA\_Q03}}
\vspace{2pt}\hrule height 0.3pt\vspace{4pt}

Let $r$ be a root of the equation $x^2 + 2x + 6 = 0$.
Then the value of the expression $(r+2)(r+3)(r+4)(r+5)$ is

\par\vspace{6pt}\noindent
\begin{tabular}{@{}p{0.48\linewidth}@{\hspace{0.04\linewidth}}p{0.48\linewidth}@{}}
$\square$\ \textbf{(A)} 51 & $\square$\ \textbf{(B)} $-51$ \\
\addlinespace[3pt]
$\square$\ \textbf{(C)} 126 & $\square$\ \textbf{(D)} $-126$
\end{tabular}

\tcblower

\textbf{Official Key:} \textbf{\large D} \hfill \textit{\small [Medium]}
\vspace{2pt}\hrule height 0.3pt\vspace{3pt}
\textbf{Topic:} Quantitative Aptitude \\
\textbf{Subtopic:} \textit{Numerical computation and estimation: ratios, percentages, powers, exponents and logarithms, permutations and combinations, series}
\vspace{2pt}\hrule height 0.3pt\vspace{3pt}
\textbf{Guide / Hint:} \\
{\small Since $r$ is a root of $x^2 + 2x + 6 = 0$, we have $r^2 + 2r = -6$. Pairing terms in the expression: $[(r+2)(r+5)][(r+3)(r+4)] = [r^2 + 7r + 10][r^2 + 7r + 12]$. Substituting $r^2 = -2r - 6$ yields $(5r + 4)(5r + 6) = 25r^2 + 50r + 24 = 25(r^2 + 2r) + 24 = 25(-6) + 24 = -126$.}

\end{tcolorbox}
\vspace{6pt}\hrule height 0.35pt\vspace{6pt}
\noindent
\begin{tcolorbox}[
    blanker,
    sidebyside,
    sidebyside align=top,
    lefthand width=0.55\textwidth,
    top=4pt, bottom=4pt, left=0pt, right=0pt,
    middle=16pt,
    underlay={
        \draw[dashed, black!80, line width=0.8pt, dash pattern=on 4pt off 3pt]
            ([xshift=0.55\textwidth+8pt]frame.north west) --
            ([xshift=0.55\textwidth+8pt]frame.south west);
    }
]
\textbf{Q.12} \quad {\footnotesize\textsc{[MCQ $\cdot$ 1 Mark(s) $\cdot$ Neg: 0.33]}} \hfill {\footnotesize\textbf{GATE 2022 CS} \quad \tiny\texttt{2022\_CS\_GA\_Q04}}
\vspace{2pt}\hrule height 0.3pt\vspace{4pt}

Given below are four statements:
\begin{itemize}[itemsep=0pt,topsep=2pt]
\item Statement 1: All students are inquisitive.
\item Statement 2: Some students are inquisitive.
\item Statement 3: No student is inquisitive.
\item Statement 4: Some students are not inquisitive.
\end{itemize}
From the given four statements, find the two statements that CANNOT BE TRUE simultaneously, assuming that there is at least one student in the class.

\par\vspace{6pt}\noindent
\begin{tabular}{@{}p{0.48\linewidth}@{\hspace{0.04\linewidth}}p{0.48\linewidth}@{}}
$\square$\ \textbf{(A)} Statement 1 and Statement 3 & $\square$\ \textbf{(B)} Statement 1 and Statement 2 \\
\addlinespace[3pt]
$\square$\ \textbf{(C)} Statement 2 and Statement 4 & $\square$\ \textbf{(D)} Statement 3 and Statement 4
\end{tabular}

\tcblower

\textbf{Official Key:} \textbf{\large A} \hfill \textit{\small [Easy]}
\vspace{2pt}\hrule height 0.3pt\vspace{3pt}
\textbf{Topic:} Analytical Aptitude \\
\textbf{Subtopic:} \textit{Logic: deduction and induction}
\vspace{2pt}\hrule height 0.3pt\vspace{3pt}
\textbf{Guide / Hint:} \\
{\small Statement 1 ("All students are inquisitive") and Statement 3 ("No student is inquisitive") are universal affirmative and universal negative statements (contraries). In classical categorical logic with a non-empty class, both statements cannot be simultaneously true.}

\end{tcolorbox}
\vspace{6pt}\hrule height 0.35pt\vspace{6pt}
\noindent
\begin{tcolorbox}[
    blanker,
    sidebyside,
    sidebyside align=top,
    lefthand width=0.55\textwidth,
    top=4pt, bottom=4pt, left=0pt, right=0pt,
    middle=16pt,
    underlay={
        \draw[dashed, black!80, line width=0.8pt, dash pattern=on 4pt off 3pt]
            ([xshift=0.55\textwidth+8pt]frame.north west) --
            ([xshift=0.55\textwidth+8pt]frame.south west);
    }
]
\textbf{Q.13} \quad {\footnotesize\textsc{[MCQ $\cdot$ 2 Mark(s) $\cdot$ Neg: 0.66]}} \hfill {\footnotesize\textbf{GATE 2022 CS} \quad \tiny\texttt{2022\_CS\_GA\_Q06}}
\vspace{2pt}\hrule height 0.3pt\vspace{4pt}

Some people believe that ``what gets measured, improves''. Some others believe that ``what gets measured, gets gamed''. One possible reason for the difference in the beliefs is the work culture in organizations. In organizations with good work culture, metrics help improve outcomes. However, the same metrics are counterproductive in organizations with poor work culture.
Which one of the following is the CORRECT logical inference based on the information in the above passage?

\par\vspace{6pt}\noindent
\begin{tabular}{@{}p{0.98\linewidth}@{}}
$\square$\ \textbf{(A)} Metrics are useful in organizations with poor work culture \\
\addlinespace[3pt]
$\square$\ \textbf{(B)} Metrics are useful in organizations with good work culture \\
\addlinespace[3pt]
$\square$\ \textbf{(C)} Metrics are always counterproductive in organizations with good work culture \\
\addlinespace[3pt]
$\square$\ \textbf{(D)} Metrics are never useful in organizations with good work culture
\end{tabular}

\tcblower

\textbf{Official Key:} \textbf{\large B} \hfill \textit{\small [Easy]}
\vspace{2pt}\hrule height 0.3pt\vspace{3pt}
\textbf{Topic:} Verbal Aptitude \\
\textbf{Subtopic:} \textit{Reading and comprehension}
\vspace{2pt}\hrule height 0.3pt\vspace{3pt}
\textbf{Guide / Hint:} \\
{\small The passage explicitly notes that "In organizations with good work culture, metrics help improve outcomes." This directly implies that metrics are useful in organizations with good work culture.}

\end{tcolorbox}
\vspace{6pt}\hrule height 0.35pt\vspace{6pt}
\noindent
\begin{tcolorbox}[
    blanker,
    sidebyside,
    sidebyside align=top,
    lefthand width=0.55\textwidth,
    top=4pt, bottom=4pt, left=0pt, right=0pt,
    middle=16pt,
    underlay={
        \draw[dashed, black!80, line width=0.8pt, dash pattern=on 4pt off 3pt]
            ([xshift=0.55\textwidth+8pt]frame.north west) --
            ([xshift=0.55\textwidth+8pt]frame.south west);
    }
]
\textbf{Q.14} \quad {\footnotesize\textsc{[MCQ $\cdot$ 2 Mark(s) $\cdot$ Neg: 0.66]}} \hfill {\footnotesize\textbf{GATE 2022 CS} \quad \tiny\texttt{2022\_CS\_GA\_Q07}}
\vspace{2pt}\hrule height 0.3pt\vspace{4pt}

In a recently conducted national entrance test, boys constituted 65\% of those who appeared for the test. Girls constituted the remaining candidates and they accounted for 60\% of the qualified candidates.
Which one of the following is the correct logical inference based on the information provided in the above passage?

\par\vspace{6pt}\noindent
\begin{tabular}{@{}p{0.98\linewidth}@{}}
$\square$\ \textbf{(A)} Equal number of boys and girls qualified \\
\addlinespace[3pt]
$\square$\ \textbf{(B)} Equal number of boys and girls appeared for the test \\
\addlinespace[3pt]
$\square$\ \textbf{(C)} The number of boys who appeared for the test is less than the number of girls who appeared \\
\addlinespace[3pt]
$\square$\ \textbf{(D)} The number of boys who qualified the test is less than the number of girls who qualified
\end{tabular}

\tcblower

\textbf{Official Key:} \textbf{\large D} \hfill \textit{\small [Easy]}
\vspace{2pt}\hrule height 0.3pt\vspace{3pt}
\textbf{Topic:} Quantitative Aptitude \\
\textbf{Subtopic:} \textit{Numerical computation and estimation: ratios, percentages, powers, exponents and logarithms, permutations and combinations, series}
\vspace{2pt}\hrule height 0.3pt\vspace{3pt}
\textbf{Guide / Hint:} \\
{\small Among all candidates who qualified, girls accounted for 60\%, meaning boys accounted for the remaining $100\% - 60\% = 40\%$. Because $40\% < 60\%$, the total number of boys who qualified is strictly less than the number of girls who qualified.}

\end{tcolorbox}
\vspace{6pt}\hrule height 0.35pt\vspace{6pt}
\noindent
\begin{tcolorbox}[
    blanker,
    sidebyside,
    sidebyside align=top,
    lefthand width=0.55\textwidth,
    top=4pt, bottom=4pt, left=0pt, right=0pt,
    middle=16pt,
    underlay={
        \draw[dashed, black!80, line width=0.8pt, dash pattern=on 4pt off 3pt]
            ([xshift=0.55\textwidth+8pt]frame.north west) --
            ([xshift=0.55\textwidth+8pt]frame.south west);
    }
]
\textbf{Q.15} \quad {\footnotesize\textsc{[MCQ $\cdot$ 2 Mark(s) $\cdot$ Neg: 0.66]}} \hfill {\footnotesize\textbf{GATE 2022 CS} \quad \tiny\texttt{2022\_CS\_GA\_Q08}}
\vspace{2pt}\hrule height 0.3pt\vspace{4pt}

A box contains five balls of same size and shape. Three of them are green coloured balls and two of them are orange coloured balls. Balls are drawn from the box one at a time. If a green ball is drawn, it is not replaced. If an orange ball is drawn, it is replaced with another orange ball.
First ball is drawn. What is the probability of getting an orange ball in the next draw?

\par\vspace{6pt}\noindent
\begin{tabular}{@{}p{0.48\linewidth}@{\hspace{0.04\linewidth}}p{0.48\linewidth}@{}}
$\square$\ \textbf{(A)} $1/2$ & $\square$\ \textbf{(B)} $8/25$ \\
\addlinespace[3pt]
$\square$\ \textbf{(C)} $19/50$ & $\square$\ \textbf{(D)} $23/50$
\end{tabular}

\tcblower

\textbf{Official Key:} \textbf{\large D} \hfill \textit{\small [Medium]}
\vspace{2pt}\hrule height 0.3pt\vspace{3pt}
\textbf{Topic:} Quantitative Aptitude \\
\textbf{Subtopic:} \textit{Elementary statistics and probability}
\vspace{2pt}\hrule height 0.3pt\vspace{3pt}
\textbf{Guide / Hint:} \\
{\small Initially, there are 3 green and 2 orange balls (total 5). If a green ball is drawn ($P(G_1) = 3/5$), 2 green and 2 orange remain, so $P(O_2|G_1) = 2/4 = 1/2$. If an orange ball is drawn ($P(O_1) = 2/5$), it is replaced, so $P(O_2|O_1) = 2/5$. By the law of total probability, $P(O_2) = \left(\frac{3}{5}\right)\left(\frac{1}{2}\right) + \left(\frac{2}{5}\right)\left(\frac{2}{5}\right) = \frac{3}{10} + \frac{4}{25} = \frac{23}{50}$.}

\end{tcolorbox}
\vspace{6pt}\hrule height 0.35pt\vspace{6pt}
\noindent
\begin{tcolorbox}[
    blanker,
    sidebyside,
    sidebyside align=top,
    lefthand width=0.55\textwidth,
    top=4pt, bottom=4pt, left=0pt, right=0pt,
    middle=16pt,
    underlay={
        \draw[dashed, black!80, line width=0.8pt, dash pattern=on 4pt off 3pt]
            ([xshift=0.55\textwidth+8pt]frame.north west) --
            ([xshift=0.55\textwidth+8pt]frame.south west);
    }
]
\textbf{Q.16} \quad {\footnotesize\textsc{[MCQ $\cdot$ 2 Mark(s) $\cdot$ Neg: 0.66]}} \hfill {\footnotesize\textbf{GATE 2022 CS} \quad \tiny\texttt{2022\_CS\_GA\_Q09}}
\vspace{2pt}\hrule height 0.3pt\vspace{4pt}

The corners and mid-points of the sides of a triangle are named using the distinct letters P, Q, R, S, T and U, but not necessarily in the same order. Consider the following statements:
\begin{itemize}[itemsep=0pt,topsep=2pt]
\item The line joining P and R is parallel to the line joining Q and S.
\item P is placed on the side opposite to the corner T.
\item S and U cannot be placed on the same side.
\end{itemize}
Which one of the following statements is correct based on the above information?

\par\vspace{6pt}\noindent
\begin{tabular}{@{}p{0.98\linewidth}@{}}
$\square$\ \textbf{(A)} P cannot be placed at a corner \\
\addlinespace[3pt]
$\square$\ \textbf{(B)} S cannot be placed at a corner \\
\addlinespace[3pt]
$\square$\ \textbf{(C)} U cannot be placed at a mid-point \\
\addlinespace[3pt]
$\square$\ \textbf{(D)} R cannot be placed at a corner
\end{tabular}

\tcblower

\textbf{Official Key:} \textbf{\large B} \hfill \textit{\small [Hard]}
\vspace{2pt}\hrule height 0.3pt\vspace{3pt}
\textbf{Topic:} Spatial Aptitude \\
\textbf{Subtopic:} \textit{Transformation of shapes: translation, rotation, scaling, mirroring, assembling, and grouping}
\vspace{2pt}\hrule height 0.3pt\vspace{3pt}
\textbf{Guide / Hint:} \\
{\small In a triangle, a segment connecting two midpoints is parallel to the opposite side joining two vertices. Between segments $PR$ and $QS$, one must connect two corners and the other two midpoints. If $S$ were a corner, then $Q$ and $S$ would both be corners, leaving $P$ and $R$ as midpoints, which violates the condition that $S$ and $U$ cannot share a side while $P$ opposes corner $T$. Thus, $S$ cannot be placed at a corner.}

\end{tcolorbox}
\vspace{6pt}\hrule height 0.35pt\vspace{6pt}
\noindent
\begin{tcolorbox}[
    blanker,
    sidebyside,
    sidebyside align=top,
    lefthand width=0.55\textwidth,
    top=4pt, bottom=4pt, left=0pt, right=0pt,
    middle=16pt,
    underlay={
        \draw[dashed, black!80, line width=0.8pt, dash pattern=on 4pt off 3pt]
            ([xshift=0.55\textwidth+8pt]frame.north west) --
            ([xshift=0.55\textwidth+8pt]frame.south west);
    }
]
\textbf{Q.17} \quad {\footnotesize\textsc{[MCQ $\cdot$ 2 Mark(s) $\cdot$ Neg: 0.66]}} \hfill {\footnotesize\textbf{GATE 2022 CS} \quad \tiny\texttt{2022\_CS\_GA\_Q10}}
\vspace{2pt}\hrule height 0.3pt\vspace{4pt}

A plot of land must be divided between four families. They want their individual plots to be similar in shape, not necessarily equal in area. The land has equally spaced poles, marked as dots in the figure. Two ropes, R1 and R2, are already present and cannot be moved.
What is the least number of additional straight ropes needed to create the desired plots? A single rope can pass through three poles that are aligned in a straight line.

\par\vspace{6pt}\noindent
\begin{tabular}{@{}p{0.48\linewidth}@{\hspace{0.04\linewidth}}p{0.48\linewidth}@{}}
$\square$\ \textbf{(A)} 2 & $\square$\ \textbf{(B)} 4 \\
\addlinespace[3pt]
$\square$\ \textbf{(C)} 5 & $\square$\ \textbf{(D)} 3
\end{tabular}

\tcblower

\textbf{Official Key:} \textbf{\large D} \hfill \textit{\small [Medium]}
\vspace{2pt}\hrule height 0.3pt\vspace{3pt}
\textbf{Topic:} Spatial Aptitude \\
\textbf{Subtopic:} \textit{Paper folding, cutting, and patterns in 2 and 3 dimensions}
\vspace{2pt}\hrule height 0.3pt\vspace{3pt}
\textbf{Guide / Hint:} \\
{\small To partition the rectangular plot into four mutually similar subplots respecting the existing fixed ropes R1 and R2, exactly 3 additional collinear straight rope segments are needed to complete the congruent triangular and trapezoidal subdivisions.}

\end{tcolorbox}
\vspace{6pt}\hrule height 0.35pt\vspace{6pt}
\noindent
\begin{tcolorbox}[
    blanker,
    sidebyside,
    sidebyside align=top,
    lefthand width=0.55\textwidth,
    top=4pt, bottom=4pt, left=0pt, right=0pt,
    middle=16pt,
    underlay={
        \draw[dashed, black!80, line width=0.8pt, dash pattern=on 4pt off 3pt]
            ([xshift=0.55\textwidth+8pt]frame.north west) --
            ([xshift=0.55\textwidth+8pt]frame.south west);
    }
]
\textbf{Q.18} \quad {\footnotesize\textsc{[MCQ $\cdot$ 1 Mark(s) $\cdot$ Neg: 0.33]}} \hfill {\footnotesize\textbf{GATE 2023 CS} \quad \tiny\texttt{2023\_CS\_GA\_Q01}}
\vspace{2pt}\hrule height 0.3pt\vspace{4pt}

We reached the station late, and \_\_\_\_\_\_\_ missed the train.

\par\vspace{6pt}\noindent
\begin{tabular}{@{}p{0.48\linewidth}@{\hspace{0.04\linewidth}}p{0.48\linewidth}@{}}
$\square$\ \textbf{(A)} near & $\square$\ \textbf{(B)} nearly \\
\addlinespace[3pt]
$\square$\ \textbf{(C)} utterly & $\square$\ \textbf{(D)} mostly
\end{tabular}

\tcblower

\textbf{Official Key:} \textbf{\large B} \hfill \textit{\small [Easy]}
\vspace{2pt}\hrule height 0.3pt\vspace{3pt}
\textbf{Topic:} Verbal Aptitude \\
\textbf{Subtopic:} \textit{Basic English grammar: tenses, articles, adjectives, prepositions, conjunctions, verb-noun agreement}
\vspace{2pt}\hrule height 0.3pt\vspace{3pt}
\textbf{Guide / Hint:} \\
{\small The adverb 'nearly' means 'almost' and correctly modifies the past-tense verb 'missed'. 'Near' is an adjective/preposition, while 'utterly' and 'mostly' do not fit the intended semantic meaning of narrowly missing a train.}

\end{tcolorbox}
\vspace{6pt}\hrule height 0.35pt\vspace{6pt}
\noindent
\begin{tcolorbox}[
    blanker,
    sidebyside,
    sidebyside align=top,
    lefthand width=0.55\textwidth,
    top=4pt, bottom=4pt, left=0pt, right=0pt,
    middle=16pt,
    underlay={
        \draw[dashed, black!80, line width=0.8pt, dash pattern=on 4pt off 3pt]
            ([xshift=0.55\textwidth+8pt]frame.north west) --
            ([xshift=0.55\textwidth+8pt]frame.south west);
    }
]
\textbf{Q.19} \quad {\footnotesize\textsc{[MCQ $\cdot$ 1 Mark(s) $\cdot$ Neg: 0.33]}} \hfill {\footnotesize\textbf{GATE 2023 CS} \quad \tiny\texttt{2023\_CS\_GA\_Q02}}
\vspace{2pt}\hrule height 0.3pt\vspace{4pt}

Kind : \_\_\_\_\_\_\_ : : Often : Frequently
(By word meaning)

\par\vspace{6pt}\noindent
\begin{tabular}{@{}p{0.48\linewidth}@{\hspace{0.04\linewidth}}p{0.48\linewidth}@{}}
$\square$\ \textbf{(A)} Mean & $\square$\ \textbf{(B)} Type \\
\addlinespace[3pt]
$\square$\ \textbf{(C)} Cruel & $\square$\ \textbf{(D)} Kindly
\end{tabular}

\tcblower

\textbf{Official Key:} \textbf{\large MTA} \hfill \textit{\small [Easy]}
\vspace{2pt}\hrule height 0.3pt\vspace{3pt}
\textbf{Topic:} Verbal Aptitude \\
\textbf{Subtopic:} \textit{Vocabulary: words, idioms, and phrases in context}
\vspace{2pt}\hrule height 0.3pt\vspace{3pt}
\textbf{Guide / Hint:} \\
{\small 'Often' and 'Frequently' are direct synonyms. 'Kind' can be a noun meaning 'Type' (category) or an adjective meaning benevolent/caring. Because of multiple ambiguous interpretations of 'Kind', GATE awarded Marks to All (MTA) for this question.}

\end{tcolorbox}
\vspace{6pt}\hrule height 0.35pt\vspace{6pt}
\noindent
\begin{tcolorbox}[
    blanker,
    sidebyside,
    sidebyside align=top,
    lefthand width=0.55\textwidth,
    top=4pt, bottom=4pt, left=0pt, right=0pt,
    middle=16pt,
    underlay={
        \draw[dashed, black!80, line width=0.8pt, dash pattern=on 4pt off 3pt]
            ([xshift=0.55\textwidth+8pt]frame.north west) --
            ([xshift=0.55\textwidth+8pt]frame.south west);
    }
]
\textbf{Q.20} \quad {\footnotesize\textsc{[MCQ $\cdot$ 1 Mark(s) $\cdot$ Neg: 0.33]}} \hfill {\footnotesize\textbf{GATE 2023 CS} \quad \tiny\texttt{2023\_CS\_GA\_Q03}}
\vspace{2pt}\hrule height 0.3pt\vspace{4pt}

A series of natural numbers $F_1, F_2, F_3, F_4, F_5, F_6, F_7, \dots$ obeys $F_{n+1} = F_n + F_{n-1}$ for all integers $n \ge 2$.
If $F_6 = 37$, and $F_7 = 60$, then what is $F_1$?

\par\vspace{6pt}\noindent
\begin{tabular}{@{}p{0.48\linewidth}@{\hspace{0.04\linewidth}}p{0.48\linewidth}@{}}
$\square$\ \textbf{(A)} 4 & $\square$\ \textbf{(B)} 5 \\
\addlinespace[3pt]
$\square$\ \textbf{(C)} 8 & $\square$\ \textbf{(D)} 9
\end{tabular}

\tcblower

\textbf{Official Key:} \textbf{\large A} \hfill \textit{\small [Easy]}
\vspace{2pt}\hrule height 0.3pt\vspace{3pt}
\textbf{Topic:} Quantitative Aptitude \\
\textbf{Subtopic:} \textit{Numerical computation and estimation: ratios, percentages, powers, exponents and logarithms, permutations and combinations, series}
\vspace{2pt}\hrule height 0.3pt\vspace{3pt}
\textbf{Guide / Hint:} \\
{\small The recurrence relation is $F_{n+1} = F_n + F_{n-1} \implies F_{n-1} = F_{n+1} - F_n$. Working backward from $F_7 = 60$ and $F_6 = 37$: $F_5 = 60 - 37 = 23$; $F_4 = 37 - 23 = 14$; $F_3 = 23 - 14 = 9$; $F_2 = 14 - 9 = 5$; and $F_1 = 9 - 5 = 4$.}

\end{tcolorbox}
\vspace{6pt}\hrule height 0.35pt\vspace{6pt}
\noindent
\begin{tcolorbox}[
    blanker,
    sidebyside,
    sidebyside align=top,
    lefthand width=0.55\textwidth,
    top=4pt, bottom=4pt, left=0pt, right=0pt,
    middle=16pt,
    underlay={
        \draw[dashed, black!80, line width=0.8pt, dash pattern=on 4pt off 3pt]
            ([xshift=0.55\textwidth+8pt]frame.north west) --
            ([xshift=0.55\textwidth+8pt]frame.south west);
    }
]
\textbf{Q.21} \quad {\footnotesize\textsc{[MCQ $\cdot$ 1 Mark(s) $\cdot$ Neg: 0.33]}} \hfill {\footnotesize\textbf{GATE 2023 CS} \quad \tiny\texttt{2023\_CS\_GA\_Q04}}
\vspace{2pt}\hrule height 0.3pt\vspace{4pt}

A survey for a certain year found that 90\% of pregnant women received medical care at least once before giving birth. Of these women, 60\% received medical care from doctors, while 40\% received medical care from other healthcare providers.
Given this information, which one of the following statements can be inferred with certainty?

\par\vspace{6pt}\noindent
\begin{tabular}{@{}p{0.98\linewidth}@{}}
$\square$\ \textbf{(A)} More than half of the pregnant women received medical care at least once from a doctor. \\
\addlinespace[3pt]
$\square$\ \textbf{(B)} Less than half of the pregnant women received medical care at least once from a doctor. \\
\addlinespace[3pt]
$\square$\ \textbf{(C)} More than half of the pregnant women received medical care at most once from a doctor. \\
\addlinespace[3pt]
$\square$\ \textbf{(D)} Less than half of the pregnant women received medical care at most once from a doctor.
\end{tabular}

\tcblower

\textbf{Official Key:} \textbf{\large A} \hfill \textit{\small [Easy]}
\vspace{2pt}\hrule height 0.3pt\vspace{3pt}
\textbf{Topic:} Quantitative Aptitude \\
\textbf{Subtopic:} \textit{Numerical computation and estimation: ratios, percentages, powers, exponents and logarithms, permutations and combinations, series}
\vspace{2pt}\hrule height 0.3pt\vspace{3pt}
\textbf{Guide / Hint:} \\
{\small Out of all pregnant women, 90\% received medical care at least once. Among those who received medical care, 60\% received care from doctors. The percentage of all pregnant women who received care at least once from a doctor is $90\% \times 60\% = 54\%$, which is strictly greater than half ($> 50\%$).}

\end{tcolorbox}
\vspace{6pt}\hrule height 0.35pt\vspace{6pt}
\noindent
\begin{tcolorbox}[
    blanker,
    sidebyside,
    sidebyside align=top,
    lefthand width=0.55\textwidth,
    top=4pt, bottom=4pt, left=0pt, right=0pt,
    middle=16pt,
    underlay={
        \draw[dashed, black!80, line width=0.8pt, dash pattern=on 4pt off 3pt]
            ([xshift=0.55\textwidth+8pt]frame.north west) --
            ([xshift=0.55\textwidth+8pt]frame.south west);
    }
]
\textbf{Q.22} \quad {\footnotesize\textsc{[MCQ $\cdot$ 1 Mark(s) $\cdot$ Neg: 0.33]}} \hfill {\footnotesize\textbf{GATE 2023 CS} \quad \tiny\texttt{2023\_CS\_GA\_Q05}}
\vspace{2pt}\hrule height 0.3pt\vspace{4pt}

Looking at the surface of a smooth 3-dimensional object from the outside, which one of the following options is TRUE?

\par\vspace{6pt}\noindent
\begin{tabular}{@{}p{0.98\linewidth}@{}}
$\square$\ \textbf{(A)} The surface of the object must be concave everywhere. \\
\addlinespace[3pt]
$\square$\ \textbf{(B)} The surface of the object must be convex everywhere. \\
\addlinespace[3pt]
$\square$\ \textbf{(C)} The surface of the object may be concave in some places and convex in other places. \\
\addlinespace[3pt]
$\square$\ \textbf{(D)} The object can have edges, but no corners.
\end{tabular}

\tcblower

\textbf{Official Key:} \textbf{\large C} \hfill \textit{\small [Easy]}
\vspace{2pt}\hrule height 0.3pt\vspace{3pt}
\textbf{Topic:} Spatial Aptitude \\
\textbf{Subtopic:} \textit{Transformation of shapes: translation, rotation, scaling, mirroring, assembling, and grouping}
\vspace{2pt}\hrule height 0.3pt\vspace{3pt}
\textbf{Guide / Hint:} \\
{\small A smooth 3-dimensional object (such as a torus or a dumbbell shape) can possess regions of positive Gaussian curvature where it is locally convex, as well as regions of negative curvature (saddle-like regions) where it is locally concave. Hence, its surface may be concave in some places and convex in others.}

\end{tcolorbox}
\vspace{6pt}\hrule height 0.35pt\vspace{6pt}
\noindent
\begin{tcolorbox}[
    blanker,
    sidebyside,
    sidebyside align=top,
    lefthand width=0.55\textwidth,
    top=4pt, bottom=4pt, left=0pt, right=0pt,
    middle=16pt,
    underlay={
        \draw[dashed, black!80, line width=0.8pt, dash pattern=on 4pt off 3pt]
            ([xshift=0.55\textwidth+8pt]frame.north west) --
            ([xshift=0.55\textwidth+8pt]frame.south west);
    }
]
\textbf{Q.23} \quad {\footnotesize\textsc{[MCQ $\cdot$ 2 Mark(s) $\cdot$ Neg: 0.66]}} \hfill {\footnotesize\textbf{GATE 2023 CS} \quad \tiny\texttt{2023\_CS\_GA\_Q06}}
\vspace{2pt}\hrule height 0.3pt\vspace{4pt}

The country of Zombieland is in distress since more than 75\% of its working population is suffering from serious health issues. Studies conducted by competent health experts concluded that a complete lack of physical exercise among its working population was one of the leading causes of their health issues. As one of the measures to address the problem, the Government of Zombieland has decided to provide monetary incentives to those who ride bicycles to work.
\par\medskip
Based only on the information provided above, which one of the following statements can be logically inferred with certainty?

\par\vspace{6pt}\noindent
\begin{tabular}{@{}p{0.98\linewidth}@{}}
$\square$\ \textbf{(A)} All the working population of Zombieland will henceforth ride bicycles to work. \\
\addlinespace[3pt]
$\square$\ \textbf{(B)} Riding bicycles will ensure that all of the working population of Zombieland is free of health issues. \\
\addlinespace[3pt]
$\square$\ \textbf{(C)} The health experts suggested to the Government of Zombieland to declare riding bicycles as mandatory. \\
\addlinespace[3pt]
$\square$\ \textbf{(D)} The Government of Zombieland believes that riding bicycles is a form of physical exercise.
\end{tabular}

\tcblower

\textbf{Official Key:} \textbf{\large D} \hfill \textit{\small [Easy]}
\vspace{2pt}\hrule height 0.3pt\vspace{3pt}
\textbf{Topic:} Verbal Aptitude \\
\textbf{Subtopic:} \textit{Reading and comprehension}
\vspace{2pt}\hrule height 0.3pt\vspace{3pt}
\textbf{Guide / Hint:} \\
{\small The passage establishes that a complete lack of physical exercise is a primary cause of health issues, and in response, the government provides financial incentives to ride bicycles to work. This direct causal measure logically implies that the government views bicycle riding as a form of physical exercise.}

\end{tcolorbox}
\vspace{6pt}\hrule height 0.35pt\vspace{6pt}
\noindent
\begin{tcolorbox}[
    blanker,
    sidebyside,
    sidebyside align=top,
    lefthand width=0.55\textwidth,
    top=4pt, bottom=4pt, left=0pt, right=0pt,
    middle=16pt,
    underlay={
        \draw[dashed, black!80, line width=0.8pt, dash pattern=on 4pt off 3pt]
            ([xshift=0.55\textwidth+8pt]frame.north west) --
            ([xshift=0.55\textwidth+8pt]frame.south west);
    }
]
\textbf{Q.24} \quad {\footnotesize\textsc{[MCQ $\cdot$ 2 Mark(s) $\cdot$ Neg: 0.66]}} \hfill {\footnotesize\textbf{GATE 2023 CS} \quad \tiny\texttt{2023\_CS\_GA\_Q07}}
\vspace{2pt}\hrule height 0.3pt\vspace{4pt}

Consider two functions of time ($t$),
\[ f(t) = 0.01 t^2, \quad g(t) = 4 t \]
where $0 < t < \infty$.
Now consider the following two statements:
\begin{enumerate}[label=(\roman*),itemsep=0pt,topsep=2pt]
\item For some $t > 0$, $g(t) > f(t)$.
\item There exists a $T$, such that $f(t) > g(t)$ for all $t > T$.
\end{enumerate}
Which one of the following options is TRUE?

\par\vspace{6pt}\noindent
\begin{tabular}{@{}p{0.48\linewidth}@{\hspace{0.04\linewidth}}p{0.48\linewidth}@{}}
$\square$\ \textbf{(A)} only (i) is correct & $\square$\ \textbf{(B)} only (ii) is correct \\
\addlinespace[3pt]
$\square$\ \textbf{(C)} both (i) and (ii) are correct & $\square$\ \textbf{(D)} neither (i) nor (ii) is correct
\end{tabular}

\tcblower

\textbf{Official Key:} \textbf{\large C} \hfill \textit{\small [Easy]}
\vspace{2pt}\hrule height 0.3pt\vspace{3pt}
\textbf{Topic:} Quantitative Aptitude \\
\textbf{Subtopic:} \textit{Numerical computation and estimation: ratios, percentages, powers, exponents and logarithms, permutations and combinations, series}
\vspace{2pt}\hrule height 0.3pt\vspace{3pt}
\textbf{Guide / Hint:} \\
{\small Setting $f(t) = g(t) \implies 0.01 t^2 = 4t \implies t = 400$. For any $0 < t < 400$ (e.g., $t = 10$), $g(10) = 40 > f(10) = 1$, making statement (i) correct. For all $t > 400$, the quadratic growth of $f(t)$ strictly dominates the linear growth of $g(t)$, so choosing $T = 400$ satisfies statement (ii). Hence, both statements are correct.}

\end{tcolorbox}
\vspace{6pt}\hrule height 0.35pt\vspace{6pt}
\noindent
\begin{tcolorbox}[
    blanker,
    sidebyside,
    sidebyside align=top,
    lefthand width=0.55\textwidth,
    top=4pt, bottom=4pt, left=0pt, right=0pt,
    middle=16pt,
    underlay={
        \draw[dashed, black!80, line width=0.8pt, dash pattern=on 4pt off 3pt]
            ([xshift=0.55\textwidth+8pt]frame.north west) --
            ([xshift=0.55\textwidth+8pt]frame.south west);
    }
]
\textbf{Q.25} \quad {\footnotesize\textsc{[MCQ $\cdot$ 2 Mark(s) $\cdot$ Neg: 0.66]}} \hfill {\footnotesize\textbf{GATE 2023 CS} \quad \tiny\texttt{2023\_CS\_GA\_Q08}}
\vspace{2pt}\hrule height 0.3pt\vspace{4pt}

Which one of the following sentence sequences creates a coherent narrative?
\begin{enumerate}[label=(\roman*),itemsep=0pt,topsep=2pt]
\item Once on the terrace, on her way to her small room in the corner, she notices the man right away.
\item She begins to pant by the time she has climbed all the stairs.
\item Mina has bought vegetables and rice at the market, so her bags are heavy.
\item He was leaning against the parapet, watching the traffic below.
\end{enumerate}

\par\vspace{6pt}\noindent
\begin{tabular}{@{}p{0.48\linewidth}@{\hspace{0.04\linewidth}}p{0.48\linewidth}@{}}
$\square$\ \textbf{(A)} (i), (ii), (iv), (iii) & $\square$\ \textbf{(B)} (ii), (iii), (i), (iv) \\
\addlinespace[3pt]
$\square$\ \textbf{(C)} (iv), (ii), (i), (iii) & $\square$\ \textbf{(D)} (iii), (ii), (i), (iv)
\end{tabular}

\tcblower

\textbf{Official Key:} \textbf{\large D} \hfill \textit{\small [Easy]}
\vspace{2pt}\hrule height 0.3pt\vspace{3pt}
\textbf{Topic:} Verbal Aptitude \\
\textbf{Subtopic:} \textit{Narrative sequencing}
\vspace{2pt}\hrule height 0.3pt\vspace{3pt}
\textbf{Guide / Hint:} \\
{\small Sentence (iii) introduces the protagonist Mina carrying heavy market purchases. Sentence (ii) describes the resulting exertion as she climbs the stairs. Sentence (i) brings her onto the terrace where she spots a man. Sentence (iv) details the man's posture against the parapet. Thus, the coherent sequence is (iii), (ii), (i), (iv).}

\end{tcolorbox}
\vspace{6pt}\hrule height 0.35pt\vspace{6pt}
\noindent
\begin{tcolorbox}[
    blanker,
    sidebyside,
    sidebyside align=top,
    lefthand width=0.55\textwidth,
    top=4pt, bottom=4pt, left=0pt, right=0pt,
    middle=16pt,
    underlay={
        \draw[dashed, black!80, line width=0.8pt, dash pattern=on 4pt off 3pt]
            ([xshift=0.55\textwidth+8pt]frame.north west) --
            ([xshift=0.55\textwidth+8pt]frame.south west);
    }
]
\textbf{Q.26} \quad {\footnotesize\textsc{[MCQ $\cdot$ 2 Mark(s) $\cdot$ Neg: 0.66]}} \hfill {\footnotesize\textbf{GATE 2023 CS} \quad \tiny\texttt{2023\_CS\_GA\_Q09}}
\vspace{2pt}\hrule height 0.3pt\vspace{4pt}

$f(x)$ and $g(y)$ are functions of $x$ and $y$, respectively, and $f(x) = g(y)$ for all real values of $x$ and $y$. Which one of the following options is necessarily TRUE for all $x$ and $y$?

\par\vspace{6pt}\noindent
\begin{tabular}{@{}p{0.98\linewidth}@{}}
$\square$\ \textbf{(A)} $f(x) = 0$ and $g(y) = 0$ \\
\addlinespace[3pt]
$\square$\ \textbf{(B)} $f(x) = g(y) = \text{constant}$ \\
\addlinespace[3pt]
$\square$\ \textbf{(C)} $f(x) \neq \text{constant}$ and $g(y) \neq \text{constant}$ \\
\addlinespace[3pt]
$\square$\ \textbf{(D)} $f(x) + g(y) = f(x) - g(y)$
\end{tabular}

\tcblower

\textbf{Official Key:} \textbf{\large B} \hfill \textit{\small [Easy]}
\vspace{2pt}\hrule height 0.3pt\vspace{3pt}
\textbf{Topic:} Quantitative Aptitude \\
\textbf{Subtopic:} \textit{Numerical computation and estimation: ratios, percentages, powers, exponents and logarithms, permutations and combinations, series}
\vspace{2pt}\hrule height 0.3pt\vspace{3pt}
\textbf{Guide / Hint:} \\
{\small We are given $f(x) = g(y)$ for all independent real variables $x$ and $y$. If $x$ varies while $y$ remains constant, $f(x)$ cannot change because $g(y)$ is fixed, meaning $f(x)$ must be a constant. Symmetrically, $g(y)$ must be identical to that same constant. Therefore, $f(x) = g(y) = \text{constant}$.}

\end{tcolorbox}
\vspace{6pt}\hrule height 0.35pt\vspace{6pt}
\noindent
\begin{tcolorbox}[
    blanker,
    sidebyside,
    sidebyside align=top,
    lefthand width=0.55\textwidth,
    top=4pt, bottom=4pt, left=0pt, right=0pt,
    middle=16pt,
    underlay={
        \draw[dashed, black!80, line width=0.8pt, dash pattern=on 4pt off 3pt]
            ([xshift=0.55\textwidth+8pt]frame.north west) --
            ([xshift=0.55\textwidth+8pt]frame.south west);
    }
]
\textbf{Q.27} \quad {\footnotesize\textsc{[MCQ $\cdot$ 2 Mark(s) $\cdot$ Neg: 0.66]}} \hfill {\footnotesize\textbf{GATE 2023 CS} \quad \tiny\texttt{2023\_CS\_GA\_Q10}}
\vspace{2pt}\hrule height 0.3pt\vspace{4pt}

Which one of the options best describes the transformation of the 2-dimensional figure P to Q, and then to R, as shown?

\par\vspace{6pt}\noindent
\begin{tabular}{@{}p{0.48\linewidth}@{\hspace{0.04\linewidth}}p{0.48\linewidth}@{}}
$\square$\ \textbf{(A)} Only S1 is TRUE. & $\square$\ \textbf{(B)} Only S1 and S2 are TRUE. \\
\addlinespace[3pt]
$\square$\ \textbf{(C)} S1, S2, and S3 are all TRUE. & $\square$\ \textbf{(D)} Only S1 and S3 are TRUE.
\end{tabular}

\tcblower

\textbf{Official Key:} \textbf{\large A} \hfill \textit{\small [Medium]}
\vspace{2pt}\hrule height 0.3pt\vspace{3pt}
\textbf{Topic:} Spatial Aptitude \\
\textbf{Subtopic:} \textit{Transformation of shapes: translation, rotation, scaling, mirroring, assembling, and grouping}
\vspace{2pt}\hrule height 0.3pt\vspace{3pt}
\textbf{Guide / Hint:} \\
{\small Analyzing the rigid 2D planar transformations from figure P to Q and Q to R shows that statement S1 accurately describes the sequence of planar rotation and reflection, while statements S2 and S3 fail to match the observed geometric orientations. Hence, only S1 is TRUE.}

\end{tcolorbox}
\vspace{6pt}\hrule height 0.35pt\vspace{6pt}
\noindent
\begin{tcolorbox}[
    blanker,
    sidebyside,
    sidebyside align=top,
    lefthand width=0.55\textwidth,
    top=4pt, bottom=4pt, left=0pt, right=0pt,
    middle=16pt,
    underlay={
        \draw[dashed, black!80, line width=0.8pt, dash pattern=on 4pt off 3pt]
            ([xshift=0.55\textwidth+8pt]frame.north west) --
            ([xshift=0.55\textwidth+8pt]frame.south west);
    }
]
\textbf{Q.28} \quad {\footnotesize\textsc{[MCQ $\cdot$ 1 Mark(s) $\cdot$ Neg: 0.33]}} \hfill {\footnotesize\textbf{GATE 2024 CS1} \quad \tiny\texttt{2024\_CS1\_GA\_Q01}}
\vspace{2pt}\hrule height 0.3pt\vspace{4pt}

If `$\to$' denotes increasing order of intensity, then the meaning of the words $[\text{dry} \to \text{arid} \to \text{parched}]$ is analogous to $[\text{diet} \to \text{fast} \to \text{\_\_\_\_\_\_\_\_}]$.
Which one of the given options is appropriate to fill the blank?

\par\vspace{6pt}\noindent
\begin{tabular}{@{}p{0.48\linewidth}@{\hspace{0.04\linewidth}}p{0.48\linewidth}@{}}
$\square$\ \textbf{(A)} starve & $\square$\ \textbf{(B)} reject \\
\addlinespace[3pt]
$\square$\ \textbf{(C)} feast & $\square$\ \textbf{(D)} deny
\end{tabular}

\tcblower

\textbf{Official Key:} \textbf{\large A} \hfill \textit{\small [Easy]}
\vspace{2pt}\hrule height 0.3pt\vspace{3pt}
\textbf{Topic:} Verbal Aptitude \\
\textbf{Subtopic:} \textit{Vocabulary: words, idioms, and phrases in context}
\vspace{2pt}\hrule height 0.3pt\vspace{3pt}
\textbf{Guide / Hint:} \\
{\small 'dry $\to$ arid $\to$ parched' depicts an increasing scale of lack of moisture. Similarly, 'diet $\to$ fast $\to$ starve' depicts an increasing scale of food restriction/deprivation, where starving represents extreme lack of nourishment.}

\end{tcolorbox}
\vspace{6pt}\hrule height 0.35pt\vspace{6pt}
\noindent
\begin{tcolorbox}[
    blanker,
    sidebyside,
    sidebyside align=top,
    lefthand width=0.55\textwidth,
    top=4pt, bottom=4pt, left=0pt, right=0pt,
    middle=16pt,
    underlay={
        \draw[dashed, black!80, line width=0.8pt, dash pattern=on 4pt off 3pt]
            ([xshift=0.55\textwidth+8pt]frame.north west) --
            ([xshift=0.55\textwidth+8pt]frame.south west);
    }
]
\textbf{Q.29} \quad {\footnotesize\textsc{[MCQ $\cdot$ 1 Mark(s) $\cdot$ Neg: 0.33]}} \hfill {\footnotesize\textbf{GATE 2024 CS1} \quad \tiny\texttt{2024\_CS1\_GA\_Q02}}
\vspace{2pt}\hrule height 0.3pt\vspace{4pt}

If two distinct non-zero real variables $x$ and $y$ are such that $(x+y)$ is proportional to $(x-y)$ then the value of $\frac{x}{y}$

\par\vspace{6pt}\noindent
\begin{tabular}{@{}p{0.98\linewidth}@{}}
$\square$\ \textbf{(A)} depends on $xy$ \\
\addlinespace[3pt]
$\square$\ \textbf{(B)} depends only on $x$ and not on $y$ \\
\addlinespace[3pt]
$\square$\ \textbf{(C)} depends only on $y$ and not on $x$ \\
\addlinespace[3pt]
$\square$\ \textbf{(D)} is a constant
\end{tabular}

\tcblower

\textbf{Official Key:} \textbf{\large D} \hfill \textit{\small [Easy]}
\vspace{2pt}\hrule height 0.3pt\vspace{3pt}
\textbf{Topic:} Quantitative Aptitude \\
\textbf{Subtopic:} \textit{Numerical computation and estimation: ratios, percentages, powers, exponents and logarithms, permutations and combinations, series}
\vspace{2pt}\hrule height 0.3pt\vspace{3pt}
\textbf{Guide / Hint:} \\
{\small If $(x+y) \propto (x-y)$, then $x+y = k(x-y)$ for some non-zero constant $k \neq 1$. Rearranging terms gives $x(1-k) = -y(1+k) \implies \frac{x}{y} = \frac{k+1}{k-1}$, which is a constant independent of $x$ and $y$.}

\end{tcolorbox}
\vspace{6pt}\hrule height 0.35pt\vspace{6pt}
\noindent
\begin{tcolorbox}[
    blanker,
    sidebyside,
    sidebyside align=top,
    lefthand width=0.55\textwidth,
    top=4pt, bottom=4pt, left=0pt, right=0pt,
    middle=16pt,
    underlay={
        \draw[dashed, black!80, line width=0.8pt, dash pattern=on 4pt off 3pt]
            ([xshift=0.55\textwidth+8pt]frame.north west) --
            ([xshift=0.55\textwidth+8pt]frame.south west);
    }
]
\textbf{Q.30} \quad {\footnotesize\textsc{[MCQ $\cdot$ 1 Mark(s) $\cdot$ Neg: 0.33]}} \hfill {\footnotesize\textbf{GATE 2024 CS1} \quad \tiny\texttt{2024\_CS1\_GA\_Q03}}
\vspace{2pt}\hrule height 0.3pt\vspace{4pt}

Consider the following sample of numbers:
\[ 9, 18, 11, 14, 15, 17, 10, 69, 11, 13 \]
The median of the sample is

\par\vspace{6pt}\noindent
\begin{tabular}{@{}p{0.48\linewidth}@{\hspace{0.04\linewidth}}p{0.48\linewidth}@{}}
$\square$\ \textbf{(A)} 13.5 & $\square$\ \textbf{(B)} 14 \\
\addlinespace[3pt]
$\square$\ \textbf{(C)} 11 & $\square$\ \textbf{(D)} 18.7
\end{tabular}

\tcblower

\textbf{Official Key:} \textbf{\large A} \hfill \textit{\small [Easy]}
\vspace{2pt}\hrule height 0.3pt\vspace{3pt}
\textbf{Topic:} Quantitative Aptitude \\
\textbf{Subtopic:} \textit{Elementary statistics and probability}
\vspace{2pt}\hrule height 0.3pt\vspace{3pt}
\textbf{Guide / Hint:} \\
{\small Arranging the 10 sample values in ascending order: $9, 10, 11, 11, 13, 14, 15, 17, 18, 69$. Since the sample size $n = 10$ is even, the median is the arithmetic mean of the 5th and 6th values: $\frac{13 + 14}{2} = 13.5$.}

\end{tcolorbox}
\vspace{6pt}\hrule height 0.35pt\vspace{6pt}
\noindent
\begin{tcolorbox}[
    blanker,
    sidebyside,
    sidebyside align=top,
    lefthand width=0.55\textwidth,
    top=4pt, bottom=4pt, left=0pt, right=0pt,
    middle=16pt,
    underlay={
        \draw[dashed, black!80, line width=0.8pt, dash pattern=on 4pt off 3pt]
            ([xshift=0.55\textwidth+8pt]frame.north west) --
            ([xshift=0.55\textwidth+8pt]frame.south west);
    }
]
\textbf{Q.31} \quad {\footnotesize\textsc{[MCQ $\cdot$ 1 Mark(s) $\cdot$ Neg: 0.33]}} \hfill {\footnotesize\textbf{GATE 2024 CS1} \quad \tiny\texttt{2024\_CS1\_GA\_Q04}}
\vspace{2pt}\hrule height 0.3pt\vspace{4pt}

The number of coins of Rs.~1, Rs.~5, and Rs.~10 denominations that a person has are in the ratio $5:3:13$. Of the total amount, the percentage of money in Rs.~5 coins is

\par\vspace{6pt}\noindent
\begin{tabular}{@{}p{0.48\linewidth}@{\hspace{0.04\linewidth}}p{0.48\linewidth}@{}}
$\square$\ \textbf{(A)} 21\% & $\square$\ \textbf{(B)} $14\frac{2}{7}\%$ \\
\addlinespace[3pt]
$\square$\ \textbf{(C)} 10\% & $\square$\ \textbf{(D)} 30\%
\end{tabular}

\tcblower

\textbf{Official Key:} \textbf{\large C} \hfill \textit{\small [Easy]}
\vspace{2pt}\hrule height 0.3pt\vspace{3pt}
\textbf{Topic:} Quantitative Aptitude \\
\textbf{Subtopic:} \textit{Numerical computation and estimation: ratios, percentages, powers, exponents and logarithms, permutations and combinations, series}
\vspace{2pt}\hrule height 0.3pt\vspace{3pt}
\textbf{Guide / Hint:} \\
{\small Let the number of Rs.~1, Rs.~5, and Rs.~10 coins be $5k$, $3k$, and $13k$ respectively. The total value is $1(5k) + 5(3k) + 10(13k) = 5k + 15k + 130k = 150k$. The percentage of money in Rs.~5 coins is $\frac{15k}{150k} \times 100\% = 10\%$.}

\end{tcolorbox}
\vspace{6pt}\hrule height 0.35pt\vspace{6pt}
\noindent
\begin{tcolorbox}[
    blanker,
    sidebyside,
    sidebyside align=top,
    lefthand width=0.55\textwidth,
    top=4pt, bottom=4pt, left=0pt, right=0pt,
    middle=16pt,
    underlay={
        \draw[dashed, black!80, line width=0.8pt, dash pattern=on 4pt off 3pt]
            ([xshift=0.55\textwidth+8pt]frame.north west) --
            ([xshift=0.55\textwidth+8pt]frame.south west);
    }
]
\textbf{Q.32} \quad {\footnotesize\textsc{[MCQ $\cdot$ 1 Mark(s) $\cdot$ Neg: 0.33]}} \hfill {\footnotesize\textbf{GATE 2024 CS1} \quad \tiny\texttt{2024\_CS1\_GA\_Q05}}
\vspace{2pt}\hrule height 0.3pt\vspace{4pt}

For positive non-zero real variables $p$ and $q$, if
\[ \log(p^2 + q^2) = \log p + \log q + 2\log 3, \]
then, the value of $\frac{p^4+q^4}{p^2q^2}$ is

\par\vspace{6pt}\noindent
\begin{tabular}{@{}p{0.48\linewidth}@{\hspace{0.04\linewidth}}p{0.48\linewidth}@{}}
$\square$\ \textbf{(A)} 79 & $\square$\ \textbf{(B)} 81 \\
\addlinespace[3pt]
$\square$\ \textbf{(C)} 9 & $\square$\ \textbf{(D)} 83
\end{tabular}

\tcblower

\textbf{Official Key:} \textbf{\large A} \hfill \textit{\small [Medium]}
\vspace{2pt}\hrule height 0.3pt\vspace{3pt}
\textbf{Topic:} Quantitative Aptitude \\
\textbf{Subtopic:} \textit{Numerical computation and estimation: ratios, percentages, powers, exponents and logarithms, permutations and combinations, series}
\vspace{2pt}\hrule height 0.3pt\vspace{3pt}
\textbf{Guide / Hint:} \\
{\small Using logarithm rules, $\log(p^2 + q^2) = \log p + \log q + 2\log 3 = \log(9pq) \implies p^2 + q^2 = 9pq$. Squaring both sides yields $(p^2+q^2)^2 = p^4 + q^4 + 2p^2q^2 = 81p^2q^2$, which gives $p^4 + q^4 = 79p^2q^2$. Hence, $\frac{p^4+q^4}{p^2q^2} = 79$.}

\end{tcolorbox}
\vspace{6pt}\hrule height 0.35pt\vspace{6pt}
\noindent
\begin{tcolorbox}[
    blanker,
    sidebyside,
    sidebyside align=top,
    lefthand width=0.55\textwidth,
    top=4pt, bottom=4pt, left=0pt, right=0pt,
    middle=16pt,
    underlay={
        \draw[dashed, black!80, line width=0.8pt, dash pattern=on 4pt off 3pt]
            ([xshift=0.55\textwidth+8pt]frame.north west) --
            ([xshift=0.55\textwidth+8pt]frame.south west);
    }
]
\textbf{Q.33} \quad {\footnotesize\textsc{[MCQ $\cdot$ 2 Mark(s) $\cdot$ Neg: 0.66]}} \hfill {\footnotesize\textbf{GATE 2024 CS1} \quad \tiny\texttt{2024\_CS1\_GA\_Q06}}
\vspace{2pt}\hrule height 0.3pt\vspace{4pt}

In the given text, the blanks are numbered (i)--(iv). Select the best match for all the blanks.
\par\medskip
Steve was advised to keep his head (i)\_\_\_\_ before heading (ii)\_\_\_\_ to bat; for, while he had a head (iii)\_\_\_\_ batting, he could only do so with a cool head (iv)\_\_\_\_ his shoulders.

\par\vspace{6pt}\noindent
\begin{tabular}{@{}p{0.98\linewidth}@{}}
$\square$\ \textbf{(A)} (i) down (ii) down (iii) on (iv) for \\
\addlinespace[3pt]
$\square$\ \textbf{(B)} (i) on (ii) down (iii) for (iv) on \\
\addlinespace[3pt]
$\square$\ \textbf{(C)} (i) down (ii) out (iii) for (iv) on \\
\addlinespace[3pt]
$\square$\ \textbf{(D)} (i) on (ii) out (iii) on (iv) for
\end{tabular}

\tcblower

\textbf{Official Key:} \textbf{\large C} \hfill \textit{\small [Easy]}
\vspace{2pt}\hrule height 0.3pt\vspace{3pt}
\textbf{Topic:} Verbal Aptitude \\
\textbf{Subtopic:} \textit{Vocabulary: words, idioms, and phrases in context}
\vspace{2pt}\hrule height 0.3pt\vspace{3pt}
\textbf{Guide / Hint:} \\
{\small 'Keep one's head down' means to stay focused and avoid trouble. 'Heading out' means going forth/departing. Having a 'head for' an activity signifies natural aptitude for it, and having a cool 'head on' one's shoulders means remaining calm and composed.}

\end{tcolorbox}
\vspace{6pt}\hrule height 0.35pt\vspace{6pt}
\noindent
\begin{tcolorbox}[
    blanker,
    sidebyside,
    sidebyside align=top,
    lefthand width=0.55\textwidth,
    top=4pt, bottom=4pt, left=0pt, right=0pt,
    middle=16pt,
    underlay={
        \draw[dashed, black!80, line width=0.8pt, dash pattern=on 4pt off 3pt]
            ([xshift=0.55\textwidth+8pt]frame.north west) --
            ([xshift=0.55\textwidth+8pt]frame.south west);
    }
]
\textbf{Q.34} \quad {\footnotesize\textsc{[MCQ $\cdot$ 2 Mark(s) $\cdot$ Neg: 0.66]}} \hfill {\footnotesize\textbf{GATE 2024 CS1} \quad \tiny\texttt{2024\_CS1\_GA\_Q07}}
\vspace{2pt}\hrule height 0.3pt\vspace{4pt}

A rectangular paper sheet of dimensions $54\text{ cm} \times 4\text{ cm}$ is taken. The two longer edges of the sheet are joined together to create a cylindrical tube. A cube whose surface area is equal to the area of the sheet is also taken.
Then, the ratio of the volume of the cylindrical tube to the volume of the cube is

\par\vspace{6pt}\noindent
\begin{tabular}{@{}p{0.48\linewidth}@{\hspace{0.04\linewidth}}p{0.48\linewidth}@{}}
$\square$\ \textbf{(A)} $1/\pi$ & $\square$\ \textbf{(B)} $2/\pi$ \\
\addlinespace[3pt]
$\square$\ \textbf{(C)} $3/\pi$ & $\square$\ \textbf{(D)} $4/\pi$
\end{tabular}

\tcblower

\textbf{Official Key:} \textbf{\large A} \hfill \textit{\small [Medium]}
\vspace{2pt}\hrule height 0.3pt\vspace{3pt}
\textbf{Topic:} Quantitative Aptitude \\
\textbf{Subtopic:} \textit{Mensuration and geometry}
\vspace{2pt}\hrule height 0.3pt\vspace{3pt}
\textbf{Guide / Hint:} \\
{\small Joining the longer edges ($54\text{ cm}$) creates a cylinder with height $h = 54\text{ cm}$ and circumference $2\pi r = 4\text{ cm} \implies r = \frac{2}{\pi}$, giving volume $V_{\text{cyl}} = \pi r^2 h = \frac{216}{\pi}$. The sheet area is $54 \times 4 = 216\text{ cm}^2$. A cube with surface area $6a^2 = 216$ has edge $a = 6\text{ cm}$ and volume $V_{\text{cube}} = 6^3 = 216$. Thus $\frac{V_{\text{cyl}}}{V_{\text{cube}}} = \frac{216/\pi}{216} = \frac{1}{\pi}$.}

\end{tcolorbox}
\vspace{6pt}\hrule height 0.35pt\vspace{6pt}
\noindent
\begin{tcolorbox}[
    blanker,
    sidebyside,
    sidebyside align=top,
    lefthand width=0.55\textwidth,
    top=4pt, bottom=4pt, left=0pt, right=0pt,
    middle=16pt,
    underlay={
        \draw[dashed, black!80, line width=0.8pt, dash pattern=on 4pt off 3pt]
            ([xshift=0.55\textwidth+8pt]frame.north west) --
            ([xshift=0.55\textwidth+8pt]frame.south west);
    }
]
\textbf{Q.35} \quad {\footnotesize\textsc{[MCQ $\cdot$ 2 Mark(s) $\cdot$ Neg: 0.66]}} \hfill {\footnotesize\textbf{GATE 2024 CS1} \quad \tiny\texttt{2024\_CS1\_GA\_Q09}}
\vspace{2pt}\hrule height 0.3pt\vspace{4pt}

A rectangular paper of $20\text{ cm} \times 8\text{ cm}$ is folded 3 times. Each fold is made along the line of symmetry, which is perpendicular to its long edge. The perimeter of the final folded sheet (in cm) is

\par\vspace{6pt}\noindent
\begin{tabular}{@{}p{0.48\linewidth}@{\hspace{0.04\linewidth}}p{0.48\linewidth}@{}}
$\square$\ \textbf{(A)} 18 & $\square$\ \textbf{(B)} 24 \\
\addlinespace[3pt]
$\square$\ \textbf{(C)} 20 & $\square$\ \textbf{(D)} 21
\end{tabular}

\tcblower

\textbf{Official Key:} \textbf{\large A} \hfill \textit{\small [Medium]}
\vspace{2pt}\hrule height 0.3pt\vspace{3pt}
\textbf{Topic:} Spatial Aptitude \\
\textbf{Subtopic:} \textit{Paper folding, cutting, and patterns in 2 and 3 dimensions}
\vspace{2pt}\hrule height 0.3pt\vspace{3pt}
\textbf{Guide / Hint:} \\
{\small Starting at $20 \times 8\text{ cm}$, the 1st fold perpendicular to the long edge yields $10 \times 8\text{ cm}$. The 2nd fold perpendicular to the new long edge ($10\text{ cm}$) yields $5 \times 8\text{ cm}$. The 3rd fold perpendicular to the new long edge ($8\text{ cm}$) yields $5 \times 4\text{ cm}$. The perimeter of this final folded sheet is $2(5 + 4) = 18\text{ cm}$.}

\end{tcolorbox}
\vspace{6pt}\hrule height 0.35pt\vspace{6pt}
\noindent
\begin{tcolorbox}[
    blanker,
    sidebyside,
    sidebyside align=top,
    lefthand width=0.55\textwidth,
    top=4pt, bottom=4pt, left=0pt, right=0pt,
    middle=16pt,
    underlay={
        \draw[dashed, black!80, line width=0.8pt, dash pattern=on 4pt off 3pt]
            ([xshift=0.55\textwidth+8pt]frame.north west) --
            ([xshift=0.55\textwidth+8pt]frame.south west);
    }
]
\textbf{Q.36} \quad {\footnotesize\textsc{[MCQ $\cdot$ 2 Mark(s) $\cdot$ Neg: 0.66]}} \hfill {\footnotesize\textbf{GATE 2024 CS1} \quad \tiny\texttt{2024\_CS1\_GA\_Q10}}
\vspace{2pt}\hrule height 0.3pt\vspace{4pt}

The least number of squares to be added in the figure to make $AB$ a line of symmetry is

\par\vspace{6pt}\noindent
\begin{tabular}{@{}p{0.48\linewidth}@{\hspace{0.04\linewidth}}p{0.48\linewidth}@{}}
$\square$\ \textbf{(A)} 6 & $\square$\ \textbf{(B)} 4 \\
\addlinespace[3pt]
$\square$\ \textbf{(C)} 5 & $\square$\ \textbf{(D)} 7
\end{tabular}

\tcblower

\textbf{Official Key:} \textbf{\large A} \hfill \textit{\small [Medium]}
\vspace{2pt}\hrule height 0.3pt\vspace{3pt}
\textbf{Topic:} Spatial Aptitude \\
\textbf{Subtopic:} \textit{Transformation of shapes: translation, rotation, scaling, mirroring, assembling, and grouping}
\vspace{2pt}\hrule height 0.3pt\vspace{3pt}
\textbf{Guide / Hint:} \\
{\small Reflecting each shaded cell across the diagonal axis $AB$ reveals that exactly 6 asymmetric squares must be shaded on the opposing side to achieve bilateral reflectional symmetry about line $AB$.}

\end{tcolorbox}
\vspace{6pt}\hrule height 0.35pt\vspace{6pt}
\noindent
\begin{tcolorbox}[
    blanker,
    sidebyside,
    sidebyside align=top,
    lefthand width=0.55\textwidth,
    top=4pt, bottom=4pt, left=0pt, right=0pt,
    middle=16pt,
    underlay={
        \draw[dashed, black!80, line width=0.8pt, dash pattern=on 4pt off 3pt]
            ([xshift=0.55\textwidth+8pt]frame.north west) --
            ([xshift=0.55\textwidth+8pt]frame.south west);
    }
]
\textbf{Q.37} \quad {\footnotesize\textsc{[MCQ $\cdot$ 1 Mark(s) $\cdot$ Neg: 0.33]}} \hfill {\footnotesize\textbf{GATE 2024 CS2} \quad \tiny\texttt{2024\_CS2\_GA\_Q01}}
\vspace{2pt}\hrule height 0.3pt\vspace{4pt}

If `$\to$' denotes increasing order of intensity, then the meaning of the words $[\text{walk} \to \text{jog} \to \text{sprint}]$ is analogous to $[\text{bothered} \to \text{\_\_\_\_\_\_\_\_} \to \text{daunted}]$.
Which one of the given options is appropriate to fill the blank?

\par\vspace{6pt}\noindent
\begin{tabular}{@{}p{0.48\linewidth}@{\hspace{0.04\linewidth}}p{0.48\linewidth}@{}}
$\square$\ \textbf{(A)} phased & $\square$\ \textbf{(B)} phrased \\
\addlinespace[3pt]
$\square$\ \textbf{(C)} fazed & $\square$\ \textbf{(D)} fused
\end{tabular}

\tcblower

\textbf{Official Key:} \textbf{\large C} \hfill \textit{\small [Easy]}
\vspace{2pt}\hrule height 0.3pt\vspace{3pt}
\textbf{Topic:} Verbal Aptitude \\
\textbf{Subtopic:} \textit{Vocabulary: words, idioms, and phrases in context}
\vspace{2pt}\hrule height 0.3pt\vspace{3pt}
\textbf{Guide / Hint:} \\
{\small 'Walk -> jog -> sprint' depicts an increasing scale of physical speed/exertion. Similarly, 'bothered -> fazed -> daunted' depicts an increasing degree of perturbation/discouragement. 'Fazed' means disconcerted or perturbed, fitting between bothered and daunted.}

\end{tcolorbox}
\vspace{6pt}\hrule height 0.35pt\vspace{6pt}
\noindent
\begin{tcolorbox}[
    blanker,
    sidebyside,
    sidebyside align=top,
    lefthand width=0.55\textwidth,
    top=4pt, bottom=4pt, left=0pt, right=0pt,
    middle=16pt,
    underlay={
        \draw[dashed, black!80, line width=0.8pt, dash pattern=on 4pt off 3pt]
            ([xshift=0.55\textwidth+8pt]frame.north west) --
            ([xshift=0.55\textwidth+8pt]frame.south west);
    }
]
\textbf{Q.38} \quad {\footnotesize\textsc{[MCQ $\cdot$ 1 Mark(s) $\cdot$ Neg: 0.33]}} \hfill {\footnotesize\textbf{GATE 2024 CS2} \quad \tiny\texttt{2024\_CS2\_GA\_Q02}}
\vspace{2pt}\hrule height 0.3pt\vspace{4pt}

Two wizards try to create a spell using all the four elements, water, air, fire, and earth. For this, they decide to mix all these elements in all possible orders. They also decide to work independently. After trying all possible combination of elements, they conclude that the spell does not work.
How many attempts does each wizard make before coming to this conclusion, independently?

\par\vspace{6pt}\noindent
\begin{tabular}{@{}p{0.48\linewidth}@{\hspace{0.04\linewidth}}p{0.48\linewidth}@{}}
$\square$\ \textbf{(A)} 24 & $\square$\ \textbf{(B)} 48 \\
\addlinespace[3pt]
$\square$\ \textbf{(C)} 16 & $\square$\ \textbf{(D)} 12
\end{tabular}

\tcblower

\textbf{Official Key:} \textbf{\large A} \hfill \textit{\small [Easy]}
\vspace{2pt}\hrule height 0.3pt\vspace{3pt}
\textbf{Topic:} Quantitative Aptitude \\
\textbf{Subtopic:} \textit{Numerical computation and estimation: ratios, percentages, powers, exponents and logarithms, permutations and combinations, series}
\vspace{2pt}\hrule height 0.3pt\vspace{3pt}
\textbf{Guide / Hint:} \\
{\small There are 4 distinct elements (water, air, fire, earth). The number of distinct permutations/orderings of these 4 elements is $4! = 4 \times 3 \times 2 \times 1 = 24$. Since each wizard tries all possible orders independently, each wizard makes 24 attempts.}

\end{tcolorbox}
\vspace{6pt}\hrule height 0.35pt\vspace{6pt}
\noindent
\begin{tcolorbox}[
    blanker,
    sidebyside,
    sidebyside align=top,
    lefthand width=0.55\textwidth,
    top=4pt, bottom=4pt, left=0pt, right=0pt,
    middle=16pt,
    underlay={
        \draw[dashed, black!80, line width=0.8pt, dash pattern=on 4pt off 3pt]
            ([xshift=0.55\textwidth+8pt]frame.north west) --
            ([xshift=0.55\textwidth+8pt]frame.south west);
    }
]
\textbf{Q.39} \quad {\footnotesize\textsc{[MCQ $\cdot$ 1 Mark(s) $\cdot$ Neg: 0.33]}} \hfill {\footnotesize\textbf{GATE 2024 CS2} \quad \tiny\texttt{2024\_CS2\_GA\_Q03}}
\vspace{2pt}\hrule height 0.3pt\vspace{4pt}

In an engineering college of 10,000 students, 1,500 like neither their core branches nor other branches. The number of students who like their core branches is 1/4th of the number of students who like other branches. The number of students who like both their core and other branches is 500.
The number of students who like their core branches is:

\par\vspace{6pt}\noindent
\begin{tabular}{@{}p{0.48\linewidth}@{\hspace{0.04\linewidth}}p{0.48\linewidth}@{}}
$\square$\ \textbf{(A)} 1,800 & $\square$\ \textbf{(B)} 3,500 \\
\addlinespace[3pt]
$\square$\ \textbf{(C)} 1,600 & $\square$\ \textbf{(D)} 1,500
\end{tabular}

\tcblower

\textbf{Official Key:} \textbf{\large A} \hfill \textit{\small [Medium]}
\vspace{2pt}\hrule height 0.3pt\vspace{3pt}
\textbf{Topic:} Quantitative Aptitude \\
\textbf{Subtopic:} \textit{Data interpretation: data graphs (bar graphs, pie charts, and other graphs representing data), 2- and 3-dimensional plots, maps, and tables}
\vspace{2pt}\hrule height 0.3pt\vspace{3pt}
\textbf{Guide / Hint:} \\
{\small Let $C$ and $O$ be sets of students liking core and other branches. Total students $N = 10000$, neither $= 1500$, so $|C \cup O| = 10000 - 1500 = 8500$. We are given $|C \cap O| = 500$ and $|C| = \frac{1}{4}|O| \implies |O| = 4|C|$. By inclusion-exclusion, $|C \cup O| = |C| + |O| - |C \cap O| \implies 8500 = |C| + 4|C| - 500 \implies 5|C| = 9000 \implies |C| = 1800$.}

\end{tcolorbox}
\vspace{6pt}\hrule height 0.35pt\vspace{6pt}
\noindent
\begin{tcolorbox}[
    blanker,
    sidebyside,
    sidebyside align=top,
    lefthand width=0.55\textwidth,
    top=4pt, bottom=4pt, left=0pt, right=0pt,
    middle=16pt,
    underlay={
        \draw[dashed, black!80, line width=0.8pt, dash pattern=on 4pt off 3pt]
            ([xshift=0.55\textwidth+8pt]frame.north west) --
            ([xshift=0.55\textwidth+8pt]frame.south west);
    }
]
\textbf{Q.40} \quad {\footnotesize\textsc{[MCQ $\cdot$ 1 Mark(s) $\cdot$ Neg: 0.33]}} \hfill {\footnotesize\textbf{GATE 2024 CS2} \quad \tiny\texttt{2024\_CS2\_GA\_Q04}}
\vspace{2pt}\hrule height 0.3pt\vspace{4pt}

For positive non-zero real variables $x$ and $y$, if
\[ \ln\left(\frac{x+y}{2}\right) = \frac{1}{2} [\ln(x) + \ln(y)] \]
then, the value of $\frac{x}{y} + \frac{y}{x}$ is:

\par\vspace{6pt}\noindent
\begin{tabular}{@{}p{0.48\linewidth}@{\hspace{0.04\linewidth}}p{0.48\linewidth}@{}}
$\square$\ \textbf{(A)} 1 & $\square$\ \textbf{(B)} 1/2 \\
\addlinespace[3pt]
$\square$\ \textbf{(C)} 2 & $\square$\ \textbf{(D)} 4
\end{tabular}

\tcblower

\textbf{Official Key:} \textbf{\large C} \hfill \textit{\small [Easy]}
\vspace{2pt}\hrule height 0.3pt\vspace{3pt}
\textbf{Topic:} Quantitative Aptitude \\
\textbf{Subtopic:} \textit{Numerical computation and estimation: ratios, percentages, powers, exponents and logarithms, permutations and combinations, series}
\vspace{2pt}\hrule height 0.3pt\vspace{3pt}
\textbf{Guide / Hint:} \\
{\small From $\ln\left(\frac{x+y}{2}\right) = \frac{1}{2}[\ln(x) + \ln(y)] = \ln(\sqrt{xy})$, we obtain $\frac{x+y}{2} = \sqrt{xy} \implies x + y = 2\sqrt{xy} \implies (\sqrt{x} - \sqrt{y})^2 = 0 \implies x = y$. Substituting $x = y$ gives $\frac{x}{y} + \frac{y}{x} = 1 + 1 = 2$.}

\end{tcolorbox}
\vspace{6pt}\hrule height 0.35pt\vspace{6pt}
\noindent
\begin{tcolorbox}[
    blanker,
    sidebyside,
    sidebyside align=top,
    lefthand width=0.55\textwidth,
    top=4pt, bottom=4pt, left=0pt, right=0pt,
    middle=16pt,
    underlay={
        \draw[dashed, black!80, line width=0.8pt, dash pattern=on 4pt off 3pt]
            ([xshift=0.55\textwidth+8pt]frame.north west) --
            ([xshift=0.55\textwidth+8pt]frame.south west);
    }
]
\textbf{Q.41} \quad {\footnotesize\textsc{[MCQ $\cdot$ 1 Mark(s) $\cdot$ Neg: 0.33]}} \hfill {\footnotesize\textbf{GATE 2024 CS2} \quad \tiny\texttt{2024\_CS2\_GA\_Q05}}
\vspace{2pt}\hrule height 0.3pt\vspace{4pt}

In the sequence 6, 9, 14, $x$, 30, 41, a possible value of $x$ is:

\par\vspace{6pt}\noindent
\begin{tabular}{@{}p{0.48\linewidth}@{\hspace{0.04\linewidth}}p{0.48\linewidth}@{}}
$\square$\ \textbf{(A)} 25 & $\square$\ \textbf{(B)} 21 \\
\addlinespace[3pt]
$\square$\ \textbf{(C)} 18 & $\square$\ \textbf{(D)} 20
\end{tabular}

\tcblower

\textbf{Official Key:} \textbf{\large B} \hfill \textit{\small [Easy]}
\vspace{2pt}\hrule height 0.3pt\vspace{3pt}
\textbf{Topic:} Quantitative Aptitude \\
\textbf{Subtopic:} \textit{Numerical computation and estimation: ratios, percentages, powers, exponents and logarithms, permutations and combinations, series}
\vspace{2pt}\hrule height 0.3pt\vspace{3pt}
\textbf{Guide / Hint:} \\
{\small Examine the consecutive differences: $9 - 6 = 3$, $14 - 9 = 5$, $x - 14 = d_3$, $30 - x = d_4$, $41 - 30 = 11$. The differences are consecutive odd numbers: $3, 5, 7, 9, 11$. For $d_3 = 7$, $x = 14 + 7 = 21$, and checking $d_4 = 30 - 21 = 9$ matches. Hence, $x = 21$.}

\end{tcolorbox}
\vspace{6pt}\hrule height 0.35pt\vspace{6pt}
\noindent
\begin{tcolorbox}[
    blanker,
    sidebyside,
    sidebyside align=top,
    lefthand width=0.55\textwidth,
    top=4pt, bottom=4pt, left=0pt, right=0pt,
    middle=16pt,
    underlay={
        \draw[dashed, black!80, line width=0.8pt, dash pattern=on 4pt off 3pt]
            ([xshift=0.55\textwidth+8pt]frame.north west) --
            ([xshift=0.55\textwidth+8pt]frame.south west);
    }
]
\textbf{Q.42} \quad {\footnotesize\textsc{[MCQ $\cdot$ 2 Mark(s) $\cdot$ Neg: 0.66]}} \hfill {\footnotesize\textbf{GATE 2024 CS2} \quad \tiny\texttt{2024\_CS2\_GA\_Q06}}
\vspace{2pt}\hrule height 0.3pt\vspace{4pt}

Sequence the following sentences in a coherent passage.
\begin{itemize}
\item[P:] This fortuitous geological event generated a colossal amount of energy and heat that resulted in the rocks rising to an average height of 4 km across the contact zone.
\item[Q:] Thus, the geophysicists tend to think of the Himalayas as an active geological event rather than as a static geological feature.
\item[R:] The natural process of the cooling of this massive edifice absorbed large quantities of atmospheric carbon dioxide, altering the earth’s atmosphere and making it better suited for life.
\item[S:] Many millennia ago, a breakaway chunk of bedrock from the Antarctic Plate collided with the massive Eurasian Plate.
\end{itemize}

\par\vspace{6pt}\noindent
\begin{tabular}{@{}p{0.48\linewidth}@{\hspace{0.04\linewidth}}p{0.48\linewidth}@{}}
$\square$\ \textbf{(A)} QPSR & $\square$\ \textbf{(B)} QSPR \\
\addlinespace[3pt]
$\square$\ \textbf{(C)} SPRQ & $\square$\ \textbf{(D)} SRPQ
\end{tabular}

\tcblower

\textbf{Official Key:} \textbf{\large C} \hfill \textit{\small [Medium]}
\vspace{2pt}\hrule height 0.3pt\vspace{3pt}
\textbf{Topic:} Verbal Aptitude \\
\textbf{Subtopic:} \textit{Narrative sequencing}
\vspace{2pt}\hrule height 0.3pt\vspace{3pt}
\textbf{Guide / Hint:} \\
{\small Sentence S introduces the foundational historical premise (many millennia ago, Antarctic plate collided with Eurasian plate). Sentence P describes the immediate effect of this event (energy/heat rising to 4 km height). Sentence R details the subsequent cooling that altered the atmosphere. Sentence Q concludes with the geophysicists' resulting viewpoint. The logical sequence is SPRQ.}

\end{tcolorbox}
\vspace{6pt}\hrule height 0.35pt\vspace{6pt}
\noindent
\begin{tcolorbox}[
    blanker,
    sidebyside,
    sidebyside align=top,
    lefthand width=0.55\textwidth,
    top=4pt, bottom=4pt, left=0pt, right=0pt,
    middle=16pt,
    underlay={
        \draw[dashed, black!80, line width=0.8pt, dash pattern=on 4pt off 3pt]
            ([xshift=0.55\textwidth+8pt]frame.north west) --
            ([xshift=0.55\textwidth+8pt]frame.south west);
    }
]
\textbf{Q.43} \quad {\footnotesize\textsc{[MCQ $\cdot$ 2 Mark(s) $\cdot$ Neg: 0.66]}} \hfill {\footnotesize\textbf{GATE 2024 CS2} \quad \tiny\texttt{2024\_CS2\_GA\_Q07}}
\vspace{2pt}\hrule height 0.3pt\vspace{4pt}

A person sold two different items at the same price. He made 10\% profit in one item, and 10\% loss in the other item. In selling these two items, the person made a total of:

\par\vspace{6pt}\noindent
\begin{tabular}{@{}p{0.48\linewidth}@{\hspace{0.04\linewidth}}p{0.48\linewidth}@{}}
$\square$\ \textbf{(A)} 1\% profit & $\square$\ \textbf{(B)} 2\% profit \\
\addlinespace[3pt]
$\square$\ \textbf{(C)} 1\% loss & $\square$\ \textbf{(D)} 2\% loss
\end{tabular}

\tcblower

\textbf{Official Key:} \textbf{\large C} \hfill \textit{\small [Easy]}
\vspace{2pt}\hrule height 0.3pt\vspace{3pt}
\textbf{Topic:} Quantitative Aptitude \\
\textbf{Subtopic:} \textit{Numerical computation and estimation: ratios, percentages, powers, exponents and logarithms, permutations and combinations, series}
\vspace{2pt}\hrule height 0.3pt\vspace{3pt}
\textbf{Guide / Hint:} \\
{\small When two items are sold at the same selling price $S$, one at a profit of $x\%$ and the other at a loss of $x\%$, there is always an overall net loss given by $\frac{x^2}{100}\%$. Here $x = 10$, so the net loss is $\frac{10^2}{100}\% = 1\%$ loss.}

\end{tcolorbox}
\vspace{6pt}\hrule height 0.35pt\vspace{6pt}
\noindent
\begin{tcolorbox}[
    blanker,
    sidebyside,
    sidebyside align=top,
    lefthand width=0.55\textwidth,
    top=4pt, bottom=4pt, left=0pt, right=0pt,
    middle=16pt,
    underlay={
        \draw[dashed, black!80, line width=0.8pt, dash pattern=on 4pt off 3pt]
            ([xshift=0.55\textwidth+8pt]frame.north west) --
            ([xshift=0.55\textwidth+8pt]frame.south west);
    }
]
\textbf{Q.44} \quad {\footnotesize\textsc{[MCQ $\cdot$ 2 Mark(s) $\cdot$ Neg: 0.66]}} \hfill {\footnotesize\textbf{GATE 2024 CS2} \quad \tiny\texttt{2024\_CS2\_GA\_Q09}}
\vspace{2pt}\hrule height 0.3pt\vspace{4pt}

A cube is to be cut into 8 pieces of equal size and shape. Here, each cut should be straight and it should not stop till it reaches the other end of the cube.
The minimum number of such cuts required is:

\par\vspace{6pt}\noindent
\begin{tabular}{@{}p{0.48\linewidth}@{\hspace{0.04\linewidth}}p{0.48\linewidth}@{}}
$\square$\ \textbf{(A)} 3 & $\square$\ \textbf{(B)} 4 \\
\addlinespace[3pt]
$\square$\ \textbf{(C)} 7 & $\square$\ \textbf{(D)} 8
\end{tabular}

\tcblower

\textbf{Official Key:} \textbf{\large A} \hfill \textit{\small [Easy]}
\vspace{2pt}\hrule height 0.3pt\vspace{3pt}
\textbf{Topic:} Spatial Aptitude \\
\textbf{Subtopic:} \textit{Paper folding, cutting, and patterns in 2 and 3 dimensions}
\vspace{2pt}\hrule height 0.3pt\vspace{3pt}
\textbf{Guide / Hint:} \\
{\small A single full planar straight cut can at most double the number of pieces. With $k$ cuts through 3 orthogonal dimensions (1 cut parallel to $xy$-plane, 1 parallel to $yz$-plane, 1 parallel to $zx$-plane), the cube is divided into $2 \times 2 \times 2 = 8$ equal pieces. Thus, a minimum of 3 cuts is required.}

\end{tcolorbox}
\vspace{6pt}\hrule height 0.35pt\vspace{6pt}
\noindent
\begin{tcolorbox}[
    blanker,
    sidebyside,
    sidebyside align=top,
    lefthand width=0.55\textwidth,
    top=4pt, bottom=4pt, left=0pt, right=0pt,
    middle=16pt,
    underlay={
        \draw[dashed, black!80, line width=0.8pt, dash pattern=on 4pt off 3pt]
            ([xshift=0.55\textwidth+8pt]frame.north west) --
            ([xshift=0.55\textwidth+8pt]frame.south west);
    }
]
\textbf{Q.45} \quad {\footnotesize\textsc{[MCQ $\cdot$ 2 Mark(s) $\cdot$ Neg: 0.66]}} \hfill {\footnotesize\textbf{GATE 2024 CS2} \quad \tiny\texttt{2024\_CS2\_GA\_Q10}}
\vspace{2pt}\hrule height 0.3pt\vspace{4pt}

In a $4 \times 4$ array, each cell of the first three rows has either a cross (X) or a number. The number in a cell represents the count of the immediate neighboring cells (left, right, top, bottom, diagonals) NOT having a cross (X). Given that the last row has no crosses (X), the sum of the four numbers to be filled in the last row is:

\par\vspace{6pt}\noindent
\begin{tabular}{@{}p{0.48\linewidth}@{\hspace{0.04\linewidth}}p{0.48\linewidth}@{}}
$\square$\ \textbf{(A)} 11 & $\square$\ \textbf{(B)} 10 \\
\addlinespace[3pt]
$\square$\ \textbf{(C)} 12 & $\square$\ \textbf{(D)} 9
\end{tabular}

\tcblower

\textbf{Official Key:} \textbf{\large A} \hfill \textit{\small [Hard]}
\vspace{2pt}\hrule height 0.3pt\vspace{3pt}
\textbf{Topic:} Analytical Aptitude \\
\textbf{Subtopic:} \textit{Logic: deduction and induction}
\vspace{2pt}\hrule height 0.3pt\vspace{3pt}
\textbf{Guide / Hint:} \\
{\small Each number in a cell denotes how many neighboring cells do NOT contain an 'X'. With crosses determined from the first three rows, the count of non-cross neighbors for each of the 4 cells in the bottom row evaluates to values whose total sum is 11.}

\end{tcolorbox}
\vspace{6pt}\hrule height 0.35pt\vspace{6pt}
\noindent
\begin{tcolorbox}[
    blanker,
    sidebyside,
    sidebyside align=top,
    lefthand width=0.55\textwidth,
    top=4pt, bottom=4pt, left=0pt, right=0pt,
    middle=16pt,
    underlay={
        \draw[dashed, black!80, line width=0.8pt, dash pattern=on 4pt off 3pt]
            ([xshift=0.55\textwidth+8pt]frame.north west) --
            ([xshift=0.55\textwidth+8pt]frame.south west);
    }
]
\textbf{Q.46} \quad {\footnotesize\textsc{[MCQ $\cdot$ 1 Mark(s) $\cdot$ Neg: 0.33]}} \hfill {\footnotesize\textbf{GATE 2025 CS1} \quad \tiny\texttt{2025\_CS1\_GA\_Q01}}
\vspace{2pt}\hrule height 0.3pt\vspace{4pt}

Ravi had \_\_\_\_\_\_ younger brother who taught at \_\_\_\_\_\_ university. He was widely regarded as \_\_\_\_\_\_ honorable man.
Select the option with the correct sequence of articles to fill in the blanks.

\par\vspace{6pt}\noindent
\begin{tabular}{@{}p{0.48\linewidth}@{\hspace{0.04\linewidth}}p{0.48\linewidth}@{}}
$\square$\ \textbf{(A)} a; a; an & $\square$\ \textbf{(B)} the; an; a \\
\addlinespace[3pt]
$\square$\ \textbf{(C)} a; an; a & $\square$\ \textbf{(D)} an; an; a
\end{tabular}

\tcblower

\textbf{Official Key:} \textbf{\large A} \hfill \textit{\small [Easy]}
\vspace{2pt}\hrule height 0.3pt\vspace{3pt}
\textbf{Topic:} Verbal Aptitude \\
\textbf{Subtopic:} \textit{Basic English grammar: tenses, articles, adjectives, prepositions, conjunctions, verb-noun agreement}
\vspace{2pt}\hrule height 0.3pt\vspace{3pt}
\textbf{Guide / Hint:} \\
{\small 'Younger' begins with a consonant sound (/j/), requiring the indefinite article 'a'. 'University' also begins with a consonant sound (/juː/), taking 'a'. 'Honorable' has a silent 'h' and begins with a vowel sound (/ɒ/), requiring 'an'. Hence, the sequence is a; a; an.}

\end{tcolorbox}
\vspace{6pt}\hrule height 0.35pt\vspace{6pt}
\noindent
\begin{tcolorbox}[
    blanker,
    sidebyside,
    sidebyside align=top,
    lefthand width=0.55\textwidth,
    top=4pt, bottom=4pt, left=0pt, right=0pt,
    middle=16pt,
    underlay={
        \draw[dashed, black!80, line width=0.8pt, dash pattern=on 4pt off 3pt]
            ([xshift=0.55\textwidth+8pt]frame.north west) --
            ([xshift=0.55\textwidth+8pt]frame.south west);
    }
]
\textbf{Q.47} \quad {\footnotesize\textsc{[MCQ $\cdot$ 1 Mark(s) $\cdot$ Neg: 0.33]}} \hfill {\footnotesize\textbf{GATE 2025 CS1} \quad \tiny\texttt{2025\_CS1\_GA\_Q02}}
\vspace{2pt}\hrule height 0.3pt\vspace{4pt}

The CEO’s decision to downsize the workforce was considered myopic because it sacrificed long-term stability to accommodate short-term gains.
Select the most appropriate option that can replace the word ``myopic'' without changing the meaning of the sentence.

\par\vspace{6pt}\noindent
\begin{tabular}{@{}p{0.48\linewidth}@{\hspace{0.04\linewidth}}p{0.48\linewidth}@{}}
$\square$\ \textbf{(A)} visionary & $\square$\ \textbf{(B)} shortsighted \\
\addlinespace[3pt]
$\square$\ \textbf{(C)} progressive & $\square$\ \textbf{(D)} innovative
\end{tabular}

\tcblower

\textbf{Official Key:} \textbf{\large B} \hfill \textit{\small [Easy]}
\vspace{2pt}\hrule height 0.3pt\vspace{3pt}
\textbf{Topic:} Verbal Aptitude \\
\textbf{Subtopic:} \textit{Vocabulary: words, idioms, and phrases in context}
\vspace{2pt}\hrule height 0.3pt\vspace{3pt}
\textbf{Guide / Hint:} \\
{\small 'Myopic' literally means nearsighted, and metaphorically refers to lacking foresight or being shortsighted by focusing only on immediate considerations rather than long-term outcomes.}

\end{tcolorbox}
\vspace{6pt}\hrule height 0.35pt\vspace{6pt}
\noindent
\begin{tcolorbox}[
    blanker,
    sidebyside,
    sidebyside align=top,
    lefthand width=0.55\textwidth,
    top=4pt, bottom=4pt, left=0pt, right=0pt,
    middle=16pt,
    underlay={
        \draw[dashed, black!80, line width=0.8pt, dash pattern=on 4pt off 3pt]
            ([xshift=0.55\textwidth+8pt]frame.north west) --
            ([xshift=0.55\textwidth+8pt]frame.south west);
    }
]
\textbf{Q.48} \quad {\footnotesize\textsc{[MCQ $\cdot$ 1 Mark(s) $\cdot$ Neg: 0.33]}} \hfill {\footnotesize\textbf{GATE 2025 CS1} \quad \tiny\texttt{2025\_CS1\_GA\_Q03}}
\vspace{2pt}\hrule height 0.3pt\vspace{4pt}

The average marks obtained by a class in an examination were calculated as 30.8. However, while checking the marks entered, the teacher found that the marks of one student were entered incorrectly as 24 instead of 42. After correcting the marks, the average becomes 31.4. How many students does the class have?

\par\vspace{6pt}\noindent
\begin{tabular}{@{}p{0.48\linewidth}@{\hspace{0.04\linewidth}}p{0.48\linewidth}@{}}
$\square$\ \textbf{(A)} 25 & $\square$\ \textbf{(B)} 28 \\
\addlinespace[3pt]
$\square$\ \textbf{(C)} 30 & $\square$\ \textbf{(D)} 32
\end{tabular}

\tcblower

\textbf{Official Key:} \textbf{\large C} \hfill \textit{\small [Easy]}
\vspace{2pt}\hrule height 0.3pt\vspace{3pt}
\textbf{Topic:} Quantitative Aptitude \\
\textbf{Subtopic:} \textit{Elementary statistics and probability}
\vspace{2pt}\hrule height 0.3pt\vspace{3pt}
\textbf{Guide / Hint:} \\
{\small Let $n$ be the number of students. The total increase in marks after correction is $42 - 24 = 18$. The increase in average is $31.4 - 30.8 = 0.6$. Therefore, $0.6 \times n = 18$, giving $n = \frac{18}{0.6} = 30$ students.}

\end{tcolorbox}
\vspace{6pt}\hrule height 0.35pt\vspace{6pt}
\noindent
\begin{tcolorbox}[
    blanker,
    sidebyside,
    sidebyside align=top,
    lefthand width=0.55\textwidth,
    top=4pt, bottom=4pt, left=0pt, right=0pt,
    middle=16pt,
    underlay={
        \draw[dashed, black!80, line width=0.8pt, dash pattern=on 4pt off 3pt]
            ([xshift=0.55\textwidth+8pt]frame.north west) --
            ([xshift=0.55\textwidth+8pt]frame.south west);
    }
]
\textbf{Q.49} \quad {\footnotesize\textsc{[MCQ $\cdot$ 1 Mark(s) $\cdot$ Neg: 0.33]}} \hfill {\footnotesize\textbf{GATE 2025 CS1} \quad \tiny\texttt{2025\_CS1\_GA\_Q04}}
\vspace{2pt}\hrule height 0.3pt\vspace{4pt}

Consider the relationships among P, Q, R, S, and T:
\begin{itemize}
\item P is the brother of Q.
\item S is the daughter of Q.
\item T is the sister of S.
\item R is the mother of Q.
\end{itemize}
The following statements are made based on the relationships given above:
(1) R is the grandmother of S.
(2) P is the uncle of S and T.
(3) R has only one son.
(4) Q has only one daughter.
Which one of the following options is correct?

\par\vspace{6pt}\noindent
\begin{tabular}{@{}p{0.48\linewidth}@{\hspace{0.04\linewidth}}p{0.48\linewidth}@{}}
$\square$\ \textbf{(A)} Both (1) and (2) are true. & $\square$\ \textbf{(B)} Both (1) and (3) are true. \\
\addlinespace[3pt]
$\square$\ \textbf{(C)} Only (3) is true. & $\square$\ \textbf{(D)} Only (4) is true.
\end{tabular}

\tcblower

\textbf{Official Key:} \textbf{\large A} \hfill \textit{\small [Easy]}
\vspace{2pt}\hrule height 0.3pt\vspace{3pt}
\textbf{Topic:} Analytical Aptitude \\
\textbf{Subtopic:} \textit{Logic: deduction and induction}
\vspace{2pt}\hrule height 0.3pt\vspace{3pt}
\textbf{Guide / Hint:} \\
{\small Since R is the mother of Q and S is the daughter of Q, R is definitely the grandmother of S (statement 1 is true). P is the brother of Q, so P is the maternal/paternal uncle of Q's children S and T (statement 2 is true). The genders and number of other children of R and Q are not specified, so (3) and (4) cannot be established.}

\end{tcolorbox}
\vspace{6pt}\hrule height 0.35pt\vspace{6pt}
\noindent
\begin{tcolorbox}[
    blanker,
    sidebyside,
    sidebyside align=top,
    lefthand width=0.55\textwidth,
    top=4pt, bottom=4pt, left=0pt, right=0pt,
    middle=16pt,
    underlay={
        \draw[dashed, black!80, line width=0.8pt, dash pattern=on 4pt off 3pt]
            ([xshift=0.55\textwidth+8pt]frame.north west) --
            ([xshift=0.55\textwidth+8pt]frame.south west);
    }
]
\textbf{Q.50} \quad {\footnotesize\textsc{[MCQ $\cdot$ 2 Mark(s) $\cdot$ Neg: 0.66]}} \hfill {\footnotesize\textbf{GATE 2025 CS1} \quad \tiny\texttt{2025\_CS1\_GA\_Q06}}
\vspace{2pt}\hrule height 0.3pt\vspace{4pt}

``I put the brown paper in my pocket along with the chalks, and possibly other things. I suppose every one must have reflected how primeval and how poetical are the things that one carries in one's pocket: the pocket-knife, for instance the type of all human tools, the infant of the sword. Once I planned to write a book of poems entirely about the things in my pocket. But I found it would be too long: and the age of the great epics is past.''
(From G.K. Chesterton's ``A Piece of Chalk'')
\par\medskip
Based only on the information provided in the above passage, which one of the following statements is true?

\par\vspace{6pt}\noindent
\begin{tabular}{@{}p{0.98\linewidth}@{}}
$\square$\ \textbf{(A)} The author of the passage carries a mirror in his pocket to reflect upon things. \\
\addlinespace[3pt]
$\square$\ \textbf{(B)} The author of the passage had decided to write a poem on epics. \\
\addlinespace[3pt]
$\square$\ \textbf{(C)} The pocket-knife is described as the infant of the sword. \\
\addlinespace[3pt]
$\square$\ \textbf{(D)} Epics are described as too inconvenient to write.
\end{tabular}

\tcblower

\textbf{Official Key:} \textbf{\large C} \hfill \textit{\small [Easy]}
\vspace{2pt}\hrule height 0.3pt\vspace{3pt}
\textbf{Topic:} Verbal Aptitude \\
\textbf{Subtopic:} \textit{Reading and comprehension}
\vspace{2pt}\hrule height 0.3pt\vspace{3pt}
\textbf{Guide / Hint:} \\
{\small The passage explicitly states: 'the pocket-knife, for instance the type of all human tools, the infant of the sword.' Statement C directly reflects this verbatim description.}

\end{tcolorbox}
\vspace{6pt}\hrule height 0.35pt\vspace{6pt}
\noindent
\begin{tcolorbox}[
    blanker,
    sidebyside,
    sidebyside align=top,
    lefthand width=0.55\textwidth,
    top=4pt, bottom=4pt, left=0pt, right=0pt,
    middle=16pt,
    underlay={
        \draw[dashed, black!80, line width=0.8pt, dash pattern=on 4pt off 3pt]
            ([xshift=0.55\textwidth+8pt]frame.north west) --
            ([xshift=0.55\textwidth+8pt]frame.south west);
    }
]
\textbf{Q.51} \quad {\footnotesize\textsc{[MCQ $\cdot$ 2 Mark(s) $\cdot$ Neg: 0.66]}} \hfill {\footnotesize\textbf{GATE 2025 CS1} \quad \tiny\texttt{2025\_CS1\_GA\_Q08}}
\vspace{2pt}\hrule height 0.3pt\vspace{4pt}

A fair six-faced dice, with the faces labelled `1', `2', `3', `4', `5', and `6', is rolled thrice. What is the probability of rolling `6' exactly once?

\par\vspace{6pt}\noindent
\begin{tabular}{@{}p{0.48\linewidth}@{\hspace{0.04\linewidth}}p{0.48\linewidth}@{}}
$\square$\ \textbf{(A)} $\frac{75}{216}$ & $\square$\ \textbf{(B)} $\frac{1}{6}$ \\
\addlinespace[3pt]
$\square$\ \textbf{(C)} $\frac{1}{18}$ & $\square$\ \textbf{(D)} $\frac{25}{216}$
\end{tabular}

\tcblower

\textbf{Official Key:} \textbf{\large A} \hfill \textit{\small [Medium]}
\vspace{2pt}\hrule height 0.3pt\vspace{3pt}
\textbf{Topic:} Quantitative Aptitude \\
\textbf{Subtopic:} \textit{Elementary statistics and probability}
\vspace{2pt}\hrule height 0.3pt\vspace{3pt}
\textbf{Guide / Hint:} \\
{\small The probability of rolling a 6 is $p = \frac{1}{6}$, and not rolling a 6 is $q = \frac{5}{6}$. Using the binomial distribution for $n = 3$ trials with $k = 1$ success, $P(X = 1) = \binom{3}{1} \left(\frac{1}{6}\right)^1 \left(\frac{5}{6}\right)^2 = 3 \times \frac{1}{6} \times \frac{25}{36} = \frac{75}{216}$.}

\end{tcolorbox}
\vspace{6pt}\hrule height 0.35pt\vspace{6pt}
\noindent
\begin{tcolorbox}[
    blanker,
    sidebyside,
    sidebyside align=top,
    lefthand width=0.55\textwidth,
    top=4pt, bottom=4pt, left=0pt, right=0pt,
    middle=16pt,
    underlay={
        \draw[dashed, black!80, line width=0.8pt, dash pattern=on 4pt off 3pt]
            ([xshift=0.55\textwidth+8pt]frame.north west) --
            ([xshift=0.55\textwidth+8pt]frame.south west);
    }
]
\textbf{Q.52} \quad {\footnotesize\textsc{[MCQ $\cdot$ 2 Mark(s) $\cdot$ Neg: 0.66]}} \hfill {\footnotesize\textbf{GATE 2025 CS1} \quad \tiny\texttt{2025\_CS1\_GA\_Q10}}
\vspace{2pt}\hrule height 0.3pt\vspace{4pt}

A shop has 4 distinct flavors of ice-cream. One can purchase any number of scoops of any flavor. The order in which the scoops are purchased is inconsequential. If one wants to purchase 3 scoops of ice-cream, in how many ways can one make that purchase?

\par\vspace{6pt}\noindent
\begin{tabular}{@{}p{0.48\linewidth}@{\hspace{0.04\linewidth}}p{0.48\linewidth}@{}}
$\square$\ \textbf{(A)} 4 & $\square$\ \textbf{(B)} 20 \\
\addlinespace[3pt]
$\square$\ \textbf{(C)} 24 & $\square$\ \textbf{(D)} 48
\end{tabular}

\tcblower

\textbf{Official Key:} \textbf{\large B} \hfill \textit{\small [Easy]}
\vspace{2pt}\hrule height 0.3pt\vspace{3pt}
\textbf{Topic:} Quantitative Aptitude \\
\textbf{Subtopic:} \textit{Numerical computation and estimation: ratios, percentages, powers, exponents and logarithms, permutations and combinations, series}
\vspace{2pt}\hrule height 0.3pt\vspace{3pt}
\textbf{Guide / Hint:} \\
{\small This is a multichoose problem (combination with repetition), where $r = 3$ scoops are chosen from $n = 4$ flavors. The number of ways is given by the stars and bars formula $\binom{n + r - 1}{r} = \binom{4 + 3 - 1}{3} = \binom{6}{3} = \frac{6 \times 5 \times 4}{3 \times 2 \times 1} = 20$.}

\end{tcolorbox}
\vspace{6pt}\hrule height 0.35pt\vspace{6pt}
\noindent
\begin{tcolorbox}[
    blanker,
    sidebyside,
    sidebyside align=top,
    lefthand width=0.55\textwidth,
    top=4pt, bottom=4pt, left=0pt, right=0pt,
    middle=16pt,
    underlay={
        \draw[dashed, black!80, line width=0.8pt, dash pattern=on 4pt off 3pt]
            ([xshift=0.55\textwidth+8pt]frame.north west) --
            ([xshift=0.55\textwidth+8pt]frame.south west);
    }
]
\textbf{Q.53} \quad {\footnotesize\textsc{[MCQ $\cdot$ 1 Mark(s) $\cdot$ Neg: 0.33]}} \hfill {\footnotesize\textbf{GATE 2025 CS2} \quad \tiny\texttt{2025\_CS2\_GA\_Q01}}
\vspace{2pt}\hrule height 0.3pt\vspace{4pt}

Despite his initial hesitation, Rehman’s \_\_\_\_\_\_\_\_\_ to contribute to the success of the project never wavered.
Select the most appropriate option to complete the above sentence.

\par\vspace{6pt}\noindent
\begin{tabular}{@{}p{0.48\linewidth}@{\hspace{0.04\linewidth}}p{0.48\linewidth}@{}}
$\square$\ \textbf{(A)} ambivalence & $\square$\ \textbf{(B)} satisfaction \\
\addlinespace[3pt]
$\square$\ \textbf{(C)} resolve & $\square$\ \textbf{(D)} revolve
\end{tabular}

\tcblower

\textbf{Official Key:} \textbf{\large C} \hfill \textit{\small [Easy]}
\vspace{2pt}\hrule height 0.3pt\vspace{3pt}
\textbf{Topic:} Verbal Aptitude \\
\textbf{Subtopic:} \textit{Vocabulary: words, idioms, and phrases in context}
\vspace{2pt}\hrule height 0.3pt\vspace{3pt}
\textbf{Guide / Hint:} \\
{\small The sentence sets up a contrast between Rehman's initial hesitation and his continuing commitment. 'Resolve' denotes firm determination and perseverance, which fits the context accurately.}

\end{tcolorbox}
\vspace{6pt}\hrule height 0.35pt\vspace{6pt}
\noindent
\begin{tcolorbox}[
    blanker,
    sidebyside,
    sidebyside align=top,
    lefthand width=0.55\textwidth,
    top=4pt, bottom=4pt, left=0pt, right=0pt,
    middle=16pt,
    underlay={
        \draw[dashed, black!80, line width=0.8pt, dash pattern=on 4pt off 3pt]
            ([xshift=0.55\textwidth+8pt]frame.north west) --
            ([xshift=0.55\textwidth+8pt]frame.south west);
    }
]
\textbf{Q.54} \quad {\footnotesize\textsc{[MCQ $\cdot$ 1 Mark(s) $\cdot$ Neg: 0.33]}} \hfill {\footnotesize\textbf{GATE 2025 CS2} \quad \tiny\texttt{2025\_CS2\_GA\_Q02}}
\vspace{2pt}\hrule height 0.3pt\vspace{4pt}

$\text{Bird} : \text{Nest} :: \text{Bee} : \text{\_\_\_\_\_\_\_}$
Select the correct option to complete the analogy.

\par\vspace{6pt}\noindent
\begin{tabular}{@{}p{0.48\linewidth}@{\hspace{0.04\linewidth}}p{0.48\linewidth}@{}}
$\square$\ \textbf{(A)} Kennel & $\square$\ \textbf{(B)} Hammock \\
\addlinespace[3pt]
$\square$\ \textbf{(C)} Hive & $\square$\ \textbf{(D)} Lair
\end{tabular}

\tcblower

\textbf{Official Key:} \textbf{\large C} \hfill \textit{\small [Easy]}
\vspace{2pt}\hrule height 0.3pt\vspace{3pt}
\textbf{Topic:} Verbal Aptitude \\
\textbf{Subtopic:} \textit{Vocabulary: word analogy and contextual relations}
\vspace{2pt}\hrule height 0.3pt\vspace{3pt}
\textbf{Guide / Hint:} \\
{\small This is a habitat/dwelling analogy: a bird naturally lives in a nest, and similarly, a bee lives in a hive.}

\end{tcolorbox}
\vspace{6pt}\hrule height 0.35pt\vspace{6pt}
\noindent
\begin{tcolorbox}[
    blanker,
    sidebyside,
    sidebyside align=top,
    lefthand width=0.55\textwidth,
    top=4pt, bottom=4pt, left=0pt, right=0pt,
    middle=16pt,
    underlay={
        \draw[dashed, black!80, line width=0.8pt, dash pattern=on 4pt off 3pt]
            ([xshift=0.55\textwidth+8pt]frame.north west) --
            ([xshift=0.55\textwidth+8pt]frame.south west);
    }
]
\textbf{Q.55} \quad {\footnotesize\textsc{[MCQ $\cdot$ 1 Mark(s) $\cdot$ Neg: 0.33]}} \hfill {\footnotesize\textbf{GATE 2025 CS2} \quad \tiny\texttt{2025\_CS2\_GA\_Q03}}
\vspace{2pt}\hrule height 0.3pt\vspace{4pt}

If $P e^x = Q e^{-x}$ for all real values of $x$, which one of the following statements is true?

\par\vspace{6pt}\noindent
\begin{tabular}{@{}p{0.48\linewidth}@{\hspace{0.04\linewidth}}p{0.48\linewidth}@{}}
$\square$\ \textbf{(A)} $P = Q = 0$ & $\square$\ \textbf{(B)} $P = Q = 1$ \\
\addlinespace[3pt]
$\square$\ \textbf{(C)} $P = 1; Q = -1$ & $\square$\ \textbf{(D)} $\frac{P}{Q} = 0$
\end{tabular}

\tcblower

\textbf{Official Key:} \textbf{\large A} \hfill \textit{\small [Medium]}
\vspace{2pt}\hrule height 0.3pt\vspace{3pt}
\textbf{Topic:} Quantitative Aptitude \\
\textbf{Subtopic:} \textit{Numerical computation and functions: exponents and equations}
\vspace{2pt}\hrule height 0.3pt\vspace{3pt}
\textbf{Guide / Hint:} \\
{\small For $P e^x = Q e^{-x}$ to hold for all real $x$, multiply both sides by $e^x$ to get $P e^{2x} = Q$. Differentiating with respect to $x$ gives $2P e^{2x} = 0$, implying $P = 0$. Substituting $P = 0$ yields $Q = 0$, so $P = Q = 0$.}

\end{tcolorbox}
\vspace{6pt}\hrule height 0.35pt\vspace{6pt}
\noindent
\begin{tcolorbox}[
    blanker,
    sidebyside,
    sidebyside align=top,
    lefthand width=0.55\textwidth,
    top=4pt, bottom=4pt, left=0pt, right=0pt,
    middle=16pt,
    underlay={
        \draw[dashed, black!80, line width=0.8pt, dash pattern=on 4pt off 3pt]
            ([xshift=0.55\textwidth+8pt]frame.north west) --
            ([xshift=0.55\textwidth+8pt]frame.south west);
    }
]
\textbf{Q.56} \quad {\footnotesize\textsc{[MCQ $\cdot$ 1 Mark(s) $\cdot$ Neg: 0.33]}} \hfill {\footnotesize\textbf{GATE 2025 CS2} \quad \tiny\texttt{2025\_CS2\_GA\_Q05}}
\vspace{2pt}\hrule height 0.3pt\vspace{4pt}

Let $p_1$ and $p_2$ denote two arbitrary prime numbers. Which one of the following statements is correct for all values of $p_1$ and $p_2$?

\par\vspace{6pt}\noindent
\begin{tabular}{@{}p{0.98\linewidth}@{}}
$\square$\ \textbf{(A)} $p_1 + p_2$ is not a prime number. \\
\addlinespace[3pt]
$\square$\ \textbf{(B)} $p_1 p_2$ is not a prime number. \\
\addlinespace[3pt]
$\square$\ \textbf{(C)} $p_1 + p_2 + 1$ is a prime number. \\
\addlinespace[3pt]
$\square$\ \textbf{(D)} $p_1 p_2 + 1$ is a prime number.
\end{tabular}

\tcblower

\textbf{Official Key:} \textbf{\large B} \hfill \textit{\small [Easy]}
\vspace{2pt}\hrule height 0.3pt\vspace{3pt}
\textbf{Topic:} Quantitative Aptitude \\
\textbf{Subtopic:} \textit{Elementary number theory and prime numbers}
\vspace{2pt}\hrule height 0.3pt\vspace{3pt}
\textbf{Guide / Hint:} \\
{\small For any two prime numbers $p_1$ and $p_2$, the product $p_1 p_2$ has at least three distinct positive divisors ($1, p_1, p_1 p_2$ when $p_1=p_2$, or $1, p_1, p_2, p_1 p_2$ when $p_1 \neq p_2$). Hence, $p_1 p_2$ can never be prime.}

\end{tcolorbox}
\vspace{6pt}\hrule height 0.35pt\vspace{6pt}
\noindent
\begin{tcolorbox}[
    blanker,
    sidebyside,
    sidebyside align=top,
    lefthand width=0.55\textwidth,
    top=4pt, bottom=4pt, left=0pt, right=0pt,
    middle=16pt,
    underlay={
        \draw[dashed, black!80, line width=0.8pt, dash pattern=on 4pt off 3pt]
            ([xshift=0.55\textwidth+8pt]frame.north west) --
            ([xshift=0.55\textwidth+8pt]frame.south west);
    }
]
\textbf{Q.57} \quad {\footnotesize\textsc{[MCQ $\cdot$ 2 Mark(s) $\cdot$ Neg: 0.66]}} \hfill {\footnotesize\textbf{GATE 2025 CS2} \quad \tiny\texttt{2025\_CS2\_GA\_Q06}}
\vspace{2pt}\hrule height 0.3pt\vspace{4pt}

Based only on the conversation below, identify the logically correct inference:
``Even if I had known that you were in the hospital, I would not have gone there to see you'', Ramya told Josephine.

\par\vspace{6pt}\noindent
\begin{tabular}{@{}p{0.98\linewidth}@{}}
$\square$\ \textbf{(A)} Ramya knew that Josephine was in the hospital. \\
\addlinespace[3pt]
$\square$\ \textbf{(B)} Ramya did not know that Josephine was in the hospital. \\
\addlinespace[3pt]
$\square$\ \textbf{(C)} Ramya and Josephine were once close friends; but now, they are not. \\
\addlinespace[3pt]
$\square$\ \textbf{(D)} Josephine was in the hospital due to an injury to her leg.
\end{tabular}

\tcblower

\textbf{Official Key:} \textbf{\large B} \hfill \textit{\small [Easy]}
\vspace{2pt}\hrule height 0.3pt\vspace{3pt}
\textbf{Topic:} Analytical Aptitude \\
\textbf{Subtopic:} \textit{Logic: deduction and conditional statements}
\vspace{2pt}\hrule height 0.3pt\vspace{3pt}
\textbf{Guide / Hint:} \\
{\small The subjunctive statement 'Even if I had known that you were in the hospital...' logically presupposes that the speaker did not know about Josephine's hospitalization at that time.}

\end{tcolorbox}
\vspace{6pt}\hrule height 0.35pt\vspace{6pt}
\noindent
\begin{tcolorbox}[
    blanker,
    sidebyside,
    sidebyside align=top,
    lefthand width=0.55\textwidth,
    top=4pt, bottom=4pt, left=0pt, right=0pt,
    middle=16pt,
    underlay={
        \draw[dashed, black!80, line width=0.8pt, dash pattern=on 4pt off 3pt]
            ([xshift=0.55\textwidth+8pt]frame.north west) --
            ([xshift=0.55\textwidth+8pt]frame.south west);
    }
]
\textbf{Q.58} \quad {\footnotesize\textsc{[MCQ $\cdot$ 2 Mark(s) $\cdot$ Neg: 0.66]}} \hfill {\footnotesize\textbf{GATE 2025 CS2} \quad \tiny\texttt{2025\_CS2\_GA\_Q07}}
\vspace{2pt}\hrule height 0.3pt\vspace{4pt}

If $\text{IMAGE}$ and $\text{FIELD}$ are coded as $\text{FHBNJ}$ and $\text{EMFJG}$ respectively then, which one among the given options is the most appropriate code for $\text{BEACH}$?

\par\vspace{6pt}\noindent
\begin{tabular}{@{}p{0.48\linewidth}@{\hspace{0.04\linewidth}}p{0.48\linewidth}@{}}
$\square$\ \textbf{(A)} CEADP & $\square$\ \textbf{(B)} IDBFC \\
\addlinespace[3pt]
$\square$\ \textbf{(C)} JGIBC & $\square$\ \textbf{(D)} IBCEC
\end{tabular}

\tcblower

\textbf{Official Key:} \textbf{\large B} \hfill \textit{\small [Medium]}
\vspace{2pt}\hrule height 0.3pt\vspace{3pt}
\textbf{Topic:} Analytical Aptitude \\
\textbf{Subtopic:} \textit{Letter patterns, sequencing and coding-decoding}
\vspace{2pt}\hrule height 0.3pt\vspace{3pt}
\textbf{Guide / Hint:} \\
{\small The coding rule applies a $+1$ alphabetical shift to each character followed by string reversal. Shifting $\text{BEACH}$ by $+1$ gives $\text{CFBDI}$, and reversing this yields $\text{IDBFC}$.}

\end{tcolorbox}
\vspace{6pt}\hrule height 0.35pt\vspace{6pt}
\noindent
\begin{tcolorbox}[
    blanker,
    sidebyside,
    sidebyside align=top,
    lefthand width=0.55\textwidth,
    top=4pt, bottom=4pt, left=0pt, right=0pt,
    middle=16pt,
    underlay={
        \draw[dashed, black!80, line width=0.8pt, dash pattern=on 4pt off 3pt]
            ([xshift=0.55\textwidth+8pt]frame.north west) --
            ([xshift=0.55\textwidth+8pt]frame.south west);
    }
]
\textbf{Q.59} \quad {\footnotesize\textsc{[MCQ $\cdot$ 2 Mark(s) $\cdot$ Neg: 0.66]}} \hfill {\footnotesize\textbf{GATE 2025 CS2} \quad \tiny\texttt{2025\_CS2\_GA\_Q08}}
\vspace{2pt}\hrule height 0.3pt\vspace{4pt}

Which one of the following options is correct for the given data in the table?
\par\medskip
\begin{center}\small
\begin{tabular}{cccc}
\toprule
Iteration ($i$) & Input ($I$) & Output ($X$) & Output ($Y$) \\
\midrule
0 & 20 & 20 & 20 \\
1 & $-4$ & 16 & $-80$ \\
2 & 10 & 26 & $-800$ \\
3 & 15 & 41 & $-12000$ \\
\bottomrule
\end{tabular}\end{center}

\par\vspace{6pt}\noindent
\begin{tabular}{@{}p{0.98\linewidth}@{}}
$\square$\ \textbf{(A)} $X(i) = X(i-1) + I(i); \; Y(i) = Y(i-1) I(i); \; i > 0$ \\
\addlinespace[3pt]
$\square$\ \textbf{(B)} $X(i) = X(i-1) I(i); \; Y(i) = Y(i-1) + I(i); \; i > 0$ \\
\addlinespace[3pt]
$\square$\ \textbf{(C)} $X(i) = X(i-1) I(i); \; Y(i) = Y(i-1) I(i); \; i > 0$ \\
\addlinespace[3pt]
$\square$\ \textbf{(D)} $X(i) = X(i-1) + I(i); \; Y(i) = Y(i-1) I(i-1); \; i > 0$
\end{tabular}

\tcblower

\textbf{Official Key:} \textbf{\large A} \hfill \textit{\small [Medium]}
\vspace{2pt}\hrule height 0.3pt\vspace{3pt}
\textbf{Topic:} Quantitative Aptitude \\
\textbf{Subtopic:} \textit{Numerical reasoning and recurrence relations}
\vspace{2pt}\hrule height 0.3pt\vspace{3pt}
\textbf{Guide / Hint:} \\
{\small Check the sequence updates: $X(i) = X(i-1) + I(i)$ yields $20 + (-4) = 16, 16 + 10 = 26, 26 + 15 = 41$, and $Y(i) = Y(i-1) \cdot I(i)$ yields $20 \times (-4) = -80, -80 \times 10 = -800, -800 \times 15 = -12000$.}

\end{tcolorbox}
\vspace{6pt}\hrule height 0.35pt\vspace{6pt}
\vspace{8pt}
\subsection*{Engineering Mathematics}
\vspace{-4pt}
\noindent
\begin{tcolorbox}[
    blanker,
    sidebyside,
    sidebyside align=top,
    lefthand width=0.55\textwidth,
    top=4pt, bottom=4pt, left=0pt, right=0pt,
    middle=16pt,
    underlay={
        \draw[dashed, black!80, line width=0.8pt, dash pattern=on 4pt off 3pt]
            ([xshift=0.55\textwidth+8pt]frame.north west) --
            ([xshift=0.55\textwidth+8pt]frame.south west);
    }
]
\textbf{Q.60} \quad {\footnotesize\textsc{[MCQ $\cdot$ 2 Mark(s) $\cdot$ Neg: 0.66]}} \hfill {\footnotesize\textbf{GATE 2018 CS} \quad \tiny\texttt{2018\_CS\_GA\_Q26}}
\vspace{2pt}\hrule height 0.3pt\vspace{4pt}

Consider a matrix $P = \begin{bmatrix} a & b \\ c & d \end{bmatrix}$. The only linearly independent eigenvector of $P$ is $\begin{bmatrix} 1 \\ 4 \end{bmatrix}$. Consider the following statements:
\begin{enumerate}[label=\Roman*.]
    \item $\det(P) = 0$
    \item $P$ has repeated eigenvalues
    \item $P$ is not diagonalizable
\end{enumerate}
Which of the statements above are necessarily TRUE?

\par\vspace{6pt}\noindent
\begin{tabular}{@{}p{0.48\linewidth}@{\hspace{0.04\linewidth}}p{0.48\linewidth}@{}}
$\square$\ \textbf{(A)} I and II only & $\square$\ \textbf{(B)} I and III only \\
\addlinespace[3pt]
$\square$\ \textbf{(C)} I, II and III & $\square$\ \textbf{(D)} II and III only
\end{tabular}

\tcblower

\textbf{Official Key:} \textbf{\large D} \hfill \textit{\small [Medium]}
\vspace{2pt}\hrule height 0.3pt\vspace{3pt}
\textbf{Topic:} Linear Algebra \\
\textbf{Subtopic:} \textit{Eigenvalues and eigenvectors}
\vspace{2pt}\hrule height 0.3pt\vspace{3pt}
\textbf{Guide / Hint:} \\
{\small A $2 \times 2$ matrix with only one linearly independent eigenvector has geometric multiplicity 1. This implies that the algebraic multiplicity of its eigenvalue must be 2 (repeated eigenvalues, so II is true) and that $P$ is not diagonalizable (so III is true). The determinant may be non-zero (e.g., $\begin{bmatrix} 2 & 1 \\ 0 & 2 \end{bmatrix}$), so I is not necessarily true.}

\end{tcolorbox}
\vspace{6pt}\hrule height 0.35pt\vspace{6pt}
\noindent
\begin{tcolorbox}[
    blanker,
    sidebyside,
    sidebyside align=top,
    lefthand width=0.55\textwidth,
    top=4pt, bottom=4pt, left=0pt, right=0pt,
    middle=16pt,
    underlay={
        \draw[dashed, black!80, line width=0.8pt, dash pattern=on 4pt off 3pt]
            ([xshift=0.55\textwidth+8pt]frame.north west) --
            ([xshift=0.55\textwidth+8pt]frame.south west);
    }
]
\textbf{Q.61} \quad {\footnotesize\textsc{[MCQ $\cdot$ 1 Mark(s) $\cdot$ Neg: 0.33]}} \hfill {\footnotesize\textbf{GATE 2022 CS} \quad \tiny\texttt{2022\_CS\_GA\_Q02}}
\vspace{2pt}\hrule height 0.3pt\vspace{4pt}

A function $y(x)$ is defined in the interval $[0, 1]$ on the $x$-axis as
\[ y(x) = \begin{cases} 2 & \text{if } 0 \le x < \frac{1}{3} \\ 3 & \text{if } \frac{1}{3} \le x < \frac{3}{4} \\ 1 & \text{if } \frac{3}{4} \le x \le 1 \end{cases} \]
Which one of the following is the area under the curve for the interval $[0, 1]$ on the $x$-axis?

\par\vspace{6pt}\noindent
\begin{tabular}{@{}p{0.48\linewidth}@{\hspace{0.04\linewidth}}p{0.48\linewidth}@{}}
$\square$\ \textbf{(A)} $5/6$ & $\square$\ \textbf{(B)} $6/5$ \\
\addlinespace[3pt]
$\square$\ \textbf{(C)} $13/6$ & $\square$\ \textbf{(D)} $6/13$
\end{tabular}

\tcblower

\textbf{Official Key:} \textbf{\large C} \hfill \textit{\small [Easy]}
\vspace{2pt}\hrule height 0.3pt\vspace{3pt}
\textbf{Topic:} Calculus \\
\textbf{Subtopic:} \textit{Integration}
\vspace{2pt}\hrule height 0.3pt\vspace{3pt}
\textbf{Guide / Hint:} \\
{\small The area under the piecewise step function is the sum of the rectangular areas: $\int_0^1 y(x) \, dx = 2\left(\frac{1}{3} - 0\right) + 3\left(\frac{3}{4} - \frac{1}{3}\right) + 1\left(1 - \frac{3}{4}\right) = \frac{2}{3} + 3\left(\frac{5}{12}\right) + \frac{1}{4} = \frac{2}{3} + \frac{5}{4} + \frac{1}{4} = \frac{2}{3} + \frac{3}{2} = \frac{13}{6}$.}

\end{tcolorbox}
\vspace{6pt}\hrule height 0.35pt\vspace{6pt}
\noindent
\begin{tcolorbox}[
    blanker,
    sidebyside,
    sidebyside align=top,
    lefthand width=0.55\textwidth,
    top=4pt, bottom=4pt, left=0pt, right=0pt,
    middle=16pt,
    underlay={
        \draw[dashed, black!80, line width=0.8pt, dash pattern=on 4pt off 3pt]
            ([xshift=0.55\textwidth+8pt]frame.north west) --
            ([xshift=0.55\textwidth+8pt]frame.south west);
    }
]
\textbf{Q.62} \quad {\footnotesize\textsc{[MCQ $\cdot$ 1 Mark(s) $\cdot$ Neg: 0.33]}} \hfill {\footnotesize\textbf{GATE 2022 CS} \quad \tiny\texttt{2022\_CS\_CS\_Q20}}
\vspace{2pt}\hrule height 0.3pt\vspace{4pt}

Consider the following two statements with respect to the matrices $A_{m \times n}, B_{n \times m}, C_{n \times n}$ and $D_{n \times n}$:
\begin{itemize}[itemsep=0pt,topsep=2pt]
\item Statement 1: $\text{tr}(AB) = \text{tr}(BA)$
\item Statement 2: $\text{tr}(CD) = \text{tr}(DC)$
\end{itemize}
where $\text{tr}()$ represents the trace of a matrix. Which one of the following holds?

\par\vspace{6pt}\noindent
\begin{tabular}{@{}p{0.98\linewidth}@{}}
$\square$\ \textbf{(A)} Statement 1 is correct and Statement 2 is wrong. \\
\addlinespace[3pt]
$\square$\ \textbf{(B)} Statement 1 is wrong and Statement 2 is correct. \\
\addlinespace[3pt]
$\square$\ \textbf{(C)} Both Statement 1 and Statement 2 are correct. \\
\addlinespace[3pt]
$\square$\ \textbf{(D)} Both Statement 1 and Statement 2 are wrong.
\end{tabular}

\tcblower

\textbf{Official Key:} \textbf{\large C} \hfill \textit{\small [Easy]}
\vspace{2pt}\hrule height 0.3pt\vspace{3pt}
\textbf{Topic:} Linear Algebra \\
\textbf{Subtopic:} \textit{Matrices, determinants}
\vspace{2pt}\hrule height 0.3pt\vspace{3pt}
\textbf{Guide / Hint:} \\
{\small The trace is cyclic: $\text{tr}(AB) = \text{tr}(BA)$ holds whenever $A$ is $m \times n$ and $B$ is $n \times m$. Similarly, for $n \times n$ square matrices $C$ and $D$, $\text{tr}(CD) = \text{tr}(DC)$. Thus, both Statement 1 and Statement 2 are correct.}

\end{tcolorbox}
\vspace{6pt}\hrule height 0.35pt\vspace{6pt}
\noindent
\begin{tcolorbox}[
    blanker,
    sidebyside,
    sidebyside align=top,
    lefthand width=0.55\textwidth,
    top=4pt, bottom=4pt, left=0pt, right=0pt,
    middle=16pt,
    underlay={
        \draw[dashed, black!80, line width=0.8pt, dash pattern=on 4pt off 3pt]
            ([xshift=0.55\textwidth+8pt]frame.north west) --
            ([xshift=0.55\textwidth+8pt]frame.south west);
    }
]
\textbf{Q.63} \quad {\footnotesize\textsc{[MCQ $\cdot$ 2 Mark(s) $\cdot$ Neg: 0.66]}} \hfill {\footnotesize\textbf{GATE 2022 CS} \quad \tiny\texttt{2022\_CS\_CS\_Q45}}
\vspace{2pt}\hrule height 0.3pt\vspace{4pt}

Consider solving the following system of simultaneous equations using LU decomposition:
\begin{align*}
x_1 + x_2 - x_3 &= 2 \\
2x_1 + 4x_2 - 3x_3 &= 7 \\
2x_1 + 5x_2 - 7x_3 &= 7
\end{align*}
where $L$ and $U$ are lower and upper triangular matrices. Which one of the following is the correct combination of values for $L_{32}$, $U_{33}$, and $x_1$?

\par\vspace{6pt}\noindent
\begin{tabular}{@{}p{0.98\linewidth}@{}}
$\square$\ \textbf{(A)} $L_{32} = 2, U_{33} = -1/2, x_1 = -1$ \\
\addlinespace[3pt]
$\square$\ \textbf{(B)} $L_{32} = 2, U_{33} = 2, x_1 = -1$ \\
\addlinespace[3pt]
$\square$\ \textbf{(C)} $L_{32} = -1/2, U_{33} = 2, x_1 = 0$ \\
\addlinespace[3pt]
$\square$\ \textbf{(D)} $L_{32} = -1/2, U_{33} = -1/2, x_1 = 0$
\end{tabular}

\tcblower

\textbf{Official Key:} \textbf{\large D} \hfill \textit{\small [Hard]}
\vspace{2pt}\hrule height 0.3pt\vspace{3pt}
\textbf{Topic:} Linear Algebra \\
\textbf{Subtopic:} \textit{LU decomposition}
\vspace{2pt}\hrule height 0.3pt\vspace{3pt}
\textbf{Guide / Hint:} \\
{\small Performing Doolittle/Crout LU factorization on matrix $A = \begin{bmatrix} 1 & 1 & -1 \\ 2 & 4 & -3 \\ 2 & 5 & -7 \end{bmatrix}$ and solving $L y = b$ and $U x = y$ reveals $L_{32} = -1/2$, $U_{33} = -1/2$, and backward substitution yields $x_1 = 0$.}

\end{tcolorbox}
\vspace{6pt}\hrule height 0.35pt\vspace{6pt}
\noindent
\begin{tcolorbox}[
    blanker,
    sidebyside,
    sidebyside align=top,
    lefthand width=0.55\textwidth,
    top=4pt, bottom=4pt, left=0pt, right=0pt,
    middle=16pt,
    underlay={
        \draw[dashed, black!80, line width=0.8pt, dash pattern=on 4pt off 3pt]
            ([xshift=0.55\textwidth+8pt]frame.north west) --
            ([xshift=0.55\textwidth+8pt]frame.south west);
    }
]
\textbf{Q.64} \quad {\footnotesize\textsc{[MCQ $\cdot$ 1 Mark(s) $\cdot$ Neg: 0.33]}} \hfill {\footnotesize\textbf{GATE 2023 CS} \quad \tiny\texttt{2023\_CS\_CS\_Q18}}
\vspace{2pt}\hrule height 0.3pt\vspace{4pt}

Let
\[ A = \begin{bmatrix} 1 & 2 & 3 & 4 \\ 4 & 1 & 2 & 3 \\ 3 & 4 & 1 & 2 \\ 2 & 3 & 4 & 1 \end{bmatrix} \quad \text{and} \quad B = \begin{bmatrix} 3 & 4 & 1 & 2 \\ 4 & 1 & 2 & 3 \\ 1 & 2 & 3 & 4 \\ 2 & 3 & 4 & 1 \end{bmatrix}. \]
Let $\det(A)$ and $\det(B)$ denote the determinants of the matrices $A$ and $B$, respectively.
Which one of the options given below is TRUE?

\par\vspace{6pt}\noindent
\begin{tabular}{@{}p{0.48\linewidth}@{\hspace{0.04\linewidth}}p{0.48\linewidth}@{}}
$\square$\ \textbf{(A)} $\det(A) = \det(B)$ & $\square$\ \textbf{(B)} $\det(B) = -\det(A)$ \\
\addlinespace[3pt]
$\square$\ \textbf{(C)} $\det(A) = 0$ & $\square$\ \textbf{(D)} $\det(AB) = \det(A) + \det(B)$
\end{tabular}

\tcblower

\textbf{Official Key:} \textbf{\large B} \hfill \textit{\small [Easy]}
\vspace{2pt}\hrule height 0.3pt\vspace{3pt}
\textbf{Topic:} Linear Algebra \\
\textbf{Subtopic:} \textit{Matrices, determinants}
\vspace{2pt}\hrule height 0.3pt\vspace{3pt}
\textbf{Guide / Hint:} \\
{\small Matrix $B$ is obtained from matrix $A$ by interchanging row 1 and row 3 ($R_1 \leftrightarrow R_3$), leaving all other rows unchanged. By the elementary row operation property of determinants, swapping any two rows of a square matrix negates its determinant: $\det(B) = -\det(A)$.}

\end{tcolorbox}
\vspace{6pt}\hrule height 0.35pt\vspace{6pt}
\noindent
\begin{tcolorbox}[
    blanker,
    sidebyside,
    sidebyside align=top,
    lefthand width=0.55\textwidth,
    top=4pt, bottom=4pt, left=0pt, right=0pt,
    middle=16pt,
    underlay={
        \draw[dashed, black!80, line width=0.8pt, dash pattern=on 4pt off 3pt]
            ([xshift=0.55\textwidth+8pt]frame.north west) --
            ([xshift=0.55\textwidth+8pt]frame.south west);
    }
]
\textbf{Q.65} \quad {\footnotesize\textsc{[MCQ $\cdot$ 1 Mark(s) $\cdot$ Neg: 0.33]}} \hfill {\footnotesize\textbf{GATE 2024 CS1} \quad \tiny\texttt{2024\_CS1\_CS\_Q11}}
\vspace{2pt}\hrule height 0.3pt\vspace{4pt}

Let $f: \mathbb{R} \to \mathbb{R}$ be a function such that $f(x) = \max\{x, x^3\}$, where $\mathbb{R}$ is the set of all real numbers. The set of all points where $f(x)$ is NOT differentiable is

\par\vspace{6pt}\noindent
\begin{tabular}{@{}p{0.48\linewidth}@{\hspace{0.04\linewidth}}p{0.48\linewidth}@{}}
$\square$\ \textbf{(A)} $\{-1, 1, 2\}$ & $\square$\ \textbf{(B)} $\{-2, -1, 1\}$ \\
\addlinespace[3pt]
$\square$\ \textbf{(C)} $\{0, 1\}$ & $\square$\ \textbf{(D)} $\{-1, 0, 1\}$
\end{tabular}

\tcblower

\textbf{Official Key:} \textbf{\large D} \hfill \textit{\small [Medium]}
\vspace{2pt}\hrule height 0.3pt\vspace{3pt}
\textbf{Topic:} Calculus \\
\textbf{Subtopic:} \textit{Limits, continuity and differentiability}
\vspace{2pt}\hrule height 0.3pt\vspace{3pt}
\textbf{Guide / Hint:} \\
{\small The functions $y = x$ and $y = x^3$ intersect at $x^3 - x = 0 \implies x \in \{-1, 0, 1\}$. The maximum function $f(x) = \max\{x, x^3\}$ switches between $x$ and $x^3$ at each of these three intersection points. Because the left and right derivatives at each transition point differ ($f'(-1^-) = 3 \neq 1 = f'(-1^+)$, $f'(0^-) = 1 \neq 0 = f'(0^+)$, $f'(1^-) = 1 \neq 3 = f'(1^+)$), $f(x)$ is not differentiable at $\{-1, 0, 1\}$.}

\end{tcolorbox}
\vspace{6pt}\hrule height 0.35pt\vspace{6pt}
\noindent
\begin{tcolorbox}[
    blanker,
    sidebyside,
    sidebyside align=top,
    lefthand width=0.55\textwidth,
    top=4pt, bottom=4pt, left=0pt, right=0pt,
    middle=16pt,
    underlay={
        \draw[dashed, black!80, line width=0.8pt, dash pattern=on 4pt off 3pt]
            ([xshift=0.55\textwidth+8pt]frame.north west) --
            ([xshift=0.55\textwidth+8pt]frame.south west);
    }
]
\textbf{Q.66} \quad {\footnotesize\textsc{[MCQ $\cdot$ 1 Mark(s) $\cdot$ Neg: 0.33]}} \hfill {\footnotesize\textbf{GATE 2024 CS1} \quad \tiny\texttt{2024\_CS1\_CS\_Q12}}
\vspace{2pt}\hrule height 0.3pt\vspace{4pt}

The product of all eigenvalues of the matrix $\begin{bmatrix} 1 & 2 & 3 \\ 4 & 5 & 6 \\ 7 & 8 & 9 \end{bmatrix}$ is

\par\vspace{6pt}\noindent
\begin{tabular}{@{}p{0.48\linewidth}@{\hspace{0.04\linewidth}}p{0.48\linewidth}@{}}
$\square$\ \textbf{(A)} $-1$ & $\square$\ \textbf{(B)} 0 \\
\addlinespace[3pt]
$\square$\ \textbf{(C)} 1 & $\square$\ \textbf{(D)} 2
\end{tabular}

\tcblower

\textbf{Official Key:} \textbf{\large B} \hfill \textit{\small [Easy]}
\vspace{2pt}\hrule height 0.3pt\vspace{3pt}
\textbf{Topic:} Linear Algebra \\
\textbf{Subtopic:} \textit{Eigenvalues and eigenvectors}
\vspace{2pt}\hrule height 0.3pt\vspace{3pt}
\textbf{Guide / Hint:} \\
{\small The product of all eigenvalues of a square matrix equals its determinant $\det(A)$. In $A = \begin{bmatrix} 1 & 2 & 3 \\ 4 & 5 & 6 \\ 7 & 8 & 9 \end{bmatrix}$, the rows are in arithmetic progression with $R_3 - R_2 = R_2 - R_1 = [3, 3, 3]$, making the rows linearly dependent and $\det(A) = 0$. Hence, the product of the eigenvalues is 0.}

\end{tcolorbox}
\vspace{6pt}\hrule height 0.35pt\vspace{6pt}
\noindent
\begin{tcolorbox}[
    blanker,
    sidebyside,
    sidebyside align=top,
    lefthand width=0.55\textwidth,
    top=4pt, bottom=4pt, left=0pt, right=0pt,
    middle=16pt,
    underlay={
        \draw[dashed, black!80, line width=0.8pt, dash pattern=on 4pt off 3pt]
            ([xshift=0.55\textwidth+8pt]frame.north west) --
            ([xshift=0.55\textwidth+8pt]frame.south west);
    }
]
\textbf{Q.67} \quad {\footnotesize\textsc{[MCQ $\cdot$ 1 Mark(s) $\cdot$ Neg: 0.33]}} \hfill {\footnotesize\textbf{GATE 2024 CS1} \quad \tiny\texttt{2024\_CS1\_CS\_Q14}}
\vspace{2pt}\hrule height 0.3pt\vspace{4pt}

Consider a permutation sampled uniformly at random from the set of all permutations of $\{1, 2, 3, \dots, n\}$ for some $n \ge 4$. Let $X$ be the event that 1 occurs before 2 in the permutation, and $Y$ the event that 3 occurs before 4. Which one of the following statements is TRUE?

\par\vspace{6pt}\noindent
\begin{tabular}{@{}p{0.98\linewidth}@{}}
$\square$\ \textbf{(A)} The events $X$ and $Y$ are mutually exclusive \\
\addlinespace[3pt]
$\square$\ \textbf{(B)} The events $X$ and $Y$ are independent \\
\addlinespace[3pt]
$\square$\ \textbf{(C)} Either event $X$ or $Y$ must occur \\
\addlinespace[3pt]
$\square$\ \textbf{(D)} Event $X$ is more likely than event $Y$
\end{tabular}

\tcblower

\textbf{Official Key:} \textbf{\large B} \hfill \textit{\small [Medium]}
\vspace{2pt}\hrule height 0.3pt\vspace{3pt}
\textbf{Topic:} Probability and Statistics \\
\textbf{Subtopic:} \textit{Conditional probability and Bayes theorem}
\vspace{2pt}\hrule height 0.3pt\vspace{3pt}
\textbf{Guide / Hint:} \\
{\small By symmetry in random permutations of $\{1, 2, \dots, n\}$, $P(X) = P(1 \text{ before } 2) = 1/2$, and $P(Y) = P(3 \text{ before } 4) = 1/2$. Because the relative order of the disjoint pair $\{1, 2\}$ has no bearing on the relative order of $\{3, 4\}$, the joint probability is $P(X \cap Y) = 1/4 = P(X)P(Y)$, proving that $X$ and $Y$ are independent.}

\end{tcolorbox}
\vspace{6pt}\hrule height 0.35pt\vspace{6pt}
\noindent
\begin{tcolorbox}[
    blanker,
    sidebyside,
    sidebyside align=top,
    lefthand width=0.55\textwidth,
    top=4pt, bottom=4pt, left=0pt, right=0pt,
    middle=16pt,
    underlay={
        \draw[dashed, black!80, line width=0.8pt, dash pattern=on 4pt off 3pt]
            ([xshift=0.55\textwidth+8pt]frame.north west) --
            ([xshift=0.55\textwidth+8pt]frame.south west);
    }
]
\textbf{Q.68} \quad {\footnotesize\textsc{[MCQ $\cdot$ 1 Mark(s) $\cdot$ Neg: 0.33]}} \hfill {\footnotesize\textbf{GATE 2024 CS2} \quad \tiny\texttt{2024\_CS2\_CS\_Q16}}
\vspace{2pt}\hrule height 0.3pt\vspace{4pt}

Let $f(x)$ be a continuous function from $\mathbb{R}$ to $\mathbb{R}$ such that
\[ f(x) = 1 - f(2 - x) \]
Which one of the following options is the CORRECT value of $\int_0^2 f(x) \, dx$?

\par\vspace{6pt}\noindent
\begin{tabular}{@{}p{0.48\linewidth}@{\hspace{0.04\linewidth}}p{0.48\linewidth}@{}}
$\square$\ \textbf{(A)} 0 & $\square$\ \textbf{(B)} 1 \\
\addlinespace[3pt]
$\square$\ \textbf{(C)} 2 & $\square$\ \textbf{(D)} -1
\end{tabular}

\tcblower

\textbf{Official Key:} \textbf{\large B} \hfill \textit{\small [Medium]}
\vspace{2pt}\hrule height 0.3pt\vspace{3pt}
\textbf{Topic:} Calculus \\
\textbf{Subtopic:} \textit{Integration}
\vspace{2pt}\hrule height 0.3pt\vspace{3pt}
\textbf{Guide / Hint:} \\
{\small Let $I = \int_0^2 f(x) \, dx$. Using the property $\int_0^a f(x) \, dx = \int_0^a f(a - x) \, dx$, $I = \int_0^2 f(2 - x) \, dx$. Adding both gives $2I = \int_0^2 [f(x) + f(2 - x)] \, dx$. Since $f(x) + f(2 - x) = 1$, we have $2I = \int_0^2 1 \, dx = 2 \implies I = 1$.}

\end{tcolorbox}
\vspace{6pt}\hrule height 0.35pt\vspace{6pt}
\noindent
\begin{tcolorbox}[
    blanker,
    sidebyside,
    sidebyside align=top,
    lefthand width=0.55\textwidth,
    top=4pt, bottom=4pt, left=0pt, right=0pt,
    middle=16pt,
    underlay={
        \draw[dashed, black!80, line width=0.8pt, dash pattern=on 4pt off 3pt]
            ([xshift=0.55\textwidth+8pt]frame.north west) --
            ([xshift=0.55\textwidth+8pt]frame.south west);
    }
]
\textbf{Q.69} \quad {\footnotesize\textsc{[MCQ $\cdot$ 1 Mark(s) $\cdot$ Neg: 0.33]}} \hfill {\footnotesize\textbf{GATE 2024 CS2} \quad \tiny\texttt{2024\_CS2\_CS\_Q18}}
\vspace{2pt}\hrule height 0.3pt\vspace{4pt}

When six unbiased dice are rolled simultaneously, the probability of getting all distinct numbers (i.e., 1, 2, 3, 4, 5, and 6) is:

\par\vspace{6pt}\noindent
\begin{tabular}{@{}p{0.48\linewidth}@{\hspace{0.04\linewidth}}p{0.48\linewidth}@{}}
$\square$\ \textbf{(A)} $\frac{1}{324}$ & $\square$\ \textbf{(B)} $\frac{5}{324}$ \\
\addlinespace[3pt]
$\square$\ \textbf{(C)} $\frac{7}{324}$ & $\square$\ \textbf{(D)} $\frac{11}{324}$
\end{tabular}

\tcblower

\textbf{Official Key:} \textbf{\large B} \hfill \textit{\small [Easy]}
\vspace{2pt}\hrule height 0.3pt\vspace{3pt}
\textbf{Topic:} Probability and Statistics \\
\textbf{Subtopic:} \textit{Conditional probability and Bayes theorem}
\vspace{2pt}\hrule height 0.3pt\vspace{3pt}
\textbf{Guide / Hint:} \\
{\small When six unbiased 6-faced dice are rolled, total outcomes are $6^6$. The number of outcomes where all 6 faces are distinct is $6! = 720$. The probability is $\frac{6!}{6^6} = \frac{720}{46656} = \frac{5}{324}$.}

\end{tcolorbox}
\vspace{6pt}\hrule height 0.35pt\vspace{6pt}
\noindent
\begin{tcolorbox}[
    blanker,
    sidebyside,
    sidebyside align=top,
    lefthand width=0.55\textwidth,
    top=4pt, bottom=4pt, left=0pt, right=0pt,
    middle=16pt,
    underlay={
        \draw[dashed, black!80, line width=0.8pt, dash pattern=on 4pt off 3pt]
            ([xshift=0.55\textwidth+8pt]frame.north west) --
            ([xshift=0.55\textwidth+8pt]frame.south west);
    }
]
\textbf{Q.70} \quad {\footnotesize\textsc{[MCQ $\cdot$ 2 Mark(s) $\cdot$ Neg: 0.66]}} \hfill {\footnotesize\textbf{GATE 2024 CS2} \quad \tiny\texttt{2024\_CS2\_CS\_Q44}}
\vspace{2pt}\hrule height 0.3pt\vspace{4pt}

Let $x$ and $y$ be random variables, not necessarily independent, that take real values in the interval $[0,1]$. Let $z = xy$ and let the mean values of $x, y, z$ be $\overline{x}, \overline{y}, \overline{z}$, respectively. Which one of the following statements is TRUE?

\par\vspace{6pt}\noindent
\begin{tabular}{@{}p{0.98\linewidth}@{}}
$\square$\ \textbf{(A)} $\overline{z} = \overline{x}\overline{y}$ \\
\addlinespace[3pt]
$\square$\ \textbf{(B)} $\overline{z} \le \overline{x}\overline{y}$ \\
\addlinespace[3pt]
$\square$\ \textbf{(C)} $\overline{z} \ge \overline{x}\overline{y}$ \\
\addlinespace[3pt]
$\square$\ \textbf{(D)} $\overline{z} \le \overline{x}$
\end{tabular}

\tcblower

\textbf{Official Key:} \textbf{\large D} \hfill \textit{\small [Medium]}
\vspace{2pt}\hrule height 0.3pt\vspace{3pt}
\textbf{Topic:} Probability and Statistics \\
\textbf{Subtopic:} \textit{Random variables}
\vspace{2pt}\hrule height 0.3pt\vspace{3pt}
\textbf{Guide / Hint:} \\
{\small Since $y \in [0, 1]$, $z = xy \le x(1) = x$ holds for every realization. Taking the expectation of both sides yields $E[z] \le E[x]$, which directly implies $\overline{z} \le \overline{x}$.}

\end{tcolorbox}
\vspace{6pt}\hrule height 0.35pt\vspace{6pt}
\noindent
\begin{tcolorbox}[
    blanker,
    sidebyside,
    sidebyside align=top,
    lefthand width=0.55\textwidth,
    top=4pt, bottom=4pt, left=0pt, right=0pt,
    middle=16pt,
    underlay={
        \draw[dashed, black!80, line width=0.8pt, dash pattern=on 4pt off 3pt]
            ([xshift=0.55\textwidth+8pt]frame.north west) --
            ([xshift=0.55\textwidth+8pt]frame.south west);
    }
]
\textbf{Q.71} \quad {\footnotesize\textsc{[MCQ $\cdot$ 2 Mark(s) $\cdot$ Neg: 0.66]}} \hfill {\footnotesize\textbf{GATE 2025 CS1} \quad \tiny\texttt{2025\_CS1\_CS\_Q41}}
\vspace{2pt}\hrule height 0.3pt\vspace{4pt}

Let $A$ be a $2 \times 2$ matrix given by:
\[ A = \begin{bmatrix} 1 & 1 \\ 1 & -1 \end{bmatrix} \]
What are the eigenvalues of the matrix $A^{13}$?

\par\vspace{6pt}\noindent
\begin{tabular}{@{}p{0.48\linewidth}@{\hspace{0.04\linewidth}}p{0.48\linewidth}@{}}
$\square$\ \textbf{(A)} $1, -1$ & $\square$\ \textbf{(B)} $2\sqrt{2}, -2\sqrt{2}$ \\
\addlinespace[3pt]
$\square$\ \textbf{(C)} $4\sqrt{2}, -4\sqrt{2}$ & $\square$\ \textbf{(D)} $64\sqrt{2}, -64\sqrt{2}$
\end{tabular}

\tcblower

\textbf{Official Key:} \textbf{\large D} \hfill \textit{\small [Medium]}
\vspace{2pt}\hrule height 0.3pt\vspace{3pt}
\textbf{Topic:} Linear Algebra \\
\textbf{Subtopic:} \textit{Eigenvalues and eigenvectors}
\vspace{2pt}\hrule height 0.3pt\vspace{3pt}
\textbf{Guide / Hint:} \\
{\small The characteristic equation of $A$ is $\det(A - \lambda I) = (1 - \lambda)(-1 - \lambda) - 1 = \lambda^2 - 2 = 0$, so eigenvalues of $A$ are $\lambda_1 = \sqrt{2}$ and $\lambda_2 = -\sqrt{2}$. The eigenvalues of $A^{13}$ are $(\sqrt{2})^{13} = 2^6 \sqrt{2} = 64\sqrt{2}$ and $(-\sqrt{2})^{13} = -64\sqrt{2}$.}

\end{tcolorbox}
\vspace{6pt}\hrule height 0.35pt\vspace{6pt}
\noindent
\begin{tcolorbox}[
    blanker,
    sidebyside,
    sidebyside align=top,
    lefthand width=0.55\textwidth,
    top=4pt, bottom=4pt, left=0pt, right=0pt,
    middle=16pt,
    underlay={
        \draw[dashed, black!80, line width=0.8pt, dash pattern=on 4pt off 3pt]
            ([xshift=0.55\textwidth+8pt]frame.north west) --
            ([xshift=0.55\textwidth+8pt]frame.south west);
    }
]
\textbf{Q.72} \quad {\footnotesize\textsc{[MCQ $\cdot$ 1 Mark(s) $\cdot$ Neg: 0.33]}} \hfill {\footnotesize\textbf{GATE 2025 CS2} \quad \tiny\texttt{2025\_CS2\_CS\_Q11}}
\vspace{2pt}\hrule height 0.3pt\vspace{4pt}

If $A = \begin{pmatrix} 1 & 2 \\ 2 & -1 \end{pmatrix}$, then which ONE of the following is $A^8$?

\par\vspace{6pt}\noindent
\begin{tabular}{@{}p{0.98\linewidth}@{}}
$\square$\ \textbf{(A)} $\begin{pmatrix} 25 & 0 \\ 0 & 25 \end{pmatrix}$ \\
\addlinespace[3pt]
$\square$\ \textbf{(B)} $\begin{pmatrix} 125 & 0 \\ 0 & 125 \end{pmatrix}$ \\
\addlinespace[3pt]
$\square$\ \textbf{(C)} $\begin{pmatrix} 625 & 0 \\ 0 & 625 \end{pmatrix}$ \\
\addlinespace[3pt]
$\square$\ \textbf{(D)} $\begin{pmatrix} 3125 & 0 \\ 0 & 3125 \end{pmatrix}$
\end{tabular}

\tcblower

\textbf{Official Key:} \textbf{\large C} \hfill \textit{\small [Medium]}
\vspace{2pt}\hrule height 0.3pt\vspace{3pt}
\textbf{Topic:} Linear Algebra \\
\textbf{Subtopic:} \textit{Matrices and matrix powers}
\vspace{2pt}\hrule height 0.3pt\vspace{3pt}
\textbf{Guide / Hint:} \\
{\small Compute $A^2 = \begin{pmatrix} 1 & 2 \\ 2 & -1 \end{pmatrix}^2 = \begin{pmatrix} 5 & 0 \\ 0 & 5 \end{pmatrix} = 5I$. Then $A^8 = (A^2)^4 = (5I)^4 = 5^4 I = 625I = \begin{pmatrix} 625 & 0 \\ 0 & 625 \end{pmatrix}$.}

\end{tcolorbox}
\vspace{6pt}\hrule height 0.35pt\vspace{6pt}
\noindent
\begin{tcolorbox}[
    blanker,
    sidebyside,
    sidebyside align=top,
    lefthand width=0.55\textwidth,
    top=4pt, bottom=4pt, left=0pt, right=0pt,
    middle=16pt,
    underlay={
        \draw[dashed, black!80, line width=0.8pt, dash pattern=on 4pt off 3pt]
            ([xshift=0.55\textwidth+8pt]frame.north west) --
            ([xshift=0.55\textwidth+8pt]frame.south west);
    }
]
\textbf{Q.73} \quad {\footnotesize\textsc{[MCQ $\cdot$ 1 Mark(s) $\cdot$ Neg: 0.33]}} \hfill {\footnotesize\textbf{GATE 2025 CS2} \quad \tiny\texttt{2025\_CS2\_CS\_Q12}}
\vspace{2pt}\hrule height 0.3pt\vspace{4pt}

The value of $x$ such that $x > 1$, satisfying the equation $\int_1^x t \ln t \, dt = \frac{1}{4}$ is

\par\vspace{6pt}\noindent
\begin{tabular}{@{}p{0.48\linewidth}@{\hspace{0.04\linewidth}}p{0.48\linewidth}@{}}
$\square$\ \textbf{(A)} $\sqrt{e}$ & $\square$\ \textbf{(B)} $e$ \\
\addlinespace[3pt]
$\square$\ \textbf{(C)} $e^2$ & $\square$\ \textbf{(D)} $e^{-1}$
\end{tabular}

\tcblower

\textbf{Official Key:} \textbf{\large A} \hfill \textit{\small [Medium]}
\vspace{2pt}\hrule height 0.3pt\vspace{3pt}
\textbf{Topic:} Calculus \\
\textbf{Subtopic:} \textit{Integration and definite integrals}
\vspace{2pt}\hrule height 0.3pt\vspace{3pt}
\textbf{Guide / Hint:} \\
{\small Using integration by parts, $\int_1^x t \ln t \, dt = \left[\frac{t^2}{2} \ln t - \frac{t^2}{4}\right]_1^x = \frac{x^2}{2} \ln x - \frac{x^2}{4} + \frac{1}{4}$. Equating this to $\frac{1}{4}$ gives $\frac{x^2}{4}(2\ln x - 1) = 0$. Since $x > 1$, we have $2\ln x = 1 \implies x = e^{1/2} = \sqrt{e}$.}

\end{tcolorbox}
\vspace{6pt}\hrule height 0.35pt\vspace{6pt}
\noindent
\begin{tcolorbox}[
    blanker,
    sidebyside,
    sidebyside align=top,
    lefthand width=0.55\textwidth,
    top=4pt, bottom=4pt, left=0pt, right=0pt,
    middle=16pt,
    underlay={
        \draw[dashed, black!80, line width=0.8pt, dash pattern=on 4pt off 3pt]
            ([xshift=0.55\textwidth+8pt]frame.north west) --
            ([xshift=0.55\textwidth+8pt]frame.south west);
    }
]
\textbf{Q.74} \quad {\footnotesize\textsc{[MCQ $\cdot$ 1 Mark(s) $\cdot$ Neg: 0.33]}} \hfill {\footnotesize\textbf{GATE 2025 CS2} \quad \tiny\texttt{2025\_CS2\_CS\_Q14}}
\vspace{2pt}\hrule height 0.3pt\vspace{4pt}

Let $L$, $M$, and $N$ be non-singular matrices of order 3 satisfying the equations $L^2 = L^{-1}$, $M = L^8$ and $N = L^2$. Which ONE of the following is the value of the determinant of $(M - N)$?

\par\vspace{6pt}\noindent
\begin{tabular}{@{}p{0.48\linewidth}@{\hspace{0.04\linewidth}}p{0.48\linewidth}@{}}
$\square$\ \textbf{(A)} 0 & $\square$\ \textbf{(B)} 1 \\
\addlinespace[3pt]
$\square$\ \textbf{(C)} 2 & $\square$\ \textbf{(D)} 3
\end{tabular}

\tcblower

\textbf{Official Key:} \textbf{\large A} \hfill \textit{\small [Medium]}
\vspace{2pt}\hrule height 0.3pt\vspace{3pt}
\textbf{Topic:} Linear Algebra \\
\textbf{Subtopic:} \textit{Matrices and determinants}
\vspace{2pt}\hrule height 0.3pt\vspace{3pt}
\textbf{Guide / Hint:} \\
{\small From $L^2 = L^{-1}$, multiplying by $L$ gives $L^3 = I$. Then $M = L^8 = (L^3)^2 L^2 = L^2$, and $N = L^2$. Hence $M - N = L^2 - L^2 = O$, whose determinant is $\det(O) = 0$.}

\end{tcolorbox}
\vspace{6pt}\hrule height 0.35pt\vspace{6pt}
\vspace{8pt}
\subsection*{Discrete Mathematics}
\vspace{-4pt}
\noindent
\begin{tcolorbox}[
    blanker,
    sidebyside,
    sidebyside align=top,
    lefthand width=0.55\textwidth,
    top=4pt, bottom=4pt, left=0pt, right=0pt,
    middle=16pt,
    underlay={
        \draw[dashed, black!80, line width=0.8pt, dash pattern=on 4pt off 3pt]
            ([xshift=0.55\textwidth+8pt]frame.north west) --
            ([xshift=0.55\textwidth+8pt]frame.south west);
    }
]
\textbf{Q.75} \quad {\footnotesize\textsc{[MCQ $\cdot$ 1 Mark(s) $\cdot$ Neg: 0.33]}} \hfill {\footnotesize\textbf{GATE 2018 CS} \quad \tiny\texttt{2018\_CS\_GA\_Q01\_v2}}
\vspace{2pt}\hrule height 0.3pt\vspace{4pt}

Which one of the following is a closed form expression for the generating function of the sequence $\{a_n\}$, where $a_n = 2n + 3$ for all $n = 0, 1, 2, \dots$?

\par\vspace{6pt}\noindent
\begin{tabular}{@{}p{0.48\linewidth}@{\hspace{0.04\linewidth}}p{0.48\linewidth}@{}}
$\square$\ \textbf{(A)} $\frac{3}{(1-x)^2}$ & $\square$\ \textbf{(B)} $\frac{3x}{(1-x)^2}$ \\
\addlinespace[3pt]
$\square$\ \textbf{(C)} $\frac{2-x}{(1-x)^2}$ & $\square$\ \textbf{(D)} $\frac{3-x}{(1-x)^2}$
\end{tabular}

\tcblower

\textbf{Official Key:} \textbf{\large D} \hfill \textit{\small [Easy]}
\vspace{2pt}\hrule height 0.3pt\vspace{3pt}
\textbf{Topic:} Discrete Mathematics \\
\textbf{Subtopic:} \textit{Combinatorics: counting, recurrence relations, generating functions}
\vspace{2pt}\hrule height 0.3pt\vspace{3pt}
\textbf{Guide / Hint:} \\
{\small The generating function is $G(x) = \sum_{n=0}^\infty (2n + 3)x^n = 2\sum_{n=0}^\infty n x^n + 3\sum_{n=0}^\infty x^n$. Using the standard power series $\sum_{n=0}^\infty x^n = \frac{1}{1-x}$ and $\sum_{n=0}^\infty n x^n = \frac{x}{(1-x)^2}$, we obtain $G(x) = \frac{2x}{(1-x)^2} + \frac{3(1-x)}{(1-x)^2} = \frac{3 - x}{(1-x)^2}$.}

\end{tcolorbox}
\vspace{6pt}\hrule height 0.35pt\vspace{6pt}
\noindent
\begin{tcolorbox}[
    blanker,
    sidebyside,
    sidebyside align=top,
    lefthand width=0.55\textwidth,
    top=4pt, bottom=4pt, left=0pt, right=0pt,
    middle=16pt,
    underlay={
        \draw[dashed, black!80, line width=0.8pt, dash pattern=on 4pt off 3pt]
            ([xshift=0.55\textwidth+8pt]frame.north west) --
            ([xshift=0.55\textwidth+8pt]frame.south west);
    }
]
\textbf{Q.76} \quad {\footnotesize\textsc{[MCQ $\cdot$ 2 Mark(s) $\cdot$ Neg: 0.66]}} \hfill {\footnotesize\textbf{GATE 2018 CS} \quad \tiny\texttt{2018\_CS\_GA\_Q27}}
\vspace{2pt}\hrule height 0.3pt\vspace{4pt}

Let $\mathbb{N}$ denote the set of all natural numbers and $\mathbb{Q}$ denote the set of all rational numbers. Consider the following sets:
\begin{itemize}[leftmargin=*, nosep]
    \item[$P$:] The set of all rational numbers between 0 and 1.
    \item[$Q$:] The set of all functions from $\{0, 1\}$ to $\mathbb{N}$.
    \item[$R$:] The set of all functions from $\mathbb{N}$ to $\{0, 1\}$.
    \item[$S$:] The set of all finite subsets of $\mathbb{N}$.
\end{itemize}
Which of the sets above are countable?

\par\vspace{6pt}\noindent
\begin{tabular}{@{}p{0.48\linewidth}@{\hspace{0.04\linewidth}}p{0.48\linewidth}@{}}
$\square$\ \textbf{(A)} $P$ and $Q$ only & $\square$\ \textbf{(B)} $Q$ and $R$ only \\
\addlinespace[3pt]
$\square$\ \textbf{(C)} $P$, $Q$ and $R$ only & $\square$\ \textbf{(D)} $P$, $Q$ and $S$ only
\end{tabular}

\tcblower

\textbf{Official Key:} \textbf{\large D} \hfill \textit{\small [Medium]}
\vspace{2pt}\hrule height 0.3pt\vspace{3pt}
\textbf{Topic:} Discrete Mathematics \\
\textbf{Subtopic:} \textit{Sets, relations, functions, partial orders and lattices}
\vspace{2pt}\hrule height 0.3pt\vspace{3pt}
\textbf{Guide / Hint:} \\
{\small Set $P$ (rationals) is countably infinite. Set $Q$ (functions $\{0, 1\} \to \mathbb{N}$) is equivalent to $\mathbb{N} \times \mathbb{N}$, which is countable. Set $S$ (finite subsets of $\mathbb{N}$) is countable. Set $R$ (functions $\mathbb{N} \to \{0, 1\}$) is equivalent to the power set $\mathcal{P}(\mathbb{N})$ and has cardinality $2^{\aleph_0}$, which is uncountable. Hence, $P$, $Q$, and $S$ are countable.}

\end{tcolorbox}
\vspace{6pt}\hrule height 0.35pt\vspace{6pt}
\noindent
\begin{tcolorbox}[
    blanker,
    sidebyside,
    sidebyside align=top,
    lefthand width=0.55\textwidth,
    top=4pt, bottom=4pt, left=0pt, right=0pt,
    middle=16pt,
    underlay={
        \draw[dashed, black!80, line width=0.8pt, dash pattern=on 4pt off 3pt]
            ([xshift=0.55\textwidth+8pt]frame.north west) --
            ([xshift=0.55\textwidth+8pt]frame.south west);
    }
]
\textbf{Q.77} \quad {\footnotesize\textsc{[MCQ $\cdot$ 2 Mark(s) $\cdot$ Neg: 0.66]}} \hfill {\footnotesize\textbf{GATE 2018 CS} \quad \tiny\texttt{2018\_CS\_GA\_Q28}}
\vspace{2pt}\hrule height 0.3pt\vspace{4pt}

Consider the first-order logic sentence
\[ \varphi \equiv \exists s \, \exists t \, \exists u \, \forall v \, \forall w \, \forall x \, \forall y \; \psi(s, t, u, v, w, x, y) \]
where $\psi(s, t, u, v, w, x, y)$ is a quantifier-free first-order logic formula using only predicate symbols and variables (no functions or constants). If $\varphi$ has a model with a universe of elements, then $\varphi$ is guaranteed to have a model of size at most

\par\vspace{6pt}\noindent
\begin{tabular}{@{}p{0.48\linewidth}@{\hspace{0.04\linewidth}}p{0.48\linewidth}@{}}
$\square$\ \textbf{(A)} 3 & $\square$\ \textbf{(B)} 4 \\
\addlinespace[3pt]
$\square$\ \textbf{(C)} 7 & $\square$\ \textbf{(D)} $\infty$
\end{tabular}

\tcblower

\textbf{Official Key:} \textbf{\large A} \hfill \textit{\small [Hard]}
\vspace{2pt}\hrule height 0.3pt\vspace{3pt}
\textbf{Topic:} Discrete Mathematics \\
\textbf{Subtopic:} \textit{Propositional and first order logic}
\vspace{2pt}\hrule height 0.3pt\vspace{3pt}
\textbf{Guide / Hint:} \\
{\small The sentence $\varphi$ has existential quantifiers for 3 variables ($s, t, u$) followed exclusively by universal quantifiers for the remaining variables. If a structure satisfies $\varphi$, the submodel formed by taking the domain as the elements corresponding to $s, t, u$ will itself satisfy the sentence because any elements evaluated under universal quantifiers are chosen from this 3-element domain.}

\end{tcolorbox}
\vspace{6pt}\hrule height 0.35pt\vspace{6pt}
\noindent
\begin{tcolorbox}[
    blanker,
    sidebyside,
    sidebyside align=top,
    lefthand width=0.55\textwidth,
    top=4pt, bottom=4pt, left=0pt, right=0pt,
    middle=16pt,
    underlay={
        \draw[dashed, black!80, line width=0.8pt, dash pattern=on 4pt off 3pt]
            ([xshift=0.55\textwidth+8pt]frame.north west) --
            ([xshift=0.55\textwidth+8pt]frame.south west);
    }
]
\textbf{Q.78} \quad {\footnotesize\textsc{[MCQ $\cdot$ 2 Mark(s) $\cdot$ Neg: 0.66]}} \hfill {\footnotesize\textbf{GATE 2022 CS} \quad \tiny\texttt{2022\_CS\_CS\_Q36}}
\vspace{2pt}\hrule height 0.3pt\vspace{4pt}

Which one of the following is the closed form for the generating function of the sequence $\{a_n\}_{n \ge 0}$ defined below?
\[ a\_n = \begin{cases} n+1, & n \text{ is odd} \\ 1, & \text{otherwise} \end{cases} \]

\par\vspace{6pt}\noindent
\begin{tabular}{@{}p{0.98\linewidth}@{}}
$\square$\ \textbf{(A)} $\frac{x(1+x^2)}{(1-x^2)^2} + \frac{1}{1-x}$ \\
\addlinespace[3pt]
$\square$\ \textbf{(B)} $\frac{x(3-x^2)}{(1-x^2)^2} + \frac{1}{1-x}$ \\
\addlinespace[3pt]
$\square$\ \textbf{(C)} $\frac{2x}{(1-x^2)^2} + \frac{1}{1-x}$ \\
\addlinespace[3pt]
$\square$\ \textbf{(D)} $\frac{x}{(1-x^2)^2} + \frac{1}{1-x}$
\end{tabular}

\tcblower

\textbf{Official Key:} \textbf{\large A} \hfill \textit{\small [Hard]}
\vspace{2pt}\hrule height 0.3pt\vspace{3pt}
\textbf{Topic:} Discrete Mathematics \\
\textbf{Subtopic:} \textit{Combinatorics: counting, recurrence relations, generating functions}
\vspace{2pt}\hrule height 0.3pt\vspace{3pt}
\textbf{Guide / Hint:} \\
{\small The sequence $a_n$ is $1, 2, 1, 4, 1, 6, \dots$. Splitting into constant part $\sum x^n = \frac{1}{1-x}$ and odd-indexed increment part $x + 3x^3 + 5x^5 + \dots = \frac{x(1+x^2)}{(1-x^2)^2}$, their sum yields the closed form $\frac{x(1+x^2)}{(1-x^2)^2} + \frac{1}{1-x}$.}

\end{tcolorbox}
\vspace{6pt}\hrule height 0.35pt\vspace{6pt}
\noindent
\begin{tcolorbox}[
    blanker,
    sidebyside,
    sidebyside align=top,
    lefthand width=0.55\textwidth,
    top=4pt, bottom=4pt, left=0pt, right=0pt,
    middle=16pt,
    underlay={
        \draw[dashed, black!80, line width=0.8pt, dash pattern=on 4pt off 3pt]
            ([xshift=0.55\textwidth+8pt]frame.north west) --
            ([xshift=0.55\textwidth+8pt]frame.south west);
    }
]
\textbf{Q.79} \quad {\footnotesize\textsc{[MCQ $\cdot$ 1 Mark(s) $\cdot$ Neg: 0.33]}} \hfill {\footnotesize\textbf{GATE 2023 CS} \quad \tiny\texttt{2023\_CS\_CS\_Q15}}
\vspace{2pt}\hrule height 0.3pt\vspace{4pt}

The Lucas sequence $L_n$ is defined by the recurrence relation:
\[ L\_n = L\_{n-1} + L\_{n-2}, \quad \text{for } n \ge 3, \]
with $L_1 = 1$ and $L_2 = 3$.
Which one of the options given is TRUE?

\par\vspace{6pt}\noindent
\begin{tabular}{@{}p{0.98\linewidth}@{}}
$\square$\ \textbf{(A)} $L_n = \left(\frac{1+\sqrt{5}}{2}\right)^n + \left(\frac{1-\sqrt{5}}{2}\right)^n$ \\
\addlinespace[3pt]
$\square$\ \textbf{(B)} $L_n = \left(\frac{1+\sqrt{5}}{2}\right)^n - \left(\frac{1-\sqrt{5}}{3}\right)^n$ \\
\addlinespace[3pt]
$\square$\ \textbf{(C)} $L_n = \left(\frac{1+\sqrt{5}}{2}\right)^n + \left(\frac{1-\sqrt{5}}{3}\right)^n$ \\
\addlinespace[3pt]
$\square$\ \textbf{(D)} $L_n = \left(\frac{1+\sqrt{5}}{2}\right)^n - \left(\frac{1-\sqrt{5}}{2}\right)^n$
\end{tabular}

\tcblower

\textbf{Official Key:} \textbf{\large A} \hfill \textit{\small [Medium]}
\vspace{2pt}\hrule height 0.3pt\vspace{3pt}
\textbf{Topic:} Discrete Mathematics \\
\textbf{Subtopic:} \textit{Combinatorics: counting, recurrence relations, generating functions}
\vspace{2pt}\hrule height 0.3pt\vspace{3pt}
\textbf{Guide / Hint:} \\
{\small The recurrence $L_n = L_{n-1} + L_{n-2}$ has characteristic equation $r^2 - r - 1 = 0$, giving roots $r_1, r_2 = \frac{1 \pm \sqrt{5}}{2}$. The general solution is $L_n = c_1 \left(\frac{1+\sqrt{5}}{2}\right)^n + c_2 \left(\frac{1-\sqrt{5}}{2}\right)^n$. Substituting initial conditions $L_1 = 1$ and $L_2 = 3$ yields $c_1 = 1$ and $c_2 = 1$.}

\end{tcolorbox}
\vspace{6pt}\hrule height 0.35pt\vspace{6pt}
\noindent
\begin{tcolorbox}[
    blanker,
    sidebyside,
    sidebyside align=top,
    lefthand width=0.55\textwidth,
    top=4pt, bottom=4pt, left=0pt, right=0pt,
    middle=16pt,
    underlay={
        \draw[dashed, black!80, line width=0.8pt, dash pattern=on 4pt off 3pt]
            ([xshift=0.55\textwidth+8pt]frame.north west) --
            ([xshift=0.55\textwidth+8pt]frame.south west);
    }
]
\textbf{Q.80} \quad {\footnotesize\textsc{[MCQ $\cdot$ 2 Mark(s) $\cdot$ Neg: 0.66]}} \hfill {\footnotesize\textbf{GATE 2023 CS} \quad \tiny\texttt{2023\_CS\_CS\_Q48}}
\vspace{2pt}\hrule height 0.3pt\vspace{4pt}

Let $U = \{1, 2, \dots, n\}$, where $n$ is a large positive integer greater than 1000. Let $k$ be a positive integer less than $n$. Let $A, B$ be subsets of $U$ with $|A| = |B| = k$ and $A \cap B = \emptyset$. We say that a permutation of $U$ separates $A$ from $B$ if one of the following is true:
\begin{itemize}[itemsep=0pt,topsep=2pt]
\item All members of $A$ appear in the permutation before any of the members of $B$.
\item All members of $B$ appear in the permutation before any of the members of $A$.
\end{itemize}
How many permutations of $U$ separate $A$ from $B$?

\par\vspace{6pt}\noindent
\begin{tabular}{@{}p{0.98\linewidth}@{}}
$\square$\ \textbf{(A)} $n!$ \\
\addlinespace[3pt]
$\square$\ \textbf{(B)} $\binom{n}{2k} (n - 2k)!$ \\
\addlinespace[3pt]
$\square$\ \textbf{(C)} $\binom{n}{2k} (n - 2k)! (k!)^2$ \\
\addlinespace[3pt]
$\square$\ \textbf{(D)} $2 \binom{n}{2k} (n - 2k)! (k!)^2$
\end{tabular}

\tcblower

\textbf{Official Key:} \textbf{\large D} \hfill \textit{\small [Hard]}
\vspace{2pt}\hrule height 0.3pt\vspace{3pt}
\textbf{Topic:} Discrete Mathematics \\
\textbf{Subtopic:} \textit{Combinatorics: counting, recurrence relations, generating functions}
\vspace{2pt}\hrule height 0.3pt\vspace{3pt}
\textbf{Guide / Hint:} \\
{\small There are 2 orderings for which block appears first ($A$ before $B$ or $B$ before $A$). There are $\binom{n}{2k}$ ways to choose the positions in the permutation for elements of $A \cup B$. The first $k$ chosen positions must contain elements of the first set ($k!$ permutations) and the next $k$ positions contain elements of the second set ($k!$ permutations). The remaining $n - 2k$ elements can be arranged in any of the $(n - 2k)!$ remaining positions. Total valid permutations $= 2 \binom{n}{2k} (n - 2k)! (k!)^2$.}

\end{tcolorbox}
\vspace{6pt}\hrule height 0.35pt\vspace{6pt}
\noindent
\begin{tcolorbox}[
    blanker,
    sidebyside,
    sidebyside align=top,
    lefthand width=0.55\textwidth,
    top=4pt, bottom=4pt, left=0pt, right=0pt,
    middle=16pt,
    underlay={
        \draw[dashed, black!80, line width=0.8pt, dash pattern=on 4pt off 3pt]
            ([xshift=0.55\textwidth+8pt]frame.north west) --
            ([xshift=0.55\textwidth+8pt]frame.south west);
    }
]
\textbf{Q.81} \quad {\footnotesize\textsc{[MCQ $\cdot$ 1 Mark(s) $\cdot$ Neg: 0.33]}} \hfill {\footnotesize\textbf{GATE 2024 CS2} \quad \tiny\texttt{2024\_CS2\_CS\_Q12}}
\vspace{2pt}\hrule height 0.3pt\vspace{4pt}

Let $p$ and $q$ be the following propositions:
$p$: Fail grade can be given.
$q$: Student scores more than 50\% marks.
Consider the statement: ``Fail grade cannot be given when student scores more than 50\% marks.''
Which one of the following is the CORRECT representation of the above statement in propositional logic?

\par\vspace{6pt}\noindent
\begin{tabular}{@{}p{0.48\linewidth}@{\hspace{0.04\linewidth}}p{0.48\linewidth}@{}}
$\square$\ \textbf{(A)} $q \to \neg p$ & $\square$\ \textbf{(B)} $q \to p$ \\
\addlinespace[3pt]
$\square$\ \textbf{(C)} $p \to q$ & $\square$\ \textbf{(D)} $\neg p \to q$
\end{tabular}

\tcblower

\textbf{Official Key:} \textbf{\large A} \hfill \textit{\small [Easy]}
\vspace{2pt}\hrule height 0.3pt\vspace{3pt}
\textbf{Topic:} Discrete Mathematics \\
\textbf{Subtopic:} \textit{Propositional and first order logic}
\vspace{2pt}\hrule height 0.3pt\vspace{3pt}
\textbf{Guide / Hint:} \\
{\small The statement is: 'Fail grade cannot be given ($\neg p$) when student scores more than 50\% marks ($q$)', which means 'If $q$, then $\neg p$'. In propositional logic, this is directly formalized as $q \to \neg p$.}

\end{tcolorbox}
\vspace{6pt}\hrule height 0.35pt\vspace{6pt}
\noindent
\begin{tcolorbox}[
    blanker,
    sidebyside,
    sidebyside align=top,
    lefthand width=0.55\textwidth,
    top=4pt, bottom=4pt, left=0pt, right=0pt,
    middle=16pt,
    underlay={
        \draw[dashed, black!80, line width=0.8pt, dash pattern=on 4pt off 3pt]
            ([xshift=0.55\textwidth+8pt]frame.north west) --
            ([xshift=0.55\textwidth+8pt]frame.south west);
    }
]
\textbf{Q.82} \quad {\footnotesize\textsc{[MCQ $\cdot$ 1 Mark(s) $\cdot$ Neg: 0.33]}} \hfill {\footnotesize\textbf{GATE 2024 CS2} \quad \tiny\texttt{2024\_CS2\_CS\_Q15}}
\vspace{2pt}\hrule height 0.3pt\vspace{4pt}

Let $T(n)$ be the recurrence relation defined as follows:
\[ T(0) = 1, \quad T(1) = 2, \quad \text{and} \quad T(n) = 5T(n-1) - 6T(n-2) \quad \text{for } n \ge 2 \]
Which one of the following statements is TRUE?

\par\vspace{6pt}\noindent
\begin{tabular}{@{}p{0.48\linewidth}@{\hspace{0.04\linewidth}}p{0.48\linewidth}@{}}
$\square$\ \textbf{(A)} $T(n) = \Theta(2^n)$ & $\square$\ \textbf{(B)} $T(n) = \Theta(n 2^n)$ \\
\addlinespace[3pt]
$\square$\ \textbf{(C)} $T(n) = \Theta(3^n)$ & $\square$\ \textbf{(D)} $T(n) = \Theta(n 3^n)$
\end{tabular}

\tcblower

\textbf{Official Key:} \textbf{\large A} \hfill \textit{\small [Easy]}
\vspace{2pt}\hrule height 0.3pt\vspace{3pt}
\textbf{Topic:} Discrete Mathematics \\
\textbf{Subtopic:} \textit{Combinatorics: counting, recurrence relations, generating functions}
\vspace{2pt}\hrule height 0.3pt\vspace{3pt}
\textbf{Guide / Hint:} \\
{\small The characteristic equation is $r^2 - 5r + 6 = 0 \implies (r - 2)(r - 3) = 0$, with roots $r = 2, 3$. The general solution is $T(n) = c_1 2^n + c_2 3^n$. Using initial conditions $T(0) = 1, T(1) = 2$: $c_1 + c_2 = 1$ and $2c_1 + 3c_2 = 2 \implies c_1 = 1, c_2 = 0$. Thus, $T(n) = 2^n = \Theta(2^n)$.}

\end{tcolorbox}
\vspace{6pt}\hrule height 0.35pt\vspace{6pt}
\noindent
\begin{tcolorbox}[
    blanker,
    sidebyside,
    sidebyside align=top,
    lefthand width=0.55\textwidth,
    top=4pt, bottom=4pt, left=0pt, right=0pt,
    middle=16pt,
    underlay={
        \draw[dashed, black!80, line width=0.8pt, dash pattern=on 4pt off 3pt]
            ([xshift=0.55\textwidth+8pt]frame.north west) --
            ([xshift=0.55\textwidth+8pt]frame.south west);
    }
]
\textbf{Q.83} \quad {\footnotesize\textsc{[MCQ $\cdot$ 1 Mark(s) $\cdot$ Neg: 0.33]}} \hfill {\footnotesize\textbf{GATE 2024 CS2} \quad \tiny\texttt{2024\_CS2\_CS\_Q17}}
\vspace{2pt}\hrule height 0.3pt\vspace{4pt}

Let $A$ be the adjacency matrix of a simple undirected graph $G$. Suppose $A$ is its own inverse. Which one of the following statements is always TRUE?

\par\vspace{6pt}\noindent
\begin{tabular}{@{}p{0.48\linewidth}@{\hspace{0.04\linewidth}}p{0.48\linewidth}@{}}
$\square$\ \textbf{(A)} $G$ is a cycle & $\square$\ \textbf{(B)} $G$ is a perfect matching \\
\addlinespace[3pt]
$\square$\ \textbf{(C)} $G$ is a complete graph & $\square$\ \textbf{(D)} There is no such graph $G$
\end{tabular}

\tcblower

\textbf{Official Key:} \textbf{\large B} \hfill \textit{\small [Medium]}
\vspace{2pt}\hrule height 0.3pt\vspace{3pt}
\textbf{Topic:} Discrete Mathematics \\
\textbf{Subtopic:} \textit{Graphs: connectivity, matching, colouring}
\vspace{2pt}\hrule height 0.3pt\vspace{3pt}
\textbf{Guide / Hint:} \\
{\small If $A = A^{-1}$, then $A^2 = I$. The $(i, j)$-th entry of $A^2$ is the number of walks of length 2 between vertex $i$ and vertex $j$. Thus $(A^2)_{ii} = \text{deg}(i) = 1$ for every vertex $i$, and $(A^2)_{ij} = 0$ for $i \neq j$. Every vertex has degree exactly 1, which defines a 1-regular graph, i.e., a perfect matching.}

\end{tcolorbox}
\vspace{6pt}\hrule height 0.35pt\vspace{6pt}
\noindent
\begin{tcolorbox}[
    blanker,
    sidebyside,
    sidebyside align=top,
    lefthand width=0.55\textwidth,
    top=4pt, bottom=4pt, left=0pt, right=0pt,
    middle=16pt,
    underlay={
        \draw[dashed, black!80, line width=0.8pt, dash pattern=on 4pt off 3pt]
            ([xshift=0.55\textwidth+8pt]frame.north west) --
            ([xshift=0.55\textwidth+8pt]frame.south west);
    }
]
\textbf{Q.84} \quad {\footnotesize\textsc{[MCQ $\cdot$ 1 Mark(s) $\cdot$ Neg: 0.33]}} \hfill {\footnotesize\textbf{GATE 2025 CS1} \quad \tiny\texttt{2025\_CS1\_CS\_Q17}}
\vspace{2pt}\hrule height 0.3pt\vspace{4pt}

$g(\cdot)$ is a function from $A$ to $B$, $f(\cdot)$ is a function from $B$ to $C$, and their composition defined as $f(g(\cdot))$ is a mapping from $A$ to $C$.
If $f(\cdot)$ and $f(g(\cdot))$ are onto (surjective) functions, which ONE of the following is TRUE about the function $g(\cdot)$?

\par\vspace{6pt}\noindent
\begin{tabular}{@{}p{0.98\linewidth}@{}}
$\square$\ \textbf{(A)} $g(\cdot)$ must be an onto (surjective) function. \\
\addlinespace[3pt]
$\square$\ \textbf{(B)} $g(\cdot)$ must be a one-to-one (injective) function. \\
\addlinespace[3pt]
$\square$\ \textbf{(C)} $g(\cdot)$ must be a bijective function, that is, both one-to-one and onto. \\
\addlinespace[3pt]
$\square$\ \textbf{(D)} $g(\cdot)$ is not required to be a one-to-one or onto function.
\end{tabular}

\tcblower

\textbf{Official Key:} \textbf{\large D} \hfill \textit{\small [Medium]}
\vspace{2pt}\hrule height 0.3pt\vspace{3pt}
\textbf{Topic:} Discrete Mathematics \\
\textbf{Subtopic:} \textit{Sets, relations, functions, partial orders and lattices}
\vspace{2pt}\hrule height 0.3pt\vspace{3pt}
\textbf{Guide / Hint:} \\
{\small If the composition $f \circ g: A \to C$ is surjective, $f$ must be surjective, but $g$ is not required to be either injective or surjective. For example, if $C$ is a singleton $\{c\}$ and $B = \{b_1, b_2\}$, $g$ can map all of $A$ to $b_1$ where $f(b_1) = c$, leaving $b_2$ uncovered (not onto) and having many collisions (not one-to-one).}

\end{tcolorbox}
\vspace{6pt}\hrule height 0.35pt\vspace{6pt}
\noindent
\begin{tcolorbox}[
    blanker,
    sidebyside,
    sidebyside align=top,
    lefthand width=0.55\textwidth,
    top=4pt, bottom=4pt, left=0pt, right=0pt,
    middle=16pt,
    underlay={
        \draw[dashed, black!80, line width=0.8pt, dash pattern=on 4pt off 3pt]
            ([xshift=0.55\textwidth+8pt]frame.north west) --
            ([xshift=0.55\textwidth+8pt]frame.south west);
    }
]
\textbf{Q.85} \quad {\footnotesize\textsc{[MCQ $\cdot$ 1 Mark(s) $\cdot$ Neg: 0.33]}} \hfill {\footnotesize\textbf{GATE 2025 CS2} \quad \tiny\texttt{2025\_CS2\_CS\_Q15}}
\vspace{2pt}\hrule height 0.3pt\vspace{4pt}

Let $P(x)$ be an arbitrary predicate over the domain of natural numbers. Which ONE of the following statements is TRUE?

\par\vspace{6pt}\noindent
\begin{tabular}{@{}p{0.98\linewidth}@{}}
$\square$\ \textbf{(A)} $(P(0) \land (\forall x [P(x) \implies P(x+1)])) \implies (\forall x P(x))$ \\
\addlinespace[3pt]
$\square$\ \textbf{(B)} $(P(0) \land (\forall x [P(x) \implies P(x-1)])) \implies (\forall x P(x))$ \\
\addlinespace[3pt]
$\square$\ \textbf{(C)} $(P(1000) \land (\forall x [P(x) \implies P(x-1)])) \implies (\forall x P(x))$ \\
\addlinespace[3pt]
$\square$\ \textbf{(D)} $(P(1000) \land (\forall x [P(x) \implies P(x+1)])) \implies (\forall x P(x))$
\end{tabular}

\tcblower

\textbf{Official Key:} \textbf{\large A} \hfill \textit{\small [Easy]}
\vspace{2pt}\hrule height 0.3pt\vspace{3pt}
\textbf{Topic:} Discrete Mathematics \\
\textbf{Subtopic:} \textit{Propositional and first order logic: mathematical induction}
\vspace{2pt}\hrule height 0.3pt\vspace{3pt}
\textbf{Guide / Hint:} \\
{\small Option A is the formal statement of the Principle of Mathematical Induction over $\mathbb{N}$: if the base case $P(0)$ holds and for all $x$, $P(x) \implies P(x+1)$, then $\forall x, P(x)$ is true.}

\end{tcolorbox}
\vspace{6pt}\hrule height 0.35pt\vspace{6pt}
\vspace{8pt}
\subsection*{Digital Logic}
\vspace{-4pt}
\noindent
\begin{tcolorbox}[
    blanker,
    sidebyside,
    sidebyside align=top,
    lefthand width=0.55\textwidth,
    top=4pt, bottom=4pt, left=0pt, right=0pt,
    middle=16pt,
    underlay={
        \draw[dashed, black!80, line width=0.8pt, dash pattern=on 4pt off 3pt]
            ([xshift=0.55\textwidth+8pt]frame.north west) --
            ([xshift=0.55\textwidth+8pt]frame.south west);
    }
]
\textbf{Q.86} \quad {\footnotesize\textsc{[MCQ $\cdot$ 1 Mark(s) $\cdot$ Neg: 0.33]}} \hfill {\footnotesize\textbf{GATE 2018 CS} \quad \tiny\texttt{2018\_CS\_GA\_Q04\_v2}}
\vspace{2pt}\hrule height 0.3pt\vspace{4pt}

Let $\oplus$ and $\odot$ denote the Exclusive OR and Exclusive NOR operations, respectively. Which one of the following is NOT CORRECT?

\par\vspace{6pt}\noindent
\begin{tabular}{@{}p{0.98\linewidth}@{}}
$\square$\ \textbf{(A)} $\overline{P \oplus Q} = P \odot Q$ \\
\addlinespace[3pt]
$\square$\ \textbf{(B)} $\overline{P \odot Q} = P \oplus Q$ \\
\addlinespace[3pt]
$\square$\ \textbf{(C)} $\overline{P} \oplus Q = P \odot Q$ \\
\addlinespace[3pt]
$\square$\ \textbf{(D)} $(P \oplus \overline{P}) \oplus Q = (P \odot \overline{P}) \odot \overline{Q}$
\end{tabular}

\tcblower

\textbf{Official Key:} \textbf{\large D} \hfill \textit{\small [Easy]}
\vspace{2pt}\hrule height 0.3pt\vspace{3pt}
\textbf{Topic:} Digital Logic \\
\textbf{Subtopic:} \textit{Boolean algebra}
\vspace{2pt}\hrule height 0.3pt\vspace{3pt}
\textbf{Guide / Hint:} \\
{\small Recall that $P \oplus \overline{P} = 1$ and $1 \oplus Q = \overline{Q}$. On the right-hand side of (D), $P \odot \overline{P} = 0$, and $0 \odot \overline{Q} = \overline{0 \oplus \overline{Q}} = \overline{\overline{Q}} = Q$. Since $\overline{Q} \neq Q$, option (D) is NOT correct.}

\end{tcolorbox}
\vspace{6pt}\hrule height 0.35pt\vspace{6pt}
\noindent
\begin{tcolorbox}[
    blanker,
    sidebyside,
    sidebyside align=top,
    lefthand width=0.55\textwidth,
    top=4pt, bottom=4pt, left=0pt, right=0pt,
    middle=16pt,
    underlay={
        \draw[dashed, black!80, line width=0.8pt, dash pattern=on 4pt off 3pt]
            ([xshift=0.55\textwidth+8pt]frame.north west) --
            ([xshift=0.55\textwidth+8pt]frame.south west);
    }
]
\textbf{Q.87} \quad {\footnotesize\textsc{[MCQ $\cdot$ 2 Mark(s) $\cdot$ Neg: 0.66]}} \hfill {\footnotesize\textbf{GATE 2018 CS} \quad \tiny\texttt{2018\_CS\_GA\_Q33}}
\vspace{2pt}\hrule height 0.3pt\vspace{4pt}

Consider the unsigned 8-bit fixed point binary number representation below:
\[ b\_7 b\_6 b\_5 b\_4 b\_3 \, . \, b\_2 b\_1 b\_0 \]
where the position of the binary point is between $b_3$ and $b_2$. Which of the following decimal numbers can be represented exactly in this format?
\begin{enumerate}[label=(\roman*)]
    \item 0.500
    \item 0.875
    \item 0.100
    \item 0.001
\end{enumerate}

\par\vspace{6pt}\noindent
\begin{tabular}{@{}p{0.48\linewidth}@{\hspace{0.04\linewidth}}p{0.48\linewidth}@{}}
$\square$\ \textbf{(A)} (i) and (ii) only & $\square$\ \textbf{(B)} (i) and (iii) only \\
\addlinespace[3pt]
$\square$\ \textbf{(C)} (i) and (ii) only & $\square$\ \textbf{(D)} (i), (ii) and (iii)
\end{tabular}

\tcblower

\textbf{Official Key:} \textbf{\large C} \hfill \textit{\small [Medium]}
\vspace{2pt}\hrule height 0.3pt\vspace{3pt}
\textbf{Topic:} Digital Logic \\
\textbf{Subtopic:} \textit{Number representations and computer arithmetic (fixed and floating point)}
\vspace{2pt}\hrule height 0.3pt\vspace{3pt}
\textbf{Guide / Hint:} \\
{\small The fractional portion uses 3 bits ($b_2 b_1 b_0$), so the fractional precision is $2^{-3} = 0.125$. Any representable fraction must be an integer multiple of $0.125$: $0.500 = 4 \times 0.125$ (i) and $0.875 = 7 \times 0.125$ (ii) are exact multiples, while $0.100$ and $0.001$ are not.}

\end{tcolorbox}
\vspace{6pt}\hrule height 0.35pt\vspace{6pt}
\noindent
\begin{tcolorbox}[
    blanker,
    sidebyside,
    sidebyside align=top,
    lefthand width=0.55\textwidth,
    top=4pt, bottom=4pt, left=0pt, right=0pt,
    middle=16pt,
    underlay={
        \draw[dashed, black!80, line width=0.8pt, dash pattern=on 4pt off 3pt]
            ([xshift=0.55\textwidth+8pt]frame.north west) --
            ([xshift=0.55\textwidth+8pt]frame.south west);
    }
]
\textbf{Q.88} \quad {\footnotesize\textsc{[MCQ $\cdot$ 1 Mark(s) $\cdot$ Neg: 0.33]}} \hfill {\footnotesize\textbf{GATE 2022 CS} \quad \tiny\texttt{2022\_CS\_CS\_Q18}}
\vspace{2pt}\hrule height 0.3pt\vspace{4pt}

Let $R_1$ and $R_2$ be two 4-bit registers that store numbers in 2's complement form. For the operation $R_1 + R_2$, which one of the following values of $R_1$ and $R_2$ gives an arithmetic overflow?

\par\vspace{6pt}\noindent
\begin{tabular}{@{}p{0.48\linewidth}@{\hspace{0.04\linewidth}}p{0.48\linewidth}@{}}
$\square$\ \textbf{(A)} $R_1 = 1011$ and $R_2 = 1110$ & $\square$\ \textbf{(B)} $R_1 = 1100$ and $R_2 = 1010$ \\
\addlinespace[3pt]
$\square$\ \textbf{(C)} $R_1 = 0011$ and $R_2 = 0100$ & $\square$\ \textbf{(D)} $R_1 = 1001$ and $R_2 = 1111$
\end{tabular}

\tcblower

\textbf{Official Key:} \textbf{\large B} \hfill \textit{\small [Medium]}
\vspace{2pt}\hrule height 0.3pt\vspace{3pt}
\textbf{Topic:} Digital Logic \\
\textbf{Subtopic:} \textit{Number representations and computer arithmetic (fixed and floating point)}
\vspace{2pt}\hrule height 0.3pt\vspace{3pt}
\textbf{Guide / Hint:} \\
{\small In a 4-bit 2's complement system, representable integers span $[-2^3, 2^3 - 1] = [-8, +7]$. For $R_1 = 1100_2 = -4$ and $R_2 = 1010_2 = -6$, their sum is $-4 + (-6) = -10$. Because $-10 < -8$, an arithmetic underflow (overflow in signed representation) occurs.}

\end{tcolorbox}
\vspace{6pt}\hrule height 0.35pt\vspace{6pt}
\noindent
\begin{tcolorbox}[
    blanker,
    sidebyside,
    sidebyside align=top,
    lefthand width=0.55\textwidth,
    top=4pt, bottom=4pt, left=0pt, right=0pt,
    middle=16pt,
    underlay={
        \draw[dashed, black!80, line width=0.8pt, dash pattern=on 4pt off 3pt]
            ([xshift=0.55\textwidth+8pt]frame.north west) --
            ([xshift=0.55\textwidth+8pt]frame.south west);
    }
]
\textbf{Q.89} \quad {\footnotesize\textsc{[MCQ $\cdot$ 2 Mark(s) $\cdot$ Neg: 0.66]}} \hfill {\footnotesize\textbf{GATE 2022 CS} \quad \tiny\texttt{2022\_CS\_CS\_Q41}}
\vspace{2pt}\hrule height 0.3pt\vspace{4pt}

Consider three floating point numbers A, B and C stored in registers RA, RB and RC, respectively as per IEEE-754 single precision floating point format. The 32-bit content stored in these registers (in hexadecimal form) are as follows:
\[ RA = \text{0xC1400000}, \quad RB = \text{0x42100000}, \quad RC = \text{0x41400000} \]
Which one of the following is FALSE?

\par\vspace{6pt}\noindent
\begin{tabular}{@{}p{0.48\linewidth}@{\hspace{0.04\linewidth}}p{0.48\linewidth}@{}}
$\square$\ \textbf{(A)} $A + C = 0$ & $\square$\ \textbf{(B)} $C = A + B$ \\
\addlinespace[3pt]
$\square$\ \textbf{(C)} $B = 3C$ & $\square$\ \textbf{(D)} $(B - C) > 0$
\end{tabular}

\tcblower

\textbf{Official Key:} \textbf{\large B} \hfill \textit{\small [Medium]}
\vspace{2pt}\hrule height 0.3pt\vspace{3pt}
\textbf{Topic:} Digital Logic \\
\textbf{Subtopic:} \textit{Number representations and computer arithmetic (fixed and floating point)}
\vspace{2pt}\hrule height 0.3pt\vspace{3pt}
\textbf{Guide / Hint:} \\
{\small Decoding 32-bit IEEE-754: $RA = \text{0xC1400000} = -1.5 \times 2^{130-127} = -12.0$; $RB = \text{0x42100000} = +1.125 \times 2^{132-127} = +36.0$; $RC = \text{0x41400000} = +1.5 \times 2^{130-127} = +12.0$. Evaluating option B: $A + B = -12 + 36 = 24 \neq C = 12.0$, so statement B is FALSE.}

\end{tcolorbox}
\vspace{6pt}\hrule height 0.35pt\vspace{6pt}
\noindent
\begin{tcolorbox}[
    blanker,
    sidebyside,
    sidebyside align=top,
    lefthand width=0.55\textwidth,
    top=4pt, bottom=4pt, left=0pt, right=0pt,
    middle=16pt,
    underlay={
        \draw[dashed, black!80, line width=0.8pt, dash pattern=on 4pt off 3pt]
            ([xshift=0.55\textwidth+8pt]frame.north west) --
            ([xshift=0.55\textwidth+8pt]frame.south west);
    }
]
\textbf{Q.90} \quad {\footnotesize\textsc{[MCQ $\cdot$ 2 Mark(s) $\cdot$ Neg: 0.66]}} \hfill {\footnotesize\textbf{GATE 2023 CS} \quad \tiny\texttt{2023\_CS\_CS\_Q45}}
\vspace{2pt}\hrule height 0.3pt\vspace{4pt}

Consider the IEEE-754 single precision floating point numbers $P = \text{0xC1800000}$ and $Q = \text{0x3F5C2EF4}$.
Which one of the following corresponds to the product of these numbers (i.e., $P \times Q$), represented in the IEEE-754 single precision format?

\par\vspace{6pt}\noindent
\begin{tabular}{@{}p{0.48\linewidth}@{\hspace{0.04\linewidth}}p{0.48\linewidth}@{}}
$\square$\ \textbf{(A)} 0x404C2EF4 & $\square$\ \textbf{(B)} 0x405C2EF4 \\
\addlinespace[3pt]
$\square$\ \textbf{(C)} 0xC15C2EF4 & $\square$\ \textbf{(D)} 0xC14C2EF4
\end{tabular}

\tcblower

\textbf{Official Key:} \textbf{\large C} \hfill \textit{\small [Medium]}
\vspace{2pt}\hrule height 0.3pt\vspace{3pt}
\textbf{Topic:} Digital Logic \\
\textbf{Subtopic:} \textit{Number representations and computer arithmetic (fixed and floating point)}
\vspace{2pt}\hrule height 0.3pt\vspace{3pt}
\textbf{Guide / Hint:} \\
{\small Number $P = \text{0xC1800000}$ has sign bit 1, biased exponent $131$, actual exponent $131 - 127 = 4$, and mantissa $1.0$, so $P = -16.0 = -2^4$. Number $Q = \text{0x3F5C2EF4}$ is positive with biased exponent $126$ (exponent $-1$) and mantissa fraction `5C2EF4`. Multiplying $P \times Q$ flips the sign to negative (1), adds 4 to the exponent (biased exponent becomes $126 + 4 = 130 = \text{0x82}$), and retains the mantissa fraction, resulting in `0xC15C2EF4`.}

\end{tcolorbox}
\vspace{6pt}\hrule height 0.35pt\vspace{6pt}
\noindent
\begin{tcolorbox}[
    blanker,
    sidebyside,
    sidebyside align=top,
    lefthand width=0.55\textwidth,
    top=4pt, bottom=4pt, left=0pt, right=0pt,
    middle=16pt,
    underlay={
        \draw[dashed, black!80, line width=0.8pt, dash pattern=on 4pt off 3pt]
            ([xshift=0.55\textwidth+8pt]frame.north west) --
            ([xshift=0.55\textwidth+8pt]frame.south west);
    }
]
\textbf{Q.91} \quad {\footnotesize\textsc{[MCQ $\cdot$ 1 Mark(s) $\cdot$ Neg: 0.33]}} \hfill {\footnotesize\textbf{GATE 2024 CS1} \quad \tiny\texttt{2024\_CS1\_CS\_Q13}}
\vspace{2pt}\hrule height 0.3pt\vspace{4pt}

Consider a system that uses 5 bits for representing signed integers in 2's complement format. In this system, two integers $A$ and $B$ are represented as $A = 01010$ and $B = 11010$. Which one of the following operations will result in either an arithmetic overflow or an arithmetic underflow?

\par\vspace{6pt}\noindent
\begin{tabular}{@{}p{0.48\linewidth}@{\hspace{0.04\linewidth}}p{0.48\linewidth}@{}}
$\square$\ \textbf{(A)} $A + B$ & $\square$\ \textbf{(B)} $A - B$ \\
\addlinespace[3pt]
$\square$\ \textbf{(C)} $B - A$ & $\square$\ \textbf{(D)} $2 * B$
\end{tabular}

\tcblower

\textbf{Official Key:} \textbf{\large B} \hfill \textit{\small [Medium]}
\vspace{2pt}\hrule height 0.3pt\vspace{3pt}
\textbf{Topic:} Digital Logic \\
\textbf{Subtopic:} \textit{Number representations and computer arithmetic (fixed and floating point)}
\vspace{2pt}\hrule height 0.3pt\vspace{3pt}
\textbf{Guide / Hint:} \\
{\small In a 5-bit 2's complement representation, the range is $[-2^4, 2^4 - 1] = [-16, +15]$. Here, $A = 01010_2 = +10$ and $B = 11010_2 = -6$. The operation $A - B = 10 - (-6) = +16$, which exceeds $+15$, resulting in an arithmetic overflow.}

\end{tcolorbox}
\vspace{6pt}\hrule height 0.35pt\vspace{6pt}
\noindent
\begin{tcolorbox}[
    blanker,
    sidebyside,
    sidebyside align=top,
    lefthand width=0.55\textwidth,
    top=4pt, bottom=4pt, left=0pt, right=0pt,
    middle=16pt,
    underlay={
        \draw[dashed, black!80, line width=0.8pt, dash pattern=on 4pt off 3pt]
            ([xshift=0.55\textwidth+8pt]frame.north west) --
            ([xshift=0.55\textwidth+8pt]frame.south west);
    }
]
\textbf{Q.92} \quad {\footnotesize\textsc{[MCQ $\cdot$ 1 Mark(s) $\cdot$ Neg: 0.33]}} \hfill {\footnotesize\textbf{GATE 2024 CS2} \quad \tiny\texttt{2024\_CS2\_CS\_Q14}}
\vspace{2pt}\hrule height 0.3pt\vspace{4pt}

The format of a single-precision floating-point number as per the IEEE 754 standard is:
Sign (1 bit), Exponent (8 bits), Mantissa (23 bits).
Choose the largest floating-point number among the following options.

\par\vspace{6pt}\noindent
\begin{tabular}{@{}p{0.98\linewidth}@{}}
$\square$\ \textbf{(A)} Sign: 0, Exponent: 0111 1111, Mantissa: 1111 1111 1111 1111 1111 111 \\
\addlinespace[3pt]
$\square$\ \textbf{(B)} Sign: 0, Exponent: 1111 1110, Mantissa: 1111 1111 1111 1111 1111 111 \\
\addlinespace[3pt]
$\square$\ \textbf{(C)} Sign: 0, Exponent: 1111 1111, Mantissa: 1111 1111 1111 1111 1111 111 \\
\addlinespace[3pt]
$\square$\ \textbf{(D)} Sign: 0, Exponent: 0111 1111, Mantissa: 0000 0000 0000 0000 0000 000
\end{tabular}

\tcblower

\textbf{Official Key:} \textbf{\large B} \hfill \textit{\small [Easy]}
\vspace{2pt}\hrule height 0.3pt\vspace{3pt}
\textbf{Topic:} Digital Logic \\
\textbf{Subtopic:} \textit{Number representations and computer arithmetic (fixed and floating point)}
\vspace{2pt}\hrule height 0.3pt\vspace{3pt}
\textbf{Guide / Hint:} \\
{\small In IEEE 754 single precision: sign bit 0 means positive. The exponent field 11111111 is reserved for $\pm\infty$ and NaN. The largest finite representable number has the maximum non-reserved biased exponent, which is $254 = 11111110_2$, with all mantissa bits set to 1 ($111\dots1$), corresponding to option B.}

\end{tcolorbox}
\vspace{6pt}\hrule height 0.35pt\vspace{6pt}
\noindent
\begin{tcolorbox}[
    blanker,
    sidebyside,
    sidebyside align=top,
    lefthand width=0.55\textwidth,
    top=4pt, bottom=4pt, left=0pt, right=0pt,
    middle=16pt,
    underlay={
        \draw[dashed, black!80, line width=0.8pt, dash pattern=on 4pt off 3pt]
            ([xshift=0.55\textwidth+8pt]frame.north west) --
            ([xshift=0.55\textwidth+8pt]frame.south west);
    }
]
\textbf{Q.93} \quad {\footnotesize\textsc{[MCQ $\cdot$ 2 Mark(s) $\cdot$ Neg: 0.66]}} \hfill {\footnotesize\textbf{GATE 2025 CS1} \quad \tiny\texttt{2025\_CS1\_CS\_Q42}}
\vspace{2pt}\hrule height 0.3pt\vspace{4pt}

Consider the following four variable Boolean function in sum-of-product form:
\[ F(b\_3, b\_2, b\_1, b\_0) = \sum(0, 2, 4, 8, 10, 11, 12) \]
where the value of the function is computed by considering $b_3b_2b_1b_0$ as a 4-bit binary number, where $b_3$ denotes the most significant bit and $b_0$ denotes the least significant bit. Note that there are no don't care terms. Which ONE of the following options is the CORRECT minimized expression?

\par\vspace{6pt}\noindent
\begin{tabular}{@{}p{0.98\linewidth}@{}}
$\square$\ \textbf{(A)} $\overline{b}_1\overline{b}_0 + \overline{b}_2\overline{b}_0 + b_1\overline{b}_2b_3$ \\
\addlinespace[3pt]
$\square$\ \textbf{(B)} $\overline{b}_1\overline{b}_0 + \overline{b}_2\overline{b}_0$ \\
\addlinespace[3pt]
$\square$\ \textbf{(C)} $\overline{b}_2\overline{b}_0 + b_1b_2b_3$ \\
\addlinespace[3pt]
$\square$\ \textbf{(D)} $\overline{b}_0\overline{b}_2 + \overline{b}_3$
\end{tabular}

\tcblower

\textbf{Official Key:} \textbf{\large A} \hfill \textit{\small [Medium]}
\vspace{2pt}\hrule height 0.3pt\vspace{3pt}
\textbf{Topic:} Digital Logic \\
\textbf{Subtopic:} \textit{Minimization}
\vspace{2pt}\hrule height 0.3pt\vspace{3pt}
\textbf{Guide / Hint:} \\
{\small Grouping minterms on a 4-variable K-map: minterms $\{0, 2, 8, 10\}$ form a 4-cell group giving $\overline{b}_2\overline{b}_0$; minterms $\{0, 4, 8, 12\}$ form a 4-cell group giving $\overline{b}_1\overline{b}_0$; remaining minterm 11 combines with 10 to give $b_3\overline{b}_2b_1$. Combining these essential prime implicants yields $\overline{b}_1\overline{b}_0 + \overline{b}_2\overline{b}_0 + b_1\overline{b}_2b_3$.}

\end{tcolorbox}
\vspace{6pt}\hrule height 0.35pt\vspace{6pt}
\vspace{8pt}
\subsection*{Computer Organization and Architecture}
\vspace{-4pt}
\noindent
\begin{tcolorbox}[
    blanker,
    sidebyside,
    sidebyside align=top,
    lefthand width=0.55\textwidth,
    top=4pt, bottom=4pt, left=0pt, right=0pt,
    middle=16pt,
    underlay={
        \draw[dashed, black!80, line width=0.8pt, dash pattern=on 4pt off 3pt]
            ([xshift=0.55\textwidth+8pt]frame.north west) --
            ([xshift=0.55\textwidth+8pt]frame.south west);
    }
]
\textbf{Q.94} \quad {\footnotesize\textsc{[MCQ $\cdot$ 1 Mark(s) $\cdot$ Neg: 0.33]}} \hfill {\footnotesize\textbf{GATE 2018 CS} \quad \tiny\texttt{2018\_CS\_GA\_Q05\_v2}}
\vspace{2pt}\hrule height 0.3pt\vspace{4pt}

Consider the following processor design characteristics.
\begin{enumerate}[label=\Roman*.]
    \item Register-to-register arithmetic operations only
    \item Fixed-length instruction format
    \item Hardwired control unit
\end{enumerate}
Which of the characteristics above are used in the design of a RISC processor?

\par\vspace{6pt}\noindent
\begin{tabular}{@{}p{0.48\linewidth}@{\hspace{0.04\linewidth}}p{0.48\linewidth}@{}}
$\square$\ \textbf{(A)} I and II only & $\square$\ \textbf{(B)} II and III only \\
\addlinespace[3pt]
$\square$\ \textbf{(C)} I and III only & $\square$\ \textbf{(D)} I, II and III
\end{tabular}

\tcblower

\textbf{Official Key:} \textbf{\large D} \hfill \textit{\small [Easy]}
\vspace{2pt}\hrule height 0.3pt\vspace{3pt}
\textbf{Topic:} Computer Organization and Architecture \\
\textbf{Subtopic:} \textit{Machine instructions and addressing modes}
\vspace{2pt}\hrule height 0.3pt\vspace{3pt}
\textbf{Guide / Hint:} \\
{\small RISC (Reduced Instruction Set Computer) architectures feature load-store designs where arithmetic operations operate only on registers (I), fixed-length instruction encodings for streamlined pipelining (II), and hardwired control units for single-cycle instruction execution (III). Hence, all three characteristics I, II, and III are typical of RISC processors.}

\end{tcolorbox}
\vspace{6pt}\hrule height 0.35pt\vspace{6pt}
\noindent
\begin{tcolorbox}[
    blanker,
    sidebyside,
    sidebyside align=top,
    lefthand width=0.55\textwidth,
    top=4pt, bottom=4pt, left=0pt, right=0pt,
    middle=16pt,
    underlay={
        \draw[dashed, black!80, line width=0.8pt, dash pattern=on 4pt off 3pt]
            ([xshift=0.55\textwidth+8pt]frame.north west) --
            ([xshift=0.55\textwidth+8pt]frame.south west);
    }
]
\textbf{Q.95} \quad {\footnotesize\textsc{[MCQ $\cdot$ 1 Mark(s) $\cdot$ Neg: 0.33]}} \hfill {\footnotesize\textbf{GATE 2018 CS} \quad \tiny\texttt{2018\_CS\_GA\_Q09}}
\vspace{2pt}\hrule height 0.3pt\vspace{4pt}

The following are some events that occur after a device controller issues an interrupt while process $L$ is under execution.
\begin{itemize}[leftmargin=*, nosep]
    \item[P.] The processor pushes the process status of $L$ on to the control stack.
    \item[Q.] The processor finishes the execution of the current instruction.
    \item[R.] The processor executes the interrupt service routine.
    \item[S.] The processor pops the process status of $L$ from the control stack.
    \item[T.] The processor fetches the new PC value from the interrupt vector.
\end{itemize}
Which one of the following is the correct order of events?

\par\vspace{6pt}\noindent
\begin{tabular}{@{}p{0.48\linewidth}@{\hspace{0.04\linewidth}}p{0.48\linewidth}@{}}
$\square$\ \textbf{(A)} Q, P, T, R, S & $\square$\ \textbf{(B)} P, Q, T, R, S \\
\addlinespace[3pt]
$\square$\ \textbf{(C)} Q, P, R, T, S & $\square$\ \textbf{(D)} P, Q, R, T, S
\end{tabular}

\tcblower

\textbf{Official Key:} \textbf{\large A} \hfill \textit{\small [Easy]}
\vspace{2pt}\hrule height 0.3pt\vspace{3pt}
\textbf{Topic:} Computer Organization and Architecture \\
\textbf{Subtopic:} \textit{I/O interface: interrupt and DMA mode}
\vspace{2pt}\hrule height 0.3pt\vspace{3pt}
\textbf{Guide / Hint:} \\
{\small Upon receiving an interrupt, the processor finishes the execution of the current instruction (Q), saves the processor status and return PC onto the stack (P), loads the starting address of the ISR from the interrupt vector table (T), executes the ISR (R), and finally restores the saved status and resumes process $L$ (S). Hence the order is QPTRS.}

\end{tcolorbox}
\vspace{6pt}\hrule height 0.35pt\vspace{6pt}
\noindent
\begin{tcolorbox}[
    blanker,
    sidebyside,
    sidebyside align=top,
    lefthand width=0.55\textwidth,
    top=4pt, bottom=4pt, left=0pt, right=0pt,
    middle=16pt,
    underlay={
        \draw[dashed, black!80, line width=0.8pt, dash pattern=on 4pt off 3pt]
            ([xshift=0.55\textwidth+8pt]frame.north west) --
            ([xshift=0.55\textwidth+8pt]frame.south west);
    }
]
\textbf{Q.96} \quad {\footnotesize\textsc{[MCQ $\cdot$ 2 Mark(s) $\cdot$ Neg: 0.66]}} \hfill {\footnotesize\textbf{GATE 2018 CS} \quad \tiny\texttt{2018\_CS\_GA\_Q34}}
\vspace{2pt}\hrule height 0.3pt\vspace{4pt}

The size of the physical address space of a processor is $2^P$ bytes. The word length is $2^W$ bytes. The capacity of cache memory is $2^N$ bytes. The size of each cache block is $2^M$ words. For a $K$-way set-associative cache memory, the length (in number of bits) of the tag field is

\par\vspace{6pt}\noindent
\begin{tabular}{@{}p{0.48\linewidth}@{\hspace{0.04\linewidth}}p{0.48\linewidth}@{}}
$\square$\ \textbf{(A)} $P - N - \log_2 K$ & $\square$\ \textbf{(B)} $P - N + \log_2 K$ \\
\addlinespace[3pt]
$\square$\ \textbf{(C)} $P - N - M - W - \log_2 K$ & $\square$\ \textbf{(D)} $P - N - M - W + \log_2 K$
\end{tabular}

\tcblower

\textbf{Official Key:} \textbf{\large B} \hfill \textit{\small [Medium]}
\vspace{2pt}\hrule height 0.3pt\vspace{3pt}
\textbf{Topic:} Computer Organization and Architecture \\
\textbf{Subtopic:} \textit{Memory hierarchy: cache, main memory and secondary storage}
\vspace{2pt}\hrule height 0.3pt\vspace{3pt}
\textbf{Guide / Hint:} \\
{\small Block size $= 2^{M+W}$ bytes $\implies M+W$ offset bits. Cache size $= 2^N$ bytes, so number of lines $= 2^N / 2^{M+W} = 2^{N-M-W}$. With $K = 2^k$ lines per set, number of sets $= 2^{N-M-W-k} \implies N-M-W-k$ set index bits. Tag bits $= P - (N-M-W-k) - (M+W) = P - N + k = P - N + \log_2 K$.}

\end{tcolorbox}
\vspace{6pt}\hrule height 0.35pt\vspace{6pt}
\noindent
\begin{tcolorbox}[
    blanker,
    sidebyside,
    sidebyside align=top,
    lefthand width=0.55\textwidth,
    top=4pt, bottom=4pt, left=0pt, right=0pt,
    middle=16pt,
    underlay={
        \draw[dashed, black!80, line width=0.8pt, dash pattern=on 4pt off 3pt]
            ([xshift=0.55\textwidth+8pt]frame.north west) --
            ([xshift=0.55\textwidth+8pt]frame.south west);
    }
]
\textbf{Q.97} \quad {\footnotesize\textsc{[MCQ $\cdot$ 1 Mark(s) $\cdot$ Neg: 0.33]}} \hfill {\footnotesize\textbf{GATE 2022 CS} \quad \tiny\texttt{2022\_CS\_CS\_Q17}}
\vspace{2pt}\hrule height 0.3pt\vspace{4pt}

Which one of the following facilitates transfer of bulk data from hard disk to main memory with the highest throughput?

\par\vspace{6pt}\noindent
\begin{tabular}{@{}p{0.48\linewidth}@{\hspace{0.04\linewidth}}p{0.48\linewidth}@{}}
$\square$\ \textbf{(A)} DMA based I/O transfer & $\square$\ \textbf{(B)} Interrupt driven I/O transfer \\
\addlinespace[3pt]
$\square$\ \textbf{(C)} Polling based I/O transfer & $\square$\ \textbf{(D)} Programmed I/O transfer
\end{tabular}

\tcblower

\textbf{Official Key:} \textbf{\large A} \hfill \textit{\small [Easy]}
\vspace{2pt}\hrule height 0.3pt\vspace{3pt}
\textbf{Topic:} Computer Organization and Architecture \\
\textbf{Subtopic:} \textit{I/O interface: interrupt and DMA mode}
\vspace{2pt}\hrule height 0.3pt\vspace{3pt}
\textbf{Guide / Hint:} \\
{\small Direct Memory Access (DMA) enables high-speed bulk data transfer directly between the disk controller and main memory without continuous CPU intervention, achieving the highest possible data throughput compared to programmed or interrupt-driven I/O.}

\end{tcolorbox}
\vspace{6pt}\hrule height 0.35pt\vspace{6pt}
\noindent
\begin{tcolorbox}[
    blanker,
    sidebyside,
    sidebyside align=top,
    lefthand width=0.55\textwidth,
    top=4pt, bottom=4pt, left=0pt, right=0pt,
    middle=16pt,
    underlay={
        \draw[dashed, black!80, line width=0.8pt, dash pattern=on 4pt off 3pt]
            ([xshift=0.55\textwidth+8pt]frame.north west) --
            ([xshift=0.55\textwidth+8pt]frame.south west);
    }
]
\textbf{Q.98} \quad {\footnotesize\textsc{[MCQ $\cdot$ 2 Mark(s) $\cdot$ Neg: 0.66]}} \hfill {\footnotesize\textbf{GATE 2023 CS} \quad \tiny\texttt{2023\_CS\_CS\_Q41}}
\vspace{2pt}\hrule height 0.3pt\vspace{4pt}

Consider the given C-code and its corresponding assembly code, with a few operands U1--U4 being unknown. Some useful information as well as the semantics of each unique assembly instruction is annotated as inline comments in the code. The memory is byte-addressable.
\begin{lstlisting}[language=C]
// C-code
int a[10], b[10], i;    // int is 32-bit
for (i = 0; i < 10; i++)
    a[i] = b[i] * 8;
\end{lstlisting}
\begin{lstlisting}
; assembly-code (; indicates comments)
; r1-r5 are 32-bit integer registers
; initialize r1=0, r2=10
; initialize r3, r4 with base address of a, b
L01: jeq r1, r2, end    ; if(r1==r2) goto end
L02: lw  r5, 0(r4)      ; r5 <- Memory[r4+0]
L03: shl r5, r5, U1     ; r5 <- r5 << U1
L04: sw  r5, 0(r3)      ; Memory[r3+0] <- r5
L05: add r3, r3, U2     ; r3 <- r3+U2
L06: add r4, r4, U3
L07: add r1, r1, 1
L08: jmp U4             ; goto U4
L09: end
\end{lstlisting}
Which one of the following options is a CORRECT replacement for operands in the position (U1, U2, U3, U4) in the above assembly code?

\par\vspace{6pt}\noindent
\begin{tabular}{@{}p{0.48\linewidth}@{\hspace{0.04\linewidth}}p{0.48\linewidth}@{}}
$\square$\ \textbf{(A)} (8, 4, 1, L02) & $\square$\ \textbf{(B)} (3, 4, 4, L01) \\
\addlinespace[3pt]
$\square$\ \textbf{(C)} (8, 1, 1, L02) & $\square$\ \textbf{(D)} (3, 1, 1, L01)
\end{tabular}

\tcblower

\textbf{Official Key:} \textbf{\large B} \hfill \textit{\small [Medium]}
\vspace{2pt}\hrule height 0.3pt\vspace{3pt}
\textbf{Topic:} Computer Organization and Architecture \\
\textbf{Subtopic:} \textit{Machine instructions and addressing modes}
\vspace{2pt}\hrule height 0.3pt\vspace{3pt}
\textbf{Guide / Hint:} \\
{\small Multiplying by 8 corresponds to shifting left by 3 bits ($2^3 = 8$), so $U_1 = 3$. Since elements are 32-bit (4-byte) integers in byte-addressable memory, the pointers in registers `r3` and `r4` must advance by 4 bytes each iteration, so $U_2 = 4$ and $U_3 = 4$. Finally, the loop branches back to the condition check at label $L_{01}$, so $U_4 = L_{01}$.}

\end{tcolorbox}
\vspace{6pt}\hrule height 0.35pt\vspace{6pt}
\noindent
\begin{tcolorbox}[
    blanker,
    sidebyside,
    sidebyside align=top,
    lefthand width=0.55\textwidth,
    top=4pt, bottom=4pt, left=0pt, right=0pt,
    middle=16pt,
    underlay={
        \draw[dashed, black!80, line width=0.8pt, dash pattern=on 4pt off 3pt]
            ([xshift=0.55\textwidth+8pt]frame.north west) --
            ([xshift=0.55\textwidth+8pt]frame.south west);
    }
]
\textbf{Q.99} \quad {\footnotesize\textsc{[MCQ $\cdot$ 1 Mark(s) $\cdot$ Neg: 0.33]}} \hfill {\footnotesize\textbf{GATE 2024 CS1} \quad \tiny\texttt{2024\_CS1\_CS\_Q15}}
\vspace{2pt}\hrule height 0.3pt\vspace{4pt}

Which one of the following statements is FALSE?

\par\vspace{6pt}\noindent
\begin{tabular}{@{}p{0.98\linewidth}@{}}
$\square$\ \textbf{(A)} In the cycle stealing mode of DMA, one word of data is transferred between an I/O device and main memory in a stolen cycle \\
\addlinespace[3pt]
$\square$\ \textbf{(B)} For bulk data transfer, the burst mode of DMA has a higher throughput than the cycle stealing mode \\
\addlinespace[3pt]
$\square$\ \textbf{(C)} Programmed I/O mechanism has a better CPU utilization than the interrupt driven I/O mechanism \\
\addlinespace[3pt]
$\square$\ \textbf{(D)} The CPU can start executing an interrupt service routine faster with vectored interrupts than with non-vectored interrupts
\end{tabular}

\tcblower

\textbf{Official Key:} \textbf{\large C} \hfill \textit{\small [Easy]}
\vspace{2pt}\hrule height 0.3pt\vspace{3pt}
\textbf{Topic:} Computer Organization and Architecture \\
\textbf{Subtopic:} \textit{I/O interface: interrupt and DMA mode}
\vspace{2pt}\hrule height 0.3pt\vspace{3pt}
\textbf{Guide / Hint:} \\
{\small Programmed I/O requires the CPU to stay in a busy-waiting polling loop to check device readiness, wasting processor cycles and yielding poor CPU utilization. Interrupt-driven I/O frees the CPU to execute other tasks while I/O operations are underway, providing significantly better CPU utilization. Thus, statement (C) is FALSE.}

\end{tcolorbox}
\vspace{6pt}\hrule height 0.35pt\vspace{6pt}
\noindent
\begin{tcolorbox}[
    blanker,
    sidebyside,
    sidebyside align=top,
    lefthand width=0.55\textwidth,
    top=4pt, bottom=4pt, left=0pt, right=0pt,
    middle=16pt,
    underlay={
        \draw[dashed, black!80, line width=0.8pt, dash pattern=on 4pt off 3pt]
            ([xshift=0.55\textwidth+8pt]frame.north west) --
            ([xshift=0.55\textwidth+8pt]frame.south west);
    }
]
\textbf{Q.100} \quad {\footnotesize\textsc{[MCQ $\cdot$ 1 Mark(s) $\cdot$ Neg: 0.33]}} \hfill {\footnotesize\textbf{GATE 2024 CS2} \quad \tiny\texttt{2024\_CS2\_CS\_Q11}}
\vspace{2pt}\hrule height 0.3pt\vspace{4pt}

Consider a computer with a 4 MHz processor. Its DMA controller can transfer 8 bytes in 1 cycle from a device to main memory through cycle stealing at regular intervals. Which one of the following is the data transfer rate (in bits per second) of the DMA controller if 1\% of the processor cycles are used for DMA?

\par\vspace{6pt}\noindent
\begin{tabular}{@{}p{0.48\linewidth}@{\hspace{0.04\linewidth}}p{0.48\linewidth}@{}}
$\square$\ \textbf{(A)} 2,56,000 & $\square$\ \textbf{(B)} 3,200 \\
\addlinespace[3pt]
$\square$\ \textbf{(C)} 25,60,000 & $\square$\ \textbf{(D)} 32,000
\end{tabular}

\tcblower

\textbf{Official Key:} \textbf{\large C} \hfill \textit{\small [Medium]}
\vspace{2pt}\hrule height 0.3pt\vspace{3pt}
\textbf{Topic:} Computer Organization and Architecture \\
\textbf{Subtopic:} \textit{I/O interface: interrupt and DMA mode}
\vspace{2pt}\hrule height 0.3pt\vspace{3pt}
\textbf{Guide / Hint:} \\
{\small Processor clock frequency is $4\text{ MHz} = 4 \times 10^6$ cycles/s. DMA steals 1\% of cycles $= 0.01 \times 4 \times 10^6 = 40,000$ cycles/s. In each stolen cycle, 8 bytes ($8 \times 8 = 64$ bits) are transferred. Transfer rate $= 40,000 \times 64 = 2,560,000$ bps $= 25,60,000$ bps.}

\end{tcolorbox}
\vspace{6pt}\hrule height 0.35pt\vspace{6pt}
\noindent
\begin{tcolorbox}[
    blanker,
    sidebyside,
    sidebyside align=top,
    lefthand width=0.55\textwidth,
    top=4pt, bottom=4pt, left=0pt, right=0pt,
    middle=16pt,
    underlay={
        \draw[dashed, black!80, line width=0.8pt, dash pattern=on 4pt off 3pt]
            ([xshift=0.55\textwidth+8pt]frame.north west) --
            ([xshift=0.55\textwidth+8pt]frame.south west);
    }
]
\textbf{Q.101} \quad {\footnotesize\textsc{[MCQ $\cdot$ 1 Mark(s) $\cdot$ Neg: 0.33]}} \hfill {\footnotesize\textbf{GATE 2025 CS1} \quad \tiny\texttt{2025\_CS1\_CS\_Q11}}
\vspace{2pt}\hrule height 0.3pt\vspace{4pt}

Suppose a program is running on a non-pipelined single processor computer system. The computer is connected to an external device that can interrupt the processor asynchronously. The processor needs to execute the interrupt service routine (ISR) to serve this interrupt. The following steps (not necessarily in order) are taken by the processor when the interrupt arrives:
\begin{enumerate}[label=(\roman*)]
\item The processor saves the content of the program counter.
\item The program counter is loaded with the start address of the ISR.
\item The processor finishes the present instruction.
\end{enumerate}
Which ONE of the following is the CORRECT sequence of steps?

\par\vspace{6pt}\noindent
\begin{tabular}{@{}p{0.48\linewidth}@{\hspace{0.04\linewidth}}p{0.48\linewidth}@{}}
$\square$\ \textbf{(A)} (iii), (i), (ii) & $\square$\ \textbf{(B)} (i), (iii), (ii) \\
\addlinespace[3pt]
$\square$\ \textbf{(C)} (i), (ii), (iii) & $\square$\ \textbf{(D)} (iii), (ii), (i)
\end{tabular}

\tcblower

\textbf{Official Key:} \textbf{\large A} \hfill \textit{\small [Easy]}
\vspace{2pt}\hrule height 0.3pt\vspace{3pt}
\textbf{Topic:} Computer Organization and Architecture \\
\textbf{Subtopic:} \textit{I/O interface: interrupt and DMA mode}
\vspace{2pt}\hrule height 0.3pt\vspace{3pt}
\textbf{Guide / Hint:} \\
{\small When an asynchronous interrupt arrives, the processor first finishes the execution of the current instruction (iii). Next, it saves the current program counter (and processor state) onto the stack (i). Finally, it loads the starting address of the interrupt service routine into the program counter (ii).}

\end{tcolorbox}
\vspace{6pt}\hrule height 0.35pt\vspace{6pt}
\noindent
\begin{tcolorbox}[
    blanker,
    sidebyside,
    sidebyside align=top,
    lefthand width=0.55\textwidth,
    top=4pt, bottom=4pt, left=0pt, right=0pt,
    middle=16pt,
    underlay={
        \draw[dashed, black!80, line width=0.8pt, dash pattern=on 4pt off 3pt]
            ([xshift=0.55\textwidth+8pt]frame.north west) --
            ([xshift=0.55\textwidth+8pt]frame.south west);
    }
]
\textbf{Q.102} \quad {\footnotesize\textsc{[MCQ $\cdot$ 2 Mark(s) $\cdot$ Neg: 0.66]}} \hfill {\footnotesize\textbf{GATE 2025 CS1} \quad \tiny\texttt{2025\_CS1\_CS\_Q36}}
\vspace{2pt}\hrule height 0.3pt\vspace{4pt}

Consider a memory system with 1M bytes of main memory and 16K bytes of cache memory. Assume that the processor generates 20-bit memory address, and the cache block size is 16 bytes. If the cache uses direct mapping, how many bits will be required to store all the tag values? [Assume memory is byte addressable, $1\text{K}=2^{10}, 1\text{M}=2^{20}$.]

\par\vspace{6pt}\noindent
\begin{tabular}{@{}p{0.48\linewidth}@{\hspace{0.04\linewidth}}p{0.48\linewidth}@{}}
$\square$\ \textbf{(A)} $6 \times 2^{10}$ & $\square$\ \textbf{(B)} $8 \times 2^{10}$ \\
\addlinespace[3pt]
$\square$\ \textbf{(C)} $2^{12}$ & $\square$\ \textbf{(D)} $2^{14}$
\end{tabular}

\tcblower

\textbf{Official Key:} \textbf{\large A} \hfill \textit{\small [Medium]}
\vspace{2pt}\hrule height 0.3pt\vspace{3pt}
\textbf{Topic:} Computer Organization and Architecture \\
\textbf{Subtopic:} \textit{Memory hierarchy: cache, main memory and secondary storage}
\vspace{2pt}\hrule height 0.3pt\vspace{3pt}
\textbf{Guide / Hint:} \\
{\small With block size $16 = 2^4$ bytes, block offset is 4 bits. Cache size is $16\text{ KB} = 2^{14}$ bytes, so number of cache lines is $2^{14}/2^4 = 2^{10} = 1024$, requiring 10 line index bits. The tag bits per line are $20 - (10 + 4) = 6$ bits. For all $2^{10}$ lines, total tag bits required are $6 \times 2^{10}$.}

\end{tcolorbox}
\vspace{6pt}\hrule height 0.35pt\vspace{6pt}
\noindent
\begin{tcolorbox}[
    blanker,
    sidebyside,
    sidebyside align=top,
    lefthand width=0.55\textwidth,
    top=4pt, bottom=4pt, left=0pt, right=0pt,
    middle=16pt,
    underlay={
        \draw[dashed, black!80, line width=0.8pt, dash pattern=on 4pt off 3pt]
            ([xshift=0.55\textwidth+8pt]frame.north west) --
            ([xshift=0.55\textwidth+8pt]frame.south west);
    }
]
\textbf{Q.103} \quad {\footnotesize\textsc{[MCQ $\cdot$ 2 Mark(s) $\cdot$ Neg: 0.66]}} \hfill {\footnotesize\textbf{GATE 2025 CS1} \quad \tiny\texttt{2025\_CS1\_CS\_Q37}}
\vspace{2pt}\hrule height 0.3pt\vspace{4pt}

A processor has 64 general-purpose registers and 50 distinct instruction types. An instruction is encoded in 32-bits. What is the maximum number of bits that can be used to store the immediate operand for the given instruction?
\begin{verbatim}
ADD R1, #25 // R1 = R1 + 25
\end{verbatim}

\par\vspace{6pt}\noindent
\begin{tabular}{@{}p{0.48\linewidth}@{\hspace{0.04\linewidth}}p{0.48\linewidth}@{}}
$\square$\ \textbf{(A)} 16 & $\square$\ \textbf{(B)} 20 \\
\addlinespace[3pt]
$\square$\ \textbf{(C)} 22 & $\square$\ \textbf{(D)} 24
\end{tabular}

\tcblower

\textbf{Official Key:} \textbf{\large B} \hfill \textit{\small [Medium]}
\vspace{2pt}\hrule height 0.3pt\vspace{3pt}
\textbf{Topic:} Computer Organization and Architecture \\
\textbf{Subtopic:} \textit{Machine instructions and addressing modes}
\vspace{2pt}\hrule height 0.3pt\vspace{3pt}
\textbf{Guide / Hint:} \\
{\small To represent 50 distinct instruction opcodes, $\lceil \log_2 50 \rceil = 6$ opcode bits are required. Specifying one destination register among 64 general-purpose registers requires $\lceil \log_2 64 \rceil = 6$ bits. In a 32-bit instruction word, the remaining bits for the immediate field are $32 - (6 + 6) = 20$ bits.}

\end{tcolorbox}
\vspace{6pt}\hrule height 0.35pt\vspace{6pt}
\noindent
\begin{tcolorbox}[
    blanker,
    sidebyside,
    sidebyside align=top,
    lefthand width=0.55\textwidth,
    top=4pt, bottom=4pt, left=0pt, right=0pt,
    middle=16pt,
    underlay={
        \draw[dashed, black!80, line width=0.8pt, dash pattern=on 4pt off 3pt]
            ([xshift=0.55\textwidth+8pt]frame.north west) --
            ([xshift=0.55\textwidth+8pt]frame.south west);
    }
]
\textbf{Q.104} \quad {\footnotesize\textsc{[MCQ $\cdot$ 2 Mark(s) $\cdot$ Neg: 0.66]}} \hfill {\footnotesize\textbf{GATE 2025 CS2} \quad \tiny\texttt{2025\_CS2\_CS\_Q39}}
\vspace{2pt}\hrule height 0.3pt\vspace{4pt}

For a direct-mapped cache, 4 bits are used for the tag field and 12 bits are used to index into a cache block. What are the cache data size and tag directory size respectively?

\par\vspace{6pt}\noindent
\begin{tabular}{@{}p{0.48\linewidth}@{\hspace{0.04\linewidth}}p{0.48\linewidth}@{}}
$\square$\ \textbf{(A)} 64 KB and 4 KB & $\square$\ \textbf{(B)} 128 KB and 16 KB \\
\addlinespace[3pt]
$\square$\ \textbf{(C)} 64 KB and 8 KB & $\square$\ \textbf{(D)} 128 KB and 6 KB
\end{tabular}

\tcblower

\textbf{Official Key:} \textbf{\large A} \hfill \textit{\small [Medium]}
\vspace{2pt}\hrule height 0.3pt\vspace{3pt}
\textbf{Topic:} Computer Organization and Architecture \\
\textbf{Subtopic:} \textit{Memory hierarchy: cache memory mapping}
\vspace{2pt}\hrule height 0.3pt\vspace{3pt}
\textbf{Guide / Hint:} \\
{\small Number of cache lines is $2^{12} = 4096$. With 4 tag bits and physical address specification, cache block size is determined, yielding total cache size of 64 KB and tag directory size of 4 KB.}

\end{tcolorbox}
\vspace{6pt}\hrule height 0.35pt\vspace{6pt}
\vspace{8pt}
\subsection*{Programming and Data Structures}
\vspace{-4pt}
\noindent
\begin{tcolorbox}[
    blanker,
    sidebyside,
    sidebyside align=top,
    lefthand width=0.55\textwidth,
    top=4pt, bottom=4pt, left=0pt, right=0pt,
    middle=16pt,
    underlay={
        \draw[dashed, black!80, line width=0.8pt, dash pattern=on 4pt off 3pt]
            ([xshift=0.55\textwidth+8pt]frame.north west) --
            ([xshift=0.55\textwidth+8pt]frame.south west);
    }
]
\textbf{Q.105} \quad {\footnotesize\textsc{[MCQ $\cdot$ 1 Mark(s) $\cdot$ Neg: 0.33]}} \hfill {\footnotesize\textbf{GATE 2018 CS} \quad \tiny\texttt{2018\_CS\_GA\_Q02\_v2}}
\vspace{2pt}\hrule height 0.3pt\vspace{4pt}

Consider the following C program.
\begin{lstlisting}[language=C]
#include <stdio.h>
struct Ournode {
    char x, y, z;
};
\par\medskip
int main() {
    struct Ournode p = {'1', '0', 'a' + 2};
    struct Ournode *q = &p;
    printf("%c, %c", *((char*)q + 1), *((char*)q + 2));
    return 0;
}
\end{lstlisting}
The output of this program is

\par\vspace{6pt}\noindent
\begin{tabular}{@{}p{0.48\linewidth}@{\hspace{0.04\linewidth}}p{0.48\linewidth}@{}}
$\square$\ \textbf{(A)} 0, c & $\square$\ \textbf{(B)} 0, a+2 \\
\addlinespace[3pt]
$\square$\ \textbf{(C)} '0', 'p+2' & $\square$\ \textbf{(D)} 1, 0
\end{tabular}

\tcblower

\textbf{Official Key:} \textbf{\large A} \hfill \textit{\small [Easy]}
\vspace{2pt}\hrule height 0.3pt\vspace{3pt}
\textbf{Topic:} Programming and Data Structures \\
\textbf{Subtopic:} \textit{Programming in C}
\vspace{2pt}\hrule height 0.3pt\vspace{3pt}
\textbf{Guide / Hint:} \\
{\small In `struct Ournode`, `char x, y, z;` are stored contiguously in memory. Since `p = {'1', '0', 'a'+2}`, `p.x = '1'`, `p.y = '0'`, and `p.z = 'c'`. Pointer `q` points to `p`; casting `q` to `(char*)` gives byte-level access where `*((char*)q + 1)` accesses `p.y` ('0') and `*((char*)q + 2)` accesses `p.z` ('c').}

\end{tcolorbox}
\vspace{6pt}\hrule height 0.35pt\vspace{6pt}
\noindent
\begin{tcolorbox}[
    blanker,
    sidebyside,
    sidebyside align=top,
    lefthand width=0.55\textwidth,
    top=4pt, bottom=4pt, left=0pt, right=0pt,
    middle=16pt,
    underlay={
        \draw[dashed, black!80, line width=0.8pt, dash pattern=on 4pt off 3pt]
            ([xshift=0.55\textwidth+8pt]frame.north west) --
            ([xshift=0.55\textwidth+8pt]frame.south west);
    }
]
\textbf{Q.106} \quad {\footnotesize\textsc{[MCQ $\cdot$ 2 Mark(s) $\cdot$ Neg: 0.66]}} \hfill {\footnotesize\textbf{GATE 2018 CS} \quad \tiny\texttt{2018\_CS\_GA\_Q29}}
\vspace{2pt}\hrule height 0.3pt\vspace{4pt}

Consider the following C program:
\begin{lstlisting}[language=C]
#include <stdio.h>
\par\medskip
void fun1(char *s1, char *s2) {
    char *tmp;
    tmp = s1;
    s1 = s2;
    s2 = tmp;
}
\par\medskip
void fun2(char **s1, char **s2) {
    char *tmp;
    tmp = *s1;
    *s1 = *s2;
    *s2 = tmp;
}
\par\medskip
int main() {
    char *str1 = "Hi", *str2 = "Bye";
    fun1(str1, str2);
    printf("%s %s ", str1, str2);
    fun2(&str1, &str2);
    printf("%s %s", str1, str2);
    return 0;
}
\end{lstlisting}
The output of the program is

\par\vspace{6pt}\noindent
\begin{tabular}{@{}p{0.48\linewidth}@{\hspace{0.04\linewidth}}p{0.48\linewidth}@{}}
$\square$\ \textbf{(A)} Hi Bye Bye Hi & $\square$\ \textbf{(B)} Hi Bye Hi Bye \\
\addlinespace[3pt]
$\square$\ \textbf{(C)} Bye Hi Bye Hi & $\square$\ \textbf{(D)} Bye Hi Hi Bye
\end{tabular}

\tcblower

\textbf{Official Key:} \textbf{\large A} \hfill \textit{\small [Medium]}
\vspace{2pt}\hrule height 0.3pt\vspace{3pt}
\textbf{Topic:} Programming and Data Structures \\
\textbf{Subtopic:} \textit{Programming in C}
\vspace{2pt}\hrule height 0.3pt\vspace{3pt}
\textbf{Guide / Hint:} \\
{\small In `fun1`, pointers `s1` and `s2` are passed by value, so the pointer swap only alters local variables and has no effect on `str1` and `str2` in `main` (printing "Hi Bye"). In `fun2`, the addresses of the pointers (`char**`) are passed, so dereferencing swaps the actual pointer values `str1` and `str2` (printing "Bye Hi").}

\end{tcolorbox}
\vspace{6pt}\hrule height 0.35pt\vspace{6pt}
\noindent
\begin{tcolorbox}[
    blanker,
    sidebyside,
    sidebyside align=top,
    lefthand width=0.55\textwidth,
    top=4pt, bottom=4pt, left=0pt, right=0pt,
    middle=16pt,
    underlay={
        \draw[dashed, black!80, line width=0.8pt, dash pattern=on 4pt off 3pt]
            ([xshift=0.55\textwidth+8pt]frame.north west) --
            ([xshift=0.55\textwidth+8pt]frame.south west);
    }
]
\textbf{Q.107} \quad {\footnotesize\textsc{[MCQ $\cdot$ 1 Mark(s) $\cdot$ Neg: 0.33]}} \hfill {\footnotesize\textbf{GATE 2022 CS} \quad \tiny\texttt{2022\_CS\_CS\_Q15}}
\vspace{2pt}\hrule height 0.3pt\vspace{4pt}

Consider the problem of reversing a singly linked list. Which one of the following statements is TRUE about the time complexity of algorithms that solve the above problem in $\mathcal{O}(1)$ space?

\par\vspace{6pt}\noindent
\begin{tabular}{@{}p{0.98\linewidth}@{}}
$\square$\ \textbf{(A)} The best algorithm for the problem takes $\Theta(n)$ time in the worst case. \\
\addlinespace[3pt]
$\square$\ \textbf{(B)} The best algorithm for the problem takes $\Theta(n \log n)$ time in the worst case. \\
\addlinespace[3pt]
$\square$\ \textbf{(C)} The best algorithm for the problem takes $\Theta(n^2)$ time in the worst case. \\
\addlinespace[3pt]
$\square$\ \textbf{(D)} It is not possible to reverse a singly linked list in $\mathcal{O}(1)$ space.
\end{tabular}

\tcblower

\textbf{Official Key:} \textbf{\large A} \hfill \textit{\small [Easy]}
\vspace{2pt}\hrule height 0.3pt\vspace{3pt}
\textbf{Topic:} Programming and Data Structures \\
\textbf{Subtopic:} \textit{Arrays, stacks, queues, linked lists}
\vspace{2pt}\hrule height 0.3pt\vspace{3pt}
\textbf{Guide / Hint:} \\
{\small A singly linked list can be reversed in-place in $\mathcal{O}(1)$ auxiliary space using three moving pointers (`prev`, `curr`, `next`) by reversing each pointer sequentially in a single pass of $n$ nodes, achieving optimal $\Theta(n)$ worst-case time.}

\end{tcolorbox}
\vspace{6pt}\hrule height 0.35pt\vspace{6pt}
\noindent
\begin{tcolorbox}[
    blanker,
    sidebyside,
    sidebyside align=top,
    lefthand width=0.55\textwidth,
    top=4pt, bottom=4pt, left=0pt, right=0pt,
    middle=16pt,
    underlay={
        \draw[dashed, black!80, line width=0.8pt, dash pattern=on 4pt off 3pt]
            ([xshift=0.55\textwidth+8pt]frame.north west) --
            ([xshift=0.55\textwidth+8pt]frame.south west);
    }
]
\textbf{Q.108} \quad {\footnotesize\textsc{[MCQ $\cdot$ 1 Mark(s) $\cdot$ Neg: 0.33]}} \hfill {\footnotesize\textbf{GATE 2022 CS} \quad \tiny\texttt{2022\_CS\_CS\_Q21}}
\vspace{2pt}\hrule height 0.3pt\vspace{4pt}

What is printed by the following ANSI C program?
\begin{lstlisting}[language=C]
#include <stdio.h>
int main(int argc, char *argv[]) {
    int x = 1, z[2] = {10, 11};
    int *p = NULL;
    p = &x;
    *p = 10;
    p = &z[1];
    *(&z[0] + 1) += 3;
    printf("%d, %d, %d\n", x, z[0], z[1]);
    return 0;
}
\end{lstlisting}

\par\vspace{6pt}\noindent
\begin{tabular}{@{}p{0.48\linewidth}@{\hspace{0.04\linewidth}}p{0.48\linewidth}@{}}
$\square$\ \textbf{(A)} 1, 10, 11 & $\square$\ \textbf{(B)} 1, 10, 14 \\
\addlinespace[3pt]
$\square$\ \textbf{(C)} 10, 14, 11 & $\square$\ \textbf{(D)} 10, 10, 14
\end{tabular}

\tcblower

\textbf{Official Key:} \textbf{\large D} \hfill \textit{\small [Easy]}
\vspace{2pt}\hrule height 0.3pt\vspace{3pt}
\textbf{Topic:} Programming and Data Structures \\
\textbf{Subtopic:} \textit{Programming in C}
\vspace{2pt}\hrule height 0.3pt\vspace{3pt}
\textbf{Guide / Hint:} \\
{\small Initially $x = 1, z[0] = 10, z[1] = 11$. The assignment `p = \&x; *p = 10;` sets $x = 10$. The line `*(\&z[0] + 1) += 3;` targets `z[1]`, updating $z[1] = 11 + 3 = 14$. The array element $z[0]$ remains 10. The output printed is `10, 10, 14`.}

\end{tcolorbox}
\vspace{6pt}\hrule height 0.35pt\vspace{6pt}
\noindent
\begin{tcolorbox}[
    blanker,
    sidebyside,
    sidebyside align=top,
    lefthand width=0.55\textwidth,
    top=4pt, bottom=4pt, left=0pt, right=0pt,
    middle=16pt,
    underlay={
        \draw[dashed, black!80, line width=0.8pt, dash pattern=on 4pt off 3pt]
            ([xshift=0.55\textwidth+8pt]frame.north west) --
            ([xshift=0.55\textwidth+8pt]frame.south west);
    }
]
\textbf{Q.109} \quad {\footnotesize\textsc{[MCQ $\cdot$ 2 Mark(s) $\cdot$ Neg: 0.66]}} \hfill {\footnotesize\textbf{GATE 2022 CS} \quad \tiny\texttt{2022\_CS\_CS\_Q43}}
\vspace{2pt}\hrule height 0.3pt\vspace{4pt}

What is printed by the following ANSI C program?
\begin{lstlisting}[language=C]
#include <stdio.h>
int main(int argc, char *argv[]) {
    int a[3][3][3] = {
        {1, 2, 3, 4, 5, 6, 7, 8, 9},
        {10, 11, 12, 13, 14, 15, 16, 17, 18},
        {19, 20, 21, 22, 23, 24, 25, 26, 27}
    };
    int i = 0, j = 0, k = 0;
    for(i = 0; i < 3; i++) {
        for(k = 0; k < 3; k++)
            printf("%d ", a[i][j][k]);
        printf("\n");
    }
    return 0;
}
\end{lstlisting}

\par\vspace{6pt}\noindent
\begin{tabular}{@{}p{0.48\linewidth}@{\hspace{0.04\linewidth}}p{0.48\linewidth}@{}}
$\square$\ \textbf{(A)} 1 2 3 \\ 10 11 12 \\ 19 20 21 & $\square$\ \textbf{(B)} 1 4 7 \\ 10 13 16 \\ 19 22 25 \\
\addlinespace[3pt]
$\square$\ \textbf{(C)} 1 2 3 \\ 4 5 6 \\ 7 8 9 & $\square$\ \textbf{(D)} 1 2 3 \\ 13 14 15 \\ 25 26 27
\end{tabular}

\tcblower

\textbf{Official Key:} \textbf{\large A} \hfill \textit{\small [Easy]}
\vspace{2pt}\hrule height 0.3pt\vspace{3pt}
\textbf{Topic:} Programming and Data Structures \\
\textbf{Subtopic:} \textit{Programming in C}
\vspace{2pt}\hrule height 0.3pt\vspace{3pt}
\textbf{Guide / Hint:} \\
{\small In the nested loops, index $j = 0$ is fixed throughout. For $i = 0$, $k$ iterates from 0 to 2, printing `a[0][0][0], a[0][0][1], a[0][0][2]` ($1\ 2\ 3$). For $i = 1$, it prints `a[1][0][0], a[1][0][1], a[1][0][2]` ($10\ 11\ 12$). For $i = 2$, it prints `a[2][0][0], a[2][0][1], a[2][0][2]` ($19\ 20\ 21$).}

\end{tcolorbox}
\vspace{6pt}\hrule height 0.35pt\vspace{6pt}
\noindent
\begin{tcolorbox}[
    blanker,
    sidebyside,
    sidebyside align=top,
    lefthand width=0.55\textwidth,
    top=4pt, bottom=4pt, left=0pt, right=0pt,
    middle=16pt,
    underlay={
        \draw[dashed, black!80, line width=0.8pt, dash pattern=on 4pt off 3pt]
            ([xshift=0.55\textwidth+8pt]frame.north west) --
            ([xshift=0.55\textwidth+8pt]frame.south west);
    }
]
\textbf{Q.110} \quad {\footnotesize\textsc{[MCQ $\cdot$ 2 Mark(s) $\cdot$ Neg: 0.66]}} \hfill {\footnotesize\textbf{GATE 2022 CS} \quad \tiny\texttt{2022\_CS\_CS\_Q44}}
\vspace{2pt}\hrule height 0.3pt\vspace{4pt}

What is printed by the following ANSI C program?
\begin{lstlisting}[language=C]
#include <stdio.h>
int main(int argc, char *argv[]) {
    char a = 'P';
    char b = 'x';
    char c = (a & b) + '*';
    char d = (a | b) - '-';
    char e = (a ^ b) + '+';
    printf("%c %c %c\n", c, d, e);
    return 0;
}
\end{lstlisting}

\par\vspace{6pt}\noindent
\begin{tabular}{@{}p{0.48\linewidth}@{\hspace{0.04\linewidth}}p{0.48\linewidth}@{}}
$\square$\ \textbf{(A)} z K S & $\square$\ \textbf{(B)} 122 75 83 \\
\addlinespace[3pt]
$\square$\ \textbf{(C)} * - + & $\square$\ \textbf{(D)} P x +
\end{tabular}

\tcblower

\textbf{Official Key:} \textbf{\large A} \hfill \textit{\small [Medium]}
\vspace{2pt}\hrule height 0.3pt\vspace{3pt}
\textbf{Topic:} Programming and Data Structures \\
\textbf{Subtopic:} \textit{Programming in C}
\vspace{2pt}\hrule height 0.3pt\vspace{3pt}
\textbf{Guide / Hint:} \\
{\small Character values: `a = 'P'` (80), `b = 'x'` (120). $c = (80 \& 120) + '*' = 80 + 42 = 122$, which is ASCII for 'z'. $d = (80 \mid 120) - '-' = 120 - 45 = 75$, which is ASCII for 'K'. $e = (80 \oplus 120) + '+' = 40 + 43 = 83$, which is ASCII for 'S'. The printed characters are `z K S`.}

\end{tcolorbox}
\vspace{6pt}\hrule height 0.35pt\vspace{6pt}
\noindent
\begin{tcolorbox}[
    blanker,
    sidebyside,
    sidebyside align=top,
    lefthand width=0.55\textwidth,
    top=4pt, bottom=4pt, left=0pt, right=0pt,
    middle=16pt,
    underlay={
        \draw[dashed, black!80, line width=0.8pt, dash pattern=on 4pt off 3pt]
            ([xshift=0.55\textwidth+8pt]frame.north west) --
            ([xshift=0.55\textwidth+8pt]frame.south west);
    }
]
\textbf{Q.111} \quad {\footnotesize\textsc{[MCQ $\cdot$ 1 Mark(s) $\cdot$ Neg: 0.33]}} \hfill {\footnotesize\textbf{GATE 2023 CS} \quad \tiny\texttt{2023\_CS\_CS\_Q12}}
\vspace{2pt}\hrule height 0.3pt\vspace{4pt}

Which one of the following sequences when stored in an array at locations $A[1], \dots, A[10]$ forms a max-heap?

\par\vspace{6pt}\noindent
\begin{tabular}{@{}p{0.98\linewidth}@{}}
$\square$\ \textbf{(A)} 23, 17, 10, 6, 13, 14, 1, 5, 7, 12 \\
\addlinespace[3pt]
$\square$\ \textbf{(B)} 23, 17, 14, 7, 13, 10, 1, 5, 6, 12 \\
\addlinespace[3pt]
$\square$\ \textbf{(C)} 23, 17, 14, 6, 13, 10, 1, 5, 7, 15 \\
\addlinespace[3pt]
$\square$\ \textbf{(D)} 23, 14, 17, 1, 10, 13, 16, 12, 7, 5
\end{tabular}

\tcblower

\textbf{Official Key:} \textbf{\large B} \hfill \textit{\small [Easy]}
\vspace{2pt}\hrule height 0.3pt\vspace{3pt}
\textbf{Topic:} Programming and Data Structures \\
\textbf{Subtopic:} \textit{Trees, binary search trees, binary heaps}
\vspace{2pt}\hrule height 0.3pt\vspace{3pt}
\textbf{Guide / Hint:} \\
{\small In a 1-indexed max-heap, every node at index $i$ must satisfy $A[i] \ge A[2i]$ and $A[i] \ge A[2i+1]$. For sequence B $[23, 17, 14, 7, 13, 10, 1, 5, 6, 12]$: $A[1]=23 \ge A[2]=17, A[3]=14$; $A[2]=17 \ge A[4]=7, A[5]=13$; $A[3]=14 \ge A[6]=10, A[7]=1$; $A[4]=7 \ge A[8]=5, A[9]=6$; $A[5]=13 \ge A[10]=12$. Every parent is greater than or equal to its children, forming a valid max-heap.}

\end{tcolorbox}
\vspace{6pt}\hrule height 0.35pt\vspace{6pt}
\noindent
\begin{tcolorbox}[
    blanker,
    sidebyside,
    sidebyside align=top,
    lefthand width=0.55\textwidth,
    top=4pt, bottom=4pt, left=0pt, right=0pt,
    middle=16pt,
    underlay={
        \draw[dashed, black!80, line width=0.8pt, dash pattern=on 4pt off 3pt]
            ([xshift=0.55\textwidth+8pt]frame.north west) --
            ([xshift=0.55\textwidth+8pt]frame.south west);
    }
]
\textbf{Q.112} \quad {\footnotesize\textsc{[MCQ $\cdot$ 1 Mark(s) $\cdot$ Neg: 0.33]}} \hfill {\footnotesize\textbf{GATE 2023 CS} \quad \tiny\texttt{2023\_CS\_CS\_Q13}}
\vspace{2pt}\hrule height 0.3pt\vspace{4pt}

Let `SLLdel` be a function that deletes a node in a singly-linked list given a pointer to the node and a pointer to the head of the list. Similarly, let `DLLdel` be another function that deletes a node in a doubly-linked list given a pointer to the node and a pointer to the head of the list.
Let $n$ denote the number of nodes in each of the linked lists. Which one of the following choices is TRUE about the worst-case time complexity of `SLLdel` and `DLLdel`?

\par\vspace{6pt}\noindent
\begin{tabular}{@{}p{0.98\linewidth}@{}}
$\square$\ \textbf{(A)} `SLLdel` is $\mathcal{O}(1)$ and `DLLdel` is $\mathcal{O}(n)$ \\
\addlinespace[3pt]
$\square$\ \textbf{(B)} Both `SLLdel` and `DLLdel` are $\mathcal{O}(\log n)$ \\
\addlinespace[3pt]
$\square$\ \textbf{(C)} Both `SLLdel` and `DLLdel` are $\mathcal{O}(1)$ \\
\addlinespace[3pt]
$\square$\ \textbf{(D)} `SLLdel` is $\mathcal{O}(n)$ and `DLLdel` is $\mathcal{O}(1)$
\end{tabular}

\tcblower

\textbf{Official Key:} \textbf{\large D} \hfill \textit{\small [Easy]}
\vspace{2pt}\hrule height 0.3pt\vspace{3pt}
\textbf{Topic:} Programming and Data Structures \\
\textbf{Subtopic:} \textit{Arrays, stacks, queues, linked lists}
\vspace{2pt}\hrule height 0.3pt\vspace{3pt}
\textbf{Guide / Hint:} \\
{\small In a singly-linked list, deleting a given node requires finding its predecessor node to bypass it, requiring an $\mathcal{O}(n)$ traversal from the head in the worst case. In a doubly-linked list, each node directly points to its predecessor (`node->prev`), allowing immediate pointer adjustments in $\mathcal{O}(1)$ time.}

\end{tcolorbox}
\vspace{6pt}\hrule height 0.35pt\vspace{6pt}
\noindent
\begin{tcolorbox}[
    blanker,
    sidebyside,
    sidebyside align=top,
    lefthand width=0.55\textwidth,
    top=4pt, bottom=4pt, left=0pt, right=0pt,
    middle=16pt,
    underlay={
        \draw[dashed, black!80, line width=0.8pt, dash pattern=on 4pt off 3pt]
            ([xshift=0.55\textwidth+8pt]frame.north west) --
            ([xshift=0.55\textwidth+8pt]frame.south west);
    }
]
\textbf{Q.113} \quad {\footnotesize\textsc{[MCQ $\cdot$ 2 Mark(s) $\cdot$ Neg: 0.66]}} \hfill {\footnotesize\textbf{GATE 2023 CS} \quad \tiny\texttt{2023\_CS\_CS\_Q46}}
\vspace{2pt}\hrule height 0.3pt\vspace{4pt}

Let $A$ be a priority queue for maintaining a set of elements. Suppose $A$ is implemented using a max-heap data structure. The operation $\text{Extract-Max}(A)$ extracts and deletes the maximum element from $A$. The operation $\text{Insert}(A, \text{key})$ inserts a new element $\text{key}$ in $A$. The properties of a max-heap are preserved at the end of each of these operations.
When $A$ contains $n$ elements, which one of the following statements about the worst case running time of these two operations is TRUE?

\par\vspace{6pt}\noindent
\begin{tabular}{@{}p{0.98\linewidth}@{}}
$\square$\ \textbf{(A)} Both `Extract-Max(A)` and `Insert(A,key)` run in $\mathcal{O}(1)$. \\
\addlinespace[3pt]
$\square$\ \textbf{(B)} Both `Extract-Max(A)` and `Insert(A,key)` run in $\mathcal{O}(\log n)$. \\
\addlinespace[3pt]
$\square$\ \textbf{(C)} `Extract-Max(A)` runs in $\mathcal{O}(1)$ whereas `Insert(A,key)` runs in $\mathcal{O}(n)$. \\
\addlinespace[3pt]
$\square$\ \textbf{(D)} `Extract-Max(A)` runs in $\mathcal{O}(1)$ whereas `Insert(A,key)` runs in $\mathcal{O}(\log n)$.
\end{tabular}

\tcblower

\textbf{Official Key:} \textbf{\large B} \hfill \textit{\small [Easy]}
\vspace{2pt}\hrule height 0.3pt\vspace{3pt}
\textbf{Topic:} Programming and Data Structures \\
\textbf{Subtopic:} \textit{Trees, binary search trees, binary heaps}
\vspace{2pt}\hrule height 0.3pt\vspace{3pt}
\textbf{Guide / Hint:} \\
{\small In a binary max-heap of $n$ elements, `Extract-Max` replaces the root with the last element and performs `maxHeapify` downward along the height $\lfloor \log_2 n \rfloor$, running in $\mathcal{O}(\log n)$ worst-case time. `Insert` places the new key at the next leaf position and sifts upward along the height, also running in $\mathcal{O}(\log n)$ worst-case time.}

\end{tcolorbox}
\vspace{6pt}\hrule height 0.35pt\vspace{6pt}
\noindent
\begin{tcolorbox}[
    blanker,
    sidebyside,
    sidebyside align=top,
    lefthand width=0.55\textwidth,
    top=4pt, bottom=4pt, left=0pt, right=0pt,
    middle=16pt,
    underlay={
        \draw[dashed, black!80, line width=0.8pt, dash pattern=on 4pt off 3pt]
            ([xshift=0.55\textwidth+8pt]frame.north west) --
            ([xshift=0.55\textwidth+8pt]frame.south west);
    }
]
\textbf{Q.114} \quad {\footnotesize\textsc{[MCQ $\cdot$ 2 Mark(s) $\cdot$ Neg: 0.66]}} \hfill {\footnotesize\textbf{GATE 2023 CS} \quad \tiny\texttt{2023\_CS\_CS\_Q47}}
\vspace{2pt}\hrule height 0.3pt\vspace{4pt}

Consider the C function `foo` and the binary tree shown:
\begin{lstlisting}[language=C]
typedef struct node {
    int val;
    struct node *left, *right;
} node;
int foo(node *p) {
    int retval;
    if (p == NULL)
        return 0;
    else {
        retval = p->val + foo(p->left) + foo(p->right);
        printf("%d ", retval);
        return retval;
    }
}
\end{lstlisting}
In the binary tree, root is 10; its left child is 5 with children 3 and 8; its right child is 11 with right child 13.
When `foo` is called with a pointer to the root node of the given binary tree, what will it print?

\par\vspace{6pt}\noindent
\begin{tabular}{@{}p{0.48\linewidth}@{\hspace{0.04\linewidth}}p{0.48\linewidth}@{}}
$\square$\ \textbf{(A)} 3 8 5 13 11 10 & $\square$\ \textbf{(B)} 3 5 8 10 11 13 \\
\addlinespace[3pt]
$\square$\ \textbf{(C)} 3 8 16 13 24 50 & $\square$\ \textbf{(D)} 3 16 8 50 24 13
\end{tabular}

\tcblower

\textbf{Official Key:} \textbf{\large C} \hfill \textit{\small [Medium]}
\vspace{2pt}\hrule height 0.3pt\vspace{3pt}
\textbf{Topic:} Programming and Data Structures \\
\textbf{Subtopic:} \textit{Recursion}
\vspace{2pt}\hrule height 0.3pt\vspace{3pt}
\textbf{Guide / Hint:} \\
{\small Function `foo` performs post-order traversal, recursively summing `p->val + foo(p->left) + foo(p->right)` and printing each subtree sum. Leaf 3 prints 3; leaf 8 prints 8; node 5 sums $5 + 3 + 8 = 16$ (prints 16); leaf 13 prints 13; node 11 sums $11 + 0 + 13 = 24$ (prints 24); root 10 sums $10 + 16 + 24 = 50$ (prints 50). The printed output sequence is 3 8 16 13 24 50.}

\end{tcolorbox}
\vspace{6pt}\hrule height 0.35pt\vspace{6pt}
\noindent
\begin{tcolorbox}[
    blanker,
    sidebyside,
    sidebyside align=top,
    lefthand width=0.55\textwidth,
    top=4pt, bottom=4pt, left=0pt, right=0pt,
    middle=16pt,
    underlay={
        \draw[dashed, black!80, line width=0.8pt, dash pattern=on 4pt off 3pt]
            ([xshift=0.55\textwidth+8pt]frame.north west) --
            ([xshift=0.55\textwidth+8pt]frame.south west);
    }
]
\textbf{Q.115} \quad {\footnotesize\textsc{[MCQ $\cdot$ 1 Mark(s) $\cdot$ Neg: 0.33]}} \hfill {\footnotesize\textbf{GATE 2024 CS1} \quad \tiny\texttt{2024\_CS1\_CS\_Q18}}
\vspace{2pt}\hrule height 0.3pt\vspace{4pt}

Consider the following C program:
\begin{lstlisting}[language=C]
#include <stdio.h>
int main() {
    int a = 6;
    int b = 0;
    while(a < 10) {
        a = a / 12 + 1;
        a += b;
    }
    printf("%d", a);
    return 0;
}
\end{lstlisting}
Which one of the following statements is CORRECT?

\par\vspace{6pt}\noindent
\begin{tabular}{@{}p{0.98\linewidth}@{}}
$\square$\ \textbf{(A)} The program prints 9 as output \\
\addlinespace[3pt]
$\square$\ \textbf{(B)} The program prints 10 as output \\
\addlinespace[3pt]
$\square$\ \textbf{(C)} The program gets stuck in an infinite loop \\
\addlinespace[3pt]
$\square$\ \textbf{(D)} The program prints 6 as output
\end{tabular}

\tcblower

\textbf{Official Key:} \textbf{\large C} \hfill \textit{\small [Easy]}
\vspace{2pt}\hrule height 0.3pt\vspace{3pt}
\textbf{Topic:} Programming and Data Structures \\
\textbf{Subtopic:} \textit{Programming in C}
\vspace{2pt}\hrule height 0.3pt\vspace{3pt}
\textbf{Guide / Hint:} \\
{\small Initially $a = 6, b = 0$. In the loop, integer division $a/12 = 6/12 = 0$, so $a = 0 + 1 = 1$, and $a += b \implies a = 1$. In subsequent iterations, $a=1 < 10$, $1/12 + 1 = 1$, so $a$ remains 1 in every iteration. The loop condition $a < 10$ remains true indefinitely, causing an infinite loop.}

\end{tcolorbox}
\vspace{6pt}\hrule height 0.35pt\vspace{6pt}
\noindent
\begin{tcolorbox}[
    blanker,
    sidebyside,
    sidebyside align=top,
    lefthand width=0.55\textwidth,
    top=4pt, bottom=4pt, left=0pt, right=0pt,
    middle=16pt,
    underlay={
        \draw[dashed, black!80, line width=0.8pt, dash pattern=on 4pt off 3pt]
            ([xshift=0.55\textwidth+8pt]frame.north west) --
            ([xshift=0.55\textwidth+8pt]frame.south west);
    }
]
\textbf{Q.116} \quad {\footnotesize\textsc{[MCQ $\cdot$ 1 Mark(s) $\cdot$ Neg: 0.33]}} \hfill {\footnotesize\textbf{GATE 2024 CS1} \quad \tiny\texttt{2024\_CS1\_CS\_Q19}}
\vspace{2pt}\hrule height 0.3pt\vspace{4pt}

Consider the following C program:
\begin{lstlisting}[language=C]
#include <stdio.h>
void fX();
int main() {
    fX();
    return 0;
}
void fX() {
    char a;
    if((a = getchar()) != '\n')
        fX();
    if(a != '\n')
        putchar(a);
}
\end{lstlisting}
Assume that the input to the program from the command line is `1234` followed by a newline character. Which one of the following statements is CORRECT?

\par\vspace{6pt}\noindent
\begin{tabular}{@{}p{0.98\linewidth}@{}}
$\square$\ \textbf{(A)} The program will not terminate \\
\addlinespace[3pt]
$\square$\ \textbf{(B)} The program will terminate with no output \\
\addlinespace[3pt]
$\square$\ \textbf{(C)} The program will terminate with 4321 as output \\
\addlinespace[3pt]
$\square$\ \textbf{(D)} The program will terminate with 1234 as output
\end{tabular}

\tcblower

\textbf{Official Key:} \textbf{\large C} \hfill \textit{\small [Easy]}
\vspace{2pt}\hrule height 0.3pt\vspace{3pt}
\textbf{Topic:} Programming and Data Structures \\
\textbf{Subtopic:} \textit{Recursion}
\vspace{2pt}\hrule height 0.3pt\vspace{3pt}
\textbf{Guide / Hint:} \\
{\small The function `fX()` recursively reads each character before printing. As each character (`'1'`, `'2'`, `'3'`, `'4'`) is read, it is saved on the call stack before `fX()` recurses. When `\n` is reached, recursion terminates and the characters are printed during stack unwinding in reverse (LIFO) order as `4321`.}

\end{tcolorbox}
\vspace{6pt}\hrule height 0.35pt\vspace{6pt}
\noindent
\begin{tcolorbox}[
    blanker,
    sidebyside,
    sidebyside align=top,
    lefthand width=0.55\textwidth,
    top=4pt, bottom=4pt, left=0pt, right=0pt,
    middle=16pt,
    underlay={
        \draw[dashed, black!80, line width=0.8pt, dash pattern=on 4pt off 3pt]
            ([xshift=0.55\textwidth+8pt]frame.north west) --
            ([xshift=0.55\textwidth+8pt]frame.south west);
    }
]
\textbf{Q.117} \quad {\footnotesize\textsc{[MCQ $\cdot$ 2 Mark(s) $\cdot$ Neg: 0.66]}} \hfill {\footnotesize\textbf{GATE 2024 CS1} \quad \tiny\texttt{2024\_CS1\_CS\_Q41}}
\vspace{2pt}\hrule height 0.3pt\vspace{4pt}

An array $[82, 101, 90, 11, 111, 75, 33, 131, 44, 93]$ is heapified. Which one of the following options represents the first three elements in the heapified array?

\par\vspace{6pt}\noindent
\begin{tabular}{@{}p{0.48\linewidth}@{\hspace{0.04\linewidth}}p{0.48\linewidth}@{}}
$\square$\ \textbf{(A)} 82, 90, 101 & $\square$\ \textbf{(B)} 82, 11, 93 \\
\addlinespace[3pt]
$\square$\ \textbf{(C)} 131, 11, 93 & $\square$\ \textbf{(D)} 131, 111, 90
\end{tabular}

\tcblower

\textbf{Official Key:} \textbf{\large D} \hfill \textit{\small [Medium]}
\vspace{2pt}\hrule height 0.3pt\vspace{3pt}
\textbf{Topic:} Programming and Data Structures \\
\textbf{Subtopic:} \textit{Trees, binary search trees, binary heaps}
\vspace{2pt}\hrule height 0.3pt\vspace{3pt}
\textbf{Guide / Hint:} \\
{\small Applying standard bottom-up max-heap construction (Floyd's algorithm) to $[82, 101, 90, 11, 111, 75, 33, 131, 44, 93]$ of length 10: non-leaf nodes are heapified from index 5 down to 1. The maximum element 131 bubbles up to index 1 (root), its left child becomes 111 (index 2), and right child becomes 90 (index 3).}

\end{tcolorbox}
\vspace{6pt}\hrule height 0.35pt\vspace{6pt}
\noindent
\begin{tcolorbox}[
    blanker,
    sidebyside,
    sidebyside align=top,
    lefthand width=0.55\textwidth,
    top=4pt, bottom=4pt, left=0pt, right=0pt,
    middle=16pt,
    underlay={
        \draw[dashed, black!80, line width=0.8pt, dash pattern=on 4pt off 3pt]
            ([xshift=0.55\textwidth+8pt]frame.north west) --
            ([xshift=0.55\textwidth+8pt]frame.south west);
    }
]
\textbf{Q.118} \quad {\footnotesize\textsc{[MCQ $\cdot$ 2 Mark(s) $\cdot$ Neg: 0.66]}} \hfill {\footnotesize\textbf{GATE 2024 CS1} \quad \tiny\texttt{2024\_CS1\_CS\_Q43}}
\vspace{2pt}\hrule height 0.3pt\vspace{4pt}

Consider a binary min-heap containing 105 distinct elements. Let $k$ be the index (in the underlying array) of the maximum element stored in the heap. The number of possible values of $k$ is

\par\vspace{6pt}\noindent
\begin{tabular}{@{}p{0.48\linewidth}@{\hspace{0.04\linewidth}}p{0.48\linewidth}@{}}
$\square$\ \textbf{(A)} 53 & $\square$\ \textbf{(B)} 52 \\
\addlinespace[3pt]
$\square$\ \textbf{(C)} 27 & $\square$\ \textbf{(D)} 1
\end{tabular}

\tcblower

\textbf{Official Key:} \textbf{\large A} \hfill \textit{\small [Medium]}
\vspace{2pt}\hrule height 0.3pt\vspace{3pt}
\textbf{Topic:} Programming and Data Structures \\
\textbf{Subtopic:} \textit{Trees, binary search trees, binary heaps}
\vspace{2pt}\hrule height 0.3pt\vspace{3pt}
\textbf{Guide / Hint:} \\
{\small In a binary min-heap, every parent is strictly smaller than its children, so the maximum element cannot be an internal node and must reside at one of the leaf nodes. In a complete binary tree of $n = 105$ nodes, the number of leaf nodes is $\lceil n/2 \rceil = \lceil 105/2 \rceil = 53$. Hence, there are 53 possible index locations for the maximum element.}

\end{tcolorbox}
\vspace{6pt}\hrule height 0.35pt\vspace{6pt}
\noindent
\begin{tcolorbox}[
    blanker,
    sidebyside,
    sidebyside align=top,
    lefthand width=0.55\textwidth,
    top=4pt, bottom=4pt, left=0pt, right=0pt,
    middle=16pt,
    underlay={
        \draw[dashed, black!80, line width=0.8pt, dash pattern=on 4pt off 3pt]
            ([xshift=0.55\textwidth+8pt]frame.north west) --
            ([xshift=0.55\textwidth+8pt]frame.south west);
    }
]
\textbf{Q.119} \quad {\footnotesize\textsc{[MCQ $\cdot$ 1 Mark(s) $\cdot$ Neg: 0.33]}} \hfill {\footnotesize\textbf{GATE 2024 CS2} \quad \tiny\texttt{2024\_CS2\_CS\_Q13}}
\vspace{2pt}\hrule height 0.3pt\vspace{4pt}

Consider the following C program. Assume parameters to a function are evaluated from right to left.
\begin{verbatim}
#include <stdio.h>
int g(int p) { printf("%d", p); return p; }
int h(int q) { printf("%d", q); return q; }
void f(int x, int y) {
    g(x);
    h(y);
}
int main() {
    f(g(10), h(20));
}
\end{verbatim}
Which one of the following is the output of the program?

\par\vspace{6pt}\noindent
\begin{tabular}{@{}p{0.48\linewidth}@{\hspace{0.04\linewidth}}p{0.48\linewidth}@{}}
$\square$\ \textbf{(A)} 20101020 & $\square$\ \textbf{(B)} 10202010 \\
\addlinespace[3pt]
$\square$\ \textbf{(C)} 20102010 & $\square$\ \textbf{(D)} 10201020
\end{tabular}

\tcblower

\textbf{Official Key:} \textbf{\large A} \hfill \textit{\small [Medium]}
\vspace{2pt}\hrule height 0.3pt\vspace{3pt}
\textbf{Topic:} Programming and Data Structures \\
\textbf{Subtopic:} \textit{Programming in C}
\vspace{2pt}\hrule height 0.3pt\vspace{3pt}
\textbf{Guide / Hint:} \\
{\small In `main()`, `f(g(10), h(20))` evaluates parameters from right to left: `h(20)` executes first, printing `20` and returning 20 ($y = 20$), then `g(10)` executes, printing `10` and returning 10 ($x = 10$). Inside `f`, `g(x)` executes printing `10`, followed by `h(y)` printing `20`. The complete printed sequence is `20101020`.}

\end{tcolorbox}
\vspace{6pt}\hrule height 0.35pt\vspace{6pt}
\noindent
\begin{tcolorbox}[
    blanker,
    sidebyside,
    sidebyside align=top,
    lefthand width=0.55\textwidth,
    top=4pt, bottom=4pt, left=0pt, right=0pt,
    middle=16pt,
    underlay={
        \draw[dashed, black!80, line width=0.8pt, dash pattern=on 4pt off 3pt]
            ([xshift=0.55\textwidth+8pt]frame.north west) --
            ([xshift=0.55\textwidth+8pt]frame.south west);
    }
]
\textbf{Q.120} \quad {\footnotesize\textsc{[MCQ $\cdot$ 2 Mark(s) $\cdot$ Neg: 0.66]}} \hfill {\footnotesize\textbf{GATE 2024 CS2} \quad \tiny\texttt{2024\_CS2\_CS\_Q36}}
\vspace{2pt}\hrule height 0.3pt\vspace{4pt}

What is the output of the following C program?
\begin{verbatim}
#include <stdio.h>
int main() {
    double a[2] = {20.0, 25.0}, *p, *q;
    p = a;
    q = p + 1;
    printf("%d,%d", (int)(q - p), (int)(*q - *p));
    return 0;
}
\end{verbatim}

\par\vspace{6pt}\noindent
\begin{tabular}{@{}p{0.48\linewidth}@{\hspace{0.04\linewidth}}p{0.48\linewidth}@{}}
$\square$\ \textbf{(A)} 4,8 & $\square$\ \textbf{(B)} 1,5 \\
\addlinespace[3pt]
$\square$\ \textbf{(C)} 8,5 & $\square$\ \textbf{(D)} 1,8
\end{tabular}

\tcblower

\textbf{Official Key:} \textbf{\large B} \hfill \textit{\small [Easy]}
\vspace{2pt}\hrule height 0.3pt\vspace{3pt}
\textbf{Topic:} Programming and Data Structures \\
\textbf{Subtopic:} \textit{Programming in C}
\vspace{2pt}\hrule height 0.3pt\vspace{3pt}
\textbf{Guide / Hint:} \\
{\small In pointer arithmetic, pointer subtraction `q - p` yields the number of elements between pointers, so $(q - p) = 1$. The expression `*q - *p` calculates the difference between pointed double values: $25.0 - 20.0 = 5.0$, cast to integer is 5. Thus, `1,5` is printed.}

\end{tcolorbox}
\vspace{6pt}\hrule height 0.35pt\vspace{6pt}
\noindent
\begin{tcolorbox}[
    blanker,
    sidebyside,
    sidebyside align=top,
    lefthand width=0.55\textwidth,
    top=4pt, bottom=4pt, left=0pt, right=0pt,
    middle=16pt,
    underlay={
        \draw[dashed, black!80, line width=0.8pt, dash pattern=on 4pt off 3pt]
            ([xshift=0.55\textwidth+8pt]frame.north west) --
            ([xshift=0.55\textwidth+8pt]frame.south west);
    }
]
\textbf{Q.121} \quad {\footnotesize\textsc{[MCQ $\cdot$ 2 Mark(s) $\cdot$ Neg: 0.66]}} \hfill {\footnotesize\textbf{GATE 2024 CS2} \quad \tiny\texttt{2024\_CS2\_CS\_Q39}}
\vspace{2pt}\hrule height 0.3pt\vspace{4pt}

You are given a set $V$ of distinct integers. A binary search tree $T$ is created by inserting all elements of $V$ one by one, starting with an empty tree. The tree $T$ follows the convention that, at each node, all values stored in the left subtree of the node are smaller than the value stored at the node. Which of the following is TRUE?

\par\vspace{6pt}\noindent
\begin{tabular}{@{}p{0.98\linewidth}@{}}
$\square$\ \textbf{(A)} Inorder traversal of $T$ can be determined from $V$ \\
\addlinespace[3pt]
$\square$\ \textbf{(B)} Root node of $T$ can be determined from $V$ \\
\addlinespace[3pt]
$\square$\ \textbf{(C)} Preorder traversal of $T$ can be determined from $V$ \\
\addlinespace[3pt]
$\square$\ \textbf{(D)} Postorder traversal of $T$ can be determined from $V$
\end{tabular}

\tcblower

\textbf{Official Key:} \textbf{\large A} \hfill \textit{\small [Easy]}
\vspace{2pt}\hrule height 0.3pt\vspace{3pt}
\textbf{Topic:} Programming and Data Structures \\
\textbf{Subtopic:} \textit{Trees, binary search trees, binary heaps}
\vspace{2pt}\hrule height 0.3pt\vspace{3pt}
\textbf{Guide / Hint:} \\
{\small For any binary search tree, the inorder traversal sequence always visits elements in strictly increasing sorted order regardless of insertion order or tree shape. Since the set of distinct elements $V$ is known, sorting $V$ directly determines the unique inorder traversal.}

\end{tcolorbox}
\vspace{6pt}\hrule height 0.35pt\vspace{6pt}
\noindent
\begin{tcolorbox}[
    blanker,
    sidebyside,
    sidebyside align=top,
    lefthand width=0.55\textwidth,
    top=4pt, bottom=4pt, left=0pt, right=0pt,
    middle=16pt,
    underlay={
        \draw[dashed, black!80, line width=0.8pt, dash pattern=on 4pt off 3pt]
            ([xshift=0.55\textwidth+8pt]frame.north west) --
            ([xshift=0.55\textwidth+8pt]frame.south west);
    }
]
\textbf{Q.122} \quad {\footnotesize\textsc{[MCQ $\cdot$ 2 Mark(s) $\cdot$ Neg: 0.66]}} \hfill {\footnotesize\textbf{GATE 2024 CS2} \quad \tiny\texttt{2024\_CS2\_CS\_Q42}}
\vspace{2pt}\hrule height 0.3pt\vspace{4pt}

Consider an array X that contains n positive integers. A subarray of X is defined to be a sequence of array locations with consecutive indices.
The C code snippet given below has been written to compute the length of the longest subarray of X that contains at most two distinct integers. The code has two missing expressions labelled $(P)$ and $(Q)$:
\begin{verbatim}
int first=0, second=0, len1=0, len2=0, maxlen=0;
for (int i=0; i < n; i++) {
    if (X[i] == first) {
        len2++; len1++;
    } else if (X[i] == second) {
        len2++;
        len1 = (P);
        second = first;
    } else {
        len2 = (Q);
        len1 = 1; second = first;
    }
    if (len2 > maxlen) maxlen = len2;
    first = X[i];
}
\end{verbatim}
Which one of the following options gives the CORRECT missing expressions?

\par\vspace{6pt}\noindent
\begin{tabular}{@{}p{0.48\linewidth}@{\hspace{0.04\linewidth}}p{0.48\linewidth}@{}}
$\square$\ \textbf{(A)} (P) len1+1 (Q) len2+1 & $\square$\ \textbf{(B)} (P) 1 (Q) len1+1 \\
\addlinespace[3pt]
$\square$\ \textbf{(C)} (P) 1 (Q) len2+1 & $\square$\ \textbf{(D)} (P) len2+1 (Q) len1+1
\end{tabular}

\tcblower

\textbf{Official Key:} \textbf{\large B} \hfill \textit{\small [Hard]}
\vspace{2pt}\hrule height 0.3pt\vspace{3pt}
\textbf{Topic:} Programming and Data Structures \\
\textbf{Subtopic:} \textit{Programming in C}
\vspace{2pt}\hrule height 0.3pt\vspace{3pt}
\textbf{Guide / Hint:} \\
{\small When $X[i] == second$, the current element matches the older distinct value, switching identity with `first`. Because $X[i]$ has just appeared once consecutively, `len1` resets to 1 (expression $P = 1$). When encountering a third distinct value (else branch), the previous window keeping at most two distinct values must keep only the contiguous run of `first` plus current element, so $Q = len1 + 1$.}

\end{tcolorbox}
\vspace{6pt}\hrule height 0.35pt\vspace{6pt}
\noindent
\begin{tcolorbox}[
    blanker,
    sidebyside,
    sidebyside align=top,
    lefthand width=0.55\textwidth,
    top=4pt, bottom=4pt, left=0pt, right=0pt,
    middle=16pt,
    underlay={
        \draw[dashed, black!80, line width=0.8pt, dash pattern=on 4pt off 3pt]
            ([xshift=0.55\textwidth+8pt]frame.north west) --
            ([xshift=0.55\textwidth+8pt]frame.south west);
    }
]
\textbf{Q.123} \quad {\footnotesize\textsc{[MCQ $\cdot$ 1 Mark(s) $\cdot$ Neg: 0.33]}} \hfill {\footnotesize\textbf{GATE 2025 CS2} \quad \tiny\texttt{2025\_CS2\_CS\_Q13}}
\vspace{2pt}\hrule height 0.3pt\vspace{4pt}

Consider a binary tree $T$ in which every node has either zero or two children. Let $n > 0$ be the number of nodes in $T$. Which ONE of the following is the number of nodes in $T$ that have exactly two children?

\par\vspace{6pt}\noindent
\begin{tabular}{@{}p{0.48\linewidth}@{\hspace{0.04\linewidth}}p{0.48\linewidth}@{}}
$\square$\ \textbf{(A)} $\frac{n-2}{2}$ & $\square$\ \textbf{(B)} $\frac{n-1}{2}$ \\
\addlinespace[3pt]
$\square$\ \textbf{(C)} $\frac{n}{2}$ & $\square$\ \textbf{(D)} $\frac{n+1}{2}$
\end{tabular}

\tcblower

\textbf{Official Key:} \textbf{\large B} \hfill \textit{\small [Easy]}
\vspace{2pt}\hrule height 0.3pt\vspace{3pt}
\textbf{Topic:} Programming and Data Structures \\
\textbf{Subtopic:} \textit{Trees and strictly binary trees}
\vspace{2pt}\hrule height 0.3pt\vspace{3pt}
\textbf{Guide / Hint:} \\
{\small In any proper (strictly) binary tree with $L$ leaves and $I$ internal nodes (nodes with 2 children), $L = I + 1$. Since $n = L + I = 2I + 1$, the number of internal nodes is $I = \frac{n-1}{2}$.}

\end{tcolorbox}
\vspace{6pt}\hrule height 0.35pt\vspace{6pt}
\noindent
\begin{tcolorbox}[
    blanker,
    sidebyside,
    sidebyside align=top,
    lefthand width=0.55\textwidth,
    top=4pt, bottom=4pt, left=0pt, right=0pt,
    middle=16pt,
    underlay={
        \draw[dashed, black!80, line width=0.8pt, dash pattern=on 4pt off 3pt]
            ([xshift=0.55\textwidth+8pt]frame.north west) --
            ([xshift=0.55\textwidth+8pt]frame.south west);
    }
]
\textbf{Q.124} \quad {\footnotesize\textsc{[MCQ $\cdot$ 1 Mark(s) $\cdot$ Neg: 0.33]}} \hfill {\footnotesize\textbf{GATE 2025 CS2} \quad \tiny\texttt{2025\_CS2\_CS\_Q19}}
\vspace{2pt}\hrule height 0.3pt\vspace{4pt}

Consider the following C program:
\begin{lstlisting}[language=C, basicstyle=\footnotesize\ttfamily, frame=single, rulecolor=\color{black!30}, tabsize=2, breaklines=true]
#include <stdio.h>
void stringcopy(char *, char *);
int main() {
    char a[30] = "@#Hello World!";
    stringcopy(a, a + 2);
    printf("%s\n", a);
    return 0;
}
void stringcopy(char *s, char *t) {
    while(*t)
        *s++ = *t++;
}
\end{lstlisting}
Which ONE of the following will be the output of the program?

\par\vspace{6pt}\noindent
\begin{tabular}{@{}p{0.48\linewidth}@{\hspace{0.04\linewidth}}p{0.48\linewidth}@{}}
$\square$\ \textbf{(A)} @\#Hello World! & $\square$\ \textbf{(B)} Hello World! \\
\addlinespace[3pt]
$\square$\ \textbf{(C)} ello World! & $\square$\ \textbf{(D)} Hello World!d!
\end{tabular}

\tcblower

\textbf{Official Key:} \textbf{\large D} \hfill \textit{\small [Medium]}
\vspace{2pt}\hrule height 0.3pt\vspace{3pt}
\textbf{Topic:} Programming and Data Structures \\
\textbf{Subtopic:} \textit{Programming in C: pointers and strings}
\vspace{2pt}\hrule height 0.3pt\vspace{3pt}
\textbf{Guide / Hint:} \\
{\small `stringcopy(a, a + 2)` copies string `"Hello World!"` into the start of array `a`, overwriting the first 12 characters. The copy stops when `*t` hits null terminator without copying the null terminator over `d!`, leaving trailing `d!\0`, resulting in `"Hello World!d!"`.}

\end{tcolorbox}
\vspace{6pt}\hrule height 0.35pt\vspace{6pt}
\noindent
\begin{tcolorbox}[
    blanker,
    sidebyside,
    sidebyside align=top,
    lefthand width=0.55\textwidth,
    top=4pt, bottom=4pt, left=0pt, right=0pt,
    middle=16pt,
    underlay={
        \draw[dashed, black!80, line width=0.8pt, dash pattern=on 4pt off 3pt]
            ([xshift=0.55\textwidth+8pt]frame.north west) --
            ([xshift=0.55\textwidth+8pt]frame.south west);
    }
]
\textbf{Q.125} \quad {\footnotesize\textsc{[MCQ $\cdot$ 2 Mark(s) $\cdot$ Neg: 0.66]}} \hfill {\footnotesize\textbf{GATE 2025 CS2} \quad \tiny\texttt{2025\_CS2\_CS\_Q38}}
\vspace{2pt}\hrule height 0.3pt\vspace{4pt}

A meld operation on two instances of a data structure combines them into one single instance of the same data structure. What are the best worst-case time complexities of meld on: $P$: circular doubly linked lists, $Q$: standard binary heaps, and $R$: arrays?

\par\vspace{6pt}\noindent
\begin{tabular}{@{}p{0.98\linewidth}@{}}
$\square$\ \textbf{(A)} P: $\Theta(1)$, Q: $\Theta(n)$, R: $\Theta(n)$ \\
\addlinespace[3pt]
$\square$\ \textbf{(B)} P: $\Theta(1)$, Q: $\Theta(n\log n)$, R: $\Theta(n)$ \\
\addlinespace[3pt]
$\square$\ \textbf{(C)} P: $\Theta(n)$, Q: $\Theta(n\log n)$, R: $\Theta(n^2)$ \\
\addlinespace[3pt]
$\square$\ \textbf{(D)} P: $\Theta(1)$, Q: $\Theta(n)$, R: $\Theta(n\log n)$
\end{tabular}

\tcblower

\textbf{Official Key:} \textbf{\large A} \hfill \textit{\small [Hard]}
\vspace{2pt}\hrule height 0.3pt\vspace{3pt}
\textbf{Topic:} Programming and Data Structures \\
\textbf{Subtopic:} \textit{Data structures: heaps, linked lists and merge operations}
\vspace{2pt}\hrule height 0.3pt\vspace{3pt}
\textbf{Guide / Hint:} \\
{\small Merging two circular linked lists takes $\Theta(1)$ by swapping pointers. Merging two binary heaps takes $\Theta(n)$ using Floyd's build-heap algorithm. Hence, time complexities are $P: \Theta(1), Q: \Theta(n), R: \Theta(n)$.}

\end{tcolorbox}
\vspace{6pt}\hrule height 0.35pt\vspace{6pt}
\vspace{8pt}
\subsection*{Algorithms}
\vspace{-4pt}
\noindent
\begin{tcolorbox}[
    blanker,
    sidebyside,
    sidebyside align=top,
    lefthand width=0.55\textwidth,
    top=4pt, bottom=4pt, left=0pt, right=0pt,
    middle=16pt,
    underlay={
        \draw[dashed, black!80, line width=0.8pt, dash pattern=on 4pt off 3pt]
            ([xshift=0.55\textwidth+8pt]frame.north west) --
            ([xshift=0.55\textwidth+8pt]frame.south west);
    }
]
\textbf{Q.126} \quad {\footnotesize\textsc{[MCQ $\cdot$ 2 Mark(s) $\cdot$ Neg: 0.66]}} \hfill {\footnotesize\textbf{GATE 2018 CS} \quad \tiny\texttt{2018\_CS\_GA\_Q30}}
\vspace{2pt}\hrule height 0.3pt\vspace{4pt}

Let $G$ be a simple undirected graph. Let $T_D$ be a depth first search tree of $G$. Let $T_B$ be a breadth first search tree of $G$. Consider the following statements:
\begin{enumerate}[label=\Roman*.]
    \item For every edge $(u, v)$ of $G$, if $u$ is at depth $i$ and $v$ is at depth $j$ in $T_D$, then $|i - j| \ge 1$.
    \item For every edge $(u, v)$ of $G$, if $u$ is at depth $i$ and $v$ is at depth $j$ in $T_B$, then $|i - j| = 1$.
\end{enumerate}
Which of the statements above is/are necessarily TRUE?

\par\vspace{6pt}\noindent
\begin{tabular}{@{}p{0.48\linewidth}@{\hspace{0.04\linewidth}}p{0.48\linewidth}@{}}
$\square$\ \textbf{(A)} I only & $\square$\ \textbf{(B)} II only \\
\addlinespace[3pt]
$\square$\ \textbf{(C)} Both I and II & $\square$\ \textbf{(D)} Neither I nor II
\end{tabular}

\tcblower

\textbf{Official Key:} \textbf{\large A} \hfill \textit{\small [Medium]}
\vspace{2pt}\hrule height 0.3pt\vspace{3pt}
\textbf{Topic:} Algorithms \\
\textbf{Subtopic:} \textit{Graph traversals, minimum spanning trees, shortest paths}
\vspace{2pt}\hrule height 0.3pt\vspace{3pt}
\textbf{Guide / Hint:} \\
{\small In a DFS tree of an undirected graph, every non-tree edge is a back edge connecting an ancestor and a descendant, so forward and cross edges do not exist; thus statement I is necessarily true. In a BFS tree, cross edges can connect vertices at the same level ($|i - j| = 0$), so statement II is false.}

\end{tcolorbox}
\vspace{6pt}\hrule height 0.35pt\vspace{6pt}
\noindent
\begin{tcolorbox}[
    blanker,
    sidebyside,
    sidebyside align=top,
    lefthand width=0.55\textwidth,
    top=4pt, bottom=4pt, left=0pt, right=0pt,
    middle=16pt,
    underlay={
        \draw[dashed, black!80, line width=0.8pt, dash pattern=on 4pt off 3pt]
            ([xshift=0.55\textwidth+8pt]frame.north west) --
            ([xshift=0.55\textwidth+8pt]frame.south west);
    }
]
\textbf{Q.127} \quad {\footnotesize\textsc{[MCQ $\cdot$ 2 Mark(s) $\cdot$ Neg: 0.66]}} \hfill {\footnotesize\textbf{GATE 2018 CS} \quad \tiny\texttt{2018\_CS\_GA\_Q31}}
\vspace{2pt}\hrule height 0.3pt\vspace{4pt}

Assume that multiplying a matrix $G_1$ of dimension $p \times q$ with another matrix $G_2$ of dimension $q \times r$ requires $pqr$ scalar multiplications. Computing the product of five matrices $F_1, F_2, F_3, F_4$ and $F_5$ of dimensions $2 \times 25$, $25 \times 3$, $3 \times 16$, $16 \times 1$ and $1 \times 10$, respectively, using the minimum number of scalar multiplications requires first computing

\par\vspace{6pt}\noindent
\begin{tabular}{@{}p{0.48\linewidth}@{\hspace{0.04\linewidth}}p{0.48\linewidth}@{}}
$\square$\ \textbf{(A)} $F_1 F_2$ and $F_3 F_4$ & $\square$\ \textbf{(B)} $F_2 F_3$ and $F_3 F_4$ \\
\addlinespace[3pt]
$\square$\ \textbf{(C)} $F_2 F_3$ and $F_4 F_5$ & $\square$\ \textbf{(D)} $F_1 F_2$ and $F_4 F_5$
\end{tabular}

\tcblower

\textbf{Official Key:} \textbf{\large C} \hfill \textit{\small [Medium]}
\vspace{2pt}\hrule height 0.3pt\vspace{3pt}
\textbf{Topic:} Algorithms \\
\textbf{Subtopic:} \textit{Algorithm design techniques: greedy, dynamic programming and divide-and-conquer}
\vspace{2pt}\hrule height 0.3pt\vspace{3pt}
\textbf{Guide / Hint:} \\
{\small Using dynamic programming for matrix chain multiplication with dimensions $d = \langle 2, 25, 3, 16, 1, 10 \rangle$: calculating minimum multiplication costs reveals optimal parenthesization as $((F_1(F_2 F_3))(F_4 F_5))$, where the intermediate pairs computed first are $F_2 F_3$ and $F_4 F_5$.}

\end{tcolorbox}
\vspace{6pt}\hrule height 0.35pt\vspace{6pt}
\noindent
\begin{tcolorbox}[
    blanker,
    sidebyside,
    sidebyside align=top,
    lefthand width=0.55\textwidth,
    top=4pt, bottom=4pt, left=0pt, right=0pt,
    middle=16pt,
    underlay={
        \draw[dashed, black!80, line width=0.8pt, dash pattern=on 4pt off 3pt]
            ([xshift=0.55\textwidth+8pt]frame.north west) --
            ([xshift=0.55\textwidth+8pt]frame.south west);
    }
]
\textbf{Q.128} \quad {\footnotesize\textsc{[MCQ $\cdot$ 2 Mark(s) $\cdot$ Neg: 0.66]}} \hfill {\footnotesize\textbf{GATE 2018 CS} \quad \tiny\texttt{2018\_CS\_GA\_Q32}}
\vspace{2pt}\hrule height 0.3pt\vspace{4pt}

Consider the following C code. Assume that \texttt{unsigned long int} type length is 64 bits.
\begin{lstlisting}[language=C]
unsigned long int fun(unsigned long int n) {
    unsigned long int i, j = 0, sum = 0;
    while (n > 1) {
        j++;
        n /= 2;
    }
    for (i = 0; j > 1; i++) {
        sum += (j % 2);
        j /= 2;
    }
    return sum;
}
\end{lstlisting}
The value returned when we call \texttt{fun(pow(2, 40))} is

\par\vspace{6pt}\noindent
\begin{tabular}{@{}p{0.48\linewidth}@{\hspace{0.04\linewidth}}p{0.48\linewidth}@{}}
$\square$\ \textbf{(A)} 4 & $\square$\ \textbf{(B)} 5 \\
\addlinespace[3pt]
$\square$\ \textbf{(C)} 6 & $\square$\ \textbf{(D)} 40
\end{tabular}

\tcblower

\textbf{Official Key:} \textbf{\large B} \hfill \textit{\small [Medium]}
\vspace{2pt}\hrule height 0.3pt\vspace{3pt}
\textbf{Topic:} Algorithms \\
\textbf{Subtopic:} \textit{Asymptotic worst case time and space complexity}
\vspace{2pt}\hrule height 0.3pt\vspace{3pt}
\textbf{Guide / Hint:} \\
{\small For $n = 2^{40}$, the first loop `n /= 2` executes until $n = 1$, which takes $\log_2(2^{40}) = 40$ steps, giving $j = 40$. The second loop starts with $j = 40$ and halves $j$ until $j \le 1$: values taken by $j$ are $40 \to 20 \to 10 \to 5 \to 2 \to 1$, incrementing `sum` exactly 5 times.}

\end{tcolorbox}
\vspace{6pt}\hrule height 0.35pt\vspace{6pt}
\noindent
\begin{tcolorbox}[
    blanker,
    sidebyside,
    sidebyside align=top,
    lefthand width=0.55\textwidth,
    top=4pt, bottom=4pt, left=0pt, right=0pt,
    middle=16pt,
    underlay={
        \draw[dashed, black!80, line width=0.8pt, dash pattern=on 4pt off 3pt]
            ([xshift=0.55\textwidth+8pt]frame.north west) --
            ([xshift=0.55\textwidth+8pt]frame.south west);
    }
]
\textbf{Q.129} \quad {\footnotesize\textsc{[MCQ $\cdot$ 1 Mark(s) $\cdot$ Neg: 0.33]}} \hfill {\footnotesize\textbf{GATE 2022 CS} \quad \tiny\texttt{2022\_CS\_CS\_Q11}}
\vspace{2pt}\hrule height 0.3pt\vspace{4pt}

Which one of the following statements is TRUE for all positive functions $f(n)$?

\par\vspace{6pt}\noindent
\begin{tabular}{@{}p{0.98\linewidth}@{}}
$\square$\ \textbf{(A)} $f(n^2) = \Theta((f(n))^2)$, when $f(n)$ is a polynomial \\
\addlinespace[3pt]
$\square$\ \textbf{(B)} $f(n^2) = o((f(n))^2)$ \\
\addlinespace[3pt]
$\square$\ \textbf{(C)} $f(n^2) = \mathcal{O}((f(n))^2)$, when $f(n)$ is an exponential function \\
\addlinespace[3pt]
$\square$\ \textbf{(D)} $f(n^2) = \Omega((f(n))^2)$
\end{tabular}

\tcblower

\textbf{Official Key:} \textbf{\large A} \hfill \textit{\small [Medium]}
\vspace{2pt}\hrule height 0.3pt\vspace{3pt}
\textbf{Topic:} Algorithms \\
\textbf{Subtopic:} \textit{Asymptotic worst case time and space complexity}
\vspace{2pt}\hrule height 0.3pt\vspace{3pt}
\textbf{Guide / Hint:} \\
{\small If $f(n)$ is a polynomial of degree $k$ ($f(n) = \Theta(n^k)$), then $f(n^2) = \Theta((n^2)^k) = \Theta(n^{2k})$, and $(f(n))^2 = (\Theta(n^k))^2 = \Theta(n^{2k})$. Since both functions grow at identical polynomial rates, $f(n^2) = \Theta((f(n))^2)$ holds.}

\end{tcolorbox}
\vspace{6pt}\hrule height 0.35pt\vspace{6pt}
\noindent
\begin{tcolorbox}[
    blanker,
    sidebyside,
    sidebyside align=top,
    lefthand width=0.55\textwidth,
    top=4pt, bottom=4pt, left=0pt, right=0pt,
    middle=16pt,
    underlay={
        \draw[dashed, black!80, line width=0.8pt, dash pattern=on 4pt off 3pt]
            ([xshift=0.55\textwidth+8pt]frame.north west) --
            ([xshift=0.55\textwidth+8pt]frame.south west);
    }
]
\textbf{Q.130} \quad {\footnotesize\textsc{[MCQ $\cdot$ 1 Mark(s) $\cdot$ Neg: 0.33]}} \hfill {\footnotesize\textbf{GATE 2022 CS} \quad \tiny\texttt{2022\_CS\_CS\_Q16}}
\vspace{2pt}\hrule height 0.3pt\vspace{4pt}

Suppose we are given $n$ keys, $m$ hash table slots, and two simple uniform hash functions $h_1$ and $h_2$. Further suppose our hashing scheme uses $h_1$ for the odd keys and $h_2$ for the even keys. What is the expected number of keys in a slot?

\par\vspace{6pt}\noindent
\begin{tabular}{@{}p{0.48\linewidth}@{\hspace{0.04\linewidth}}p{0.48\linewidth}@{}}
$\square$\ \textbf{(A)} $m/n$ & $\square$\ \textbf{(B)} $n/m$ \\
\addlinespace[3pt]
$\square$\ \textbf{(C)} $2n/m$ & $\square$\ \textbf{(D)} $n / (2m)$
\end{tabular}

\tcblower

\textbf{Official Key:} \textbf{\large B} \hfill \textit{\small [Easy]}
\vspace{2pt}\hrule height 0.3pt\vspace{3pt}
\textbf{Topic:} Algorithms \\
\textbf{Subtopic:} \textit{Searching, sorting, hashing}
\vspace{2pt}\hrule height 0.3pt\vspace{3pt}
\textbf{Guide / Hint:} \\
{\small Let $n_1$ be the number of odd keys and $n_2$ be the number of even keys ($n_1 + n_2 = n$). Under simple uniform hashing, each key independently hashes to any slot with probability $1/m$. By linearity of expectation, the expected number of keys in any slot is $\frac{n_1}{m} + \frac{n_2}{m} = \frac{n}{m}$.}

\end{tcolorbox}
\vspace{6pt}\hrule height 0.35pt\vspace{6pt}
\noindent
\begin{tcolorbox}[
    blanker,
    sidebyside,
    sidebyside align=top,
    lefthand width=0.55\textwidth,
    top=4pt, bottom=4pt, left=0pt, right=0pt,
    middle=16pt,
    underlay={
        \draw[dashed, black!80, line width=0.8pt, dash pattern=on 4pt off 3pt]
            ([xshift=0.55\textwidth+8pt]frame.north west) --
            ([xshift=0.55\textwidth+8pt]frame.south west);
    }
]
\textbf{Q.131} \quad {\footnotesize\textsc{[MCQ $\cdot$ 2 Mark(s) $\cdot$ Neg: 0.66]}} \hfill {\footnotesize\textbf{GATE 2022 CS} \quad \tiny\texttt{2022\_CS\_CS\_Q37}}
\vspace{2pt}\hrule height 0.3pt\vspace{4pt}

Consider a simple undirected unweighted graph with at least three vertices. If $A$ is the adjacency matrix of the graph, then the number of 3-cycles in the graph is given by the trace of

\par\vspace{6pt}\noindent
\begin{tabular}{@{}p{0.48\linewidth}@{\hspace{0.04\linewidth}}p{0.48\linewidth}@{}}
$\square$\ \textbf{(A)} $A^3$ & $\square$\ \textbf{(B)} $A^3$ divided by 2 \\
\addlinespace[3pt]
$\square$\ \textbf{(C)} $A^3$ divided by 3 & $\square$\ \textbf{(D)} $A^3$ divided by 6
\end{tabular}

\tcblower

\textbf{Official Key:} \textbf{\large D} \hfill \textit{\small [Medium]}
\vspace{2pt}\hrule height 0.3pt\vspace{3pt}
\textbf{Topic:} Algorithms \\
\textbf{Subtopic:} \textit{Graph traversals, minimum spanning trees, shortest paths}
\vspace{2pt}\hrule height 0.3pt\vspace{3pt}
\textbf{Guide / Hint:} \\
{\small The diagonal entry $(A^3)_{ii}$ counts the number of closed walks of length 3 starting and ending at vertex $i$, so $\text{tr}(A^3) = \sum (A^3)_{ii}$ counts all closed walks of length 3 in $G$. In a simple undirected graph, each triangle contains 3 vertices and can be traversed in 2 directions, counting each 3-cycle $3 \times 2 = 6$ times. Thus, the number of 3-cycles is $\frac{\text{tr}(A^3)}{6}$.}

\end{tcolorbox}
\vspace{6pt}\hrule height 0.35pt\vspace{6pt}
\noindent
\begin{tcolorbox}[
    blanker,
    sidebyside,
    sidebyside align=top,
    lefthand width=0.55\textwidth,
    top=4pt, bottom=4pt, left=0pt, right=0pt,
    middle=16pt,
    underlay={
        \draw[dashed, black!80, line width=0.8pt, dash pattern=on 4pt off 3pt]
            ([xshift=0.55\textwidth+8pt]frame.north west) --
            ([xshift=0.55\textwidth+8pt]frame.south west);
    }
]
\textbf{Q.132} \quad {\footnotesize\textsc{[MCQ $\cdot$ 1 Mark(s) $\cdot$ Neg: 0.33]}} \hfill {\footnotesize\textbf{GATE 2023 CS} \quad \tiny\texttt{2023\_CS\_CS\_Q20}}
\vspace{2pt}\hrule height 0.3pt\vspace{4pt}

An algorithm has to store several keys generated by an adversary in a hash table. The adversary is malicious who tries to maximize the number of collisions. Let $k$ be the number of keys, $m$ be the number of slots in the hash table, and $k > m$.
Which one of the following is the best hashing strategy to counteract the adversary?

\par\vspace{6pt}\noindent
\begin{tabular}{@{}p{0.98\linewidth}@{}}
$\square$\ \textbf{(A)} Division method, i.e., use the hash function $h(k) = k \bmod m$. \\
\addlinespace[3pt]
$\square$\ \textbf{(B)} Multiplication method, i.e., use the hash function $h(k) = \lfloor m(kA - \lfloor kA \rfloor)\rfloor$, where $A$ is a carefully chosen constant. \\
\addlinespace[3pt]
$\square$\ \textbf{(C)} Universal hashing method. \\
\addlinespace[3pt]
$\square$\ \textbf{(D)} If $k$ is a prime number, use Division method. Otherwise, use Multiplication method.
\end{tabular}

\tcblower

\textbf{Official Key:} \textbf{\large C} \hfill \textit{\small [Easy]}
\vspace{2pt}\hrule height 0.3pt\vspace{3pt}
\textbf{Topic:} Algorithms \\
\textbf{Subtopic:} \textit{Searching, sorting, hashing}
\vspace{2pt}\hrule height 0.3pt\vspace{3pt}
\textbf{Guide / Hint:} \\
{\small In universal hashing, the hash function is chosen uniformly at random from a carefully designed universal collection of hash functions independently of the keys. Because the adversary cannot predict which function will be selected at runtime, any fixed choice of malicious keys cannot guarantee worst-case collision behavior.}

\end{tcolorbox}
\vspace{6pt}\hrule height 0.35pt\vspace{6pt}
\noindent
\begin{tcolorbox}[
    blanker,
    sidebyside,
    sidebyside align=top,
    lefthand width=0.55\textwidth,
    top=4pt, bottom=4pt, left=0pt, right=0pt,
    middle=16pt,
    underlay={
        \draw[dashed, black!80, line width=0.8pt, dash pattern=on 4pt off 3pt]
            ([xshift=0.55\textwidth+8pt]frame.north west) --
            ([xshift=0.55\textwidth+8pt]frame.south west);
    }
]
\textbf{Q.133} \quad {\footnotesize\textsc{[MCQ $\cdot$ 1 Mark(s) $\cdot$ Neg: 0.33]}} \hfill {\footnotesize\textbf{GATE 2024 CS1} \quad \tiny\texttt{2024\_CS1\_CS\_Q17}}
\vspace{2pt}\hrule height 0.3pt\vspace{4pt}

Given an integer array of size $N$, we want to check if the array is sorted (in either ascending or descending order). An algorithm solves this problem by making a single pass through the array and comparing each element of the array only with its adjacent elements. The worst-case time complexity of this algorithm is

\par\vspace{6pt}\noindent
\begin{tabular}{@{}p{0.98\linewidth}@{}}
$\square$\ \textbf{(A)} both $\mathcal{O}(N)$ and $\Omega(N)$ \\
\addlinespace[3pt]
$\square$\ \textbf{(B)} $\mathcal{O}(N)$ but not $\Omega(N)$ \\
\addlinespace[3pt]
$\square$\ \textbf{(C)} $\Omega(N)$ but not $\mathcal{O}(N)$ \\
\addlinespace[3pt]
$\square$\ \textbf{(D)} neither $\mathcal{O}(N)$ nor $\Omega(N)$
\end{tabular}

\tcblower

\textbf{Official Key:} \textbf{\large A} \hfill \textit{\small [Easy]}
\vspace{2pt}\hrule height 0.3pt\vspace{3pt}
\textbf{Topic:} Algorithms \\
\textbf{Subtopic:} \textit{Asymptotic worst case time and space complexity}
\vspace{2pt}\hrule height 0.3pt\vspace{3pt}
\textbf{Guide / Hint:} \\
{\small In the worst case (such as when the array is sorted or the disorder occurs at the very end), a single-pass algorithm that compares adjacent elements must traverse the entire array, performing $N-1$ comparisons. The worst-case runtime is asymptotically tight at $\Theta(N)$, meaning it is both $\mathcal{O}(N)$ and $\Omega(N)$.}

\end{tcolorbox}
\vspace{6pt}\hrule height 0.35pt\vspace{6pt}
\noindent
\begin{tcolorbox}[
    blanker,
    sidebyside,
    sidebyside align=top,
    lefthand width=0.55\textwidth,
    top=4pt, bottom=4pt, left=0pt, right=0pt,
    middle=16pt,
    underlay={
        \draw[dashed, black!80, line width=0.8pt, dash pattern=on 4pt off 3pt]
            ([xshift=0.55\textwidth+8pt]frame.north west) --
            ([xshift=0.55\textwidth+8pt]frame.south west);
    }
]
\textbf{Q.134} \quad {\footnotesize\textsc{[MCQ $\cdot$ 2 Mark(s) $\cdot$ Neg: 0.66]}} \hfill {\footnotesize\textbf{GATE 2024 CS1} \quad \tiny\texttt{2024\_CS1\_CS\_Q42}}
\vspace{2pt}\hrule height 0.3pt\vspace{4pt}

Consider the following recurrence relation:
\[ T(n) = \begin{cases} \sqrt{n} T(\sqrt{n}) + n & \text{for } n > 1, \\ 1 & \text{for } n = 1. \end{cases} \]
Which one of the following options is CORRECT?

\par\vspace{6pt}\noindent
\begin{tabular}{@{}p{0.98\linewidth}@{}}
$\square$\ \textbf{(A)} $T(n) = \Theta(n \log \log n)$ \\
\addlinespace[3pt]
$\square$\ \textbf{(B)} $T(n) = \Theta(n \log n)$ \\
\addlinespace[3pt]
$\square$\ \textbf{(C)} $T(n) = \Theta(n^2 \log n)$ \\
\addlinespace[3pt]
$\square$\ \textbf{(D)} $T(n) = \Theta(n^2 \log \log n)$
\end{tabular}

\tcblower

\textbf{Official Key:} \textbf{\large A} \hfill \textit{\small [Medium]}
\vspace{2pt}\hrule height 0.3pt\vspace{3pt}
\textbf{Topic:} Algorithms \\
\textbf{Subtopic:} \textit{Asymptotic worst case time and space complexity}
\vspace{2pt}\hrule height 0.3pt\vspace{3pt}
\textbf{Guide / Hint:} \\
{\small Dividing $T(n) = \sqrt{n} T(\sqrt{n}) + n$ by $n$ gives $\frac{T(n)}{n} = \frac{T(\sqrt{n})}{\sqrt{n}} + 1$. Defining $S(n) = \frac{T(n)}{n}$ yields $S(n) = S(\sqrt{n}) + 1$. Substituting $n = 2^{2^k}$ gives $S(2^{2^k}) = S(2^{2^{k-1}}) + 1 = \Theta(k) = \Theta(\log \log n)$. Therefore, $T(n) = n S(n) = \Theta(n \log \log n)$.}

\end{tcolorbox}
\vspace{6pt}\hrule height 0.35pt\vspace{6pt}
\noindent
\begin{tcolorbox}[
    blanker,
    sidebyside,
    sidebyside align=top,
    lefthand width=0.55\textwidth,
    top=4pt, bottom=4pt, left=0pt, right=0pt,
    middle=16pt,
    underlay={
        \draw[dashed, black!80, line width=0.8pt, dash pattern=on 4pt off 3pt]
            ([xshift=0.55\textwidth+8pt]frame.north west) --
            ([xshift=0.55\textwidth+8pt]frame.south west);
    }
]
\textbf{Q.135} \quad {\footnotesize\textsc{[MCQ $\cdot$ 1 Mark(s) $\cdot$ Neg: 0.33]}} \hfill {\footnotesize\textbf{GATE 2025 CS1} \quad \tiny\texttt{2025\_CS1\_CS\_Q18}}
\vspace{2pt}\hrule height 0.3pt\vspace{4pt}

Let $G$ be any undirected graph with positive edge weights, and $T$ be a minimum spanning tree of $G$. For any two vertices, $u$ and $v$, let $d_1(u, v)$ and $d_2(u, v)$ be the shortest distances between $u$ and $v$ in $G$ and $T$, respectively. Which ONE of the options is CORRECT for all possible $G, T, u$ and $v$?

\par\vspace{6pt}\noindent
\begin{tabular}{@{}p{0.48\linewidth}@{\hspace{0.04\linewidth}}p{0.48\linewidth}@{}}
$\square$\ \textbf{(A)} $d_1(u, v) = d_2(u, v)$ & $\square$\ \textbf{(B)} $d_1(u, v) \le d_2(u, v)$ \\
\addlinespace[3pt]
$\square$\ \textbf{(C)} $d_1(u, v) \ge d_2(u, v)$ & $\square$\ \textbf{(D)} $d_1(u, v) \ne d_2(u, v)$
\end{tabular}

\tcblower

\textbf{Official Key:} \textbf{\large B} \hfill \textit{\small [Easy]}
\vspace{2pt}\hrule height 0.3pt\vspace{3pt}
\textbf{Topic:} Algorithms \\
\textbf{Subtopic:} \textit{Graph traversals, minimum spanning trees, shortest paths}
\vspace{2pt}\hrule height 0.3pt\vspace{3pt}
\textbf{Guide / Hint:} \\
{\small Since the minimum spanning tree $T$ is a subgraph of $G$ ($E(T) \subseteq E(G)$), every path between $u$ and $v$ in $T$ is also a valid path in $G$. Thus, the shortest path distance in $G$ cannot exceed that in $T$, guaranteeing $d_1(u, v) \le d_2(u, v)$.}

\end{tcolorbox}
\vspace{6pt}\hrule height 0.35pt\vspace{6pt}
\noindent
\begin{tcolorbox}[
    blanker,
    sidebyside,
    sidebyside align=top,
    lefthand width=0.55\textwidth,
    top=4pt, bottom=4pt, left=0pt, right=0pt,
    middle=16pt,
    underlay={
        \draw[dashed, black!80, line width=0.8pt, dash pattern=on 4pt off 3pt]
            ([xshift=0.55\textwidth+8pt]frame.north west) --
            ([xshift=0.55\textwidth+8pt]frame.south west);
    }
]
\textbf{Q.136} \quad {\footnotesize\textsc{[MCQ $\cdot$ 1 Mark(s) $\cdot$ Neg: 0.33]}} \hfill {\footnotesize\textbf{GATE 2025 CS1} \quad \tiny\texttt{2025\_CS1\_CS\_Q20}}
\vspace{2pt}\hrule height 0.3pt\vspace{4pt}

Consider the following recurrence relation:
\[ T(n) = 2T(n-1) + n2^n \quad \text{for } n > 0, \; T(0) = 1. \]
Which ONE of the following options is CORRECT?

\par\vspace{6pt}\noindent
\begin{tabular}{@{}p{0.48\linewidth}@{\hspace{0.04\linewidth}}p{0.48\linewidth}@{}}
$\square$\ \textbf{(A)} $T(n) = \Theta(n^2 2^n)$ & $\square$\ \textbf{(B)} $T(n) = \Theta(n 2^n)$ \\
\addlinespace[3pt]
$\square$\ \textbf{(C)} $T(n) = \Theta((\log n)^2 2^n)$ & $\square$\ \textbf{(D)} $T(n) = \Theta(4^n)$
\end{tabular}

\tcblower

\textbf{Official Key:} \textbf{\large A} \hfill \textit{\small [Medium]}
\vspace{2pt}\hrule height 0.3pt\vspace{3pt}
\textbf{Topic:} Algorithms \\
\textbf{Subtopic:} \textit{Asymptotic worst case time and space complexity}
\vspace{2pt}\hrule height 0.3pt\vspace{3pt}
\textbf{Guide / Hint:} \\
{\small Dividing both sides by $2^n$ yields $\frac{T(n)}{2^n} = \frac{T(n-1)}{2^{n-1}} + n$. Letting $S(n) = \frac{T(n)}{2^n}$, we have $S(n) = S(n-1) + n = S(0) + \sum_{i=1}^n i = 1 + \frac{n(n+1)}{2} = \Theta(n^2)$. Multiplying back gives $T(n) = \Theta(n^2 2^n)$.}

\end{tcolorbox}
\vspace{6pt}\hrule height 0.35pt\vspace{6pt}
\noindent
\begin{tcolorbox}[
    blanker,
    sidebyside,
    sidebyside align=top,
    lefthand width=0.55\textwidth,
    top=4pt, bottom=4pt, left=0pt, right=0pt,
    middle=16pt,
    underlay={
        \draw[dashed, black!80, line width=0.8pt, dash pattern=on 4pt off 3pt]
            ([xshift=0.55\textwidth+8pt]frame.north west) --
            ([xshift=0.55\textwidth+8pt]frame.south west);
    }
]
\textbf{Q.137} \quad {\footnotesize\textsc{[MCQ $\cdot$ 2 Mark(s) $\cdot$ Neg: 0.66]}} \hfill {\footnotesize\textbf{GATE 2025 CS1} \quad \tiny\texttt{2025\_CS1\_CS\_Q43}}
\vspace{2pt}\hrule height 0.3pt\vspace{4pt}

Let $G(V, E)$ be an undirected and unweighted graph with 100 vertices. Let $d(u, v)$ denote the number of edges in a shortest path between vertices $u$ and $v$ in $V$. Let the maximum value of $d(u, v), u, v \in V$ such that $u \neq v$, be 30. Let $T$ be any breadth-first-search tree of $G$. Which ONE of the given options is CORRECT for every such graph $G$?

\par\vspace{6pt}\noindent
\begin{tabular}{@{}p{0.98\linewidth}@{}}
$\square$\ \textbf{(A)} The height of $T$ is exactly 15. \\
\addlinespace[3pt]
$\square$\ \textbf{(B)} The height of $T$ is exactly 30. \\
\addlinespace[3pt]
$\square$\ \textbf{(C)} The height of $T$ is at least 15. \\
\addlinespace[3pt]
$\square$\ \textbf{(D)} The height of $T$ is at least 30.
\end{tabular}

\tcblower

\textbf{Official Key:} \textbf{\large C} \hfill \textit{\small [Medium]}
\vspace{2pt}\hrule height 0.3pt\vspace{3pt}
\textbf{Topic:} Algorithms \\
\textbf{Subtopic:} \textit{Graph traversals, minimum spanning trees, shortest paths}
\vspace{2pt}\hrule height 0.3pt\vspace{3pt}
\textbf{Guide / Hint:} \\
{\small The diameter of $G$ is $\text{diam}(G) = \max_{u,v} d(u, v) = 30$. For any BFS root $r$, the triangle inequality ensures that for the pair $(u, v)$ realizing the diameter, $30 = d(u, v) \le d(r, u) + d(r, v) \le 2 \times \text{height}(T)$. Hence, $\text{height}(T) \ge \lceil 30/2 \rceil = 15$.}

\end{tcolorbox}
\vspace{6pt}\hrule height 0.35pt\vspace{6pt}
\noindent
\begin{tcolorbox}[
    blanker,
    sidebyside,
    sidebyside align=top,
    lefthand width=0.55\textwidth,
    top=4pt, bottom=4pt, left=0pt, right=0pt,
    middle=16pt,
    underlay={
        \draw[dashed, black!80, line width=0.8pt, dash pattern=on 4pt off 3pt]
            ([xshift=0.55\textwidth+8pt]frame.north west) --
            ([xshift=0.55\textwidth+8pt]frame.south west);
    }
]
\textbf{Q.138} \quad {\footnotesize\textsc{[MCQ $\cdot$ 1 Mark(s) $\cdot$ Neg: 0.33]}} \hfill {\footnotesize\textbf{GATE 2025 CS2} \quad \tiny\texttt{2025\_CS2\_CS\_Q20}}
\vspace{2pt}\hrule height 0.3pt\vspace{4pt}

Consider an unordered list of $N$ distinct integers. What is the minimum number of element comparisons required to find an integer in the list that is NOT the largest in the list?

\par\vspace{6pt}\noindent
\begin{tabular}{@{}p{0.48\linewidth}@{\hspace{0.04\linewidth}}p{0.48\linewidth}@{}}
$\square$\ \textbf{(A)} 1 & $\square$\ \textbf{(B)} $N - 1$ \\
\addlinespace[3pt]
$\square$\ \textbf{(C)} $N$ & $\square$\ \textbf{(D)} $2N - 1$
\end{tabular}

\tcblower

\textbf{Official Key:} \textbf{\large A} \hfill \textit{\small [Easy]}
\vspace{2pt}\hrule height 0.3pt\vspace{3pt}
\textbf{Topic:} Algorithms \\
\textbf{Subtopic:} \textit{Searching and comparison complexity}
\vspace{2pt}\hrule height 0.3pt\vspace{3pt}
\textbf{Guide / Hint:} \\
{\small To find any element that is not the maximum, compare any two distinct elements in 1 comparison; the smaller of the two is guaranteed not to be the maximum.}

\end{tcolorbox}
\vspace{6pt}\hrule height 0.35pt\vspace{6pt}
\noindent
\begin{tcolorbox}[
    blanker,
    sidebyside,
    sidebyside align=top,
    lefthand width=0.55\textwidth,
    top=4pt, bottom=4pt, left=0pt, right=0pt,
    middle=16pt,
    underlay={
        \draw[dashed, black!80, line width=0.8pt, dash pattern=on 4pt off 3pt]
            ([xshift=0.55\textwidth+8pt]frame.north west) --
            ([xshift=0.55\textwidth+8pt]frame.south west);
    }
]
\textbf{Q.139} \quad {\footnotesize\textsc{[MCQ $\cdot$ 2 Mark(s) $\cdot$ Neg: 0.66]}} \hfill {\footnotesize\textbf{GATE 2025 CS2} \quad \tiny\texttt{2025\_CS2\_CS\_Q37}}
\vspace{2pt}\hrule height 0.3pt\vspace{4pt}

Let $G$ be an edge-weighted undirected graph with positive edge weights. Suppose a positive constant $\alpha$ is added to the weight of every edge in $G$. Which ONE of the following statements is TRUE?

\par\vspace{6pt}\noindent
\begin{tabular}{@{}p{0.98\linewidth}@{}}
$\square$\ \textbf{(A)} Every MST remains an MST, and every SP remains an SP. \\
\addlinespace[3pt]
$\square$\ \textbf{(B)} MSTs need not remain MSTs, and every SP remains an SP. \\
\addlinespace[3pt]
$\square$\ \textbf{(C)} Every MST remains an MST, and SPs need not remain SPs. \\
\addlinespace[3pt]
$\square$\ \textbf{(D)} MSTs need not remain MSTs, and SPs need not remain SPs.
\end{tabular}

\tcblower

\textbf{Official Key:} \textbf{\large C} \hfill \textit{\small [Medium]}
\vspace{2pt}\hrule height 0.3pt\vspace{3pt}
\textbf{Topic:} Algorithms \\
\textbf{Subtopic:} \textit{Graph algorithms: minimum spanning trees and shortest paths}
\vspace{2pt}\hrule height 0.3pt\vspace{3pt}
\textbf{Guide / Hint:} \\
{\small Adding a constant $\alpha$ to all edge weights preserves the relative order of all spanning trees because every spanning tree has exactly $V - 1$ edges (so all tree costs increase by $(V-1)\alpha$). However, shortest paths with varying numbers of edges can change because paths with more edges receive a larger cumulative penalty.}

\end{tcolorbox}
\vspace{6pt}\hrule height 0.35pt\vspace{6pt}
\noindent
\begin{tcolorbox}[
    blanker,
    sidebyside,
    sidebyside align=top,
    lefthand width=0.55\textwidth,
    top=4pt, bottom=4pt, left=0pt, right=0pt,
    middle=16pt,
    underlay={
        \draw[dashed, black!80, line width=0.8pt, dash pattern=on 4pt off 3pt]
            ([xshift=0.55\textwidth+8pt]frame.north west) --
            ([xshift=0.55\textwidth+8pt]frame.south west);
    }
]
\textbf{Q.140} \quad {\footnotesize\textsc{[MCQ $\cdot$ 2 Mark(s) $\cdot$ Neg: 0.66]}} \hfill {\footnotesize\textbf{GATE 2025 CS2} \quad \tiny\texttt{2025\_CS2\_CS\_Q41}}
\vspace{2pt}\hrule height 0.3pt\vspace{4pt}

An array $A$ of length $n$ with distinct elements is said to be bitonic if there is an index $1 \le i \le n$ such that $A[1..i]$ is strictly increasing and $A[i..n]$ is strictly decreasing. What is the optimal time complexity to search for a key in a bitonic array?

\par\vspace{6pt}\noindent
\begin{tabular}{@{}p{0.48\linewidth}@{\hspace{0.04\linewidth}}p{0.48\linewidth}@{}}
$\square$\ \textbf{(A)} $\Theta(n)$ & $\square$\ \textbf{(B)} $\Theta(1)$ \\
\addlinespace[3pt]
$\square$\ \textbf{(C)} $\Theta(\log^2 n)$ & $\square$\ \textbf{(D)} $\Theta(\log n)$
\end{tabular}

\tcblower

\textbf{Official Key:} \textbf{\large D} \hfill \textit{\small [Medium]}
\vspace{2pt}\hrule height 0.3pt\vspace{3pt}
\textbf{Topic:} Algorithms \\
\textbf{Subtopic:} \textit{Searching and divide-and-conquer: bitonic array search}
\vspace{2pt}\hrule height 0.3pt\vspace{3pt}
\textbf{Guide / Hint:} \\
{\small In a bitonic array, the maximum element can be found in $\Theta(\log n)$ time using binary search on mid-elements. Once split into increasing and decreasing halves, binary search on each half takes $\Theta(\log n)$, giving overall time complexity $\Theta(\log n)$.}

\end{tcolorbox}
\vspace{6pt}\hrule height 0.35pt\vspace{6pt}
\vspace{8pt}
\subsection*{Theory of Computation}
\vspace{-4pt}
\noindent
\begin{tcolorbox}[
    blanker,
    sidebyside,
    sidebyside align=top,
    lefthand width=0.55\textwidth,
    top=4pt, bottom=4pt, left=0pt, right=0pt,
    middle=16pt,
    underlay={
        \draw[dashed, black!80, line width=0.8pt, dash pattern=on 4pt off 3pt]
            ([xshift=0.55\textwidth+8pt]frame.north west) --
            ([xshift=0.55\textwidth+8pt]frame.south west);
    }
]
\textbf{Q.141} \quad {\footnotesize\textsc{[MCQ $\cdot$ 1 Mark(s) $\cdot$ Neg: 0.33]}} \hfill {\footnotesize\textbf{GATE 2018 CS} \quad \tiny\texttt{2018\_CS\_GA\_Q06\_v2}}
\vspace{2pt}\hrule height 0.3pt\vspace{4pt}

Let $N$ be an NFA with $n$ states. Let $k$ be the number of states of a minimal DFA which is equivalent to $N$. Which one of the following is necessarily true?

\par\vspace{6pt}\noindent
\begin{tabular}{@{}p{0.48\linewidth}@{\hspace{0.04\linewidth}}p{0.48\linewidth}@{}}
$\square$\ \textbf{(A)} $k \ge 2^n$ & $\square$\ \textbf{(B)} $k \ge n$ \\
\addlinespace[3pt]
$\square$\ \textbf{(C)} $k \le n^2$ & $\square$\ \textbf{(D)} $k \le 2^n$
\end{tabular}

\tcblower

\textbf{Official Key:} \textbf{\large D} \hfill \textit{\small [Easy]}
\vspace{2pt}\hrule height 0.3pt\vspace{3pt}
\textbf{Topic:} Theory of Computation \\
\textbf{Subtopic:} \textit{Regular expressions and finite automata}
\vspace{2pt}\hrule height 0.3pt\vspace{3pt}
\textbf{Guide / Hint:} \\
{\small By the subset construction algorithm, converting an $n$-state NFA to a DFA produces states that are subsets of the NFA states. Thus, the power set has $2^n$ possible states, ensuring that any equivalent minimal DFA has $k \le 2^n$ states.}

\end{tcolorbox}
\vspace{6pt}\hrule height 0.35pt\vspace{6pt}
\noindent
\begin{tcolorbox}[
    blanker,
    sidebyside,
    sidebyside align=top,
    lefthand width=0.55\textwidth,
    top=4pt, bottom=4pt, left=0pt, right=0pt,
    middle=16pt,
    underlay={
        \draw[dashed, black!80, line width=0.8pt, dash pattern=on 4pt off 3pt]
            ([xshift=0.55\textwidth+8pt]frame.north west) --
            ([xshift=0.55\textwidth+8pt]frame.south west);
    }
]
\textbf{Q.142} \quad {\footnotesize\textsc{[MCQ $\cdot$ 1 Mark(s) $\cdot$ Neg: 0.33]}} \hfill {\footnotesize\textbf{GATE 2018 CS} \quad \tiny\texttt{2018\_CS\_GA\_Q07\_v2}}
\vspace{2pt}\hrule height 0.3pt\vspace{4pt}

The set of all recursively enumerable languages is

\par\vspace{6pt}\noindent
\begin{tabular}{@{}p{0.98\linewidth}@{}}
$\square$\ \textbf{(A)} closed under complementation. \\
\addlinespace[3pt]
$\square$\ \textbf{(B)} closed under intersection. \\
\addlinespace[3pt]
$\square$\ \textbf{(C)} a subset of the set of all context-sensitive languages. \\
\addlinespace[3pt]
$\square$\ \textbf{(D)} an uncountable set.
\end{tabular}

\tcblower

\textbf{Official Key:} \textbf{\large B} \hfill \textit{\small [Easy]}
\vspace{2pt}\hrule height 0.3pt\vspace{3pt}
\textbf{Topic:} Theory of Computation \\
\textbf{Subtopic:} \textit{Regular and context-free languages, pumping lemma}
\vspace{2pt}\hrule height 0.3pt\vspace{3pt}
\textbf{Guide / Hint:} \\
{\small The set of recursively enumerable (RE) languages is closed under union, intersection, concatenation, and Kleene star, but not closed under complementation or set difference.}

\end{tcolorbox}
\vspace{6pt}\hrule height 0.35pt\vspace{6pt}
\noindent
\begin{tcolorbox}[
    blanker,
    sidebyside,
    sidebyside align=top,
    lefthand width=0.55\textwidth,
    top=4pt, bottom=4pt, left=0pt, right=0pt,
    middle=16pt,
    underlay={
        \draw[dashed, black!80, line width=0.8pt, dash pattern=on 4pt off 3pt]
            ([xshift=0.55\textwidth+8pt]frame.north west) --
            ([xshift=0.55\textwidth+8pt]frame.south west);
    }
]
\textbf{Q.143} \quad {\footnotesize\textsc{[MCQ $\cdot$ 2 Mark(s) $\cdot$ Neg: 0.66]}} \hfill {\footnotesize\textbf{GATE 2018 CS} \quad \tiny\texttt{2018\_CS\_GA\_Q35}}
\vspace{2pt}\hrule height 0.3pt\vspace{4pt}

Consider the following languages:
\begin{itemize}[leftmargin=*, nosep]
    \item[I.] $\{a^m b^n c^p d^q \mid m + p = n + q, \text{ where } m, n, p, q \ge 0\}$
    \item[II.] $\{a^m b^n c^p d^q \mid m = n \text{ and } p = q, \text{ where } m, n, p, q \ge 0\}$
    \item[III.] $\{a^m b^n c^p d^q \mid m = p \text{ and } n = q, \text{ where } m, n, p, q \ge 0\}$
    \item[IV.] $\{a^m b^n c^p d^q \mid m = q \text{ and } n = p, \text{ where } m, n, p, q \ge 0\}$
\end{itemize}
Which of the languages above are context-free?

\par\vspace{6pt}\noindent
\begin{tabular}{@{}p{0.48\linewidth}@{\hspace{0.04\linewidth}}p{0.48\linewidth}@{}}
$\square$\ \textbf{(A)} I and IV only & $\square$\ \textbf{(B)} I and II only \\
\addlinespace[3pt]
$\square$\ \textbf{(C)} II and III only & $\square$\ \textbf{(D)} II and IV only
\end{tabular}

\tcblower

\textbf{Official Key:} \textbf{\large B} \hfill \textit{\small [Medium]}
\vspace{2pt}\hrule height 0.3pt\vspace{3pt}
\textbf{Topic:} Theory of Computation \\
\textbf{Subtopic:} \textit{Context-free grammars and push-down automata}
\vspace{2pt}\hrule height 0.3pt\vspace{3pt}
\textbf{Guide / Hint:} \\
{\small Language I has condition $m - n = q - p$, which can be checked by pushing $a$'s, popping for $b$'s, then pushing/popping for $c$ and $d$ on a single PDA stack (CFL). Language II is the concatenation of CFLs $L_1 = \{a^m b^m\}$ and $L_2 = \{c^p d^p\}$, hence a CFL. Languages III and IV require multi-variable dependencies across non-adjacent symbols simultaneously, which cannot be handled by a single stack (non-CFL).}

\end{tcolorbox}
\vspace{6pt}\hrule height 0.35pt\vspace{6pt}
\noindent
\begin{tcolorbox}[
    blanker,
    sidebyside,
    sidebyside align=top,
    lefthand width=0.55\textwidth,
    top=4pt, bottom=4pt, left=0pt, right=0pt,
    middle=16pt,
    underlay={
        \draw[dashed, black!80, line width=0.8pt, dash pattern=on 4pt off 3pt]
            ([xshift=0.55\textwidth+8pt]frame.north west) --
            ([xshift=0.55\textwidth+8pt]frame.south west);
    }
]
\textbf{Q.144} \quad {\footnotesize\textsc{[MCQ $\cdot$ 2 Mark(s) $\cdot$ Neg: 0.66]}} \hfill {\footnotesize\textbf{GATE 2018 CS} \quad \tiny\texttt{2018\_CS\_GA\_Q36}}
\vspace{2pt}\hrule height 0.3pt\vspace{4pt}

Consider the following problems. $L(G)$ denotes the language generated by a grammar $G$. $L(M)$ denotes the language accepted by a machine $M$.
\begin{enumerate}[label=\Roman*.]
    \item For an unrestricted grammar $G$ and a string $w$, whether $w \in L(G)$
    \item Given a Turing machine $M$, whether $L(M)$ is regular
    \item Given two grammars $G_1$ and $G_2$, whether $L(G_1) = L(G_2)$
    \item Given an NFA $N$ and a DPDA $D$, whether $L(N) = L(D)$
\end{enumerate}
Which of the above problems are UNDECIDABLE?

\par\vspace{6pt}\noindent
\begin{tabular}{@{}p{0.48\linewidth}@{\hspace{0.04\linewidth}}p{0.48\linewidth}@{}}
$\square$\ \textbf{(A)} I and IV only & $\square$\ \textbf{(B)} II and III only \\
\addlinespace[3pt]
$\square$\ \textbf{(C)} III and IV only & $\square$\ \textbf{(D)} I, II and III only
\end{tabular}

\tcblower

\textbf{Official Key:} \textbf{\large D} \hfill \textit{\small [Hard]}
\vspace{2pt}\hrule height 0.3pt\vspace{3pt}
\textbf{Topic:} Theory of Computation \\
\textbf{Subtopic:} \textit{Turing machines and undecidability}
\vspace{2pt}\hrule height 0.3pt\vspace{3pt}
\textbf{Guide / Hint:} \\
{\small Problem I (membership for unrestricted grammars) is undecidable. Problem II (testing if a TM accepts a regular language) is undecidable by Rice's theorem. Problem III (equivalence of CFGs) is undecidable. Problem IV asks whether a regular language (NFA) equals a deterministic context-free language (DPDA), which is decidable. Hence, I, II, and III are undecidable.}

\end{tcolorbox}
\vspace{6pt}\hrule height 0.35pt\vspace{6pt}
\noindent
\begin{tcolorbox}[
    blanker,
    sidebyside,
    sidebyside align=top,
    lefthand width=0.55\textwidth,
    top=4pt, bottom=4pt, left=0pt, right=0pt,
    middle=16pt,
    underlay={
        \draw[dashed, black!80, line width=0.8pt, dash pattern=on 4pt off 3pt]
            ([xshift=0.55\textwidth+8pt]frame.north west) --
            ([xshift=0.55\textwidth+8pt]frame.south west);
    }
]
\textbf{Q.145} \quad {\footnotesize\textsc{[MCQ $\cdot$ 1 Mark(s) $\cdot$ Neg: 0.33]}} \hfill {\footnotesize\textbf{GATE 2022 CS} \quad \tiny\texttt{2022\_CS\_CS\_Q12}}
\vspace{2pt}\hrule height 0.3pt\vspace{4pt}

Which one of the following regular expressions correctly represents the language of the finite automaton given below?

\par\vspace{6pt}\noindent
\begin{tabular}{@{}p{0.98\linewidth}@{}}
$\square$\ \textbf{(A)} $ab^*ba^* + ba^*ab^*$ \\
\addlinespace[3pt]
$\square$\ \textbf{(B)} $(ab^*b + ba^*a)^*(a + b)$ \\
\addlinespace[3pt]
$\square$\ \textbf{(C)} $(ab^*b + ba^*a)^*(a^*b + b^*a)$ \\
\addlinespace[3pt]
$\square$\ \textbf{(D)} $(ba^*a + ab^*b)^*(ab^* + ba^*)$
\end{tabular}

\tcblower

\textbf{Official Key:} \textbf{\large D} \hfill \textit{\small [Medium]}
\vspace{2pt}\hrule height 0.3pt\vspace{3pt}
\textbf{Topic:} Theory of Computation \\
\textbf{Subtopic:} \textit{Regular expressions and finite automata}
\vspace{2pt}\hrule height 0.3pt\vspace{3pt}
\textbf{Guide / Hint:} \\
{\small From the start state, loops on intermediate states allow repeating patterns of $ba^*a$ and $ab^*b$ before moving to an accepting state via either $ab^*$ or $ba^*$. Factoring loop paths and exit transitions yields the regular expression $(ba^*a + ab^*b)^*(ab^* + ba^*)$.}

\end{tcolorbox}
\vspace{6pt}\hrule height 0.35pt\vspace{6pt}
\noindent
\begin{tcolorbox}[
    blanker,
    sidebyside,
    sidebyside align=top,
    lefthand width=0.55\textwidth,
    top=4pt, bottom=4pt, left=0pt, right=0pt,
    middle=16pt,
    underlay={
        \draw[dashed, black!80, line width=0.8pt, dash pattern=on 4pt off 3pt]
            ([xshift=0.55\textwidth+8pt]frame.north west) --
            ([xshift=0.55\textwidth+8pt]frame.south west);
    }
]
\textbf{Q.146} \quad {\footnotesize\textsc{[MCQ $\cdot$ 2 Mark(s) $\cdot$ Neg: 0.66]}} \hfill {\footnotesize\textbf{GATE 2023 CS} \quad \tiny\texttt{2023\_CS\_CS\_Q39}}
\vspace{2pt}\hrule height 0.3pt\vspace{4pt}

Consider the context-free grammar $G$ below:
\begin{align*}
S &\to aSb \mid X \\
X &\to aX \mid Xb \mid a \mid b,
\end{align*}
where $S$ and $X$ are non-terminals, and $a$ and $b$ are terminal symbols. The starting non-terminal is $S$.
Which one of the following statements is CORRECT?

\par\vspace{6pt}\noindent
\begin{tabular}{@{}p{0.98\linewidth}@{}}
$\square$\ \textbf{(A)} The language generated by $G$ is $(a + b)^*$ \\
\addlinespace[3pt]
$\square$\ \textbf{(B)} The language generated by $G$ is $a^*(a + b)b^*$ \\
\addlinespace[3pt]
$\square$\ \textbf{(C)} The language generated by $G$ is $a^*b^*(a + b)$ \\
\addlinespace[3pt]
$\square$\ \textbf{(D)} The language generated by $G$ is not a regular language
\end{tabular}

\tcblower

\textbf{Official Key:} \textbf{\large B} \hfill \textit{\small [Medium]}
\vspace{2pt}\hrule height 0.3pt\vspace{3pt}
\textbf{Topic:} Theory of Computation \\
\textbf{Subtopic:} \textit{Context-free grammars and push-down automata}
\vspace{2pt}\hrule height 0.3pt\vspace{3pt}
\textbf{Guide / Hint:} \\
{\small The subgrammar $X \to aX \mid Xb \mid a \mid b$ derives all non-empty strings with any number of $a$'s followed by any number of $b$'s having at least one $a$ or $b$, described by $a^*(a+b)b^*$. Rule $S \to aSb \mid X$ prefixes $k$ $a$'s and suffixes $k$ $b$'s ($k \ge 0$), yielding $a^k a^*(a+b)b^* b^k = a^*(a+b)b^*$. Because this language is represented by a regular expression, it is a regular language.}

\end{tcolorbox}
\vspace{6pt}\hrule height 0.35pt\vspace{6pt}
\noindent
\begin{tcolorbox}[
    blanker,
    sidebyside,
    sidebyside align=top,
    lefthand width=0.55\textwidth,
    top=4pt, bottom=4pt, left=0pt, right=0pt,
    middle=16pt,
    underlay={
        \draw[dashed, black!80, line width=0.8pt, dash pattern=on 4pt off 3pt]
            ([xshift=0.55\textwidth+8pt]frame.north west) --
            ([xshift=0.55\textwidth+8pt]frame.south west);
    }
]
\textbf{Q.147} \quad {\footnotesize\textsc{[MCQ $\cdot$ 1 Mark(s) $\cdot$ Neg: 0.33]}} \hfill {\footnotesize\textbf{GATE 2025 CS1} \quad \tiny\texttt{2025\_CS1\_CS\_Q19}}
\vspace{2pt}\hrule height 0.3pt\vspace{4pt}

Consider the following context-free grammar $G$, where $S, A$, and $B$ are the variables (non-terminals), $a$ and $b$ are the terminal symbols, $S$ is the start variable, and the rules of $G$ are described as:
\[ S \to aaB \mid Abb \]
\[ A \to a \mid aA \]
\[ B \to b \mid bB \]
Which ONE of the languages $L(G)$ is accepted by $G$?

\par\vspace{6pt}\noindent
\begin{tabular}{@{}p{0.98\linewidth}@{}}
$\square$\ \textbf{(A)} $L(G) = \{a^2b^n \mid n \ge 1\} \cup \{a^nb^2 \mid n \ge 1\}$ \\
\addlinespace[3pt]
$\square$\ \textbf{(B)} $L(G) = \{a^nb^{2n} \mid n \ge 1\} \cup \{a^{2n}b^n \mid n \ge 1\}$ \\
\addlinespace[3pt]
$\square$\ \textbf{(C)} $L(G) = \{a^nb^n \mid n \ge 1\}$ \\
\addlinespace[3pt]
$\square$\ \textbf{(D)} $L(G) = \{a^{2n}b^{2n} \mid n \ge 1\}$
\end{tabular}

\tcblower

\textbf{Official Key:} \textbf{\large A} \hfill \textit{\small [Easy]}
\vspace{2pt}\hrule height 0.3pt\vspace{3pt}
\textbf{Topic:} Theory of Computation \\
\textbf{Subtopic:} \textit{Context-free grammars and push-down automata}
\vspace{2pt}\hrule height 0.3pt\vspace{3pt}
\textbf{Guide / Hint:} \\
{\small From the rules, $A$ generates $a^+$ (any $a^n$ with $n \ge 1$) and $B$ generates $b^+$ (any $b^n$ with $n \ge 1$). Thus, $aaB$ produces $a^2b^n$ ($n \ge 1$), and $Abb$ produces $a^nb^2$ ($n \ge 1$). Therefore, $L(G) = \{a^2b^n \mid n \ge 1\} \cup \{a^nb^2 \mid n \ge 1\}$.}

\end{tcolorbox}
\vspace{6pt}\hrule height 0.35pt\vspace{6pt}
\noindent
\begin{tcolorbox}[
    blanker,
    sidebyside,
    sidebyside align=top,
    lefthand width=0.55\textwidth,
    top=4pt, bottom=4pt, left=0pt, right=0pt,
    middle=16pt,
    underlay={
        \draw[dashed, black!80, line width=0.8pt, dash pattern=on 4pt off 3pt]
            ([xshift=0.55\textwidth+8pt]frame.north west) --
            ([xshift=0.55\textwidth+8pt]frame.south west);
    }
]
\textbf{Q.148} \quad {\footnotesize\textsc{[MCQ $\cdot$ 2 Mark(s) $\cdot$ Neg: 0.66]}} \hfill {\footnotesize\textbf{GATE 2025 CS1} \quad \tiny\texttt{2025\_CS1\_CS\_Q44}}
\vspace{2pt}\hrule height 0.3pt\vspace{4pt}

Consider the following two languages over the alphabet $\{a, b\}$:
\[ L\_1 = \{ \alpha\beta\alpha \mid \alpha \in \{a, b\}^+ \text{ AND } \beta \in \{a, b\}^+ \} \]
\[ L\_2 = \{ \alpha\beta\alpha \mid \alpha \in \{a\}^+ \text{ AND } \beta \in \{a, b\}^+ \} \]
Which ONE of the following statements is CORRECT?

\par\vspace{6pt}\noindent
\begin{tabular}{@{}p{0.98\linewidth}@{}}
$\square$\ \textbf{(A)} Both $L_1$ and $L_2$ are regular languages. \\
\addlinespace[3pt]
$\square$\ \textbf{(B)} $L_1$ is a regular language but $L_2$ is not a regular language. \\
\addlinespace[3pt]
$\square$\ \textbf{(C)} $L_1$ is not a regular language but $L_2$ is a regular language. \\
\addlinespace[3pt]
$\square$\ \textbf{(D)} Neither $L_1$ nor $L_2$ is a regular language.
\end{tabular}

\tcblower

\textbf{Official Key:} \textbf{\large C} \hfill \textit{\small [Medium]}
\vspace{2pt}\hrule height 0.3pt\vspace{3pt}
\textbf{Topic:} Theory of Computation \\
\textbf{Subtopic:} \textit{Regular and context-free languages, pumping lemma}
\vspace{2pt}\hrule height 0.3pt\vspace{3pt}
\textbf{Guide / Hint:} \\
{\small For $L_2$, $\alpha = a^k$ ($k \ge 1$). A string belongs to $L_2$ if and only if it starts with $a$, ends with $a$, and has length at least 3 ($a(a+b)^+a$), which is clearly regular. For $L_1$, the prefix $\alpha$ and suffix $\alpha$ must match over $\{a, b\}^+$, which requires unbounded memory to verify matching arbitrary binary strings, making $L_1$ non-regular.}

\end{tcolorbox}
\vspace{6pt}\hrule height 0.35pt\vspace{6pt}
\noindent
\begin{tcolorbox}[
    blanker,
    sidebyside,
    sidebyside align=top,
    lefthand width=0.55\textwidth,
    top=4pt, bottom=4pt, left=0pt, right=0pt,
    middle=16pt,
    underlay={
        \draw[dashed, black!80, line width=0.8pt, dash pattern=on 4pt off 3pt]
            ([xshift=0.55\textwidth+8pt]frame.north west) --
            ([xshift=0.55\textwidth+8pt]frame.south west);
    }
]
\textbf{Q.149} \quad {\footnotesize\textsc{[MCQ $\cdot$ 2 Mark(s) $\cdot$ Neg: 0.66]}} \hfill {\footnotesize\textbf{GATE 2025 CS1} \quad \tiny\texttt{2025\_CS1\_CS\_Q45}}
\vspace{2pt}\hrule height 0.3pt\vspace{4pt}

Consider the following two languages over the alphabet $\{a, b, c\}$, where $m$ and $n$ are natural numbers.
\[ L\_1 = \{a^mb^mc^{m+n} \mid m, n \ge 1\} \]
\[ L\_2 = \{a^mb^nc^{m+n} \mid m, n \ge 1\} \]
Which ONE of the following statements is CORRECT?

\par\vspace{6pt}\noindent
\begin{tabular}{@{}p{0.98\linewidth}@{}}
$\square$\ \textbf{(A)} Both $L_1$ and $L_2$ are context-free languages. \\
\addlinespace[3pt]
$\square$\ \textbf{(B)} $L_1$ is a context-free language but $L_2$ is not a context-free language. \\
\addlinespace[3pt]
$\square$\ \textbf{(C)} $L_1$ is not a context-free language but $L_2$ is a context-free language. \\
\addlinespace[3pt]
$\square$\ \textbf{(D)} Neither $L_1$ nor $L_2$ are context-free languages.
\end{tabular}

\tcblower

\textbf{Official Key:} \textbf{\large C} \hfill \textit{\small [Medium]}
\vspace{2pt}\hrule height 0.3pt\vspace{3pt}
\textbf{Topic:} Theory of Computation \\
\textbf{Subtopic:} \textit{Context-free grammars and push-down automata}
\vspace{2pt}\hrule height 0.3pt\vspace{3pt}
\textbf{Guide / Hint:} \\
{\small In $L_2$, rewrite $c^{m+n} = c^n c^m$ or generate using PDA by pushing $a$'s and $b$'s and popping on $c$'s ($c$'s match total count of $a$'s and $b$'s), which requires a single stack, so $L_2$ is context-free. In $L_1$, $a^m b^m c^m c^n$ simultaneously requires equality between powers of $a$ and $b$, plus matching with $c$, requiring two independent stack comparisons, making $L_1$ non-CFL.}

\end{tcolorbox}
\vspace{6pt}\hrule height 0.35pt\vspace{6pt}
\noindent
\begin{tcolorbox}[
    blanker,
    sidebyside,
    sidebyside align=top,
    lefthand width=0.55\textwidth,
    top=4pt, bottom=4pt, left=0pt, right=0pt,
    middle=16pt,
    underlay={
        \draw[dashed, black!80, line width=0.8pt, dash pattern=on 4pt off 3pt]
            ([xshift=0.55\textwidth+8pt]frame.north west) --
            ([xshift=0.55\textwidth+8pt]frame.south west);
    }
]
\textbf{Q.150} \quad {\footnotesize\textsc{[MCQ $\cdot$ 1 Mark(s) $\cdot$ Neg: 0.33]}} \hfill {\footnotesize\textbf{GATE 2025 CS2} \quad \tiny\texttt{2025\_CS2\_CS\_Q24}}
\vspace{2pt}\hrule height 0.3pt\vspace{4pt}

Which ONE of the following languages is accepted by a deterministic pushdown automaton?

\par\vspace{6pt}\noindent
\begin{tabular}{@{}p{0.98\linewidth}@{}}
$\square$\ \textbf{(A)} Any regular language. \\
\addlinespace[3pt]
$\square$\ \textbf{(B)} Any context-free language. \\
\addlinespace[3pt]
$\square$\ \textbf{(C)} Any language accepted by a non-deterministic pushdown automaton. \\
\addlinespace[3pt]
$\square$\ \textbf{(D)} Any decidable language.
\end{tabular}

\tcblower

\textbf{Official Key:} \textbf{\large A} \hfill \textit{\small [Easy]}
\vspace{2pt}\hrule height 0.3pt\vspace{3pt}
\textbf{Topic:} Theory of Computation \\
\textbf{Subtopic:} \textit{Context-free grammars and push-down automata}
\vspace{2pt}\hrule height 0.3pt\vspace{3pt}
\textbf{Guide / Hint:} \\
{\small Every regular language is recognized by a Deterministic Finite Automaton (DFA), which can be directly simulated by a DPDA without using stack operations.}

\end{tcolorbox}
\vspace{6pt}\hrule height 0.35pt\vspace{6pt}
\noindent
\begin{tcolorbox}[
    blanker,
    sidebyside,
    sidebyside align=top,
    lefthand width=0.55\textwidth,
    top=4pt, bottom=4pt, left=0pt, right=0pt,
    middle=16pt,
    underlay={
        \draw[dashed, black!80, line width=0.8pt, dash pattern=on 4pt off 3pt]
            ([xshift=0.55\textwidth+8pt]frame.north west) --
            ([xshift=0.55\textwidth+8pt]frame.south west);
    }
]
\textbf{Q.151} \quad {\footnotesize\textsc{[MCQ $\cdot$ 1 Mark(s) $\cdot$ Neg: 0.33]}} \hfill {\footnotesize\textbf{GATE 2025 CS2} \quad \tiny\texttt{2025\_CS2\_CS\_Q25}}
\vspace{2pt}\hrule height 0.3pt\vspace{4pt}

Let $G_1, G_2$ be Context Free Grammars (CFGs) and $R$ be a regular expression. For a grammar $G$, let $L(G)$ denote the language generated by $G$. Which ONE among the following questions is decidable?

\par\vspace{6pt}\noindent
\begin{tabular}{@{}p{0.98\linewidth}@{}}
$\square$\ \textbf{(A)} Is $L(G_1) = L(G_2)$? \\
\addlinespace[3pt]
$\square$\ \textbf{(B)} Is $L(G_1) \cap L(G_2) = \emptyset$? \\
\addlinespace[3pt]
$\square$\ \textbf{(C)} Is $L(G_1) = L(R)$? \\
\addlinespace[3pt]
$\square$\ \textbf{(D)} Is $L(G_1) = \emptyset$?
\end{tabular}

\tcblower

\textbf{Official Key:} \textbf{\large D} \hfill \textit{\small [Easy]}
\vspace{2pt}\hrule height 0.3pt\vspace{3pt}
\textbf{Topic:} Theory of Computation \\
\textbf{Subtopic:} \textit{Turing machines and undecidability}
\vspace{2pt}\hrule height 0.3pt\vspace{3pt}
\textbf{Guide / Hint:} \\
{\small For Context-Free Grammars, emptiness ($L(G) = \emptyset$) is decidable via finding generating non-terminals, whereas equivalence ($L(G_1) = L(G_2)$) and disjointness ($L(G_1) \cap L(G_2) = \emptyset$) are undecidable.}

\end{tcolorbox}
\vspace{6pt}\hrule height 0.35pt\vspace{6pt}
\vspace{8pt}
\subsection*{Compiler Design}
\vspace{-4pt}
\noindent
\begin{tcolorbox}[
    blanker,
    sidebyside,
    sidebyside align=top,
    lefthand width=0.55\textwidth,
    top=4pt, bottom=4pt, left=0pt, right=0pt,
    middle=16pt,
    underlay={
        \draw[dashed, black!80, line width=0.8pt, dash pattern=on 4pt off 3pt]
            ([xshift=0.55\textwidth+8pt]frame.north west) --
            ([xshift=0.55\textwidth+8pt]frame.south west);
    }
]
\textbf{Q.152} \quad {\footnotesize\textsc{[MCQ $\cdot$ 1 Mark(s) $\cdot$ Neg: 0.33]}} \hfill {\footnotesize\textbf{GATE 2018 CS} \quad \tiny\texttt{2018\_CS\_GA\_Q08\_v2}}
\vspace{2pt}\hrule height 0.3pt\vspace{4pt}

Which one of the following statements is FALSE?

\par\vspace{6pt}\noindent
\begin{tabular}{@{}p{0.98\linewidth}@{}}
$\square$\ \textbf{(A)} Context-free grammar can be used to specify both lexical and syntax rules. \\
\addlinespace[3pt]
$\square$\ \textbf{(B)} Type checking is done before parsing. \\
\addlinespace[3pt]
$\square$\ \textbf{(C)} High-level language programs can be converted to intermediate code directly. \\
\addlinespace[3pt]
$\square$\ \textbf{(D)} Arguments to a function can be passed using the program stack.
\end{tabular}

\tcblower

\textbf{Official Key:} \textbf{\large B} \hfill \textit{\small [Easy]}
\vspace{2pt}\hrule height 0.3pt\vspace{3pt}
\textbf{Topic:} Compiler Design \\
\textbf{Subtopic:} \textit{Lexical analysis, parsing, syntax-directed translation}
\vspace{2pt}\hrule height 0.3pt\vspace{3pt}
\textbf{Guide / Hint:} \\
{\small Type checking verifies semantic consistency according to language type rules and is performed during semantic analysis, which occurs after syntax analysis (parsing). Therefore, the statement that type checking is done during syntax analysis is FALSE.}

\end{tcolorbox}
\vspace{6pt}\hrule height 0.35pt\vspace{6pt}
\noindent
\begin{tcolorbox}[
    blanker,
    sidebyside,
    sidebyside align=top,
    lefthand width=0.55\textwidth,
    top=4pt, bottom=4pt, left=0pt, right=0pt,
    middle=16pt,
    underlay={
        \draw[dashed, black!80, line width=0.8pt, dash pattern=on 4pt off 3pt]
            ([xshift=0.55\textwidth+8pt]frame.north west) --
            ([xshift=0.55\textwidth+8pt]frame.south west);
    }
]
\textbf{Q.153} \quad {\footnotesize\textsc{[MCQ $\cdot$ 2 Mark(s) $\cdot$ Neg: 0.66]}} \hfill {\footnotesize\textbf{GATE 2018 CS} \quad \tiny\texttt{2018\_CS\_GA\_Q37}}
\vspace{2pt}\hrule height 0.3pt\vspace{4pt}

A lexical analyzer uses the following patterns to recognize three tokens $T_1, T_2$, and $T_3$ over the alphabet $\{a, b, c\}$:
\[ T\_1: a?(b|c)^*a \qquad T\_2: b?(a|c)^*b \qquad T\_3: c?(b|a)^*c \]
Note that '?' means 0 or 1 occurrence of the preceding symbol. The analyzer matches the longest prefix of the remaining input. For the input string `bbaacabc`, the token sequence generated is

\par\vspace{6pt}\noindent
\begin{tabular}{@{}p{0.48\linewidth}@{\hspace{0.04\linewidth}}p{0.48\linewidth}@{}}
$\square$\ \textbf{(A)} $T_1 T_2 T_3$ & $\square$\ \textbf{(B)} $T_1 T_3$ \\
\addlinespace[3pt]
$\square$\ \textbf{(C)} $T_2 T_1 T_3$ & $\square$\ \textbf{(D)} $T_3 T_3$
\end{tabular}

\tcblower

\textbf{Official Key:} \textbf{\large D} \hfill \textit{\small [Medium]}
\vspace{2pt}\hrule height 0.3pt\vspace{3pt}
\textbf{Topic:} Compiler Design \\
\textbf{Subtopic:} \textit{Lexical analysis, parsing, syntax-directed translation}
\vspace{2pt}\hrule height 0.3pt\vspace{3pt}
\textbf{Guide / Hint:} \\
{\small A lexical analyzer applies the longest match rule. For input `bbaacabc`, pattern $T_3 = c?(b|a)^*c$ matches prefix `bbaac` (starts without $c$, matches $(b|a)^* = bbaa$, ends with $c$), consuming `bbaac` as token $T_3$. The remaining string is `abc`, which also matches $T_3$ (starts without $c$, $(b|a)^* = ab$, ends with $c$). Hence the recognized sequence of tokens is $T_3 T_3$.}

\end{tcolorbox}
\vspace{6pt}\hrule height 0.35pt\vspace{6pt}
\noindent
\begin{tcolorbox}[
    blanker,
    sidebyside,
    sidebyside align=top,
    lefthand width=0.55\textwidth,
    top=4pt, bottom=4pt, left=0pt, right=0pt,
    middle=16pt,
    underlay={
        \draw[dashed, black!80, line width=0.8pt, dash pattern=on 4pt off 3pt]
            ([xshift=0.55\textwidth+8pt]frame.north west) --
            ([xshift=0.55\textwidth+8pt]frame.south west);
    }
]
\textbf{Q.154} \quad {\footnotesize\textsc{[MCQ $\cdot$ 1 Mark(s) $\cdot$ Neg: 0.33]}} \hfill {\footnotesize\textbf{GATE 2022 CS} \quad \tiny\texttt{2022\_CS\_CS\_Q13}}
\vspace{2pt}\hrule height 0.3pt\vspace{4pt}

Which one of the following statements is TRUE?

\par\vspace{6pt}\noindent
\begin{tabular}{@{}p{0.98\linewidth}@{}}
$\square$\ \textbf{(A)} The LALR(1) parser for a grammar G cannot have reduce-reduce conflict if the LR(1) parser for G does not have reduce-reduce conflict. \\
\addlinespace[3pt]
$\square$\ \textbf{(B)} Symbol table is accessed only during the lexical analysis phase. \\
\addlinespace[3pt]
$\square$\ \textbf{(C)} Data flow analysis is necessary for run-time memory management. \\
\addlinespace[3pt]
$\square$\ \textbf{(D)} LR(1) parsing is sufficient for deterministic context-free languages.
\end{tabular}

\tcblower

\textbf{Official Key:} \textbf{\large D} \hfill \textit{\small [Easy]}
\vspace{2pt}\hrule height 0.3pt\vspace{3pt}
\textbf{Topic:} Compiler Design \\
\textbf{Subtopic:} \textit{Lexical analysis, parsing, syntax-directed translation}
\vspace{2pt}\hrule height 0.3pt\vspace{3pt}
\textbf{Guide / Hint:} \\
{\small LR(1) grammars describe exactly the family of Deterministic Context-Free Languages (DCFLs), meaning LR(1) parsing is both necessary and sufficient for any DCFL. LALR(1) can introduce reduce-reduce conflicts not present in LR(1), and symbol tables are accessed throughout compilation.}

\end{tcolorbox}
\vspace{6pt}\hrule height 0.35pt\vspace{6pt}
\noindent
\begin{tcolorbox}[
    blanker,
    sidebyside,
    sidebyside align=top,
    lefthand width=0.55\textwidth,
    top=4pt, bottom=4pt, left=0pt, right=0pt,
    middle=16pt,
    underlay={
        \draw[dashed, black!80, line width=0.8pt, dash pattern=on 4pt off 3pt]
            ([xshift=0.55\textwidth+8pt]frame.north west) --
            ([xshift=0.55\textwidth+8pt]frame.south west);
    }
]
\textbf{Q.155} \quad {\footnotesize\textsc{[MCQ $\cdot$ 1 Mark(s) $\cdot$ Neg: 0.33]}} \hfill {\footnotesize\textbf{GATE 2023 CS} \quad \tiny\texttt{2023\_CS\_CS\_Q36}}
\vspace{2pt}\hrule height 0.3pt\vspace{4pt}

Consider the following program:
\begin{lstlisting}[language=C]
int main() {
    f1();
    f2(2);
    f3();
    return(0);
}
int f1() {
    return(1);
}
int f2(int X) {
    f3();
    if (X == 1)
        return f1();
    else
        return (X * f2(X - 1));
}
int f3() {
    return(5);
}
\end{lstlisting}
Which one of the following options represents the activation tree corresponding to the main function?

\par\vspace{6pt}\noindent
\begin{tabular}{@{}p{0.98\linewidth}@{}}
$\square$\ \textbf{(A)} main -> [f1, f2, f3]; where f2(2) -> [f3, f2(1)]; and f2(1) -> [f3, f1] \\
\addlinespace[3pt]
$\square$\ \textbf{(B)} main -> [f2, f3]; where f2 -> [f1, f1, f3] \\
\addlinespace[3pt]
$\square$\ \textbf{(C)} main -> [f2]; where f2 -> [f1, f3, f1] \\
\addlinespace[3pt]
$\square$\ \textbf{(D)} main -> [f1, f2, f3]; where f2 -> [f3, f2, f1]
\end{tabular}

\tcblower

\textbf{Official Key:} \textbf{\large A} \hfill \textit{\small [Medium]}
\vspace{2pt}\hrule height 0.3pt\vspace{3pt}
\textbf{Topic:} Compiler Design \\
\textbf{Subtopic:} \textit{Runtime environments}
\vspace{2pt}\hrule height 0.3pt\vspace{3pt}
\textbf{Guide / Hint:} \\
{\small In `main()`, the calls are `f1()`, `f2(2)`, and `f3()` as direct children. Inside `f2(2)`, it invokes `f3()` first, then because $X=2 \neq 1$, recursively invokes `f2(1)`. Inside `f2(1)`, it invokes `f3()` first, and then since $X=1$, invokes `f1()`. This call sequence maps directly to the tree structure in option A.}

\end{tcolorbox}
\vspace{6pt}\hrule height 0.35pt\vspace{6pt}
\noindent
\begin{tcolorbox}[
    blanker,
    sidebyside,
    sidebyside align=top,
    lefthand width=0.55\textwidth,
    top=4pt, bottom=4pt, left=0pt, right=0pt,
    middle=16pt,
    underlay={
        \draw[dashed, black!80, line width=0.8pt, dash pattern=on 4pt off 3pt]
            ([xshift=0.55\textwidth+8pt]frame.north west) --
            ([xshift=0.55\textwidth+8pt]frame.south west);
    }
]
\textbf{Q.156} \quad {\footnotesize\textsc{[MCQ $\cdot$ 2 Mark(s) $\cdot$ Neg: 0.66]}} \hfill {\footnotesize\textbf{GATE 2024 CS1} \quad \tiny\texttt{2024\_CS1\_CS\_Q37}}
\vspace{2pt}\hrule height 0.3pt\vspace{4pt}

Consider the following syntax-directed definition (SDD):
\begin{align*}
S &\to DHTU & \{ S.val = D.val + H.val + T.val + U.val; \} \\
D &\to \text{``M''} D_1 & \{ D.val = 5 + D_1.val; \} \\
D &\to \epsilon & \{ D.val = -5; \} \\
H &\to \text{``L''} H_1 & \{ H.val = 5 * 10 + H_1.val; \} \\
H &\to \epsilon & \{ H.val = -10; \} \\
T &\to \text{``C''} T_1 & \{ T.val = 5 * 100 + T_1.val; \} \\
T &\to \epsilon & \{ T.val = -5; \} \\
U &\to \text{``K''} & \{ U.val = 5; \}
\end{align*}
Given ``MMLK'' as the input, which one of the following options is the CORRECT value computed by the SDD (in the attribute $S.val$)?

\par\vspace{6pt}\noindent
\begin{tabular}{@{}p{0.48\linewidth}@{\hspace{0.04\linewidth}}p{0.48\linewidth}@{}}
$\square$\ \textbf{(A)} 45 & $\square$\ \textbf{(B)} 50 \\
\addlinespace[3pt]
$\square$\ \textbf{(C)} 55 & $\square$\ \textbf{(D)} 65
\end{tabular}

\tcblower

\textbf{Official Key:} \textbf{\large A} \hfill \textit{\small [Medium]}
\vspace{2pt}\hrule height 0.3pt\vspace{3pt}
\textbf{Topic:} Compiler Design \\
\textbf{Subtopic:} \textit{Lexical analysis, parsing, syntax-directed translation}
\vspace{2pt}\hrule height 0.3pt\vspace{3pt}
\textbf{Guide / Hint:} \\
{\small For input "MMLK" with production $S \to DHTU$: $D$ derives "MM" via $D \to \text{"M"} D_1 \to \text{"M"}\text{"M"} D_2 \to \text{"M"}\text{"M"}\epsilon$, giving $D.val = 5 + 5 - 5 = 5$. $H$ derives "L" via $H \to \text{"L"} H_1 \to \text{"L"}\epsilon$, giving $H.val = 50 - 10 = 40$. $T$ derives $\epsilon$, giving $T.val = -5$. $U$ derives "K", giving $U.val = 5$. Summing attribute values yields $S.val = 5 + 40 - 5 + 5 = 45$.}

\end{tcolorbox}
\vspace{6pt}\hrule height 0.35pt\vspace{6pt}
\noindent
\begin{tcolorbox}[
    blanker,
    sidebyside,
    sidebyside align=top,
    lefthand width=0.55\textwidth,
    top=4pt, bottom=4pt, left=0pt, right=0pt,
    middle=16pt,
    underlay={
        \draw[dashed, black!80, line width=0.8pt, dash pattern=on 4pt off 3pt]
            ([xshift=0.55\textwidth+8pt]frame.north west) --
            ([xshift=0.55\textwidth+8pt]frame.south west);
    }
]
\textbf{Q.157} \quad {\footnotesize\textsc{[MCQ $\cdot$ 2 Mark(s) $\cdot$ Neg: 0.66]}} \hfill {\footnotesize\textbf{GATE 2024 CS1} \quad \tiny\texttt{2024\_CS1\_CS\_Q38}}
\vspace{2pt}\hrule height 0.3pt\vspace{4pt}

Consider the following grammar $G$, with $S$ as the start symbol. The grammar $G$ has three incomplete productions denoted by (1), (2), and (3).
\begin{align*}
S &\to daT \mid (1) \\
T &\to aS \mid bT \mid (2) \\
R &\to (3) \mid \epsilon
\end{align*}
The set of terminals is $\{a, b, c, d, f\}$. The FIRST and FOLLOW sets of the different non-terminals are as follows:
$$\text{FIRST}(S) = \{c, d, f\}, \quad \text{FIRST}(T) = \{a, b, \epsilon\}, \quad \text{FIRST}(R) = \{c, \epsilon\}$$
$$\text{FOLLOW}(S) = \text{FOLLOW}(T) = \{c, f, \$\}, \quad \text{FOLLOW}(R) = \{f\}$$
Which one of the following options CORRECTLY fills in the incomplete productions?

\par\vspace{6pt}\noindent
\begin{tabular}{@{}p{0.98\linewidth}@{}}
$\square$\ \textbf{(A)} (1) $S \to Rf$ (2) $T \to \epsilon$ (3) $R \to cTR$ \\
\addlinespace[3pt]
$\square$\ \textbf{(B)} (1) $S \to fR$ (2) $T \to \epsilon$ (3) $R \to cTR$ \\
\addlinespace[3pt]
$\square$\ \textbf{(C)} (1) $S \to fR$ (2) $T \to cT$ (3) $R \to cR$ \\
\addlinespace[3pt]
$\square$\ \textbf{(D)} (1) $S \to Rf$ (2) $T \to cT$ (3) $R \to cR$
\end{tabular}

\tcblower

\textbf{Official Key:} \textbf{\large A} \hfill \textit{\small [Medium]}
\vspace{2pt}\hrule height 0.3pt\vspace{3pt}
\textbf{Topic:} Compiler Design \\
\textbf{Subtopic:} \textit{Lexical analysis, parsing, syntax-directed translation}
\vspace{2pt}\hrule height 0.3pt\vspace{3pt}
\textbf{Guide / Hint:} \\
{\small Since $\text{FIRST}(S) = \{c, d, f\}$ and $S \to daT$ generates $d$, production (1) must generate $c$ and $f$. With (1) as $S \to Rf$ where $\text{FIRST}(R) = \{c, \epsilon\}$, $\text{FIRST}(Rf) = \{c, f\}$. Since $\epsilon \in \text{FIRST}(T)$, production (2) must be $T \to \epsilon$. Production (3) as $R \to cTR$ matches $\text{FIRST}(R) = \{c, \epsilon\}$ and $\text{FOLLOW}(R) = \{f\}$, matching option (A).}

\end{tcolorbox}
\vspace{6pt}\hrule height 0.35pt\vspace{6pt}
\noindent
\begin{tcolorbox}[
    blanker,
    sidebyside,
    sidebyside align=top,
    lefthand width=0.55\textwidth,
    top=4pt, bottom=4pt, left=0pt, right=0pt,
    middle=16pt,
    underlay={
        \draw[dashed, black!80, line width=0.8pt, dash pattern=on 4pt off 3pt]
            ([xshift=0.55\textwidth+8pt]frame.north west) --
            ([xshift=0.55\textwidth+8pt]frame.south west);
    }
]
\textbf{Q.158} \quad {\footnotesize\textsc{[MCQ $\cdot$ 2 Mark(s) $\cdot$ Neg: 0.66]}} \hfill {\footnotesize\textbf{GATE 2024 CS1} \quad \tiny\texttt{2024\_CS1\_CS\_Q39}}
\vspace{2pt}\hrule height 0.3pt\vspace{4pt}

Consider the following pseudo-code:
\begin{lstlisting}
L1:  t1 = -1
L2:  t2 = 0
L3:  t3 = 0
L4:  t4 = 4 * t3
L5:  t5 = 4 * t2
L6:  t6 = t5 * M
L7:  t7 = t4 + t6
L8:  t8 = a[t7]
L9:  if t8 <= max goto L11
L10: t1 = t8
L11: t3 = t3 + 1
L12: if t3 < M goto L4
L13: t2 = t2 + 1
L14: if t2 < N goto L3
L15: max = t1
\end{lstlisting}
Which one of the following options CORRECTLY specifies the number of basic blocks and the number of instructions in the largest basic block, respectively?

\par\vspace{6pt}\noindent
\begin{tabular}{@{}p{0.48\linewidth}@{\hspace{0.04\linewidth}}p{0.48\linewidth}@{}}
$\square$\ \textbf{(A)} 6 and 6 & $\square$\ \textbf{(B)} 6 and 7 \\
\addlinespace[3pt]
$\square$\ \textbf{(C)} 7 and 7 & $\square$\ \textbf{(D)} 7 and 6
\end{tabular}

\tcblower

\textbf{Official Key:} \textbf{\large D} \hfill \textit{\small [Medium]}
\vspace{2pt}\hrule height 0.3pt\vspace{3pt}
\textbf{Topic:} Compiler Design \\
\textbf{Subtopic:} \textit{Intermediate code generation}
\vspace{2pt}\hrule height 0.3pt\vspace{3pt}
\textbf{Guide / Hint:} \\
{\small Leaders are instructions that start a basic block: $L_1$ (first instruction), targets of conditional/unconditional jumps ($L_4, L_{11}, L_3$), and instructions immediately following jumps ($L_{10}, L_{13}, L_{15}$). This gives 7 leaders: $\{L_1, L_4, L_{10}, L_{11}, L_3, L_{13}, L_{15}\}$, forming 7 basic blocks. The block starting at $L_4$ contains lines $L_4$ through $L_9$, which has 6 instructions, the maximum in any block.}

\end{tcolorbox}
\vspace{6pt}\hrule height 0.35pt\vspace{6pt}
\noindent
\begin{tcolorbox}[
    blanker,
    sidebyside,
    sidebyside align=top,
    lefthand width=0.55\textwidth,
    top=4pt, bottom=4pt, left=0pt, right=0pt,
    middle=16pt,
    underlay={
        \draw[dashed, black!80, line width=0.8pt, dash pattern=on 4pt off 3pt]
            ([xshift=0.55\textwidth+8pt]frame.north west) --
            ([xshift=0.55\textwidth+8pt]frame.south west);
    }
]
\textbf{Q.159} \quad {\footnotesize\textsc{[MCQ $\cdot$ 1 Mark(s) $\cdot$ Neg: 0.33]}} \hfill {\footnotesize\textbf{GATE 2024 CS2} \quad \tiny\texttt{2024\_CS2\_CS\_Q21}}
\vspace{2pt}\hrule height 0.3pt\vspace{4pt}

Consider the following two sets:
\par\medskip
\begin{tabular}{ll}
\textbf{Set X} & \textbf{Set Y} \\
P. Lexical Analyzer & 1. Abstract Syntax Tree \\
Q. Syntax Analyzer & 2. Token \\
R. Intermediate Code Generator & 3. Parse Tree \\
S. Code Optimizer & 4. Constant Folding
\end{tabular}
\par\medskip
Which one of the following options is the CORRECT match from Set X to Set Y?

\par\vspace{6pt}\noindent
\begin{tabular}{@{}p{0.48\linewidth}@{\hspace{0.04\linewidth}}p{0.48\linewidth}@{}}
$\square$\ \textbf{(A)} P -- 4; Q -- 1; R -- 3; S -- 2 & $\square$\ \textbf{(B)} P -- 2; Q -- 3; R -- 1; S -- 4 \\
\addlinespace[3pt]
$\square$\ \textbf{(C)} P -- 2; Q -- 1; R -- 3; S -- 4 & $\square$\ \textbf{(D)} P -- 4; Q -- 3; R -- 2; S -- 1
\end{tabular}

\tcblower

\textbf{Official Key:} \textbf{\large B} \hfill \textit{\small [Easy]}
\vspace{2pt}\hrule height 0.3pt\vspace{3pt}
\textbf{Topic:} Compiler Design \\
\textbf{Subtopic:} \textit{Lexical analysis, parsing, syntax-directed translation}
\vspace{2pt}\hrule height 0.3pt\vspace{3pt}
\textbf{Guide / Hint:} \\
{\small Lexical Analyzer outputs tokens (P - 2). Syntax Analyzer constructs parse trees (Q - 3). Intermediate Code Generator generates abstract syntax trees or intermediate representations (R - 1). Code Optimizer performs constant folding (S - 4). Thus, the correct matching is P - 2; Q - 3; R - 1; S - 4.}

\end{tcolorbox}
\vspace{6pt}\hrule height 0.35pt\vspace{6pt}
\noindent
\begin{tcolorbox}[
    blanker,
    sidebyside,
    sidebyside align=top,
    lefthand width=0.55\textwidth,
    top=4pt, bottom=4pt, left=0pt, right=0pt,
    middle=16pt,
    underlay={
        \draw[dashed, black!80, line width=0.8pt, dash pattern=on 4pt off 3pt]
            ([xshift=0.55\textwidth+8pt]frame.north west) --
            ([xshift=0.55\textwidth+8pt]frame.south west);
    }
]
\textbf{Q.160} \quad {\footnotesize\textsc{[MCQ $\cdot$ 2 Mark(s) $\cdot$ Neg: 0.66]}} \hfill {\footnotesize\textbf{GATE 2024 CS2} \quad \tiny\texttt{2024\_CS2\_CS\_Q40}}
\vspace{2pt}\hrule height 0.3pt\vspace{4pt}

Consider the following context-free grammar where the start symbol is $S$ and the set of terminals is $\{a,b,c,d\}$:
\[ S \to AaAb \mid BbBa \]
\[ A \to cS \mid \epsilon \]
\[ B \to dS \mid \epsilon \]
The following is a partially-filled LL(1) parsing table:
\par\medskip
\begin{tabular}{|c|c|c|c|c|c|}
\hline
& $a$ & $b$ & $c$ & $d$ & \$ \\
\hline
$S$ & $S \to AaAb$ & $S \to BbBa$ & (1) & (2) & \\
$A$ & $A \to \epsilon$ & (3) & $A \to cS$ & & \\
$B$ & (4) & $B \to \epsilon$ & & $B \to dS$ & \\
\hline
\end{tabular}
\par\medskip
Which one of the following options represents the CORRECT combination for the numbered cells in the parsing table?

\par\vspace{6pt}\noindent
\begin{tabular}{@{}p{0.98\linewidth}@{}}
$\square$\ \textbf{(A)} (1) $S \to AaAb$ (2) $S \to BbBa$ (3) $A \to \epsilon$ (4) $B \to \epsilon$ \\
\addlinespace[3pt]
$\square$\ \textbf{(B)} (1) $S \to BbBa$ (2) $S \to AaAb$ (3) $A \to \epsilon$ (4) $B \to \epsilon$ \\
\addlinespace[3pt]
$\square$\ \textbf{(C)} (1) $S \to AaAb$ (2) $S \to BbBa$ (3) blank (4) blank \\
\addlinespace[3pt]
$\square$\ \textbf{(D)} (1) $S \to BbBa$ (2) $S \to AaAb$ (3) blank (4) blank
\end{tabular}

\tcblower

\textbf{Official Key:} \textbf{\large A} \hfill \textit{\small [Medium]}
\vspace{2pt}\hrule height 0.3pt\vspace{3pt}
\textbf{Topic:} Compiler Design \\
\textbf{Subtopic:} \textit{Lexical analysis, parsing, syntax-directed translation}
\vspace{2pt}\hrule height 0.3pt\vspace{3pt}
\textbf{Guide / Hint:} \\
{\small FIRST(A) = {c, $\epsilon$} and FIRST(B) = {d, $\epsilon$}. For $S \to AaAb$, FIRST($AaAb$) = {c, a}, so entry $M[S, c]$ gets $S \to AaAb$ (cell 1). For $S \to BbBa$, FIRST($BbBa$) = {d, b}, so entry $M[S, d]$ gets $S \to BbBa$ (cell 2). FOLLOW(A) = {a, b}, so $A \to \epsilon$ appears on $a$ and $b$ (cell 3). FOLLOW(B) = {b, a}, so $B \to \epsilon$ appears on $b$ and $a$ (cell 4).}

\end{tcolorbox}
\vspace{6pt}\hrule height 0.35pt\vspace{6pt}
\noindent
\begin{tcolorbox}[
    blanker,
    sidebyside,
    sidebyside align=top,
    lefthand width=0.55\textwidth,
    top=4pt, bottom=4pt, left=0pt, right=0pt,
    middle=16pt,
    underlay={
        \draw[dashed, black!80, line width=0.8pt, dash pattern=on 4pt off 3pt]
            ([xshift=0.55\textwidth+8pt]frame.north west) --
            ([xshift=0.55\textwidth+8pt]frame.south west);
    }
]
\textbf{Q.161} \quad {\footnotesize\textsc{[MCQ $\cdot$ 2 Mark(s) $\cdot$ Neg: 0.66]}} \hfill {\footnotesize\textbf{GATE 2024 CS2} \quad \tiny\texttt{2024\_CS2\_CS\_Q43}}
\vspace{2pt}\hrule height 0.3pt\vspace{4pt}

Consider the following expression: $x[i] = (p + r) * -s[i] + u / w$. The following sequence shows the list of triples representing the given expression, with entries missing for triples (1), (3), and (6):
\begin{verbatim}
(0) + p r
(1)
(2) uminus (1)
(3)
(4) / u w
(5) + (3) (4)
(6)
(7) = (6) (5)
\end{verbatim}
Which one of the following options fills in the missing entries CORRECTLY?

\par\vspace{6pt}\noindent
\begin{tabular}{@{}p{0.98\linewidth}@{}}
$\square$\ \textbf{(A)} (1) =[] s i; (3) * (0) (2); (6) []= x i \\
\addlinespace[3pt]
$\square$\ \textbf{(B)} (1) []= s i; (3) - (0) (2); (6) =[] x (5) \\
\addlinespace[3pt]
$\square$\ \textbf{(C)} (1) =[] s i; (3) * (0) (2); (6) []= x (5) \\
\addlinespace[3pt]
$\square$\ \textbf{(D)} (1) []= s i; (3) - (0) (2); (6) =[] x i
\end{tabular}

\tcblower

\textbf{Official Key:} \textbf{\large A} \hfill \textit{\small [Medium]}
\vspace{2pt}\hrule height 0.3pt\vspace{3pt}
\textbf{Topic:} Compiler Design \\
\textbf{Subtopic:} \textit{Intermediate code generation}
\vspace{2pt}\hrule height 0.3pt\vspace{3pt}
\textbf{Guide / Hint:} \\
{\small In compiler triples for array references: array indexing $s[i]$ uses index operation `=[] s i` at triple (1). Triple (2) applies unary minus to (1): `uminus (1)`. Triple (3) multiplies $(p + r)$ at (0) with $-s[i]$ at (2): `* (0) (2)`. For assignment to array $x[i]$, triple (6) performs indexed target `[]= x i`, and triple (7) stores value (5) into (6).}

\end{tcolorbox}
\vspace{6pt}\hrule height 0.35pt\vspace{6pt}
\noindent
\begin{tcolorbox}[
    blanker,
    sidebyside,
    sidebyside align=top,
    lefthand width=0.55\textwidth,
    top=4pt, bottom=4pt, left=0pt, right=0pt,
    middle=16pt,
    underlay={
        \draw[dashed, black!80, line width=0.8pt, dash pattern=on 4pt off 3pt]
            ([xshift=0.55\textwidth+8pt]frame.north west) --
            ([xshift=0.55\textwidth+8pt]frame.south west);
    }
]
\textbf{Q.162} \quad {\footnotesize\textsc{[MCQ $\cdot$ 1 Mark(s) $\cdot$ Neg: 0.33]}} \hfill {\footnotesize\textbf{GATE 2025 CS1} \quad \tiny\texttt{2025\_CS1\_CS\_Q12}}
\vspace{2pt}\hrule height 0.3pt\vspace{4pt}

Which ONE of the following statements is FALSE regarding the symbol table?

\par\vspace{6pt}\noindent
\begin{tabular}{@{}p{0.98\linewidth}@{}}
$\square$\ \textbf{(A)} Symbol table is responsible for keeping track of the scope of variables. \\
\addlinespace[3pt]
$\square$\ \textbf{(B)} Symbol table can be implemented using a binary search tree. \\
\addlinespace[3pt]
$\square$\ \textbf{(C)} Symbol table is not required after the parsing phase. \\
\addlinespace[3pt]
$\square$\ \textbf{(D)} Symbol table is created during the lexical analysis phase.
\end{tabular}

\tcblower

\textbf{Official Key:} \textbf{\large C} \hfill \textit{\small [Easy]}
\vspace{2pt}\hrule height 0.3pt\vspace{3pt}
\textbf{Topic:} Compiler Design \\
\textbf{Subtopic:} \textit{Lexical analysis, parsing, syntax-directed translation}
\vspace{2pt}\hrule height 0.3pt\vspace{3pt}
\textbf{Guide / Hint:} \\
{\small The symbol table is actively used throughout all phases of the compiler, including semantic analysis, intermediate code generation, code optimization, and target code generation for type checking, memory allocation, and address resolution. Thus, claiming it is not required after parsing is false.}

\end{tcolorbox}
\vspace{6pt}\hrule height 0.35pt\vspace{6pt}
\noindent
\begin{tcolorbox}[
    blanker,
    sidebyside,
    sidebyside align=top,
    lefthand width=0.55\textwidth,
    top=4pt, bottom=4pt, left=0pt, right=0pt,
    middle=16pt,
    underlay={
        \draw[dashed, black!80, line width=0.8pt, dash pattern=on 4pt off 3pt]
            ([xshift=0.55\textwidth+8pt]frame.north west) --
            ([xshift=0.55\textwidth+8pt]frame.south west);
    }
]
\textbf{Q.163} \quad {\footnotesize\textsc{[MCQ $\cdot$ 1 Mark(s) $\cdot$ Neg: 0.33]}} \hfill {\footnotesize\textbf{GATE 2025 CS1} \quad \tiny\texttt{2025\_CS1\_CS\_Q13}}
\vspace{2pt}\hrule height 0.3pt\vspace{4pt}

Which ONE of the following techniques used in compiler code optimization uses live variable analysis?

\par\vspace{6pt}\noindent
\begin{tabular}{@{}p{0.98\linewidth}@{}}
$\square$\ \textbf{(A)} Run-time function call management \\
\addlinespace[3pt]
$\square$\ \textbf{(B)} Register assignment to variables \\
\addlinespace[3pt]
$\square$\ \textbf{(C)} Strength reduction \\
\addlinespace[3pt]
$\square$\ \textbf{(D)} Constant folding
\end{tabular}

\tcblower

\textbf{Official Key:} \textbf{\large B} \hfill \textit{\small [Easy]}
\vspace{2pt}\hrule height 0.3pt\vspace{3pt}
\textbf{Topic:} Compiler Design \\
\textbf{Subtopic:} \textit{Data flow analyses: constant propagation, liveness analysis, common sub expression elimination}
\vspace{2pt}\hrule height 0.3pt\vspace{3pt}
\textbf{Guide / Hint:} \\
{\small Live variable analysis determines whether the value in a variable could be read along some execution path before being redefined. This information is directly used in register allocation and assignment to construct interference graphs and allocate registers efficiently.}

\end{tcolorbox}
\vspace{6pt}\hrule height 0.35pt\vspace{6pt}
\noindent
\begin{tcolorbox}[
    blanker,
    sidebyside,
    sidebyside align=top,
    lefthand width=0.55\textwidth,
    top=4pt, bottom=4pt, left=0pt, right=0pt,
    middle=16pt,
    underlay={
        \draw[dashed, black!80, line width=0.8pt, dash pattern=on 4pt off 3pt]
            ([xshift=0.55\textwidth+8pt]frame.north west) --
            ([xshift=0.55\textwidth+8pt]frame.south west);
    }
]
\textbf{Q.164} \quad {\footnotesize\textsc{[MCQ $\cdot$ 1 Mark(s) $\cdot$ Neg: 0.33]}} \hfill {\footnotesize\textbf{GATE 2025 CS2} \quad \tiny\texttt{2025\_CS2\_CS\_Q21}}
\vspace{2pt}\hrule height 0.3pt\vspace{4pt}

Consider the following statements about the use of backpatching in a compiler for intermediate code generation:
(I) Backpatching can be used to generate code for Boolean expression in one pass.
(II) Backpatching can be used to generate code for flow-of-control statements in one pass.
Which ONE of the following options is CORRECT?

\par\vspace{6pt}\noindent
\begin{tabular}{@{}p{0.98\linewidth}@{}}
$\square$\ \textbf{(A)} Only (I) is correct. \\
\addlinespace[3pt]
$\square$\ \textbf{(B)} Only (II) is correct. \\
\addlinespace[3pt]
$\square$\ \textbf{(C)} Both (I) and (II) are correct. \\
\addlinespace[3pt]
$\square$\ \textbf{(D)} Neither (I) nor (II) is correct.
\end{tabular}

\tcblower

\textbf{Official Key:} \textbf{\large C} \hfill \textit{\small [Easy]}
\vspace{2pt}\hrule height 0.3pt\vspace{3pt}
\textbf{Topic:} Compiler Design \\
\textbf{Subtopic:} \textit{Intermediate code generation and backpatching}
\vspace{2pt}\hrule height 0.3pt\vspace{3pt}
\textbf{Guide / Hint:} \\
{\small Backpatching is a standard technique used in single-pass compilers to synthesize target labels for forward jumps in both Boolean expressions and flow-of-control statements.}

\end{tcolorbox}
\vspace{6pt}\hrule height 0.35pt\vspace{6pt}
\noindent
\begin{tcolorbox}[
    blanker,
    sidebyside,
    sidebyside align=top,
    lefthand width=0.55\textwidth,
    top=4pt, bottom=4pt, left=0pt, right=0pt,
    middle=16pt,
    underlay={
        \draw[dashed, black!80, line width=0.8pt, dash pattern=on 4pt off 3pt]
            ([xshift=0.55\textwidth+8pt]frame.north west) --
            ([xshift=0.55\textwidth+8pt]frame.south west);
    }
]
\textbf{Q.165} \quad {\footnotesize\textsc{[MCQ $\cdot$ 1 Mark(s) $\cdot$ Neg: 0.33]}} \hfill {\footnotesize\textbf{GATE 2025 CS2} \quad \tiny\texttt{2025\_CS2\_CS\_Q22}}
\vspace{2pt}\hrule height 0.3pt\vspace{4pt}

Given the following syntax directed translation rules:
- Rule 1: $R \to AB \;\{ B.i = R.i - 1; \; A.i = B.i; \; R.i = A.i + 1; \}$
- Rule 2: $P \to CD \;\{ P.i = C.i + D.i; \; D.i = C.i + 2; \}$
- Rule 3: $Q \to EF \;\{ Q.i = E.i + F.i; \}$
\par\medskip
Which ONE is the CORRECT option among the following?

\par\vspace{6pt}\noindent
\begin{tabular}{@{}p{0.98\linewidth}@{}}
$\square$\ \textbf{(A)} Rule 1 is S-attributed and L-attributed; Rule 2 is S-attributed and not L-attributed; Rule 3 is neither S-attributed nor L-attributed \\
\addlinespace[3pt]
$\square$\ \textbf{(B)} Rule 1 is neither S-attributed nor L-attributed; Rule 2 is S-attributed and L-attributed; Rule 3 is S-attributed and L-attributed \\
\addlinespace[3pt]
$\square$\ \textbf{(C)} Rule 1 is neither S-attributed nor L-attributed; Rule 2 is not S-attributed and is L-attributed; Rule 3 is S-attributed and L-attributed \\
\addlinespace[3pt]
$\square$\ \textbf{(D)} Rule 1 is S-attributed and not L-attributed; Rule 2 is not S-attributed and is L-attributed; Rule 3 is S-attributed and L-attributed
\end{tabular}

\tcblower

\textbf{Official Key:} \textbf{\large C} \hfill \textit{\small [Hard]}
\vspace{2pt}\hrule height 0.3pt\vspace{3pt}
\textbf{Topic:} Compiler Design \\
\textbf{Subtopic:} \textit{Syntax-directed translation: S-attributed and L-attributed definitions}
\vspace{2pt}\hrule height 0.3pt\vspace{3pt}
\textbf{Guide / Hint:} \\
{\small In Rule 1, $R.i$ is computed from children (synthesized-like) but defined circularly with inherited attributes, making it neither. In Rule 2, attributes depend only on left siblings, satisfying L-attributed definition. In Rule 3, $Q.i$ depends only on children, which is S-attributed (and hence L-attributed).}

\end{tcolorbox}
\vspace{6pt}\hrule height 0.35pt\vspace{6pt}
\noindent
\begin{tcolorbox}[
    blanker,
    sidebyside,
    sidebyside align=top,
    lefthand width=0.55\textwidth,
    top=4pt, bottom=4pt, left=0pt, right=0pt,
    middle=16pt,
    underlay={
        \draw[dashed, black!80, line width=0.8pt, dash pattern=on 4pt off 3pt]
            ([xshift=0.55\textwidth+8pt]frame.north west) --
            ([xshift=0.55\textwidth+8pt]frame.south west);
    }
]
\textbf{Q.166} \quad {\footnotesize\textsc{[MCQ $\cdot$ 2 Mark(s) $\cdot$ Neg: 0.66]}} \hfill {\footnotesize\textbf{GATE 2025 CS2} \quad \tiny\texttt{2025\_CS2\_CS\_Q40}}
\vspace{2pt}\hrule height 0.3pt\vspace{4pt}

Given a Context-Free Grammar $G$:
$$\begin{aligned} S &\to Aa \mid bAc \mid dc \mid bda \\ A &\to d \end{aligned}$$
Which ONE of the following statements is CORRECT?

\par\vspace{6pt}\noindent
\begin{tabular}{@{}p{0.98\linewidth}@{}}
$\square$\ \textbf{(A)} $G$ is neither LALR(1) nor SLR(1) \\
\addlinespace[3pt]
$\square$\ \textbf{(B)} $G$ is CLR(1), not LALR(1) \\
\addlinespace[3pt]
$\square$\ \textbf{(C)} $G$ is LALR(1), not SLR(1) \\
\addlinespace[3pt]
$\square$\ \textbf{(D)} $G$ is LALR(1), also SLR(1)
\end{tabular}

\tcblower

\textbf{Official Key:} \textbf{\large C} \hfill \textit{\small [Hard]}
\vspace{2pt}\hrule height 0.3pt\vspace{3pt}
\textbf{Topic:} Compiler Design \\
\textbf{Subtopic:} \textit{Parsing: LR, LALR, and SLR parsers}
\vspace{2pt}\hrule height 0.3pt\vspace{3pt}
\textbf{Guide / Hint:} \\
{\small In SLR(1), FOLLOW sets create reduce-reduce conflicts on lookahead $c$. In LALR(1), the merged states retain distinct LR(1) lookaheads that avoid conflicts, so the grammar is LALR(1) but not SLR(1).}

\end{tcolorbox}
\vspace{6pt}\hrule height 0.35pt\vspace{6pt}
\vspace{8pt}
\subsection*{Operating Systems}
\vspace{-4pt}
\noindent
\begin{tcolorbox}[
    blanker,
    sidebyside,
    sidebyside align=top,
    lefthand width=0.55\textwidth,
    top=4pt, bottom=4pt, left=0pt, right=0pt,
    middle=16pt,
    underlay={
        \draw[dashed, black!80, line width=0.8pt, dash pattern=on 4pt off 3pt]
            ([xshift=0.55\textwidth+8pt]frame.north west) --
            ([xshift=0.55\textwidth+8pt]frame.south west);
    }
]
\textbf{Q.167} \quad {\footnotesize\textsc{[MCQ $\cdot$ 1 Mark(s) $\cdot$ Neg: 0.33]}} \hfill {\footnotesize\textbf{GATE 2018 CS} \quad \tiny\texttt{2018\_CS\_GA\_Q10\_v2}}
\vspace{2pt}\hrule height 0.3pt\vspace{4pt}

Consider a process executing on an operating system that uses demand paging. The average time for a memory access in the system is $M$ units if the corresponding memory page is available in memory, and $D$ units $(D \gg M)$ if the page needs to be fetched from the disk. If $X$ is the effective memory access time, the page fault rate $p$ is

\par\vspace{6pt}\noindent
\begin{tabular}{@{}p{0.48\linewidth}@{\hspace{0.04\linewidth}}p{0.48\linewidth}@{}}
$\square$\ \textbf{(A)} $\frac{X - M}{D}$ & $\square$\ \textbf{(B)} $\frac{X - M}{D - M}$ \\
\addlinespace[3pt]
$\square$\ \textbf{(C)} $\frac{D - X}{D - M}$ & $\square$\ \textbf{(D)} $\frac{X - M}{D + M}$
\end{tabular}

\tcblower

\textbf{Official Key:} \textbf{\large B} \hfill \textit{\small [Medium]}
\vspace{2pt}\hrule height 0.3pt\vspace{3pt}
\textbf{Topic:} Operating System \\
\textbf{Subtopic:} \textit{Memory management and virtual memory}
\vspace{2pt}\hrule height 0.3pt\vspace{3pt}
\textbf{Guide / Hint:} \\
{\small The effective access time is $X = (1 - p)M + pD = M + p(D - M)$. Solving for $p$ yields $p(D - M) = X - M \implies p = \frac{X - M}{D - M}$.}

\end{tcolorbox}
\vspace{6pt}\hrule height 0.35pt\vspace{6pt}
\noindent
\begin{tcolorbox}[
    blanker,
    sidebyside,
    sidebyside align=top,
    lefthand width=0.55\textwidth,
    top=4pt, bottom=4pt, left=0pt, right=0pt,
    middle=16pt,
    underlay={
        \draw[dashed, black!80, line width=0.8pt, dash pattern=on 4pt off 3pt]
            ([xshift=0.55\textwidth+8pt]frame.north west) --
            ([xshift=0.55\textwidth+8pt]frame.south west);
    }
]
\textbf{Q.168} \quad {\footnotesize\textsc{[MCQ $\cdot$ 2 Mark(s) $\cdot$ Neg: 0.66]}} \hfill {\footnotesize\textbf{GATE 2018 CS} \quad \tiny\texttt{2018\_CS\_GA\_Q39}}
\vspace{2pt}\hrule height 0.3pt\vspace{4pt}

In a system, there are three types of resources: $E, F$ and $G$. Four processes $P_0, P_1, P_2$ and $P_3$ execute concurrently. At time $t_0$, the maximum need and current allocation of resources are:
\begin{center}
\begin{tabular}{c|ccc|ccc}
\textbf{Process} & \multicolumn{3}{c|}{\textbf{Allocation}} & \multicolumn{3}{c}{\textbf{Max Need}} \\
& $E$ & $F$ & $G$ & $E$ & $F$ & $G$ \\
\hline
$P_0$ & 2 & 0 & 1 & 5 & 3 & 1 \\
$P_1$ & 0 & 1 & 0 & 1 & 2 & 1 \\
$P_2$ & 0 & 0 & 1 & 0 & 3 & 1 \\
$P_3$ & 1 & 0 & 0 & 3 & 2 & 1 \\
\end{tabular}
\end{center}
Total resources are: $E = 6, F = 4, G = 2$. The system at time $t_0$ is in

\par\vspace{6pt}\noindent
\begin{tabular}{@{}p{0.98\linewidth}@{}}
$\square$\ \textbf{(A)} a safe state. \\
\addlinespace[3pt]
$\square$\ \textbf{(B)} an unsafe state, but not in a deadlock. \\
\addlinespace[3pt]
$\square$\ \textbf{(C)} a deadlock. \\
\addlinespace[3pt]
$\square$\ \textbf{(D)} none of the above.
\end{tabular}

\tcblower

\textbf{Official Key:} \textbf{\large A} \hfill \textit{\small [Medium]}
\vspace{2pt}\hrule height 0.3pt\vspace{3pt}
\textbf{Topic:} Operating System \\
\textbf{Subtopic:} \textit{Deadlock}
\vspace{2pt}\hrule height 0.3pt\vspace{3pt}
\textbf{Guide / Hint:} \\
{\small Currently Allocated: $E = 2+0+0+1 = 3$, $F = 0+1+0+0 = 1$, $G = 1+0+1+0 = 2$. Available resources $= (6-3, 4-1, 2-2) = (3, 3, 0)$. Need matrix: $P_0(3, 3, 0)$, $P_1(1, 1, 1)$, $P_2(0, 3, 0)$, $P_3(2, 2, 1)$. Since $\text{Need}(P_0) \le \text{Available}$, $P_0$ can finish, releasing its allocation $(2, 0, 1)$ to make Available $(5, 3, 1)$. Next $P_2, P_1, P_3$ can all execute. Thus, the system is in a safe state.}

\end{tcolorbox}
\vspace{6pt}\hrule height 0.35pt\vspace{6pt}
\noindent
\begin{tcolorbox}[
    blanker,
    sidebyside,
    sidebyside align=top,
    lefthand width=0.55\textwidth,
    top=4pt, bottom=4pt, left=0pt, right=0pt,
    middle=16pt,
    underlay={
        \draw[dashed, black!80, line width=0.8pt, dash pattern=on 4pt off 3pt]
            ([xshift=0.55\textwidth+8pt]frame.north west) --
            ([xshift=0.55\textwidth+8pt]frame.south west);
    }
]
\textbf{Q.169} \quad {\footnotesize\textsc{[MCQ $\cdot$ 2 Mark(s) $\cdot$ Neg: 0.66]}} \hfill {\footnotesize\textbf{GATE 2018 CS} \quad \tiny\texttt{2018\_CS\_GA\_Q40}}
\vspace{2pt}\hrule height 0.3pt\vspace{4pt}

Consider the following solution to the producer-consumer synchronization problem. The shared buffer size is $N$. Three semaphores $empty, full$ and $mutex$ are used with initial values $N, 0$ and $1$, respectively.
\begin{lstlisting}[language=C]
// Producer:
while (1) {
    produce();
    wait(P);
    wait(mutex);
    put_item();
    signal(mutex);
    signal(Q);
}
\par\medskip
// Consumer:
while (1) {
    wait(R);
    wait(mutex);
    remove_item();
    signal(mutex);
    signal(S);
    consume();
}
\end{lstlisting}
Which of the following choices for $P, Q, R, S$ makes the solution correct?

\par\vspace{6pt}\noindent
\begin{tabular}{@{}p{0.98\linewidth}@{}}
$\square$\ \textbf{(A)} $P: empty, Q: full, R: empty, S: full$ \\
\addlinespace[3pt]
$\square$\ \textbf{(B)} $P: empty, Q: full, R: full, S: empty$ \\
\addlinespace[3pt]
$\square$\ \textbf{(C)} $P: full, Q: empty, R: empty, S: full$ \\
\addlinespace[3pt]
$\square$\ \textbf{(D)} $P: full, Q: empty, R: full, S: empty$
\end{tabular}

\tcblower

\textbf{Official Key:} \textbf{\large C} \hfill \textit{\small [Medium]}
\vspace{2pt}\hrule height 0.3pt\vspace{3pt}
\textbf{Topic:} Operating System \\
\textbf{Subtopic:} \textit{Concurrency and synchronization}
\vspace{2pt}\hrule height 0.3pt\vspace{3pt}
\textbf{Guide / Hint:} \\
{\small In the standard producer-consumer problem with a buffer of size $N$, `empty` is initialized to $N$ and `full` is initialized to $0$. The producer waits on `empty` and signals `full`. Conversely, the consumer waits on `full` and signals `empty`. Comparing this with the provided code reveals that semaphore `empty` corresponds to buffer spaces and `full` corresponds to items, making $P: \text{full}, Q: \text{empty}, R: \text{empty}, S: \text{full}$.}

\end{tcolorbox}
\vspace{6pt}\hrule height 0.35pt\vspace{6pt}
\noindent
\begin{tcolorbox}[
    blanker,
    sidebyside,
    sidebyside align=top,
    lefthand width=0.55\textwidth,
    top=4pt, bottom=4pt, left=0pt, right=0pt,
    middle=16pt,
    underlay={
        \draw[dashed, black!80, line width=0.8pt, dash pattern=on 4pt off 3pt]
            ([xshift=0.55\textwidth+8pt]frame.north west) --
            ([xshift=0.55\textwidth+8pt]frame.south west);
    }
]
\textbf{Q.170} \quad {\footnotesize\textsc{[MCQ $\cdot$ 1 Mark(s) $\cdot$ Neg: 0.33]}} \hfill {\footnotesize\textbf{GATE 2022 CS} \quad \tiny\texttt{2022\_CS\_CS\_Q19}}
\vspace{2pt}\hrule height 0.3pt\vspace{4pt}

Consider the following threads, T1, T2, and T3 executing on a single processor, synchronized using three binary semaphore variables, S1, S2, and S3, operated upon using standard `wait()` and `signal()`:
\begin{center}
\begin{tabular}{p{3.5cm}p{3.5cm}p{3.5cm}}
\textbf{T1} & \textbf{T2} & \textbf{T3} \\
\texttt{while(true)\{} & \texttt{while(true)\{} & \texttt{while(true)\{} \\
\texttt{\ \ wait(S3);} & \texttt{\ \ wait(S1);} & \texttt{\ \ wait(S2);} \\
\texttt{\ \ print("C");} & \texttt{\ \ print("B");} & \texttt{\ \ print("A");} \\
\texttt{\ \ signal(S2);\}} & \texttt{\ \ signal(S3);\}} & \texttt{\ \ signal(S1);\}}
\end{tabular}
\end{center}
Which initialization of the semaphores would print the sequence BCABCABCA\dots?

\par\vspace{6pt}\noindent
\begin{tabular}{@{}p{0.48\linewidth}@{\hspace{0.04\linewidth}}p{0.48\linewidth}@{}}
$\square$\ \textbf{(A)} $S_1 = 1; S_2 = 1; S_3 = 1$ & $\square$\ \textbf{(B)} $S_1 = 1; S_2 = 1; S_3 = 0$ \\
\addlinespace[3pt]
$\square$\ \textbf{(C)} $S_1 = 1; S_2 = 0; S_3 = 0$ & $\square$\ \textbf{(D)} $S_1 = 0; S_2 = 1; S_3 = 1$
\end{tabular}

\tcblower

\textbf{Official Key:} \textbf{\large C} \hfill \textit{\small [Medium]}
\vspace{2pt}\hrule height 0.3pt\vspace{3pt}
\textbf{Topic:} Operating System \\
\textbf{Subtopic:} \textit{Concurrency and synchronization}
\vspace{2pt}\hrule height 0.3pt\vspace{3pt}
\textbf{Guide / Hint:} \\
{\small To output BCABCABCA..., the sequence begins with 'B' from T2, which requires $S_1 = 1$. After printing 'B', T2 signals $S_3$, allowing T1 to print 'C'. T1 then signals $S_2$, allowing T3 to print 'A'. Finally, T3 signals $S_1$, enabling T2 again. Therefore, the required initial semaphore values are $S_1 = 1, S_2 = 0, S_3 = 0$.}

\end{tcolorbox}
\vspace{6pt}\hrule height 0.35pt\vspace{6pt}
\noindent
\begin{tcolorbox}[
    blanker,
    sidebyside,
    sidebyside align=top,
    lefthand width=0.55\textwidth,
    top=4pt, bottom=4pt, left=0pt, right=0pt,
    middle=16pt,
    underlay={
        \draw[dashed, black!80, line width=0.8pt, dash pattern=on 4pt off 3pt]
            ([xshift=0.55\textwidth+8pt]frame.north west) --
            ([xshift=0.55\textwidth+8pt]frame.south west);
    }
]
\textbf{Q.171} \quad {\footnotesize\textsc{[MCQ $\cdot$ 2 Mark(s) $\cdot$ Neg: 0.66]}} \hfill {\footnotesize\textbf{GATE 2022 CS} \quad \tiny\texttt{2022\_CS\_CS\_Q38}}
\vspace{2pt}\hrule height 0.3pt\vspace{4pt}

Which one of the following statements is FALSE?

\par\vspace{6pt}\noindent
\begin{tabular}{@{}p{0.98\linewidth}@{}}
$\square$\ \textbf{(A)} The TLB performs an associative search in parallel on all its valid entries using page number of incoming virtual address. \\
\addlinespace[3pt]
$\square$\ \textbf{(B)} If the virtual address of a word given by CPU has a TLB hit, but the subsequent search for the word results in a cache miss, then the word will always be present in the main memory. \\
\addlinespace[3pt]
$\square$\ \textbf{(C)} The memory access time using a given inverted page table is always same for all incoming virtual addresses. \\
\addlinespace[3pt]
$\square$\ \textbf{(D)} In a system that uses hashed page tables, if two distinct virtual addresses V1 and V2 map to the same value while hashing, then the memory access time of these addresses will not be the same.
\end{tabular}

\tcblower

\textbf{Official Key:} \textbf{\large C} \hfill \textit{\small [Medium]}
\vspace{2pt}\hrule height 0.3pt\vspace{3pt}
\textbf{Topic:} Operating System \\
\textbf{Subtopic:} \textit{Memory management and virtual memory}
\vspace{2pt}\hrule height 0.3pt\vspace{3pt}
\textbf{Guide / Hint:} \\
{\small An inverted page table organizes entries indexed by physical frames using a hash table with collision chaining. Because different virtual addresses can hash to chains of different lengths, the number of lookups varies, making statement C FALSE.}

\end{tcolorbox}
\vspace{6pt}\hrule height 0.35pt\vspace{6pt}
\noindent
\begin{tcolorbox}[
    blanker,
    sidebyside,
    sidebyside align=top,
    lefthand width=0.55\textwidth,
    top=4pt, bottom=4pt, left=0pt, right=0pt,
    middle=16pt,
    underlay={
        \draw[dashed, black!80, line width=0.8pt, dash pattern=on 4pt off 3pt]
            ([xshift=0.55\textwidth+8pt]frame.north west) --
            ([xshift=0.55\textwidth+8pt]frame.south west);
    }
]
\textbf{Q.172} \quad {\footnotesize\textsc{[MCQ $\cdot$ 2 Mark(s) $\cdot$ Neg: 0.66]}} \hfill {\footnotesize\textbf{GATE 2022 CS} \quad \tiny\texttt{2022\_CS\_CS\_Q42}}
\vspace{2pt}\hrule height 0.3pt\vspace{4pt}

Consider four processes P, Q, R, and S scheduled on a CPU as per round robin algorithm with a time quantum of 4 units. The processes arrive in the order P, Q, R, S, all at time $t = 0$. There is exactly one context switch from S to Q, exactly one context switch from R to Q, and exactly two context switches from Q to R. There is no context switch from S to P. Switching to a ready process after the termination of another process is also considered a context switch. Which one of the following is NOT possible as CPU burst time (in time units) of these processes?

\par\vspace{6pt}\noindent
\begin{tabular}{@{}p{0.48\linewidth}@{\hspace{0.04\linewidth}}p{0.48\linewidth}@{}}
$\square$\ \textbf{(A)} P = 4, Q = 10, R = 6, S = 2 & $\square$\ \textbf{(B)} P = 2, Q = 9, R = 5, S = 1 \\
\addlinespace[3pt]
$\square$\ \textbf{(C)} P = 4, Q = 12, R = 5, S = 4 & $\square$\ \textbf{(D)} P = 3, Q = 7, R = 7, S = 3
\end{tabular}

\tcblower

\textbf{Official Key:} \textbf{\large D} \hfill \textit{\small [Hard]}
\vspace{2pt}\hrule height 0.3pt\vspace{3pt}
\textbf{Topic:} Operating System \\
\textbf{Subtopic:} \textit{CPU and I/O scheduling}
\vspace{2pt}\hrule height 0.3pt\vspace{3pt}
\textbf{Guide / Hint:} \\
{\small With Round Robin scheduling and time quantum 4 for processes arriving as P, Q, R, S: in option D ($P=3, Q=7, R=7, S=3$), execution proceeds $P(3) \to Q(4) \to R(4) \to S(3) \to Q(3) \to R(3)$. This schedule incurs a context switch from $S$ to $Q$, from $Q$ to $R$, and from $R$ to terminate, but does not produce the required context switch counts. Hence, burst combination D is NOT possible.}

\end{tcolorbox}
\vspace{6pt}\hrule height 0.35pt\vspace{6pt}
\noindent
\begin{tcolorbox}[
    blanker,
    sidebyside,
    sidebyside align=top,
    lefthand width=0.55\textwidth,
    top=4pt, bottom=4pt, left=0pt, right=0pt,
    middle=16pt,
    underlay={
        \draw[dashed, black!80, line width=0.8pt, dash pattern=on 4pt off 3pt]
            ([xshift=0.55\textwidth+8pt]frame.north west) --
            ([xshift=0.55\textwidth+8pt]frame.south west);
    }
]
\textbf{Q.173} \quad {\footnotesize\textsc{[MCQ $\cdot$ 2 Mark(s) $\cdot$ Neg: 0.66]}} \hfill {\footnotesize\textbf{GATE 2023 CS} \quad \tiny\texttt{2023\_CS\_CS\_Q38}}
\vspace{2pt}\hrule height 0.3pt\vspace{4pt}

Consider the two functions `incr` and `decr` shown below:
\begin{lstlisting}[language=C]
incr() {
    wait(s);
    X = X + 1;
    signal(s);
}
decr() {
    wait(s);
    X = X - 1;
    signal(s);
}
\end{lstlisting}
There are 5 threads each invoking `incr` once, and 3 threads each invoking `decr` once, on the same shared variable $X$. The initial value of $X$ is 10.
Suppose there are two implementations of the semaphore $s$, as follows:
\begin{itemize}[itemsep=0pt,topsep=2pt]
\item I-1: $s$ is a binary semaphore initialized to 1.
\item I-2: $s$ is a counting semaphore initialized to 2.
\end{itemize}
Let $V_1, V_2$ be the values of $X$ at the end of execution of all the threads with implementations I-1, I-2, respectively.
Which one of the following choices corresponds to the minimum possible values of $V_1, V_2$, respectively?

\par\vspace{6pt}\noindent
\begin{tabular}{@{}p{0.48\linewidth}@{\hspace{0.04\linewidth}}p{0.48\linewidth}@{}}
$\square$\ \textbf{(A)} 15, 7 & $\square$\ \textbf{(B)} 7, 7 \\
\addlinespace[3pt]
$\square$\ \textbf{(C)} 12, 7 & $\square$\ \textbf{(D)} 12, 8
\end{tabular}

\tcblower

\textbf{Official Key:} \textbf{\large C} \hfill \textit{\small [Medium]}
\vspace{2pt}\hrule height 0.3pt\vspace{3pt}
\textbf{Topic:} Operating System \\
\textbf{Subtopic:} \textit{Concurrency and synchronization}
\vspace{2pt}\hrule height 0.3pt\vspace{3pt}
\textbf{Guide / Hint:} \\
{\small With binary semaphore $s=1$ (I-1), mutual exclusion guarantees serialized critical section access, yielding net change $+5 - 3 = +2 \implies V_1 = 10 + 2 = 12$. With counting semaphore $s=2$ (I-2), two threads can enter simultaneously; a decrement thread can read $X=10$, be preempted while all 5 increments and 2 decrements complete, and then write back $10-1 = 9$, repeatedly clobbering increments to achieve minimum $V_2 = 10 - 3 = 7$.}

\end{tcolorbox}
\vspace{6pt}\hrule height 0.35pt\vspace{6pt}
\noindent
\begin{tcolorbox}[
    blanker,
    sidebyside,
    sidebyside align=top,
    lefthand width=0.55\textwidth,
    top=4pt, bottom=4pt, left=0pt, right=0pt,
    middle=16pt,
    underlay={
        \draw[dashed, black!80, line width=0.8pt, dash pattern=on 4pt off 3pt]
            ([xshift=0.55\textwidth+8pt]frame.north west) --
            ([xshift=0.55\textwidth+8pt]frame.south west);
    }
]
\textbf{Q.174} \quad {\footnotesize\textsc{[MCQ $\cdot$ 2 Mark(s) $\cdot$ Neg: 0.66]}} \hfill {\footnotesize\textbf{GATE 2024 CS1} \quad \tiny\texttt{2024\_CS1\_CS\_Q40}}
\vspace{2pt}\hrule height 0.3pt\vspace{4pt}

Consider the following two threads T1 and T2 that update two shared variables `a` and `b`. Assume that initially $a = b = 1$. Though context switching between threads can happen at any time, each statement of T1 or T2 is executed atomically without interruption.
\begin{center}
\begin{tabular}{p{4cm}p{4cm}}
\textbf{T1} & \textbf{T2} \\
\texttt{a = a + 1;} & \texttt{b = 2 * b;} \\
\texttt{b = b + 1;} & \texttt{a = 2 * a;}
\end{tabular}
\end{center}
Which one of the following options lists all the possible combinations of values of `a` and `b` after both T1 and T2 finish execution?

\par\vspace{6pt}\noindent
\begin{tabular}{@{}p{0.98\linewidth}@{}}
$\square$\ \textbf{(A)} $(a = 4, b = 4)$; $(a = 3, b = 3)$; $(a = 4, b = 3)$ \\
\addlinespace[3pt]
$\square$\ \textbf{(B)} $(a = 3, b = 4)$; $(a = 4, b = 3)$; $(a = 3, b = 3)$ \\
\addlinespace[3pt]
$\square$\ \textbf{(C)} $(a = 4, b = 4)$; $(a = 4, b = 3)$; $(a = 3, b = 4)$ \\
\addlinespace[3pt]
$\square$\ \textbf{(D)} $(a = 2, b = 2)$; $(a = 2, b = 3)$; $(a = 3, b = 4)$
\end{tabular}

\tcblower

\textbf{Official Key:} \textbf{\large A} \hfill \textit{\small [Medium]}
\vspace{2pt}\hrule height 0.3pt\vspace{3pt}
\textbf{Topic:} Operating System \\
\textbf{Subtopic:} \textit{Concurrency and synchronization}
\vspace{2pt}\hrule height 0.3pt\vspace{3pt}
\textbf{Guide / Hint:} \\
{\small Initially $a=b=1$. Let statements be $T_1: [S_1: a=a+1, S_2: b=b+1]$ and $T_2: [S_3: b=2b, S_4: a=2a]$. Executing $T_1$ then $T_2$ yields $a=2, b=2 \to b=4, a=4 \implies (4, 4)$. Executing $T_2$ then $T_1$ yields $b=2, a=2 \to a=3, b=3 \implies (3, 3)$. Executing $S_1, S_3, S_4, S_2$ yields $a=2 \to b=2 \to a=4 \to b=3 \implies (4, 3)$. No valid interleaving can produce $(a=3, b=4)$.}

\end{tcolorbox}
\vspace{6pt}\hrule height 0.35pt\vspace{6pt}
\noindent
\begin{tcolorbox}[
    blanker,
    sidebyside,
    sidebyside align=top,
    lefthand width=0.55\textwidth,
    top=4pt, bottom=4pt, left=0pt, right=0pt,
    middle=16pt,
    underlay={
        \draw[dashed, black!80, line width=0.8pt, dash pattern=on 4pt off 3pt]
            ([xshift=0.55\textwidth+8pt]frame.north west) --
            ([xshift=0.55\textwidth+8pt]frame.south west);
    }
]
\textbf{Q.175} \quad {\footnotesize\textsc{[MCQ $\cdot$ 2 Mark(s) $\cdot$ Neg: 0.66]}} \hfill {\footnotesize\textbf{GATE 2024 CS2} \quad \tiny\texttt{2024\_CS2\_CS\_Q37}}
\vspace{2pt}\hrule height 0.3pt\vspace{4pt}

Consider a single processor system with four processes A, B, C, and D, represented as given below, where for each process the first value is its arrival time, and the second value is its CPU burst time:
A (0, 10), B (2, 6), C (4, 3), and D (6, 7).
Which one of the following options gives the average waiting time for Shortest Remaining Time First (SRTF) and Non-Preemptive Shortest Job First (NP-SJF) scheduling?

\par\vspace{6pt}\noindent
\begin{tabular}{@{}p{0.48\linewidth}@{\hspace{0.04\linewidth}}p{0.48\linewidth}@{}}
$\square$\ \textbf{(A)} SRTF = 6, NP-SJF = 7 & $\square$\ \textbf{(B)} SRTF = 6, NP-SJF = 7.5 \\
\addlinespace[3pt]
$\square$\ \textbf{(C)} SRTF = 7, NP-SJF = 7.5 & $\square$\ \textbf{(D)} SRTF = 7, NP-SJF = 8.5
\end{tabular}

\tcblower

\textbf{Official Key:} \textbf{\large B} \hfill \textit{\small [Medium]}
\vspace{2pt}\hrule height 0.3pt\vspace{3pt}
\textbf{Topic:} Operating System \\
\textbf{Subtopic:} \textit{CPU and I/O scheduling}
\vspace{2pt}\hrule height 0.3pt\vspace{3pt}
\textbf{Guide / Hint:} \\
{\small Under SRTF: execution order is A(0-2), B(2-4), C(4-7), B(7-11), D(11-18), A(18-26). Turnaround times: A=26, B=9, C=3, D=12; Waiting times: A=16, B=3, C=0, D=5. Average waiting time $= 24/4 = 6.0$. Under NP-SJF: A runs non-preemptively (0-10), then C(10-13), B(13-19), D(19-26). Waiting times: A=0, C=6, B=11, D=13. Average waiting time $= 30/4 = 7.5$.}

\end{tcolorbox}
\vspace{6pt}\hrule height 0.35pt\vspace{6pt}
\noindent
\begin{tcolorbox}[
    blanker,
    sidebyside,
    sidebyside align=top,
    lefthand width=0.55\textwidth,
    top=4pt, bottom=4pt, left=0pt, right=0pt,
    middle=16pt,
    underlay={
        \draw[dashed, black!80, line width=0.8pt, dash pattern=on 4pt off 3pt]
            ([xshift=0.55\textwidth+8pt]frame.north west) --
            ([xshift=0.55\textwidth+8pt]frame.south west);
    }
]
\textbf{Q.176} \quad {\footnotesize\textsc{[MCQ $\cdot$ 1 Mark(s) $\cdot$ Neg: 0.33]}} \hfill {\footnotesize\textbf{GATE 2025 CS1} \quad \tiny\texttt{2025\_CS1\_CS\_Q14}}
\vspace{2pt}\hrule height 0.3pt\vspace{4pt}

Consider a demand paging memory management system with 32-bit logical address, 20-bit physical address, and page size of 2048 bytes. Assuming that the memory is byte addressable, what is the maximum number of entries in the page table?

\par\vspace{6pt}\noindent
\begin{tabular}{@{}p{0.48\linewidth}@{\hspace{0.04\linewidth}}p{0.48\linewidth}@{}}
$\square$\ \textbf{(A)} $2^{21}$ & $\square$\ \textbf{(B)} $2^{20}$ \\
\addlinespace[3pt]
$\square$\ \textbf{(C)} $2^{22}$ & $\square$\ \textbf{(D)} $2^{24}$
\end{tabular}

\tcblower

\textbf{Official Key:} \textbf{\large A} \hfill \textit{\small [Medium]}
\vspace{2pt}\hrule height 0.3pt\vspace{3pt}
\textbf{Topic:} Operating System \\
\textbf{Subtopic:} \textit{Memory management and virtual memory}
\vspace{2pt}\hrule height 0.3pt\vspace{3pt}
\textbf{Guide / Hint:} \\
{\small With a page size of $2048 = 2^{11}$ bytes, the page offset requires 11 bits. The 32-bit logical address is split into page number and page offset, so the page number has $32 - 11 = 21$ bits. Hence, the page table contains at most $2^{21}$ entries.}

\end{tcolorbox}
\vspace{6pt}\hrule height 0.35pt\vspace{6pt}
\noindent
\begin{tcolorbox}[
    blanker,
    sidebyside,
    sidebyside align=top,
    lefthand width=0.55\textwidth,
    top=4pt, bottom=4pt, left=0pt, right=0pt,
    middle=16pt,
    underlay={
        \draw[dashed, black!80, line width=0.8pt, dash pattern=on 4pt off 3pt]
            ([xshift=0.55\textwidth+8pt]frame.north west) --
            ([xshift=0.55\textwidth+8pt]frame.south west);
    }
]
\textbf{Q.177} \quad {\footnotesize\textsc{[MCQ $\cdot$ 2 Mark(s) $\cdot$ Neg: 0.66]}} \hfill {\footnotesize\textbf{GATE 2025 CS1} \quad \tiny\texttt{2025\_CS1\_CS\_Q38}}
\vspace{2pt}\hrule height 0.3pt\vspace{4pt}

A computer has two processors, $M_1$ and $M_2$. Four processes $P_1, P_2, P_3, P_4$ with CPU bursts of 20, 16, 25, and 10 milliseconds, respectively, arrive at the same time and these are the only processes in the system. The scheduler uses non-preemptive priority scheduling, with priorities decided as follows:
\begin{itemize}
\item $M_1$ uses priority of execution for the processes as, $P_1 > P_3 > P_2 > P_4$, i.e., $P_1$ and $P_4$ have highest and lowest priorities, respectively.
\item $M_2$ uses priority of execution for the processes as, $P_2 > P_3 > P_4 > P_1$, i.e., $P_2$ and $P_1$ have highest and lowest priorities, respectively.
\end{itemize}
A process $P_i$ is scheduled to a processor $M_k$, if the processor is free and no other process $P_j$ is waiting with higher priority. At any given point of time, a process can be allocated to any one of the free processors without violating the execution priority rules. Ignore the context switch time. What will be the average waiting time of the processes in milliseconds?

\par\vspace{6pt}\noindent
\begin{tabular}{@{}p{0.48\linewidth}@{\hspace{0.04\linewidth}}p{0.48\linewidth}@{}}
$\square$\ \textbf{(A)} 9.00 & $\square$\ \textbf{(B)} 8.75 \\
\addlinespace[3pt]
$\square$\ \textbf{(C)} 6.50 & $\square$\ \textbf{(D)} 7.50
\end{tabular}

\tcblower

\textbf{Official Key:} \textbf{\large A} \hfill \textit{\small [Hard]}
\vspace{2pt}\hrule height 0.3pt\vspace{3pt}
\textbf{Topic:} Operating System \\
\textbf{Subtopic:} \textit{CPU and I/O scheduling}
\vspace{2pt}\hrule height 0.3pt\vspace{3pt}
\textbf{Guide / Hint:} \\
{\small At $t = 0$, $M_1$ takes highest priority process $P_1$ (burst 20) and $M_2$ takes $P_2$ (burst 16). At $t = 16$, $M_2$ frees up; among remaining processes $\{P_3, P_4\}$, $P_3$ has higher priority on $M_2$, so $P_3$ runs on $M_2$ ($t = 16$ to 41). At $t = 20$, $M_1$ frees up and executes $P_4$ ($t = 20$ to 30). Waiting times are: $P_1 = 0$, $P_2 = 0$, $P_3 = 16$, $P_4 = 20$. Average waiting time is $(0 + 0 + 16 + 20) / 4 = 36 / 4 = 9.00$ ms.}

\end{tcolorbox}
\vspace{6pt}\hrule height 0.35pt\vspace{6pt}
\noindent
\begin{tcolorbox}[
    blanker,
    sidebyside,
    sidebyside align=top,
    lefthand width=0.55\textwidth,
    top=4pt, bottom=4pt, left=0pt, right=0pt,
    middle=16pt,
    underlay={
        \draw[dashed, black!80, line width=0.8pt, dash pattern=on 4pt off 3pt]
            ([xshift=0.55\textwidth+8pt]frame.north west) --
            ([xshift=0.55\textwidth+8pt]frame.south west);
    }
]
\textbf{Q.178} \quad {\footnotesize\textsc{[MCQ $\cdot$ 1 Mark(s) $\cdot$ Neg: 0.33]}} \hfill {\footnotesize\textbf{GATE 2025 CS2} \quad \tiny\texttt{2025\_CS2\_CS\_Q26}}
\vspace{2pt}\hrule height 0.3pt\vspace{4pt}

Processes $P_1, P_2, P_3, P_4$ arrive in that order at times 0, 1, 2, and 8 milliseconds respectively, and have execution times of 10, 13, 6, and 9 milliseconds respectively. Shortest Remaining Time First (SRTF) algorithm is used as the CPU scheduling policy. Ignore context switching times. Which ONE of the following correctly gives the average turnaround time of the four processes in milliseconds?

\par\vspace{6pt}\noindent
\begin{tabular}{@{}p{0.48\linewidth}@{\hspace{0.04\linewidth}}p{0.48\linewidth}@{}}
$\square$\ \textbf{(A)} 22 & $\square$\ \textbf{(B)} 15 \\
\addlinespace[3pt]
$\square$\ \textbf{(C)} 37 & $\square$\ \textbf{(D)} 19
\end{tabular}

\tcblower

\textbf{Official Key:} \textbf{\large D} \hfill \textit{\small [Medium]}
\vspace{2pt}\hrule height 0.3pt\vspace{3pt}
\textbf{Topic:} Operating System \\
\textbf{Subtopic:} \textit{CPU and I/O scheduling: SRTF}
\vspace{2pt}\hrule height 0.3pt\vspace{3pt}
\textbf{Guide / Hint:} \\
{\small Using preemptive SRTF: $P_1$ runs [0,2], $P_3$ runs [2,8] (finishes at 8, TAT = $8-2=6$). Then $P_4$ (remaining 9) runs [8,17] (finishes at 17, TAT = $17-8=9$). $P_1$ (remaining 8) runs [17,25] (finishes at 25, TAT = $25-0=25$). $P_2$ runs [25,38] (finishes at 38, TAT = $38-1=37$). Average TAT = $\frac{25 + 37 + 6 + 9}{4} = \frac{76}{4} = 19$ ms.}

\end{tcolorbox}
\vspace{6pt}\hrule height 0.35pt\vspace{6pt}
\vspace{8pt}
\subsection*{Databases}
\vspace{-4pt}
\noindent
\begin{tcolorbox}[
    blanker,
    sidebyside,
    sidebyside align=top,
    lefthand width=0.55\textwidth,
    top=4pt, bottom=4pt, left=0pt, right=0pt,
    middle=16pt,
    underlay={
        \draw[dashed, black!80, line width=0.8pt, dash pattern=on 4pt off 3pt]
            ([xshift=0.55\textwidth+8pt]frame.north west) --
            ([xshift=0.55\textwidth+8pt]frame.south west);
    }
]
\textbf{Q.179} \quad {\footnotesize\textsc{[MCQ $\cdot$ 1 Mark(s) $\cdot$ Neg: 0.33]}} \hfill {\footnotesize\textbf{GATE 2018 CS} \quad \tiny\texttt{2018\_CS\_GA\_Q11}}
\vspace{2pt}\hrule height 0.3pt\vspace{4pt}

In an Entity-Relationship (ER) model, suppose $R$ is a many-to-one relationship from entity set $E_1$ to entity set $E_2$. Assume that $E_1$ has total participation in $R$, but $E_2$ has partial participation in $R$. Which one of the following statements is necessarily TRUE?

\par\vspace{6pt}\noindent
\begin{tabular}{@{}p{0.98\linewidth}@{}}
$\square$\ \textbf{(A)} Every entity in $E_1$ is associated with exactly one entity in $E_2$. \\
\addlinespace[3pt]
$\square$\ \textbf{(B)} Some entity in $E_1$ is associated with more than one entity in $E_2$. \\
\addlinespace[3pt]
$\square$\ \textbf{(C)} Every entity in $E_2$ is associated with at least one entity in $E_1$. \\
\addlinespace[3pt]
$\square$\ \textbf{(D)} Every entity in $E_2$ is associated with at most one entity in $E_1$.
\end{tabular}

\tcblower

\textbf{Official Key:} \textbf{\large A} \hfill \textit{\small [Easy]}
\vspace{2pt}\hrule height 0.3pt\vspace{3pt}
\textbf{Topic:} Databases \\
\textbf{Subtopic:} \textit{ER-model}
\vspace{2pt}\hrule height 0.3pt\vspace{3pt}
\textbf{Guide / Hint:} \\
{\small In an ER diagram, a many-to-one relationship from $E_1$ to $E_2$ means each entity in $E_1$ is related to at most one entity in $E_2$. Total participation of $E_1$ requires that every entity in $E_1$ participates in at least one relationship. Combining these, every entity in $E_1$ is associated with exactly one entity in $E_2$.}

\end{tcolorbox}
\vspace{6pt}\hrule height 0.35pt\vspace{6pt}
\noindent
\begin{tcolorbox}[
    blanker,
    sidebyside,
    sidebyside align=top,
    lefthand width=0.55\textwidth,
    top=4pt, bottom=4pt, left=0pt, right=0pt,
    middle=16pt,
    underlay={
        \draw[dashed, black!80, line width=0.8pt, dash pattern=on 4pt off 3pt]
            ([xshift=0.55\textwidth+8pt]frame.north west) --
            ([xshift=0.55\textwidth+8pt]frame.south west);
    }
]
\textbf{Q.180} \quad {\footnotesize\textsc{[MCQ $\cdot$ 1 Mark(s) $\cdot$ Neg: 0.33]}} \hfill {\footnotesize\textbf{GATE 2018 CS} \quad \tiny\texttt{2018\_CS\_GA\_Q12}}
\vspace{2pt}\hrule height 0.3pt\vspace{4pt}

Consider the following two tables and four queries in SQL.
\texttt{Book (isbn, bname)}
\texttt{Stock (isbn, copies)}
\par\medskip
Query 1: \texttt{SELECT B.isbn, S.copies FROM Book B INNER JOIN Stock S ON B.isbn = S.isbn;}
Query 2: \texttt{SELECT B.isbn, S.copies FROM Book B LEFT OUTER JOIN Stock S ON B.isbn = S.isbn;}
Query 3: \texttt{SELECT B.isbn, S.copies FROM Book B RIGHT OUTER JOIN Stock S ON B.isbn = S.isbn;}
Query 4: \texttt{SELECT B.isbn, S.copies FROM Book B FULL OUTER JOIN Stock S ON B.isbn = S.isbn;}
\par\medskip
Which one of the following statements is CORRECT?

\par\vspace{6pt}\noindent
\begin{tabular}{@{}p{0.98\linewidth}@{}}
$\square$\ \textbf{(A)} Query 1 and Query 2 always produce the same number of rows. \\
\addlinespace[3pt]
$\square$\ \textbf{(B)} Query 2 and Query 3 always produce the same number of rows. \\
\addlinespace[3pt]
$\square$\ \textbf{(C)} Query 1 always produces the maximum number of rows among all queries. \\
\addlinespace[3pt]
$\square$\ \textbf{(D)} All rows from Query 1, Query 2, and Query 3 are always present in the result of Query 4.
\end{tabular}

\tcblower

\textbf{Official Key:} \textbf{\large D} \hfill \textit{\small [Easy]}
\vspace{2pt}\hrule height 0.3pt\vspace{3pt}
\textbf{Topic:} Databases \\
\textbf{Subtopic:} \textit{Relational model: relational algebra, tuple calculus, SQL}
\vspace{2pt}\hrule height 0.3pt\vspace{3pt}
\textbf{Guide / Hint:} \\
{\small Query 1 uses natural join (inner join), Query 2 uses left outer join, Query 3 uses right outer join, and Query 4 uses full outer join. A full outer join preserves all tuples from both relation $Book$ and relation $Stock$ (filling missing values with NULLs). Therefore, every tuple output by Query 1, Query 2, and Query 3 is guaranteed to appear in Query 4.}

\end{tcolorbox}
\vspace{6pt}\hrule height 0.35pt\vspace{6pt}
\noindent
\begin{tcolorbox}[
    blanker,
    sidebyside,
    sidebyside align=top,
    lefthand width=0.55\textwidth,
    top=4pt, bottom=4pt, left=0pt, right=0pt,
    middle=16pt,
    underlay={
        \draw[dashed, black!80, line width=0.8pt, dash pattern=on 4pt off 3pt]
            ([xshift=0.55\textwidth+8pt]frame.north west) --
            ([xshift=0.55\textwidth+8pt]frame.south west);
    }
]
\textbf{Q.181} \quad {\footnotesize\textsc{[MCQ $\cdot$ 2 Mark(s) $\cdot$ Neg: 0.66]}} \hfill {\footnotesize\textbf{GATE 2018 CS} \quad \tiny\texttt{2018\_CS\_GA\_Q41}}
\vspace{2pt}\hrule height 0.3pt\vspace{4pt}

Consider the relations $r(A, B)$ and $s(B, C)$, where $s.B$ is a primary key and $r.B$ is a foreign key referencing $s.B$. Consider the query:
\[ Q: \Pi\_{A, B, C}(r \bowtie \sigma\_{B < 5}(s)) \]
Which one of the following relational algebra queries is NOT equivalent to $Q$?

\par\vspace{6pt}\noindent
\begin{tabular}{@{}p{0.98\linewidth}@{}}
$\square$\ \textbf{(A)} $\Pi_{A, B, C}(\sigma_{B < 5}(r) \bowtie s)$ \\
\addlinespace[3pt]
$\square$\ \textbf{(B)} $\sigma_{B < 5}(\Pi_{A, B, C}(r \bowtie s))$ \\
\addlinespace[3pt]
$\square$\ \textbf{(C)} $\Pi_{A, B, C}(r \,\text{⟕}\, \sigma_{B < 5}(s))$ \\
\addlinespace[3pt]
$\square$\ \textbf{(D)} $\Pi_{A, B, C}(\sigma_{B < 5}(r \bowtie s))$
\end{tabular}

\tcblower

\textbf{Official Key:} \textbf{\large C} \hfill \textit{\small [Medium]}
\vspace{2pt}\hrule height 0.3pt\vspace{3pt}
\textbf{Topic:} Databases \\
\textbf{Subtopic:} \textit{Relational model: relational algebra, tuple calculus, SQL}
\vspace{2pt}\hrule height 0.3pt\vspace{3pt}
\textbf{Guide / Hint:} \\
{\small Query $Q$ inner joins $r$ with $\sigma_{B < 5}(s)$, which discards all tuples from $r$ where $B \ge 5$. Option (C) performs a left outer join between $r$ and $\sigma_{B < 5}(s)$, which retains all tuples from $r$ (filling missing $C$ values with NULL for $B \ge 5$), so its result is NOT equivalent to $Q$.}

\end{tcolorbox}
\vspace{6pt}\hrule height 0.35pt\vspace{6pt}
\noindent
\begin{tcolorbox}[
    blanker,
    sidebyside,
    sidebyside align=top,
    lefthand width=0.55\textwidth,
    top=4pt, bottom=4pt, left=0pt, right=0pt,
    middle=16pt,
    underlay={
        \draw[dashed, black!80, line width=0.8pt, dash pattern=on 4pt off 3pt]
            ([xshift=0.55\textwidth+8pt]frame.north west) --
            ([xshift=0.55\textwidth+8pt]frame.south west);
    }
]
\textbf{Q.182} \quad {\footnotesize\textsc{[MCQ $\cdot$ 2 Mark(s) $\cdot$ Neg: 0.66]}} \hfill {\footnotesize\textbf{GATE 2018 CS} \quad \tiny\texttt{2018\_CS\_GA\_Q42}}
\vspace{2pt}\hrule height 0.3pt\vspace{4pt}

Consider the following four relational schemas. For each schema, all non-trivial functional dependencies are listed. The underlined attributes denote candidate keys.
\begin{itemize}[leftmargin=*, nosep]
    \item[Schema I:] $R_1(A, B, C, D)$ with FDs: $A \to B, B \to C, C \to D$
    \item[Schema II:] $R_2(\underline{rollno}, \underline{courseid}, email)$ with FDs: $rollno \to email, email \to rollno$
    \item[Schema III:] $R_3(A, B, C)$ with FDs: $A \to B, B \to C$
    \item[Schema IV:] $R_4(A, B, C)$ with FDs: $AB \to C, C \to A$
\end{itemize}
Which one of the relational schemas is in 3NF but not in BCNF?

\par\vspace{6pt}\noindent
\begin{tabular}{@{}p{0.48\linewidth}@{\hspace{0.04\linewidth}}p{0.48\linewidth}@{}}
$\square$\ \textbf{(A)} Schema I & $\square$\ \textbf{(B)} Schema II \\
\addlinespace[3pt]
$\square$\ \textbf{(C)} Schema III & $\square$\ \textbf{(D)} Schema IV
\end{tabular}

\tcblower

\textbf{Official Key:} \textbf{\large B} \hfill \textit{\small [Medium]}
\vspace{2pt}\hrule height 0.3pt\vspace{3pt}
\textbf{Topic:} Databases \\
\textbf{Subtopic:} \textit{Integrity constraints, normal forms}
\vspace{2pt}\hrule height 0.3pt\vspace{3pt}
\textbf{Guide / Hint:} \\
{\small In Schema II, the functional dependencies are $rollno \to email$ and $email \to rollno$. The candidate keys are $\{rollno, courseid\}$ and $\{email, courseid\}$. In $email \to rollno$, the left side is not a superkey (violating BCNF), but the right side $rollno$ is a prime attribute (satisfying 3NF). Hence Schema II is in 3NF but not in BCNF.}

\end{tcolorbox}
\vspace{6pt}\hrule height 0.35pt\vspace{6pt}
\noindent
\begin{tcolorbox}[
    blanker,
    sidebyside,
    sidebyside align=top,
    lefthand width=0.55\textwidth,
    top=4pt, bottom=4pt, left=0pt, right=0pt,
    middle=16pt,
    underlay={
        \draw[dashed, black!80, line width=0.8pt, dash pattern=on 4pt off 3pt]
            ([xshift=0.55\textwidth+8pt]frame.north west) --
            ([xshift=0.55\textwidth+8pt]frame.south west);
    }
]
\textbf{Q.183} \quad {\footnotesize\textsc{[MCQ $\cdot$ 1 Mark(s) $\cdot$ Neg: 0.33]}} \hfill {\footnotesize\textbf{GATE 2022 CS} \quad \tiny\texttt{2022\_CS\_CS\_Q14}}
\vspace{2pt}\hrule height 0.3pt\vspace{4pt}

In a relational data model, which one of the following statements is TRUE?

\par\vspace{6pt}\noindent
\begin{tabular}{@{}p{0.98\linewidth}@{}}
$\square$\ \textbf{(A)} A relation with only two attributes is always in BCNF. \\
\addlinespace[3pt]
$\square$\ \textbf{(B)} If all attributes of a relation are prime attributes, then the relation is in BCNF. \\
\addlinespace[3pt]
$\square$\ \textbf{(C)} Every relation has at least one non-prime attribute. \\
\addlinespace[3pt]
$\square$\ \textbf{(D)} BCNF decompositions preserve functional dependencies.
\end{tabular}

\tcblower

\textbf{Official Key:} \textbf{\large A} \hfill \textit{\small [Easy]}
\vspace{2pt}\hrule height 0.3pt\vspace{3pt}
\textbf{Topic:} Databases \\
\textbf{Subtopic:} \textit{Integrity constraints, normal forms}
\vspace{2pt}\hrule height 0.3pt\vspace{3pt}
\textbf{Guide / Hint:} \\
{\small For any binary relation $R(A, B)$, the only possible non-trivial functional dependencies are $A \to B$ or $B \to A$. In either case, the left-hand side is a superkey of $R$, satisfying the definition of Boyce-Codd Normal Form (BCNF).}

\end{tcolorbox}
\vspace{6pt}\hrule height 0.35pt\vspace{6pt}
\noindent
\begin{tcolorbox}[
    blanker,
    sidebyside,
    sidebyside align=top,
    lefthand width=0.55\textwidth,
    top=4pt, bottom=4pt, left=0pt, right=0pt,
    middle=16pt,
    underlay={
        \draw[dashed, black!80, line width=0.8pt, dash pattern=on 4pt off 3pt]
            ([xshift=0.55\textwidth+8pt]frame.north west) --
            ([xshift=0.55\textwidth+8pt]frame.south west);
    }
]
\textbf{Q.184} \quad {\footnotesize\textsc{[MCQ $\cdot$ 2 Mark(s) $\cdot$ Neg: 0.66]}} \hfill {\footnotesize\textbf{GATE 2022 CS} \quad \tiny\texttt{2022\_CS\_CS\_Q39}}
\vspace{2pt}\hrule height 0.3pt\vspace{4pt}

Let $R_i(z)$ and $W_i(z)$ denote read and write operations on a data element $z$ by a transaction $T_i$, respectively. Consider the schedule $S$ with four transactions:
\[ S: R\_4(x) R\_2(x) R\_3(x) R\_1(y) W\_1(y) W\_2(x) W\_3(y) R\_4(y) \]
Which one of the following serial schedules is conflict equivalent to $S$?

\par\vspace{6pt}\noindent
\begin{tabular}{@{}p{0.48\linewidth}@{\hspace{0.04\linewidth}}p{0.48\linewidth}@{}}
$\square$\ \textbf{(A)} $T_1 \to T_3 \to T_4 \to T_2$ & $\square$\ \textbf{(B)} $T_1 \to T_4 \to T_3 \to T_2$ \\
\addlinespace[3pt]
$\square$\ \textbf{(C)} $T_4 \to T_1 \to T_3 \to T_2$ & $\square$\ \textbf{(D)} $T_3 \to T_1 \to T_4 \to T_2$
\end{tabular}

\tcblower

\textbf{Official Key:} \textbf{\large A} \hfill \textit{\small [Medium]}
\vspace{2pt}\hrule height 0.3pt\vspace{3pt}
\textbf{Topic:} Databases \\
\textbf{Subtopic:} \textit{Transactions and concurrency control}
\vspace{2pt}\hrule height 0.3pt\vspace{3pt}
\textbf{Guide / Hint:} \\
{\small Constructing the serialization graph from conflict operations in $S$: on $y$, $R_1(y), W_1(y) \to W_3(y) \to R_4(y)$ gives edges $T_1 \to T_3 \to T_4$. On $x$, $R_4(x), R_3(x) \to W_2(x)$ gives edges $T_4 \to T_2$ and $T_3 \to T_2$. The unique topological order is $T_1 \to T_3 \to T_4 \to T_2$.}

\end{tcolorbox}
\vspace{6pt}\hrule height 0.35pt\vspace{6pt}
\noindent
\begin{tcolorbox}[
    blanker,
    sidebyside,
    sidebyside align=top,
    lefthand width=0.55\textwidth,
    top=4pt, bottom=4pt, left=0pt, right=0pt,
    middle=16pt,
    underlay={
        \draw[dashed, black!80, line width=0.8pt, dash pattern=on 4pt off 3pt]
            ([xshift=0.55\textwidth+8pt]frame.north west) --
            ([xshift=0.55\textwidth+8pt]frame.south west);
    }
]
\textbf{Q.185} \quad {\footnotesize\textsc{[MCQ $\cdot$ 1 Mark(s) $\cdot$ Neg: 0.33]}} \hfill {\footnotesize\textbf{GATE 2023 CS} \quad \tiny\texttt{2023\_CS\_CS\_Q16}}
\vspace{2pt}\hrule height 0.3pt\vspace{4pt}

Which one of the options given below refers to the degree (or arity) of a relation in relational database systems?

\par\vspace{6pt}\noindent
\begin{tabular}{@{}p{0.98\linewidth}@{}}
$\square$\ \textbf{(A)} Number of attributes of its relation schema. \\
\addlinespace[3pt]
$\square$\ \textbf{(B)} Number of tuples stored in the relation. \\
\addlinespace[3pt]
$\square$\ \textbf{(C)} Number of entries in the relation. \\
\addlinespace[3pt]
$\square$\ \textbf{(D)} Number of distinct domains of its relation schema.
\end{tabular}

\tcblower

\textbf{Official Key:} \textbf{\large A} \hfill \textit{\small [Easy]}
\vspace{2pt}\hrule height 0.3pt\vspace{3pt}
\textbf{Topic:} Databases \\
\textbf{Subtopic:} \textit{Relational model: relational algebra, tuple calculus, SQL}
\vspace{2pt}\hrule height 0.3pt\vspace{3pt}
\textbf{Guide / Hint:} \\
{\small In relational database theory, the degree (or arity) of a relation is defined as the number of attributes (columns) in its relation schema, whereas the cardinality of a relation is the number of tuples (rows).}

\end{tcolorbox}
\vspace{6pt}\hrule height 0.35pt\vspace{6pt}
\noindent
\begin{tcolorbox}[
    blanker,
    sidebyside,
    sidebyside align=top,
    lefthand width=0.55\textwidth,
    top=4pt, bottom=4pt, left=0pt, right=0pt,
    middle=16pt,
    underlay={
        \draw[dashed, black!80, line width=0.8pt, dash pattern=on 4pt off 3pt]
            ([xshift=0.55\textwidth+8pt]frame.north west) --
            ([xshift=0.55\textwidth+8pt]frame.south west);
    }
]
\textbf{Q.186} \quad {\footnotesize\textsc{[MCQ $\cdot$ 1 Mark(s) $\cdot$ Neg: 0.33]}} \hfill {\footnotesize\textbf{GATE 2024 CS1} \quad \tiny\texttt{2024\_CS1\_CS\_Q20}}
\vspace{2pt}\hrule height 0.3pt\vspace{4pt}

Let $S$ be the specification: ``Instructors teach courses. Students register for courses. Courses are allocated classrooms. Instructors guide students.'' Which one of the following ER diagrams CORRECTLY represents $S$?

\par\vspace{6pt}\noindent
\begin{tabular}{@{}p{0.48\linewidth}@{\hspace{0.04\linewidth}}p{0.48\linewidth}@{}}
$\square$\ \textbf{(A)} (i) & $\square$\ \textbf{(B)} (ii) \\
\addlinespace[3pt]
$\square$\ \textbf{(C)} (iii) & $\square$\ \textbf{(D)} (iv)
\end{tabular}

\tcblower

\textbf{Official Key:} \textbf{\large D} \hfill \textit{\small [Medium]}
\vspace{2pt}\hrule height 0.3pt\vspace{3pt}
\textbf{Topic:} Databases \\
\textbf{Subtopic:} \textit{ER-model}
\vspace{2pt}\hrule height 0.3pt\vspace{3pt}
\textbf{Guide / Hint:} \\
{\small Specification $S$ describes four distinct binary relationships among entities: Instructors teach courses, Students register for courses, Courses are allocated classrooms, and Instructors guide students. In standard ER modeling, entity sets are represented by rectangles and relationships by diamond diamonds connecting them; diagram (iv) correctly captures all four relationship sets.}

\end{tcolorbox}
\vspace{6pt}\hrule height 0.35pt\vspace{6pt}
\noindent
\begin{tcolorbox}[
    blanker,
    sidebyside,
    sidebyside align=top,
    lefthand width=0.55\textwidth,
    top=4pt, bottom=4pt, left=0pt, right=0pt,
    middle=16pt,
    underlay={
        \draw[dashed, black!80, line width=0.8pt, dash pattern=on 4pt off 3pt]
            ([xshift=0.55\textwidth+8pt]frame.north west) --
            ([xshift=0.55\textwidth+8pt]frame.south west);
    }
]
\textbf{Q.187} \quad {\footnotesize\textsc{[MCQ $\cdot$ 1 Mark(s) $\cdot$ Neg: 0.33]}} \hfill {\footnotesize\textbf{GATE 2024 CS1} \quad \tiny\texttt{2024\_CS1\_CS\_Q21}}
\vspace{2pt}\hrule height 0.3pt\vspace{4pt}

In a $\text{B}^+$ tree, the requirement of at least half-full (50\%) node occupancy is relaxed for which one of the following cases?

\par\vspace{6pt}\noindent
\begin{tabular}{@{}p{0.48\linewidth}@{\hspace{0.04\linewidth}}p{0.48\linewidth}@{}}
$\square$\ \textbf{(A)} Only the root node & $\square$\ \textbf{(B)} All leaf nodes \\
\addlinespace[3pt]
$\square$\ \textbf{(C)} All internal nodes & $\square$\ \textbf{(D)} Only the leftmost leaf node
\end{tabular}

\tcblower

\textbf{Official Key:} \textbf{\large A} \hfill \textit{\small [Easy]}
\vspace{2pt}\hrule height 0.3pt\vspace{3pt}
\textbf{Topic:} Databases \\
\textbf{Subtopic:} \textit{File organization, indexing (e.g., B and B+ trees)}
\vspace{2pt}\hrule height 0.3pt\vspace{3pt}
\textbf{Guide / Hint:} \\
{\small In a $\text{B}^+$ tree, every node except the root node must satisfy the minimum occupancy constraint of being at least half full ($	ext{at least } \lceil p/2 \rceil$ child pointers). The root node is the sole exception, and is permitted to hold as few as two pointers (or one key) when non-leaf, and one key when leaf.}

\end{tcolorbox}
\vspace{6pt}\hrule height 0.35pt\vspace{6pt}
\noindent
\begin{tcolorbox}[
    blanker,
    sidebyside,
    sidebyside align=top,
    lefthand width=0.55\textwidth,
    top=4pt, bottom=4pt, left=0pt, right=0pt,
    middle=16pt,
    underlay={
        \draw[dashed, black!80, line width=0.8pt, dash pattern=on 4pt off 3pt]
            ([xshift=0.55\textwidth+8pt]frame.north west) --
            ([xshift=0.55\textwidth+8pt]frame.south west);
    }
]
\textbf{Q.188} \quad {\footnotesize\textsc{[MCQ $\cdot$ 1 Mark(s) $\cdot$ Neg: 0.33]}} \hfill {\footnotesize\textbf{GATE 2024 CS2} \quad \tiny\texttt{2024\_CS2\_CS\_Q19}}
\vspace{2pt}\hrule height 0.3pt\vspace{4pt}

Once the DBMS informs the user that a transaction has been successfully completed, its effect should persist even if the system crashes before all its changes are reflected on disk. This property is called:

\par\vspace{6pt}\noindent
\begin{tabular}{@{}p{0.48\linewidth}@{\hspace{0.04\linewidth}}p{0.48\linewidth}@{}}
$\square$\ \textbf{(A)} durability & $\square$\ \textbf{(B)} atomicity \\
\addlinespace[3pt]
$\square$\ \textbf{(C)} consistency & $\square$\ \textbf{(D)} isolation
\end{tabular}

\tcblower

\textbf{Official Key:} \textbf{\large A} \hfill \textit{\small [Easy]}
\vspace{2pt}\hrule height 0.3pt\vspace{3pt}
\textbf{Topic:} Databases \\
\textbf{Subtopic:} \textit{Transactions and concurrency control}
\vspace{2pt}\hrule height 0.3pt\vspace{3pt}
\textbf{Guide / Hint:} \\
{\small The Durability property of the ACID paradigm guarantees that once a transaction has committed successfully, its updates and effects survive permanently in non-volatile storage, even in the event of a subsequent system crash.}

\end{tcolorbox}
\vspace{6pt}\hrule height 0.35pt\vspace{6pt}
\noindent
\begin{tcolorbox}[
    blanker,
    sidebyside,
    sidebyside align=top,
    lefthand width=0.55\textwidth,
    top=4pt, bottom=4pt, left=0pt, right=0pt,
    middle=16pt,
    underlay={
        \draw[dashed, black!80, line width=0.8pt, dash pattern=on 4pt off 3pt]
            ([xshift=0.55\textwidth+8pt]frame.north west) --
            ([xshift=0.55\textwidth+8pt]frame.south west);
    }
]
\textbf{Q.189} \quad {\footnotesize\textsc{[MCQ $\cdot$ 1 Mark(s) $\cdot$ Neg: 0.33]}} \hfill {\footnotesize\textbf{GATE 2024 CS2} \quad \tiny\texttt{2024\_CS2\_CS\_Q20}}
\vspace{2pt}\hrule height 0.3pt\vspace{4pt}

In the context of owner and weak entity sets in the ER (Entity-Relationship) data model, which one of the following statements is TRUE?

\par\vspace{6pt}\noindent
\begin{tabular}{@{}p{0.98\linewidth}@{}}
$\square$\ \textbf{(A)} The weak entity set MUST have total participation in the identifying relationship \\
\addlinespace[3pt]
$\square$\ \textbf{(B)} The owner entity set MUST have total participation in the identifying relationship \\
\addlinespace[3pt]
$\square$\ \textbf{(C)} Both weak and owner entity sets MUST have total participation in the identifying relationship \\
\addlinespace[3pt]
$\square$\ \textbf{(D)} Neither weak entity set nor owner entity set MUST have total participation in the identifying relationship
\end{tabular}

\tcblower

\textbf{Official Key:} \textbf{\large A} \hfill \textit{\small [Easy]}
\vspace{2pt}\hrule height 0.3pt\vspace{3pt}
\textbf{Topic:} Databases \\
\textbf{Subtopic:} \textit{ER-model}
\vspace{2pt}\hrule height 0.3pt\vspace{3pt}
\textbf{Guide / Hint:} \\
{\small A weak entity cannot exist without its identifying (owner) entity. Therefore, every instance of the weak entity set must be associated with an owner entity instance via the identifying relationship, mandating total participation of the weak entity set.}

\end{tcolorbox}
\vspace{6pt}\hrule height 0.35pt\vspace{6pt}
\noindent
\begin{tcolorbox}[
    blanker,
    sidebyside,
    sidebyside align=top,
    lefthand width=0.55\textwidth,
    top=4pt, bottom=4pt, left=0pt, right=0pt,
    middle=16pt,
    underlay={
        \draw[dashed, black!80, line width=0.8pt, dash pattern=on 4pt off 3pt]
            ([xshift=0.55\textwidth+8pt]frame.north west) --
            ([xshift=0.55\textwidth+8pt]frame.south west);
    }
]
\textbf{Q.190} \quad {\footnotesize\textsc{[MCQ $\cdot$ 2 Mark(s) $\cdot$ Neg: 0.66]}} \hfill {\footnotesize\textbf{GATE 2024 CS2} \quad \tiny\texttt{2024\_CS2\_CS\_Q45}}
\vspace{2pt}\hrule height 0.3pt\vspace{4pt}

The relation schema, Person(pid, city), describes the city of residence for every person uniquely identified by pid. The following relational algebra operators are available: selection, projection, cross product, and rename.
To find the list of cities where at least 3 persons reside, what is the minimum number of cross product operations required?

\par\vspace{6pt}\noindent
\begin{tabular}{@{}p{0.48\linewidth}@{\hspace{0.04\linewidth}}p{0.48\linewidth}@{}}
$\square$\ \textbf{(A)} 1 & $\square$\ \textbf{(B)} 2 \\
\addlinespace[3pt]
$\square$\ \textbf{(C)} 3 & $\square$\ \textbf{(D)} 4
\end{tabular}

\tcblower

\textbf{Official Key:} \textbf{\large B} \hfill \textit{\small [Medium]}
\vspace{2pt}\hrule height 0.3pt\vspace{3pt}
\textbf{Topic:} Databases \\
\textbf{Subtopic:} \textit{Relational model: relational algebra, tuple calculus, SQL}
\vspace{2pt}\hrule height 0.3pt\vspace{3pt}
\textbf{Guide / Hint:} \\
{\small To find cities with at least 3 distinct persons residing in them using selection, projection, cross-product, and rename, we need 3 copies of the Person table ($P_1, P_2, P_3$). Combining 3 relations requires 2 cross-product operations: $(P_1 \times P_2) \times P_3$.}

\end{tcolorbox}
\vspace{6pt}\hrule height 0.35pt\vspace{6pt}
\noindent
\begin{tcolorbox}[
    blanker,
    sidebyside,
    sidebyside align=top,
    lefthand width=0.55\textwidth,
    top=4pt, bottom=4pt, left=0pt, right=0pt,
    middle=16pt,
    underlay={
        \draw[dashed, black!80, line width=0.8pt, dash pattern=on 4pt off 3pt]
            ([xshift=0.55\textwidth+8pt]frame.north west) --
            ([xshift=0.55\textwidth+8pt]frame.south west);
    }
]
\textbf{Q.191} \quad {\footnotesize\textsc{[MCQ $\cdot$ 1 Mark(s) $\cdot$ Neg: 0.33]}} \hfill {\footnotesize\textbf{GATE 2025 CS1} \quad \tiny\texttt{2025\_CS1\_CS\_Q15}}
\vspace{2pt}\hrule height 0.3pt\vspace{4pt}

A schedule of three database transactions $T_1, T_2$, and $T_3$ is shown. $R_i(A)$ and $W_i(A)$ denote read and write of data item $A$ by transaction $T_i$, $i= 1,2,3$. The transaction $T_1$ aborts at the end. Which other transaction(s) will be required to be rolled back?
\[ R\_1(X) \; W\_1(Y) \; R\_2(X) \; R\_2(Y) \; R\_3(Y) \; ABORT(T\_1) \]

\par\vspace{6pt}\noindent
\begin{tabular}{@{}p{0.48\linewidth}@{\hspace{0.04\linewidth}}p{0.48\linewidth}@{}}
$\square$\ \textbf{(A)} Only $T_2$ & $\square$\ \textbf{(B)} Only $T_3$ \\
\addlinespace[3pt]
$\square$\ \textbf{(C)} Both $T_2$ and $T_3$ & $\square$\ \textbf{(D)} Neither $T_2$ nor $T_3$
\end{tabular}

\tcblower

\textbf{Official Key:} \textbf{\large C} \hfill \textit{\small [Medium]}
\vspace{2pt}\hrule height 0.3pt\vspace{3pt}
\textbf{Topic:} Databases \\
\textbf{Subtopic:} \textit{Transactions and concurrency control}
\vspace{2pt}\hrule height 0.3pt\vspace{3pt}
\textbf{Guide / Hint:} \\
{\small Transaction $T_1$ writes to item $Y$ ($W_1(Y)$). Subsequently, both $T_2$ and $T_3$ read uncommitted data item $Y$ ($R_2(Y)$ and $R_3(Y)$) written by $T_1$ before $T_1$ commits. When $T_1$ aborts, a cascading rollback is triggered, requiring both $T_2$ and $T_3$ to be rolled back.}

\end{tcolorbox}
\vspace{6pt}\hrule height 0.35pt\vspace{6pt}
\noindent
\begin{tcolorbox}[
    blanker,
    sidebyside,
    sidebyside align=top,
    lefthand width=0.55\textwidth,
    top=4pt, bottom=4pt, left=0pt, right=0pt,
    middle=16pt,
    underlay={
        \draw[dashed, black!80, line width=0.8pt, dash pattern=on 4pt off 3pt]
            ([xshift=0.55\textwidth+8pt]frame.north west) --
            ([xshift=0.55\textwidth+8pt]frame.south west);
    }
]
\textbf{Q.192} \quad {\footnotesize\textsc{[MCQ $\cdot$ 2 Mark(s) $\cdot$ Neg: 0.66]}} \hfill {\footnotesize\textbf{GATE 2025 CS1} \quad \tiny\texttt{2025\_CS1\_CS\_Q39}}
\vspace{2pt}\hrule height 0.3pt\vspace{4pt}

Consider two relations describing $teams$ and $players$ in a sports league:
\begin{itemize}
\item $teams(tid, tname)$: $tid, tname$ are team-id and team-name, respectively
\item $players(pid, pname, tid)$: $pid, pname$, and $tid$ denote player-id, player-name and the team-id of the player, respectively
\end{itemize}
Which ONE of the following tuple relational calculus queries returns the name of the players who play for the team having $tname$ as 'MI'?

\par\vspace{6pt}\noindent
\begin{tabular}{@{}p{0.98\linewidth}@{}}
$\square$\ \textbf{(A)} $\{ p.pname \mid p \in players \land \exists t (t \in teams \land p.tid = t.tid \land t.tname = \text{'MI'}) \}$ \\
\addlinespace[3pt]
$\square$\ \textbf{(B)} $\{ p.pname \mid p \in teams \land \exists t (t \in players \land p.tid = t.tid \land t.tname = \text{'MI'}) \}$ \\
\addlinespace[3pt]
$\square$\ \textbf{(C)} $\{ p.pname \mid p \in players \land \exists t (t \in teams \land t.tname = \text{'MI'}) \}$ \\
\addlinespace[3pt]
$\square$\ \textbf{(D)} $\{ p.pname \mid p \in teams \land \exists t (t \in players \land t.tname = \text{'MI'}) \}$
\end{tabular}

\tcblower

\textbf{Official Key:} \textbf{\large A} \hfill \textit{\small [Medium]}
\vspace{2pt}\hrule height 0.3pt\vspace{3pt}
\textbf{Topic:} Databases \\
\textbf{Subtopic:} \textit{Relational model: relational algebra, tuple calculus, SQL}
\vspace{2pt}\hrule height 0.3pt\vspace{3pt}
\textbf{Guide / Hint:} \\
{\small The query must select `p.pname` from the `players` relation ($p \in players$) and join with a tuple $t \in teams$ on the foreign key condition $p.tid = t.tid$, while filtering for $t.tname = \text{'MI'}$. Option A accurately expresses this quantified condition in Tuple Relational Calculus.}

\end{tcolorbox}
\vspace{6pt}\hrule height 0.35pt\vspace{6pt}
\vspace{8pt}
\subsection*{Computer Networks}
\vspace{-4pt}
\noindent
\begin{tcolorbox}[
    blanker,
    sidebyside,
    sidebyside align=top,
    lefthand width=0.55\textwidth,
    top=4pt, bottom=4pt, left=0pt, right=0pt,
    middle=16pt,
    underlay={
        \draw[dashed, black!80, line width=0.8pt, dash pattern=on 4pt off 3pt]
            ([xshift=0.55\textwidth+8pt]frame.north west) --
            ([xshift=0.55\textwidth+8pt]frame.south west);
    }
]
\textbf{Q.193} \quad {\footnotesize\textsc{[MCQ $\cdot$ 1 Mark(s) $\cdot$ Neg: 0.33]}} \hfill {\footnotesize\textbf{GATE 2018 CS} \quad \tiny\texttt{2018\_CS\_GA\_Q13}}
\vspace{2pt}\hrule height 0.3pt\vspace{4pt}

Match the following:
\begin{center}
\begin{tabular}{ll@{\hspace{2cm}}ll}
\textbf{Field} & & \textbf{Length in bits} & \\
P. & UDP Header's Port Number & I. & 48 \\
Q. & Ethernet MAC Address & II. & 8 \\
R. & IPv6 Next Header & III. & 32 \\
S. & TCP Header's Sequence Number & IV. & 16
\end{tabular}
\end{center}

\par\vspace{6pt}\noindent
\begin{tabular}{@{}p{0.48\linewidth}@{\hspace{0.04\linewidth}}p{0.48\linewidth}@{}}
$\square$\ \textbf{(A)} P-III, Q-I, R-II, S-IV & $\square$\ \textbf{(B)} P-IV, Q-II, R-I, S-III \\
\addlinespace[3pt]
$\square$\ \textbf{(C)} P-IV, Q-I, R-II, S-III & $\square$\ \textbf{(D)} P-I, Q-IV, R-II, S-III
\end{tabular}

\tcblower

\textbf{Official Key:} \textbf{\large C} \hfill \textit{\small [Easy]}
\vspace{2pt}\hrule height 0.3pt\vspace{3pt}
\textbf{Topic:} Computer Networks \\
\textbf{Subtopic:} \textit{Concept of layering: OSI and TCP/IP Protocol Stacks}
\vspace{2pt}\hrule height 0.3pt\vspace{3pt}
\textbf{Guide / Hint:} \\
{\small UDP port numbers are 16 bits long (P-IV). Ethernet MAC addresses are 48 bits long (Q-I). The IPv6 Next Header field is 8 bits long (R-II). TCP sequence numbers are 32 bits long (S-III). Hence the correct matching is P-IV, Q-I, R-II, S-III.}

\end{tcolorbox}
\vspace{6pt}\hrule height 0.35pt\vspace{6pt}
\noindent
\begin{tcolorbox}[
    blanker,
    sidebyside,
    sidebyside align=top,
    lefthand width=0.55\textwidth,
    top=4pt, bottom=4pt, left=0pt, right=0pt,
    middle=16pt,
    underlay={
        \draw[dashed, black!80, line width=0.8pt, dash pattern=on 4pt off 3pt]
            ([xshift=0.55\textwidth+8pt]frame.north west) --
            ([xshift=0.55\textwidth+8pt]frame.south west);
    }
]
\textbf{Q.194} \quad {\footnotesize\textsc{[MCQ $\cdot$ 1 Mark(s) $\cdot$ Neg: 0.33]}} \hfill {\footnotesize\textbf{GATE 2018 CS} \quad \tiny\texttt{2018\_CS\_GA\_Q14}}
\vspace{2pt}\hrule height 0.3pt\vspace{4pt}

Consider the following statements regarding the slow start phase of the TCP congestion control algorithm. Note that $cwnd$ stands for the congestion window size and $MSS$ denotes the Maximum Segment Size.
\begin{itemize}[leftmargin=*, nosep]
    \item[(i)] The $cwnd$ increases by 2 $MSS$ on every successful acknowledgement.
    \item[(ii)] The $cwnd$ increases by 1 $MSS$ on every successful acknowledgement.
    \item[(iii)] The $cwnd$ increases by 1 $MSS$ every round trip time.
    \item[(iv)] The $cwnd$ approximately doubles every round trip time.
\end{itemize}
Which of the following statements is/are CORRECT?

\par\vspace{6pt}\noindent
\begin{tabular}{@{}p{0.48\linewidth}@{\hspace{0.04\linewidth}}p{0.48\linewidth}@{}}
$\square$\ \textbf{(A)} (i) and (iii) only & $\square$\ \textbf{(B)} (i) and (iv) only \\
\addlinespace[3pt]
$\square$\ \textbf{(C)} (ii) and (iv) only & $\square$\ \textbf{(D)} (ii) and (iii) only
\end{tabular}

\tcblower

\textbf{Official Key:} \textbf{\large C} \hfill \textit{\small [Medium]}
\vspace{2pt}\hrule height 0.3pt\vspace{3pt}
\textbf{Topic:} Computer Networks \\
\textbf{Subtopic:} \textit{Transport layer: flow control and congestion control, UDP, TCP, sockets}
\vspace{2pt}\hrule height 0.3pt\vspace{3pt}
\textbf{Guide / Hint:} \\
{\small In TCP slow start, the congestion window (`cwnd`) increases by 1 MSS for each acknowledged segment. Consequently, when all segments in a flight are acknowledged, `cwnd` approximately doubles every round-trip time (RTT). Thus, statement (iv) is TRUE.}

\end{tcolorbox}
\vspace{6pt}\hrule height 0.35pt\vspace{6pt}
\noindent
\begin{tcolorbox}[
    blanker,
    sidebyside,
    sidebyside align=top,
    lefthand width=0.55\textwidth,
    top=4pt, bottom=4pt, left=0pt, right=0pt,
    middle=16pt,
    underlay={
        \draw[dashed, black!80, line width=0.8pt, dash pattern=on 4pt off 3pt]
            ([xshift=0.55\textwidth+8pt]frame.north west) --
            ([xshift=0.55\textwidth+8pt]frame.south west);
    }
]
\textbf{Q.195} \quad {\footnotesize\textsc{[MCQ $\cdot$ 1 Mark(s) $\cdot$ Neg: 0.33]}} \hfill {\footnotesize\textbf{GATE 2022 CS} \quad \tiny\texttt{2022\_CS\_CS\_Q22}}
\vspace{2pt}\hrule height 0.3pt\vspace{4pt}

Consider an enterprise network with two Ethernet segments, a web server and a firewall, connected via three routers. What is the number of subnets inside the enterprise network?

\par\vspace{6pt}\noindent
\begin{tabular}{@{}p{0.48\linewidth}@{\hspace{0.04\linewidth}}p{0.48\linewidth}@{}}
$\square$\ \textbf{(A)} 3 & $\square$\ \textbf{(B)} 12 \\
\addlinespace[3pt]
$\square$\ \textbf{(C)} 6 & $\square$\ \textbf{(D)} 8
\end{tabular}

\tcblower

\textbf{Official Key:} \textbf{\large C} \hfill \textit{\small [Medium]}
\vspace{2pt}\hrule height 0.3pt\vspace{3pt}
\textbf{Topic:} Computer Networks \\
\textbf{Subtopic:} \textit{Concept of layering: OSI and TCP/IP Protocol Stacks}
\vspace{2pt}\hrule height 0.3pt\vspace{3pt}
\textbf{Guide / Hint:} \\
{\small In the enterprise topology, subnets are formed by distinct broadcast domains separated by router interfaces: 2 Ethernet LAN segments (2 subnets), 1 web server subnet (1 subnet), 1 internal firewall segment (1 subnet), and 2 point-to-point links between the routers (2 subnets), yielding a total of 6 subnets.}

\end{tcolorbox}
\vspace{6pt}\hrule height 0.35pt\vspace{6pt}
\noindent
\begin{tcolorbox}[
    blanker,
    sidebyside,
    sidebyside align=top,
    lefthand width=0.55\textwidth,
    top=4pt, bottom=4pt, left=0pt, right=0pt,
    middle=16pt,
    underlay={
        \draw[dashed, black!80, line width=0.8pt, dash pattern=on 4pt off 3pt]
            ([xshift=0.55\textwidth+8pt]frame.north west) --
            ([xshift=0.55\textwidth+8pt]frame.south west);
    }
]
\textbf{Q.196} \quad {\footnotesize\textsc{[MCQ $\cdot$ 1 Mark(s) $\cdot$ Neg: 0.33]}} \hfill {\footnotesize\textbf{GATE 2023 CS} \quad \tiny\texttt{2023\_CS\_CS\_Q17}}
\vspace{2pt}\hrule height 0.3pt\vspace{4pt}

Suppose two hosts are connected by a point-to-point link and they are configured to use Stop-and-Wait protocol for reliable data transfer. Identify in which one of the following scenarios, the utilization of the link is the lowest.

\par\vspace{6pt}\noindent
\begin{tabular}{@{}p{0.98\linewidth}@{}}
$\square$\ \textbf{(A)} Longer link length and lower transmission rate \\
\addlinespace[3pt]
$\square$\ \textbf{(B)} Longer link length and higher transmission rate \\
\addlinespace[3pt]
$\square$\ \textbf{(C)} Shorter link length and lower transmission rate \\
\addlinespace[3pt]
$\square$\ \textbf{(D)} Shorter link length and higher transmission rate
\end{tabular}

\tcblower

\textbf{Official Key:} \textbf{\large B} \hfill \textit{\small [Medium]}
\vspace{2pt}\hrule height 0.3pt\vspace{3pt}
\textbf{Topic:} Computer Networks \\
\textbf{Subtopic:} \textit{Transport layer: flow control and congestion control, UDP, TCP, sockets}
\vspace{2pt}\hrule height 0.3pt\vspace{3pt}
\textbf{Guide / Hint:} \\
{\small Link utilization in Stop-and-Wait is $\eta = \frac{1}{1 + 2a}$, where $a = \frac{T_{\text{prop}}}{T_{\text{trans}}} = \frac{d/v}{L/B} = \frac{d \cdot B}{v \cdot L}$. Utilization $\eta$ is lowest when parameter $a$ is maximized, which occurs when the link length $d$ is longer (increasing propagation delay) and the transmission rate $B$ is higher (decreasing packet transmission time).}

\end{tcolorbox}
\vspace{6pt}\hrule height 0.35pt\vspace{6pt}
\noindent
\begin{tcolorbox}[
    blanker,
    sidebyside,
    sidebyside align=top,
    lefthand width=0.55\textwidth,
    top=4pt, bottom=4pt, left=0pt, right=0pt,
    middle=16pt,
    underlay={
        \draw[dashed, black!80, line width=0.8pt, dash pattern=on 4pt off 3pt]
            ([xshift=0.55\textwidth+8pt]frame.north west) --
            ([xshift=0.55\textwidth+8pt]frame.south west);
    }
]
\textbf{Q.197} \quad {\footnotesize\textsc{[MCQ $\cdot$ 1 Mark(s) $\cdot$ Neg: 0.33]}} \hfill {\footnotesize\textbf{GATE 2024 CS1} \quad \tiny\texttt{2024\_CS1\_CS\_Q16}}
\vspace{2pt}\hrule height 0.3pt\vspace{4pt}

A user starts browsing a webpage hosted at a remote server. The browser opens a single TCP connection to fetch the entire webpage from the server. The webpage consists of a top-level index page with multiple embedded image objects. Assume that all caches (e.g., DNS cache, browser cache) are all initially empty. The following packets leave the user's computer in some order:
\begin{enumerate}[label=(\roman*),itemsep=0pt,topsep=2pt]
\item HTTP GET request for the index page
\item DNS request to resolve the web server's name to its IP address
\item HTTP GET request for an image object
\item TCP SYN to open a connection to the web server
\end{enumerate}
Which one of the following is the CORRECT chronological order (earliest in time to latest) of the packets leaving the computer?

\par\vspace{6pt}\noindent
\begin{tabular}{@{}p{0.48\linewidth}@{\hspace{0.04\linewidth}}p{0.48\linewidth}@{}}
$\square$\ \textbf{(A)} (iv), (ii), (iii), (i) & $\square$\ \textbf{(B)} (ii), (iv), (iii), (i) \\
\addlinespace[3pt]
$\square$\ \textbf{(C)} (ii), (iv), (i), (iii) & $\square$\ \textbf{(D)} (iv), (ii), (i), (iii)
\end{tabular}

\tcblower

\textbf{Official Key:} \textbf{\large C} \hfill \textit{\small [Easy]}
\vspace{2pt}\hrule height 0.3pt\vspace{3pt}
\textbf{Topic:} Computer Networks \\
\textbf{Subtopic:} \textit{Concept of layering: OSI and TCP/IP Protocol Stacks}
\vspace{2pt}\hrule height 0.3pt\vspace{3pt}
\textbf{Guide / Hint:} \\
{\small Initially all caches are empty. The client must first send a DNS query (ii) to resolve the server hostname to an IP. Next, it establishes a TCP connection by sending TCP SYN (iv). Once connected, it sends an HTTP GET request for the index HTML page (i). After parsing the HTML and identifying embedded images, it sends HTTP GET requests for the images (iii).}

\end{tcolorbox}
\vspace{6pt}\hrule height 0.35pt\vspace{6pt}
\noindent
\begin{tcolorbox}[
    blanker,
    sidebyside,
    sidebyside align=top,
    lefthand width=0.55\textwidth,
    top=4pt, bottom=4pt, left=0pt, right=0pt,
    middle=16pt,
    underlay={
        \draw[dashed, black!80, line width=0.8pt, dash pattern=on 4pt off 3pt]
            ([xshift=0.55\textwidth+8pt]frame.north west) --
            ([xshift=0.55\textwidth+8pt]frame.south west);
    }
]
\textbf{Q.198} \quad {\footnotesize\textsc{[MCQ $\cdot$ 2 Mark(s) $\cdot$ Neg: 0.66]}} \hfill {\footnotesize\textbf{GATE 2024 CS1} \quad \tiny\texttt{2024\_CS1\_CS\_Q36}}
\vspace{2pt}\hrule height 0.3pt\vspace{4pt}

Consider a network path P---Q---R between nodes P and R via router Q. Node P sends a file of size $10^6$ bytes to R via this path by splitting the file into chunks of $10^3$ bytes each. Node P sends these chunks one after the other without any wait time between the successive chunk transmissions. Assume that the size of extra headers added to these chunks is negligible, and that the chunk size is less than the MTU.
\par\medskip
Each of the links P---Q and Q---R has a bandwidth of $10^6$ bits/sec, and negligible propagation latency. Router Q immediately transmits every packet it receives from P to R, with negligible processing and queueing delays. Router Q can simultaneously receive on link P---Q and transmit on link Q---R.
\par\medskip
Assume P starts transmitting the chunks at time $t = 0$.
Which one of the following options gives the time (in seconds, rounded off to 3 decimal places) at which R receives all the chunks of the file?

\par\vspace{6pt}\noindent
\begin{tabular}{@{}p{0.48\linewidth}@{\hspace{0.04\linewidth}}p{0.48\linewidth}@{}}
$\square$\ \textbf{(A)} 8.000 & $\square$\ \textbf{(B)} 8.008 \\
\addlinespace[3pt]
$\square$\ \textbf{(C)} 15.992 & $\square$\ \textbf{(D)} 16.000
\end{tabular}

\tcblower

\textbf{Official Key:} \textbf{\large B} \hfill \textit{\small [Medium]}
\vspace{2pt}\hrule height 0.3pt\vspace{3pt}
\textbf{Topic:} Computer Networks \\
\textbf{Subtopic:} \textit{Concept of layering: OSI and TCP/IP Protocol Stacks}
\vspace{2pt}\hrule height 0.3pt\vspace{3pt}
\textbf{Guide / Hint:} \\
{\small The file size is $10^6$ bytes split into 1000 chunks of $10^3$ bytes ($8000$ bits) each. Link transmission time for one chunk is $\frac{8000\text{ bits}}{10^6\text{ bps}} = 0.008\text{ s}$. Node P completes sending all 1000 chunks at $t = 1000 \times 0.008 = 8.000\text{ s}$. The last chunk arrives at router Q at $t = 8.000\text{ s}$ and takes another $0.008\text{ s}$ to reach R over link Q---R, completing at $8.008\text{ s}$.}

\end{tcolorbox}
\vspace{6pt}\hrule height 0.35pt\vspace{6pt}
\noindent
\begin{tcolorbox}[
    blanker,
    sidebyside,
    sidebyside align=top,
    lefthand width=0.55\textwidth,
    top=4pt, bottom=4pt, left=0pt, right=0pt,
    middle=16pt,
    underlay={
        \draw[dashed, black!80, line width=0.8pt, dash pattern=on 4pt off 3pt]
            ([xshift=0.55\textwidth+8pt]frame.north west) --
            ([xshift=0.55\textwidth+8pt]frame.south west);
    }
]
\textbf{Q.199} \quad {\footnotesize\textsc{[MCQ $\cdot$ 2 Mark(s) $\cdot$ Neg: 0.66]}} \hfill {\footnotesize\textbf{GATE 2024 CS2} \quad \tiny\texttt{2024\_CS2\_CS\_Q38}}
\vspace{2pt}\hrule height 0.3pt\vspace{4pt}

Which one of the following CIDR prefixes exactly represents the range of IP addresses 10.12.2.0 to 10.12.3.255?

\par\vspace{6pt}\noindent
\begin{tabular}{@{}p{0.48\linewidth}@{\hspace{0.04\linewidth}}p{0.48\linewidth}@{}}
$\square$\ \textbf{(A)} 10.12.2.0/23 & $\square$\ \textbf{(B)} 10.12.2.0/24 \\
\addlinespace[3pt]
$\square$\ \textbf{(C)} 10.12.0.0/22 & $\square$\ \textbf{(D)} 10.12.2.0/22
\end{tabular}

\tcblower

\textbf{Official Key:} \textbf{\large A} \hfill \textit{\small [Easy]}
\vspace{2pt}\hrule height 0.3pt\vspace{3pt}
\textbf{Topic:} Computer Networks \\
\textbf{Subtopic:} \textit{Fragmentation and IP addressing, IPv4, CIDR notation}
\vspace{2pt}\hrule height 0.3pt\vspace{3pt}
\textbf{Guide / Hint:} \\
{\small The address range spans from $10.12.2.0$ ($00000010_2$) to $10.12.3.255$ ($00000011_2$). The first 23 bits ($8 + 8 + 7 = 23$) are identical ($10.12.2$). Total addresses in range $= 2 \times 256 = 512 = 2^{32 - 23}$, which exactly matches CIDR prefix `10.12.2.0/23`.}

\end{tcolorbox}
\vspace{6pt}\hrule height 0.35pt\vspace{6pt}
\noindent
\begin{tcolorbox}[
    blanker,
    sidebyside,
    sidebyside align=top,
    lefthand width=0.55\textwidth,
    top=4pt, bottom=4pt, left=0pt, right=0pt,
    middle=16pt,
    underlay={
        \draw[dashed, black!80, line width=0.8pt, dash pattern=on 4pt off 3pt]
            ([xshift=0.55\textwidth+8pt]frame.north west) --
            ([xshift=0.55\textwidth+8pt]frame.south west);
    }
]
\textbf{Q.200} \quad {\footnotesize\textsc{[MCQ $\cdot$ 1 Mark(s) $\cdot$ Neg: 0.33]}} \hfill {\footnotesize\textbf{GATE 2025 CS1} \quad \tiny\texttt{2025\_CS1\_CS\_Q16}}
\vspace{2pt}\hrule height 0.3pt\vspace{4pt}

Identify the ONE CORRECT matching between the OSI layers and their corresponding functionalities as shown.
\par\medskip
\begin{tabular}{ll}
\textbf{OSI Layers} & \textbf{Functionalities} \\
(a) Network layer & (I) Packet routing \\
(b) Transport layer & (II) Framing and error handling \\
(c) Datalink layer & (III) Host to host communication
\end{tabular}

\par\vspace{6pt}\noindent
\begin{tabular}{@{}p{0.48\linewidth}@{\hspace{0.04\linewidth}}p{0.48\linewidth}@{}}
$\square$\ \textbf{(A)} (a)-(I), (b)-(II), (c)-(III) & $\square$\ \textbf{(B)} (a)-(I), (b)-(III), (c)-(II) \\
\addlinespace[3pt]
$\square$\ \textbf{(C)} (a)-(II), (b)-(I), (c)-(III) & $\square$\ \textbf{(D)} (a)-(III), (b)-(II), (c)-(I)
\end{tabular}

\tcblower

\textbf{Official Key:} \textbf{\large B} \hfill \textit{\small [Easy]}
\vspace{2pt}\hrule height 0.3pt\vspace{3pt}
\textbf{Topic:} Computer Networks \\
\textbf{Subtopic:} \textit{Concept of layering: OSI and TCP/IP Protocol Stacks}
\vspace{2pt}\hrule height 0.3pt\vspace{3pt}
\textbf{Guide / Hint:} \\
{\small The Network layer is responsible for packet routing across networks ((a)-(I)). The Transport layer provides end-to-end communication services between hosts ((b)-(III)). The Data Link layer handles framing, flow control, and physical error handling ((c)-(II)).}

\end{tcolorbox}
\vspace{6pt}\hrule height 0.35pt\vspace{6pt}
\noindent
\begin{tcolorbox}[
    blanker,
    sidebyside,
    sidebyside align=top,
    lefthand width=0.55\textwidth,
    top=4pt, bottom=4pt, left=0pt, right=0pt,
    middle=16pt,
    underlay={
        \draw[dashed, black!80, line width=0.8pt, dash pattern=on 4pt off 3pt]
            ([xshift=0.55\textwidth+8pt]frame.north west) --
            ([xshift=0.55\textwidth+8pt]frame.south west);
    }
]
\textbf{Q.201} \quad {\footnotesize\textsc{[MCQ $\cdot$ 2 Mark(s) $\cdot$ Neg: 0.66]}} \hfill {\footnotesize\textbf{GATE 2025 CS1} \quad \tiny\texttt{2025\_CS1\_CS\_Q40}}
\vspace{2pt}\hrule height 0.3pt\vspace{4pt}

A packet with the destination IP address 145.36.109.70 arrives at a router whose routing table is shown below:
\par\medskip
\begin{tabular}{|l|l|c|}
\hline
\textbf{Subnet Address} & \textbf{Subnet Mask (CIDR)} & \textbf{Interface} \\
\hline
145.36.0.0 & /16 & E1 \\
145.36.128.0 & /17 & E2 \\
145.36.64.0 & /18 & E3 \\
145.36.255.0 & /24 & E4 \\
Default & -- & E5 \\
\hline
\end{tabular}
\par\medskip
Which interface will the packet be forwarded to?

\par\vspace{6pt}\noindent
\begin{tabular}{@{}p{0.48\linewidth}@{\hspace{0.04\linewidth}}p{0.48\linewidth}@{}}
$\square$\ \textbf{(A)} E3 & $\square$\ \textbf{(B)} E1 \\
\addlinespace[3pt]
$\square$\ \textbf{(C)} E2 & $\square$\ \textbf{(D)} E5
\end{tabular}

\tcblower

\textbf{Official Key:} \textbf{\large A} \hfill \textit{\small [Medium]}
\vspace{2pt}\hrule height 0.3pt\vspace{3pt}
\textbf{Topic:} Computer Networks \\
\textbf{Subtopic:} \textit{Routing protocols: shortest path, flooding, distance vector and link state routing}
\vspace{2pt}\hrule height 0.3pt\vspace{3pt}
\textbf{Guide / Hint:} \\
{\small The destination address has third octet $109 = 01101101_2$. For /16, $145.36.0.0$ matches. For /17 ($128 = 10000000_2$), the leading bit is 0, so /17 does not match. For /18 ($64 = 01000000_2$), the first two bits are $01_2$, which matches $109$ ($01101101_2$). For /24, 255 does not match. By Longest Prefix Match, /18 (interface E3) is chosen.}

\end{tcolorbox}
\vspace{6pt}\hrule height 0.35pt\vspace{6pt}
\noindent
\begin{tcolorbox}[
    blanker,
    sidebyside,
    sidebyside align=top,
    lefthand width=0.55\textwidth,
    top=4pt, bottom=4pt, left=0pt, right=0pt,
    middle=16pt,
    underlay={
        \draw[dashed, black!80, line width=0.8pt, dash pattern=on 4pt off 3pt]
            ([xshift=0.55\textwidth+8pt]frame.north west) --
            ([xshift=0.55\textwidth+8pt]frame.south west);
    }
]
\textbf{Q.202} \quad {\footnotesize\textsc{[MCQ $\cdot$ 1 Mark(s) $\cdot$ Neg: 0.33]}} \hfill {\footnotesize\textbf{GATE 2025 CS2} \quad \tiny\texttt{2025\_CS2\_CS\_Q16}}
\vspace{2pt}\hrule height 0.3pt\vspace{4pt}

Consider the following statements:
(i) Address Resolution Protocol (ARP) provides a mapping from an IP address to the corresponding hardware (link-layer) address.
(ii) A single TCP segment from a sender S to a receiver R cannot carry both data from S to R and acknowledgement for a segment from R to S.
Which ONE of the following is CORRECT?

\par\vspace{6pt}\noindent
\begin{tabular}{@{}p{0.48\linewidth}@{\hspace{0.04\linewidth}}p{0.48\linewidth}@{}}
$\square$\ \textbf{(A)} Both (i) and (ii) are TRUE & $\square$\ \textbf{(B)} (i) is TRUE and (ii) is FALSE \\
\addlinespace[3pt]
$\square$\ \textbf{(C)} (i) is FALSE and (ii) is TRUE & $\square$\ \textbf{(D)} Both (i) and (ii) are FALSE
\end{tabular}

\tcblower

\textbf{Official Key:} \textbf{\large B} \hfill \textit{\small [Easy]}
\vspace{2pt}\hrule height 0.3pt\vspace{3pt}
\textbf{Topic:} Computer Networks \\
\textbf{Subtopic:} \textit{Basics of IP support protocols: ARP and TCP headers}
\vspace{2pt}\hrule height 0.3pt\vspace{3pt}
\textbf{Guide / Hint:} \\
{\small ARP resolves an IPv4 address to a physical MAC (link-layer) address, so (i) is true. A TCP segment supports piggybacking, where data and acknowledgement flags/numbers are sent in the same segment, so (ii) is false.}

\end{tcolorbox}
\vspace{6pt}\hrule height 0.35pt\vspace{6pt}
\noindent
\begin{tcolorbox}[
    blanker,
    sidebyside,
    sidebyside align=top,
    lefthand width=0.55\textwidth,
    top=4pt, bottom=4pt, left=0pt, right=0pt,
    middle=16pt,
    underlay={
        \draw[dashed, black!80, line width=0.8pt, dash pattern=on 4pt off 3pt]
            ([xshift=0.55\textwidth+8pt]frame.north west) --
            ([xshift=0.55\textwidth+8pt]frame.south west);
    }
]
\textbf{Q.203} \quad {\footnotesize\textsc{[MCQ $\cdot$ 1 Mark(s) $\cdot$ Neg: 0.33]}} \hfill {\footnotesize\textbf{GATE 2025 CS2} \quad \tiny\texttt{2025\_CS2\_CS\_Q17}}
\vspace{2pt}\hrule height 0.3pt\vspace{4pt}

Consider the routing protocols given in List I and the names given in List II:
\par\medskip
\begin{center}\small
\begin{tabular}{ll}
\toprule
List I & List II \\
\midrule
(i) Distance vector routing & (a) Bellman-Ford \\
(ii) Link state routing & (b) Dijkstra \\
\bottomrule
\end{tabular}\end{center}
\par\medskip
For matching of items in List I with those in List II, which ONE of the following options is CORRECT?

\par\vspace{6pt}\noindent
\begin{tabular}{@{}p{0.48\linewidth}@{\hspace{0.04\linewidth}}p{0.48\linewidth}@{}}
$\square$\ \textbf{(A)} (i) – (a) and (ii) – (b) & $\square$\ \textbf{(B)} (i) – (a) and (ii) – (a) \\
\addlinespace[3pt]
$\square$\ \textbf{(C)} (i) – (b) and (ii) – (a) & $\square$\ \textbf{(D)} (i) – (b) and (ii) – (b)
\end{tabular}

\tcblower

\textbf{Official Key:} \textbf{\large A} \hfill \textit{\small [Easy]}
\vspace{2pt}\hrule height 0.3pt\vspace{3pt}
\textbf{Topic:} Computer Networks \\
\textbf{Subtopic:} \textit{Routing protocols: distance vector and link state}
\vspace{2pt}\hrule height 0.3pt\vspace{3pt}
\textbf{Guide / Hint:} \\
{\small Distance Vector routing is based on the distributed Bellman-Ford algorithm, whereas Link State routing uses Dijkstra's shortest path algorithm.}

\end{tcolorbox}
\vspace{6pt}\hrule height 0.35pt\vspace{6pt}
\noindent
\begin{tcolorbox}[
    blanker,
    sidebyside,
    sidebyside align=top,
    lefthand width=0.55\textwidth,
    top=4pt, bottom=4pt, left=0pt, right=0pt,
    middle=16pt,
    underlay={
        \draw[dashed, black!80, line width=0.8pt, dash pattern=on 4pt off 3pt]
            ([xshift=0.55\textwidth+8pt]frame.north west) --
            ([xshift=0.55\textwidth+8pt]frame.south west);
    }
]
\textbf{Q.204} \quad {\footnotesize\textsc{[MCQ $\cdot$ 1 Mark(s) $\cdot$ Neg: 0.33]}} \hfill {\footnotesize\textbf{GATE 2025 CS2} \quad \tiny\texttt{2025\_CS2\_CS\_Q18}}
\vspace{2pt}\hrule height 0.3pt\vspace{4pt}

A machine receives an IPv4 datagram. The protocol field of the IPv4 header has the protocol number of a protocol X. Which ONE of the following is NOT a possible candidate for X?

\par\vspace{6pt}\noindent
\begin{tabular}{@{}p{0.98\linewidth}@{}}
$\square$\ \textbf{(A)} Internet Control Message Protocol (ICMP) \\
\addlinespace[3pt]
$\square$\ \textbf{(B)} Internet Group Management Protocol (IGMP) \\
\addlinespace[3pt]
$\square$\ \textbf{(C)} Open Shortest Path First (OSPF) \\
\addlinespace[3pt]
$\square$\ \textbf{(D)} Routing Information Protocol (RIP)
\end{tabular}

\tcblower

\textbf{Official Key:} \textbf{\large D} \hfill \textit{\small [Medium]}
\vspace{2pt}\hrule height 0.3pt\vspace{3pt}
\textbf{Topic:} Computer Networks \\
\textbf{Subtopic:} \textit{Concept of layering: OSI and TCP/IP protocol stacks}
\vspace{2pt}\hrule height 0.3pt\vspace{3pt}
\textbf{Guide / Hint:} \\
{\small The IPv4 Protocol field identifies the payload carried directly in the IP datagram: ICMP (protocol 1), IGMP (protocol 2), and OSPF (protocol 89). In contrast, RIP runs as an application-layer process encapsulated inside UDP (port 520).}

\end{tcolorbox}
\vspace{6pt}\hrule height 0.35pt\vspace{6pt}
\noindent
\begin{tcolorbox}[
    blanker,
    sidebyside,
    sidebyside align=top,
    lefthand width=0.55\textwidth,
    top=4pt, bottom=4pt, left=0pt, right=0pt,
    middle=16pt,
    underlay={
        \draw[dashed, black!80, line width=0.8pt, dash pattern=on 4pt off 3pt]
            ([xshift=0.55\textwidth+8pt]frame.north west) --
            ([xshift=0.55\textwidth+8pt]frame.south west);
    }
]
\textbf{Q.205} \quad {\footnotesize\textsc{[MCQ $\cdot$ 1 Mark(s) $\cdot$ Neg: 0.33]}} \hfill {\footnotesize\textbf{GATE 2025 CS2} \quad \tiny\texttt{2025\_CS2\_CS\_Q23}}
\vspace{2pt}\hrule height 0.3pt\vspace{4pt}

Consider a network that uses Ethernet and IPv4. Assume that IPv4 headers do not use any options field. Each Ethernet frame can carry a maximum of 1500 bytes in its data field. A UDP segment is transmitted. The payload (data) in the UDP segment is 7488 bytes.
Which ONE of the following choices has the CORRECT total number of fragments transmitted and the size of the last fragment including IPv4 header?

\par\vspace{6pt}\noindent
\begin{tabular}{@{}p{0.48\linewidth}@{\hspace{0.04\linewidth}}p{0.48\linewidth}@{}}
$\square$\ \textbf{(A)} 5 fragments, 1488 bytes & $\square$\ \textbf{(B)} 6 fragments, 88 bytes \\
\addlinespace[3pt]
$\square$\ \textbf{(C)} 6 fragments, 108 bytes & $\square$\ \textbf{(D)} 6 fragments, 116 bytes
\end{tabular}

\tcblower

\textbf{Official Key:} \textbf{\large D} \hfill \textit{\small [Medium]}
\vspace{2pt}\hrule height 0.3pt\vspace{3pt}
\textbf{Topic:} Computer Networks \\
\textbf{Subtopic:} \textit{Fragmentation and IP addressing}
\vspace{2pt}\hrule height 0.3pt\vspace{3pt}
\textbf{Guide / Hint:} \\
{\small UDP payload is 7488 bytes + 8 bytes UDP header = 7496 bytes IP payload. Max IP payload per Ethernet fragment is $\lfloor(1500 - 20)/8\rfloor \times 8 = 1480$ bytes. Fragment sizes are 5 fragments of 1480 bytes ($5 \times 1480 = 7400$), leaving $7496 - 7400 = 96$ bytes for the 6th fragment, so total size is $96 + 20 = 116$ bytes.}

\end{tcolorbox}
\vspace{6pt}\hrule height 0.35pt\vspace{6pt}
\noindent
\begin{tcolorbox}[
    blanker,
    sidebyside,
    sidebyside align=top,
    lefthand width=0.55\textwidth,
    top=4pt, bottom=4pt, left=0pt, right=0pt,
    middle=16pt,
    underlay={
        \draw[dashed, black!80, line width=0.8pt, dash pattern=on 4pt off 3pt]
            ([xshift=0.55\textwidth+8pt]frame.north west) --
            ([xshift=0.55\textwidth+8pt]frame.south west);
    }
]
\textbf{Q.206} \quad {\footnotesize\textsc{[MCQ $\cdot$ 2 Mark(s) $\cdot$ Neg: 0.66]}} \hfill {\footnotesize\textbf{GATE 2025 CS2} \quad \tiny\texttt{2025\_CS2\_CS\_Q36}}
\vspace{2pt}\hrule height 0.3pt\vspace{4pt}

Suppose we are transmitting frames between two nodes using Stop-and-Wait protocol. The frame size is 3000 bits. Which ONE of the following gives the efficiency of the channel?

\par\vspace{6pt}\noindent
\begin{tabular}{@{}p{0.48\linewidth}@{\hspace{0.04\linewidth}}p{0.48\linewidth}@{}}
$\square$\ \textbf{(A)} 88.23 & $\square$\ \textbf{(B)} 93.75 \\
\addlinespace[3pt]
$\square$\ \textbf{(C)} 85.44 & $\square$\ \textbf{(D)} 66.67
\end{tabular}

\tcblower

\textbf{Official Key:} \textbf{\large A} \hfill \textit{\small [Medium]}
\vspace{2pt}\hrule height 0.3pt\vspace{3pt}
\textbf{Topic:} Computer Networks \\
\textbf{Subtopic:} \textit{Data link layer: flow control and Stop-and-Wait efficiency}
\vspace{2pt}\hrule height 0.3pt\vspace{3pt}
\textbf{Guide / Hint:} \\
{\small Efficiency in Stop-and-Wait is $\eta = \frac{T_t}{T_t + 2 T_p}$. Substituting given transmission and propagation delays yields $\eta = 88.23\%$.}

\end{tcolorbox}
\vspace{6pt}\hrule height 0.35pt\vspace{6pt}
\end{document}